% Generated mechanically from frozen input inputs/p01/ja/sites.tex.
% Do not edit this generated file; edit the bound locale source instead.

% OK, start here.
%

\setcounter{chapter}{6}
\chapter{サイトと層}



\phantomsection
\label{sites-section-phantom}


\section{序論}
\label{sites-section-introduction}

\noindent
サイトという概念は、スキームのエタール位相における層を研究できるようにするため、
Grothendieck によって導入された。この概念の基本的な参考文献はおそらく
\cite{SGA4} である。本章でいうサイトは \cite{SGA4} のものとは異なる。
本章でサイトと呼ぶものは、\cite[Expos\'e II, D\'efinition 1.3]{SGA4}
では前位相を備えた圏と呼ばれる。このようにする理由は、代数幾何では、
例えばスキームの性質が与えられた位相について局所的であることを定義する際に、
指定された被覆の類を用いるのが便利なことが多いからである。Descent,
Section \ref{descent-section-descending-properties} を参照せよ。
本章の叙述は \cite{ArtinTopologies} に忠実に従う。
宇宙は用いない。









\section{前層}
\label{sites-section-presheaves}

\noindent
$\mathcal{C}$ を圏とする。
{\it 集合の前層}とは、$\mathcal{C}$ から $\textit{Sets}$ への
反変関手 $\mathcal{F}$ をいう（Categories, Remark
\ref{categories-remark-functor-into-sets} を参照せよ）。
したがって、$\mathcal{C}$ の各対象 $U$ に対して集合
$\mathcal{F}(U)$ がある。この集合の元を、$U$ 上の
$\mathcal{F}$ の{\it 切断}という。各射 $f : V \to U$ に対して、写像
$\mathcal{F}(f) : \mathcal{F}(U) \to \mathcal{F}(V)$ を
{\it 制限写像}といい、しばしば
$f^\ast : \mathcal{F}(U) \to \mathcal{F}(V)$ と書く。
別の言い方をすれば、$f^*(s)$ は $f$ による $s$ の
{\it 引き戻し}である。関手性は $g^* f^* (s) = (f \circ g)^*(s)$
を意味する。ときには記法 $s|_V := f^\ast(s)$ を用いる。
この記法は、位相における関数の制限という概念と整合する。
実際、$W \to V \to U$ が $\mathcal{C}$ の射で、$s$ が
$U$ 上の $\mathcal{F}$ の切断ならば、$\mathcal{F}$ の関手性により
$s|_W = (s|_V)|_W$ である。もちろん、射 $V \to U$ が複数存在することは
十分あり得て、それらに対応する引き戻し写像が等しくなるとは限らないので、
注意しなければならない。

\begin{definition}
\label{sites-definition-presheaves-sets}
$\mathcal{C}$ 上の{\it 集合の前層}とは、$\mathcal{C}$ から
$\textit{Sets}$ への反変関手をいう。{\it 前層の射}とは関手の変換をいう。
集合の前層の圏を $\textit{PSh}(\mathcal{C})$ と書く。
\end{definition}

\noindent
$\mathcal{C}$ の任意の対象 $U$ に対して、点関手 $h_U$
（Categories, Example \ref{categories-example-hom-functor} を参照せよ）は
前層であることに注意する。これらを{\it 表現可能前層}という。
これらの前層には、任意の前層 $\mathcal{F}$ に対して
\begin{equation}
\label{sites-equation-map-representable-into-presheaf}
\Mor_{\textit{PSh}(\mathcal{C})}(h_U, \mathcal{F})
=
\mathcal{F}(U).
\end{equation}
となるという好ましい性質がある。これは米田の補題である
（Categories, Lemma \ref{categories-lemma-yoneda}）。

\medskip\noindent
同様に、アーベル群、環などの前層という概念を定義できる。
より一般に、ある圏に値をもつ前層を定義できる。

\begin{definition}
\label{sites-definition-presheaf}
$\mathcal{C}$, $\mathcal{A}$ を圏とする。
$\mathcal{A}$ に値をもつ $\mathcal{C}$ 上の{\it 前層}
$\mathcal{F}$ とは、$\mathcal{C}$ から $\mathcal{A}$ への反変関手、
すなわち $\mathcal{F} : \mathcal{C}^{opp} \to \mathcal{A}$ をいう。
$\mathcal{A}$ に値をもつ $\mathcal{C}$ 上の前層の
{\it 射} $\mathcal{F} \to \mathcal{G}$ とは、$\mathcal{F}$ から
$\mathcal{G}$ への関手の変換をいう。
\end{definition}

\noindent
これらが、$\mathcal{A}$ に値をもつ $\mathcal{C}$ 上の前層の圏の
対象と射をなす。

\begin{remark}
\label{sites-remark-big-presheaves}
すでに指摘したように、Categories, Remark
\ref{categories-remark-big-categories} に挙げられた任意の「大きな」圏に
値をもつ前層の圏を考えることができる。これらも「大きな」圏となり、
議論を進めるにつれて上記の注記に列挙される。
\end{remark}













\section{前層の単射と全射}
\label{sites-section-injective-surjective}

\begin{definition}
\label{sites-definition-presheaves-injective-surjective}
$\mathcal{C}$ を圏とし、$\varphi : \mathcal{F}
\to \mathcal{G}$ を集合の前層の射とする。
\begin{enumerate}
\item $\mathcal{C}$ の各対象 $U$ に対して写像
$\varphi_U : \mathcal{F}(U) \to \mathcal{G}(U)$ が単射であるとき、
$\varphi$ は{\it 単射}であるという。
\item $\mathcal{C}$ の各対象 $U$ に対して写像
$\varphi_U : \mathcal{F}(U) \to \mathcal{G}(U)$ が全射であるとき、
$\varphi$ は{\it 全射}であるという。
\end{enumerate}
\end{definition}

\begin{lemma}
\label{sites-lemma-mono-epi}
上で定義した単射（それぞれ全射）は、ちょうど
$\textit{PSh}(\mathcal{C})$ のモノ射（それぞれエピ射）である。
射が同型であるための必要十分条件は、それが単射かつ全射であることである。
\end{lemma}

\begin{proof}
$\varphi : \mathcal{F} \to \mathcal{G}$ が単射であることと、
$\textit{PSh}(\mathcal{C})$ のモノ射であることが同値であることを示す。
実際、「単射ならモノ射」の向きは直ちに分かるので、逆を示そう。
$\varphi$ がモノ射であると仮定する。$U \in \Ob(\mathcal{C})$ とし、
$\varphi_U$ が単射であることを示さなければならない。
$a, b \in \mathcal{F}(U)$ が $\varphi_U (a) = \varphi_U (b)$ を
満たすとし、$a = b$ を確かめる。同型
(\ref{sites-equation-map-representable-into-presheaf}) のもとで、
$\mathcal{F}(U)$ の元 $a$, $b$ は二つの自然変換
$a', b' \in \Mor_{\textit{PSh}(\mathcal{C})}(h_U, \mathcal{F})$
に対応する。同様に、類似の同型
$\Mor_{\textit{PSh}(\mathcal{C})}(h_U, \mathcal{G})
= \mathcal{G}(U)$ のもとで、$\mathcal{G}(U)$ の二つの等しい元
$\varphi_U (a)$, $\varphi_U (b)$ は二つの自然変換
$\varphi \circ a', \varphi \circ b'
\in \Mor_{\textit{PSh}(\mathcal{C})}(h_U, \mathcal{G})$
に対応し、したがってこれらも等しい。ゆえに
$\varphi \circ a' = \varphi \circ b'$ であり、$\varphi$ はモノ射なので
$a' = b'$、したがって $a = b$ である。これで (1) が終わる。

\medskip\noindent
$\varphi : \mathcal{F} \to \mathcal{G}$ が全射であることと、
$\textit{PSh}(\mathcal{C})$ のエピ射であることが同値であることを示す。
実際、「全射ならエピ射」の向きは直ちに分かるので、逆を示そう。
$\varphi$ がエピ射であると仮定する。

\medskip\noindent
圏 $\textit{Sets}$ における任意の二つの射
$f : A \to B$ と $g : A \to C$ に対して、$B$ および $C$ から
$B \coprod_A C$ への二つの標準写像をそれぞれ
$\text{inl}_{f,g}$ および $\text{inr}_{f,g}$ と書く。
（ここで押し出しは $\textit{Sets}$ において計算する。）

\medskip\noindent
いま、$\mathcal{C}$ 上の集合の前層 $\mathcal{H}$ を、各
$U \in \Ob(\mathcal{C})$ に対して
$\mathcal{H}(U) = \mathcal{G}(U) \coprod_{\mathcal{F}(U)} \mathcal{G}(U)$
（ここで押し出しは $\textit{Sets}$ において、写像
$\varphi_U : \mathcal{F}(U) \to \mathcal{G}(U)$ により計算する）
と置くことで定義する。射への作用は（押し出しの関手性により）明らかな方法で
定める。すると二つの自然変換 $i_1 : \mathcal{G} \to \mathcal{H}$ と
$i_2 : \mathcal{G} \to \mathcal{H}$ があり、対象
$U \in \Ob(\mathcal{C})$ におけるその成分はそれぞれ
$\text{inl}_{\varphi_U, \varphi_U}$ と
$\text{inr}_{\varphi_U, \varphi_U}$ で与えられる。
押し出しの定義から $i_1 \circ \varphi = i_2 \circ \varphi$ であり、
$\varphi$ はエピ射なので $i_1 = i_2$ となる。したがって、各
$U \in \Ob(\mathcal{C})$ に対して
$\text{inl}_{\varphi_U, \varphi_U}
= \text{inr}_{\varphi_U, \varphi_U}$ である。ゆえに
$\varphi_U$ は全射でなければならない（実際、簡単な組合せ論的議論により、
$f : A \to B$ が $\textit{Sets}$ の射ならば、
$\text{inl}_{f,f} = \text{inr}_{f,f}$ であるための必要十分条件は
$f$ が全射であることである）。言い換えると $\varphi$ は全射であり、
(2) が証明された。

\medskip\noindent
$\varphi : \mathcal{F} \to \mathcal{G}$ が単射かつ全射であることと、
$\textit{PSh}(\mathcal{C})$ の同型であることが同値であることを示す。
今度は「同型なら単射かつ全射」の向きが直ちに分かる。
逆を示すには、$\varphi$ が単射かつ全射ならば、各
$U \in \Ob(\mathcal{C})$ に対して $\varphi_U$ は可逆な写像であり、
すべての $U$ に対するこれらの写像の逆を組み合わせると、
$\varphi$ の逆となる自然変換 $\mathcal{G} \to \mathcal{F}$ が
得られることに注意すれば十分である。
\end{proof}


\begin{definition}
\label{sites-definition-sub-presheaf}
各対象 $U \in \Ob(\mathcal{C})$ に対して集合 $\mathcal{F}(U)$ が
$\mathcal{G}(U)$ の部分集合であり、しかも制限写像と両立するとき、
$\mathcal{F}$ は $\mathcal{G}$ の{\it 部分前層}であるという。
\end{definition}

\noindent
言い換えると、包含写像 $\mathcal{F}(U) \to \mathcal{G}(U)$ は
貼り合わさって、前層の（単射な）射
$\mathcal{F} \to \mathcal{G}$ を与える。

\begin{lemma}
\label{sites-lemma-image}
$\mathcal{C}$ を圏とする。
$\varphi : \mathcal{F} \to \mathcal{G}$ を、$\mathcal{C}$ 上の
集合の前層の射とする。$\varphi$ が
$\mathcal{F} \to \mathcal{G}' \to \mathcal{G}$ と分解し、
かつ第一の写像が全射となるような部分前層
$\mathcal{G}' \subset \mathcal{G}$ が一意に存在する。
\end{lemma}

\begin{proof}
存在を示すには、各 $U \in \Ob(C)$ に対して
$\mathcal{G}'(U) = \varphi_U \left(\mathcal{F}(U)\right)$ と置き
（射への作用は $\mathcal{G}$ から引き継ぐ）、これが
$\mathcal{G}$ の部分前層を定め、$\varphi$ が
$\mathcal{F} \to \mathcal{G}' \to \mathcal{G}$ と分解し、
第一の写像が全射となることを証明すればよい。一意性は直ちに分かる。
\end{proof}

\begin{definition}
\label{sites-definition-image}
記法は Lemma \ref{sites-lemma-image} のとおりとする。
$\mathcal{G}'$ を $\varphi$ の{\it 像}という。
\end{definition}


\section{前層の極限と余極限}
\label{sites-section-limits-colimits-PSh}

\noindent
$\mathcal{C}$ を圏とする。
圏 $\textit{PSh}(\mathcal{C})$ には極限と余極限が存在する。
さらに、任意の $U \in \Ob(\mathcal{C})$ に対して関手
$$
\textit{PSh}(\mathcal{C})
\longrightarrow
\textit{Sets}, \quad
\mathcal{F}
\longmapsto
\mathcal{F}(U)
$$
は極限および余極限と可換である。これらの主張を証明する最も簡単な方法は、
おそらく次のものである。図式
$
\mathcal{F} :
\mathcal{I}
\to
\textit{PSh}(\mathcal{C})
$
が与えられたとき、前層を
$$
\mathcal{F}_{\lim} :
U
\longmapsto
\lim_{i \in \mathcal{I}} \mathcal{F}_i(U)
\text{  and  }
\mathcal{F}_{\colim} :
U
\longmapsto
\colim_{i \in \mathcal{I}} \mathcal{F}_i(U)
$$
で定める。明らかに射影写像 $\mathcal{F}_{\lim} \to \mathcal{F}_i$ と
標準写像 $\mathcal{F}_i \to \mathcal{F}_{\colim}$ がある。これらの写像は、
Categories, Definition \ref{categories-definition-limit}
（それぞれ Categories, Definition \ref{categories-definition-colimit}）における
極限（それぞれ余極限）の写像に求められる条件を満たす。
実際、これらは明らかに $\mathcal{F}$ 上の錐（それぞれ余錐）をなす。
さらに、$(\mathcal{G}, q_i : \mathcal{G} \to \mathcal{F}_i)$ が別の系
（極限の定義におけるもの）ならば、各 $U$ に対して、適切な関手性の条件を
備えた写像の系 $\mathcal{G}(U) \to \mathcal{F}_i(U)$ が得られる。
したがって、一意な写像 $\mathcal{G}(U) \to \mathcal{F}_{\lim}(U)$ が得られる。
$U$ を動かすとき、これらが両立し、求める写像
$\mathcal{G} \to \mathcal{F}_{\lim}$ から生じることは容易に確かめられる。
余極限の場合も同様の議論が成り立つ。




















\section{前層の圏の関手性}
\label{sites-section-functoriality-PSh}

\noindent
$u : \mathcal{C} \to \mathcal{D}$ を圏の間の関手とする。このとき
$$
u^p :
\textit{PSh}(\mathcal{D})
\longrightarrow
\textit{PSh}(\mathcal{C})
$$
によって、$\mathcal{D}$ 上の $\mathcal{G}$ に前層
$u^p\mathcal{G} = \mathcal{G} \circ u$ を対応させる関手を表す。
前節により、この関手はすべての極限および余極限と可換することに注意せよ。

\medskip\noindent
$V \in \Ob(\mathcal{D})$ に対し、次を満たす圏を $\mathcal{I}^u_V$ と書く：
\begin{equation}
\label{sites-equation-colim-category}
\begin{matrix}
\Ob(\mathcal{I}^u_V)
&
=
&
\{
(U, \phi)
\mid
U \in \Ob(\mathcal{C}),
\phi : V \to u(U)
\}
\\
\Mor_{\mathcal{I}^u_V}((U, \phi), (U', \phi'))
&
=
&
\{
f : U \to U' \text{ in }\mathcal{C}
\mid
u(f) \circ \phi = \phi'
\}
\end{matrix}
\end{equation}
上付き添字 ${}^u$ を記法から省き、単に $\mathcal{I}_V$ と書くこともある。
これらの圏を用いて、関手 $u^p$ の左随伴を定義する。
その前に、いくつかの技術的な補題を証明する。

\begin{lemma}
\label{sites-lemma-almost-directed}
$u : \mathcal{C} \to \mathcal{D}$ を圏の間の関手とする。
$\mathcal{C}$ がファイバー積と等化子をもち、$u$ がそれらと可換すると仮定する。
このとき圏 $(\mathcal{I}_V)^{opp}$ は
Categories, Lemma \ref{categories-lemma-split-into-directed}
の仮定を満たす。
\end{lemma}

\begin{proof}
確かめるべき条件は二つある。

\medskip\noindent
まず、三つの対象
$\phi : V \to u(U)$、$\phi' : V \to u(U')$、
$\phi'' : V \to u(U'')$ と、
$u(a) \circ \phi' = \phi$ および $u(b) \circ \phi'' = \phi$ を満たす射
$a : U' \to U$、$b : U'' \to U$ が与えられたとする。
別の対象 $\phi''' : V \to u(U''')$ と射
$c : U''' \to U'$、$d : U''' \to U''$ で、
$u(c) \circ \phi''' = \phi'$、$u(d) \circ \phi''' = \phi''$、かつ
$a \circ c = b \circ d$ を満たすものが存在することを示さなければならない。
$U''' = U' \times_U U''$ とし、$c$ と $d$ を射影とする。
$u$ はファイバー積と可換するので、これでよい。確認は省略する。

\medskip\noindent
次に、二つの対象 $\phi : V \to u(U)$ と $\phi' : V \to u(U')$、および射
$a, b : (U, \phi) \to (U', \phi')$ が与えられたとする。
$a$ と $b$ を等化する射 $c : (U'', \phi'') \to (U, \phi)$ を
見つけなければならない。$c : U'' \to U$ を圏 $\mathcal{C}$ における
$a$ と $b$ の等化子とする。$u$ は等化子と可換し、かつ
$u(a) \circ \phi = u(b) \circ \phi = \phi'$ であるから、
射 $\phi'' : V \to u(U'')$ が得られる。
\end{proof}

\begin{lemma}
\label{sites-lemma-directed}
$u : \mathcal{C} \to \mathcal{D}$ を圏の間の関手とする。次を仮定する。
\begin{enumerate}
\item 圏 $\mathcal{C}$ は終対象 $X$ をもち、$u(X)$ は
$\mathcal{D}$ の終対象である。
\item 圏 $\mathcal{C}$ はファイバー積をもち、$u$ はそれらと可換する。
\end{enumerate}
このとき添字圏 $(\mathcal{I}^u_V)^{opp}$ はフィルター付きである
（Categories, Definition \ref{categories-definition-directed} を参照せよ）。
\end{lemma}

\begin{proof}
これらの仮定から Lemma \ref{sites-lemma-almost-directed} の仮定が満たされる
（Categories, Section \ref{categories-section-finite-limits} の議論を参照せよ）。
Categories, Lemma \ref{categories-lemma-split-into-directed} により、
$\mathcal{I}_V$ は有向圏の（空であるかもしれない）非交和である。
したがって $\mathcal{I}_V$ が連結であることを示せば十分である。

\medskip\noindent
まず、$\mathcal{I}_V$ が空でないことを示す。
仮定により存在する $\mathcal{C}$ の終対象を $X$ とする。
$u(X)$ が仮定により $\mathcal{D}$ の終対象であることから得られる射
$V \to u(X)$ を取る。これは $\mathcal{I}_V$ の対象を与える。

\medskip\noindent
次に、$\mathcal{I}_V$ が連結であることを示す。
$\phi_1 : V \to u(U_1)$ と $\phi_2 : V \to u(U_2)$ を
$\Ob(\mathcal{I}_V)$ の対象とする。仮定により $U_1\times U_2$ が存在し、
$u(U_1\times U_2) = u(U_1)\times u(U_2)$ である。
積の普遍性により $(\phi_1, \phi_2)$ に対応する射
$\phi : V \to u(U_1\times U_2)$ を考える。
対象 $\phi : V \to u(U_1\times U_2)$ が
$\phi_1 : V \to u(U_1)$ と $\phi_2 : V \to u(U_2)$ の双方へ写ることは明らかである。
\end{proof}

\noindent
$\mathcal{D}$ の射 $g : V' \to V$ が与えられたとき、対象上
$\overline{g}(U, \phi) = (U, \phi \circ g)$ と置くことにより関手
$\overline{g} : \mathcal{I}_V \to \mathcal{I}_{V'}$ が得られる。
$\mathcal{C}$ 上の前層 $\mathcal{F}$ が与えられたとき、関手
$$
\mathcal{F}_V :
\mathcal{I}_V^{opp}
\longrightarrow
\textit{Sets}, \quad
(U, \phi)
\longmapsto
\mathcal{F}(U).
$$
が得られる。言い換えると、$\mathcal{F}_V$ は $\mathcal{I}_V$ 上の
集合の前層である。$\mathcal{F}_{V'} \circ \overline{g} = \mathcal{F}_V$
であることに注意せよ。次で定める：
$$
u_p\mathcal{F}(V) =
\colim_{\mathcal{I}_V^{opp}} \mathcal{F}_V
$$
余極限であるから、各 $(U, \phi) \in \Ob(\mathcal{I}_V)$ に対して
標準写像 $\mathcal{F}(U)\xrightarrow{c(\phi)}u_p\mathcal{F}(V)$ が得られる。
上のような $g : V' \to V$ に対して、Categories, Lemma
\ref{categories-lemma-functorial-colimit} により、
$\mathcal{F}_{V'} \circ \overline{g} = \mathcal{F}_V$ と両立する標準的な制限写像
$g^* : u_p\mathcal{F}(V) \to u_p\mathcal{F}(V')$ がある。
これは、すべての $(U, \phi) \in \Ob(\mathcal{I}_V)$ に対して図式
$$
\xymatrix{
\mathcal{F}(U) \ar[r]^{c(\phi)} \ar[d]_{\text{id}}
&
u_p\mathcal{F}(V) \ar[d]^{g^*}
\\
\mathcal{F}(U) \ar[r]^{c(\phi \circ g)}
&
u_p\mathcal{F}(V')
}
$$
を可換にする一意な写像である。これらの写像の一意性により前層が得られる。
この前層を $u_p\mathcal{F}$ と書く。

\begin{lemma}
\label{sites-lemma-recover}
標準写像 $\mathcal{F}(U) \to u_p\mathcal{F}(u(U))$ があり、
これは（$\mathcal{F}$ および $u_p\mathcal{F}$ の）制限写像と両立する。
\end{lemma}

\begin{proof}
これは上で導入した写像 $c(\text{id}_{u(U)})$ にほかならない。
\end{proof}

\noindent
前層の任意の写像 $\mathcal{F} \to \mathcal{F}'$ は、関手の間の両立する写像の系
$\mathcal{F}_V \to \mathcal{F}'_V$ を、したがって前層の写像
$u_p\mathcal{F} \to u_p\mathcal{F}'$ を生じさせることに注意せよ。
言い換えると、関手
$$
u_p :
\textit{PSh}(\mathcal{C})
\longrightarrow
\textit{PSh}(\mathcal{D})
$$
を定義したことになる。

\begin{lemma}
\label{sites-lemma-adjoints-u}
関手 $u_p$ は関手 $u^p$ の左随伴である。言い換えると、公式
$$
\Mor_{\textit{PSh}(\mathcal{C})}(\mathcal{F}, u^p\mathcal{G})
=
\Mor_{\textit{PSh}(\mathcal{D})}(u_p\mathcal{F}, \mathcal{G})
$$
は $\mathcal{F}$ と $\mathcal{G}$ に関して二関手的に成り立つ。
\end{lemma}

\begin{proof}
$\mathcal{G}$ を $\mathcal{D}$ 上の前層、$\mathcal{F}$ を $\mathcal{C}$ 上の前層とする。
両方向の写像を構成することで、表示された公式が成り立つことを示す。
それらが互いに逆であることの確認は読者に委ねる。

\medskip\noindent
写像 $\alpha : u_p \mathcal{F} \to \mathcal{G}$ が与えられると、
$u^p\alpha : u^p u_p \mathcal{F} \to u^p \mathcal{G}$ が得られる。
Lemma \ref{sites-lemma-recover} により写像
$\mathcal{F} \to u^p u_p \mathcal{F}$ がある。この二つの合成が
求める写像を与える。（この構成のよい点は、現れるすべてのものについて
明らかに関手的であることである。）

\medskip\noindent
逆に、写像 $\beta : \mathcal{F} \to u^p\mathcal{G}$ が与えられると、写像
$u_p\beta : u_p\mathcal{F} \to u_p u^p\mathcal{G}$ が得られる。
$V \in \Ob(\mathcal{D})$ とする。$\mathcal{I}_V$ 上の関手
$u^p\mathcal{G}_V$ から、値 $\mathcal{G}(V)$ をもつ定値関手への標準写像が
あると主張する。実際、$\mathcal{I}_V$ の各対象 $(U, \phi)$ における
$u^p\mathcal{G}_V$ の値は $\mathcal{G}(u(U))$ であり、これは
$\mathcal{G}(\phi) = \phi^*$ によって $\mathcal{G}(V)$ へ写る。
$\mathcal{G}$ 自身が関手なので、これは関手の変換である。
これにより写像 $u_p u^p \mathcal{G}(V) \to \mathcal{G}(V)$ が得られる。
さらに容易な確認により、これは $V$ に関して関手的であり、前層の写像
$u_p u^p \mathcal{G} \to \mathcal{G}$ を与える。合成
$u_p\mathcal{F} \to u_p u^p\mathcal{G} \to \mathcal{G}$ が求める写像である。
\end{proof}

\begin{remark}
\label{sites-remark-functoriality-presheaves-values}
$\mathcal{A}$ を、任意の図式 $\mathcal{I}_Y \to \mathcal{A}$ が
$\mathcal{A}$ に余極限をもつような圏と仮定する。このとき、
$\mathcal{A}$ に値をもつ前層の圏の上で、上とまったく同じ方法により定義される
関手 $u^p$ と $u_p$ があることは明らかである。
さらに、対 $u^p$ と $u_p$ の随伴性もこの状況で引き続き成り立つ。
\end{remark}

\begin{lemma}
\label{sites-lemma-pullback-representable-presheaf}
$u : \mathcal{C} \to \mathcal{D}$ を圏の間の関手とする。
$\mathcal{C}$ の任意の対象 $U$ に対して $u_ph_U = h_{u(U)}$ である。
\end{lemma}

\begin{proof}
$u_p$ と $u^p$ の随伴性により
$$
\Mor_{\textit{PSh}(\mathcal{D})}(u_ph_U, \mathcal{G})
=
\Mor_{\textit{PSh}(\mathcal{C})}(h_U, u^p\mathcal{G})
=
u^p\mathcal{G}(U) =
\mathcal{G}(u(U))
$$
である。したがって米田の補題により、前層として
$u_ph_U = h_{u(U)}$ であることが分かる。
\end{proof}

















\section{サイト}
\label{sites-section-sites-definitions}

\noindent
本章のサイトという概念には、次の種類の構造を用いる。

\begin{definition}
\label{sites-definition-family-morphisms-fixed-target}
$\mathcal{C}$ を圏とする。Conventions, Section
\ref{conventions-section-categories} を参照せよ。
$\mathcal{C}$ における{\it 終域を固定した射の族}とは、対象
$U \in \Ob(\mathcal{C})$、集合 $I$、および各 $i\in I$ に対する
$\mathcal{C}$ の射 $U_i \to U$ で終域が $U$ であるものからなるデータをいう。
これを $\{U_i \to U\}_{i\in I}$ と表す。
\end{definition}

\noindent
集合 $I$ が空であることもあり得る！この記法は、位相における
開被覆を想起させることを意図している。

\begin{definition}
\label{sites-definition-site}
{\it サイト}\footnote{この記法は序論で説明したように
\cite{SGA4} のものとは異なる。}とは、圏 $\mathcal{C}$ と、
{\it $\mathcal{C}$ の被覆}と呼ばれる終域を固定した射の族
$\{U_i \to U\}_{i \in I}$ の集合 $\text{Cov}(\mathcal{C})$ で、
次の公理を満たすものからなるデータをいう。
\begin{enumerate}
\item $V \to U$ が同型ならば、$\{V \to U\} \in
\text{Cov}(\mathcal{C})$ である。
\item $\{U_i \to U\}_{i\in I} \in \text{Cov}(\mathcal{C})$ であり、
各 $i$ に対して $\{V_{ij} \to U_i\}_{j\in J_i} \in
\text{Cov}(\mathcal{C})$ ならば、
$\{V_{ij} \to U\}_{i \in I, j\in J_i} \in \text{Cov}(\mathcal{C})$ である。
\item $\{U_i \to U\}_{i\in I}\in \text{Cov}(\mathcal{C})$ であり、
$V \to U$ が $\mathcal{C}$ の射ならば、すべての $i$ に対して
$U_i \times_U V$ が存在し、
$\{U_i \times_U V \to V \}_{i\in I} \in \text{Cov}(\mathcal{C})$ である。
\end{enumerate}
\end{definition}

\noindent
いくつか明確にしておく。公理 (1) では、$I = \{i\}$ が一元集合で、
射 $U_i \to U$ が (1) で与えられた射 $V \to U$ と等しいような
$\mathcal{C}$ の被覆 $\{U_i \to U\}_{i \in I}$ が存在することを要求する。
以下では、添字集合が一元集合である終域を固定した射の族を、しばしば
$\{V \to U\}$ と表す。公理 (3) では、$i \in I$ に対する
ファイバー積 $U_i \times_U V$ のある選択について被覆が存在することを要求する。

\begin{remark}
\label{sites-remark-no-big-sites}
（集合論上の問題について。初読時には飛ばしてよい。）
サイトを導入する主な理由は、サイト上の層の圏を研究するためである。
それは、代数幾何学で非常に重要な役割を果たしてきた位相空間上の層の圏を
一般化するものだからである。「類の類」などを考えることを避けるため、
圏や $2$-圏の場合とは異なり、サイトが「大きな」圏であることを許さない。

\medskip\noindent
$\mathcal{C}$ を圏とし、$\text{Cov}(\mathcal{C})$ を上の (1), (2), (3) を
満たす被覆の真の類と仮定する。これもサイトとしては許さない。
その主な理由は、被覆にわたる極限を取ることになるからである。
しかし、$\text{Cov}(\mathcal{C})$ を被覆の集合または少し異なる構造に
置き換え、同じ層の圏を得る自然な方法がいくつかある。例えば次のとおりである。
\begin{enumerate}
\item Sets, Section \ref{sets-section-coverings-site} では、同じ層の圏を
与える適切な被覆の集合を選ぶ方法を示す。
\item 別の方法は、付随する位相を取ることである
（Definition \ref{sites-definition-topology-associated-site} を参照せよ）。
$\mathcal{C}$ 上に得られる位相は同じ層の圏をもつ。
二つの位相が同じ層の圏をもつための必要十分条件は、それらが等しいことである。
Theorem \ref{sites-theorem-topology-and-topos} を参照せよ。
圏上の位相は、各対象上のふるいを選ぶことによって与えられる。
あらゆるふるいの集まり、さらには $\mathcal{C}$ 上のあらゆる位相の集まりも集合である。
\item サイトの概念を少し修正することもできる。下の Remark
\ref{sites-remark-shrink-coverings} を参照せよ。この方法では標準的な被覆の集合が得られる。
\end{enumerate}
これらの解決法にはそれぞれ小さな欠点がある。第一の方法では、後に行う構成が
被覆の集合の選択に依存しないことを確かめなければならない。
第二の方法では、位相について学び、サイトに関する多くの議論をやり直さなければならない。
第三の方法については Remark \ref{sites-remark-shrink-coverings} の最後の文を参照せよ。

\medskip\noindent
本章では、上の Definition \ref{sites-definition-site} の意味でのサイトを用いる。
上のような被覆の真の類をもつ圏 $\mathcal{C}$ が与えられたとき、
Sets, Lemma \ref{sets-lemma-coverings-site} を用いて、これをサイトを与える
被覆の集合に置き換える。下の Lemma
\ref{sites-lemma-choice-set-coverings-immaterial} において、こうして得られる
層の圏（トポス）がこの選択に依存しないことを示す。
望むならば、読者はこれらの問題を扱うために残り二つの方針のいずれかを用いてよい。
\end{remark}

\begin{example}
\label{sites-example-site-topological}
$X$ を位相空間とする。$X$ のすべての開集合 $U$ を対象とし、包含写像だけを
射とする圏を $X_{Zar}$ と書く。すなわち、$X_{Zar}$ の任意の二対象の間には
高々一つの射しかない。ここで
$\{U_i \to U\}_{i \in I}\in \text{Cov}(X_{Zar})$ を
$\bigcup U_i = U$ であることと同値であるように定める。
上の条件 (1) と (2) は明らかであり、$X_{Zar}$ において
$U \times V = U \cap V$ であることに気づけば (3) も明らかである。
特に空集合は $X_{Zar}$ の元でなければならないことに注意せよ。
そうでなければ、一般にはこの構成は成り立たない。さらに、後で見るように、
{\it 空集合の空被覆を被覆として許す}ことも同じく重要である！
Sets, Lemma \ref{sets-lemma-coverings-site} に従い、適切な被覆の集合
$\text{Cov}(X_{Zar})_{\kappa, \alpha}$ を選ぶことで、$X_{Zar}$ をサイトにする。
% P01-E0058
サイト $X_{Zar}$ 上の前層と層（以下で定義する）は、Sheaves, Section
\ref{sheaves-section-introduction} で定義された位相空間上の通常の前層と層の概念に
正確に一致する。
\end{example}

\begin{example}
\label{sites-example-site-on-group}
$G$ を群とする。左 $G$-作用をもつ集合 $X$ を対象とし、$G$-同変写像を射とする
圏 $G\textit{-Sets}$ を考える。重要な例は ${}_GG$ である。
これは台集合が $G$ で、作用が左乗法で与えられる $G$-集合である。
この圏はファイバー積をもつ。Categories, Section
\ref{categories-section-example-fibre-products} を参照せよ。
$\bigcup_{i\in I} \varphi_i(U_i) = U$ であるとき、かつそのときに限り、
$\{\varphi_i : U_i \to U\}_{i\in I}$ を被覆と定める。
これは $G\textit{-Sets}$ 上の被覆の類を与え、Definition
\ref{sites-definition-site} の条件 (1), (2), (3) を満たすことは容易に分かる。
基礎となる圏の対象の集まりも被覆の集まりも真の類をなすので、得られるものは
サイトではない。
% P01-E0059
まず、Sets, Lemma \ref{sets-lemma-sets-with-group-action} に従って
$G\textit{-Sets}$ を充満部分圏 $G\textit{-Sets}_\alpha$ に置き換える。
その後、Sets, Lemma \ref{sets-lemma-coverings-site} に従って被覆の集まりを
適切に制限して得られるサイト
$(G\textit{-Sets}_\alpha,
\text{Cov}_{\kappa, \alpha'}(G\textit{-Sets}_\alpha))$
を $\mathcal{T}_G$ と書く。

\medskip\noindent
特に、群 $G$ が可算ならば、$\mathcal{T}_G$ を可算 $G$-集合の圏とし、
被覆を、$I$ が可算であるような共同全射的な射の族
$\{\varphi_i : U_i \to U\}_{i \in I}$ とすることができる。
\end{example}

\begin{example}
\label{sites-example-indiscrete}
$\mathcal{C}$ を圏とする。$U$ の被覆を
$\{f : V \to U \mid f\text{ is an isomorphism}\}$ とすることにより、
これをサイトにする標準的な方法がある。
このサイト上の層は $\mathcal{C}$ 上の前層である。
% P01-E0060
対応する位相を{\it カオス位相}または{\it 密着位相}という。
\end{example}























\section{層}
\label{sites-section-sheaves}

\noindent
$\mathcal{C}$ をサイトとする。圏に値をもつ層という概念を導入する前に、
集合の前層が層であるとはどういうことかを説明する。
$\mathcal{F}$ を $\mathcal{C}$ 上の集合の前層とし、
$\{U_i \to U\}_{i\in I}$ を $\text{Cov}(\mathcal{C})$ の元とする。
仮定により、すべてのファイバー積 $U_i \times_U U_j$ が
$\mathcal{C}$ に存在する。二つの自然な写像
$$
\xymatrix{
\prod\nolimits_{i\in I}
\mathcal{F}(U_i)
\ar@<1ex>[r]^-{\text{pr}_0^*} \ar@<-1ex>[r]_-{\text{pr}_1^*}
&
\prod\nolimits_{(i_0, i_1) \in I \times I}
\mathcal{F}(U_{i_0} \times_U U_{i_1})
}
$$
があり、表示式に示したように $\text{pr}^*_i$, $i = 0, 1$ と書く。
実際、左辺の元は族 $(s_i)_{i\in I}$ に対応し、各 $s_i$ は
$U_i$ 上の $\mathcal{F}$ の切断である。各対
$(i_0, i_1) \in I \times I$ に対して射影
$$
\text{pr}^{(i_0, i_1)}_{i_0} :
U_{i_0} \times_U U_{i_1}
\longrightarrow
U_{i_0}
\text{ and }
\text{pr}^{(i_0, i_1)}_{i_1} :
U_{i_0} \times_U U_{i_1}
\longrightarrow
U_{i_1}.
$$
がある。したがって、切断 $s_{i_0}$ を第一の写像で引き戻すことも、
切断 $s_{i_1}$ を第二の写像で引き戻すこともできる。
具体的には、上で述べた写像は
$$
\text{pr}_0^* :
(s_i)_{i\in I}
\longmapsto
\Big(
\text{pr}^{(i_0, i_1), *}_{i_0}(s_{i_0})
\Big)_{(i_0, i_1) \in I \times I}
$$
および
$$
\text{pr}_1^* :
(s_i)_{i\in I}
\longmapsto
\Big(
\text{pr}^{(i_0, i_1), *}_{i_1}(s_{i_1})
\Big)_{(i_0, i_1) \in I \times I}.
$$
である。最後に、自然な写像
$$
\mathcal{F}(U)
\longrightarrow
\prod\nolimits_{i\in I}
\mathcal{F}(U_i), \quad
s
\longmapsto
(s|_{U_i})_{i \in I}
$$
を考える。ここで $s|_{U_i}$ は写像 $U_i \to U$ による $s$ の引き戻しを表す。
$\mathcal{F}$ の関手性とファイバー積図式の可換性から、
$\text{pr}_0^*( (s|_{U_i})_{i \in I} ) =
\text{pr}_1^*( (s|_{U_i})_{i \in I} )$ であることは明らかである。

\begin{definition}
\label{sites-definition-sheaf-sets}
$\mathcal{C}$ をサイトとし、$\mathcal{F}$ を $\mathcal{C}$ 上の集合の前層とする。
すべての被覆 $\{U_i \to U\}_{i \in I} \in \text{Cov}(\mathcal{C})$
に対して図式
\begin{equation}
\label{sites-equation-sheaf-condition}
\xymatrix{
\mathcal{F}(U) \ar[r]
&
\prod\nolimits_{i\in I}
\mathcal{F}(U_i)
\ar@<1ex>[r]^-{\text{pr}_0^*} \ar@<-1ex>[r]_-{\text{pr}_1^*}
&
\prod\nolimits_{(i_0, i_1) \in I \times I}
\mathcal{F}(U_{i_0} \times_U U_{i_1})
}
\end{equation}
が第一の矢を $\text{pr}_0^*$ と $\text{pr}_1^*$ の等化子として表すとき、
$\mathcal{F}$ は{\it 層}であるという。
\end{definition}

\noindent
大まかに言えば、これは切断 $s_i \in \mathcal{F}(U_i)$ が与えられ、
すべての対 $(i, j) \in I \times I$ に対して
$$
s_i|_{U_i \times_U U_j} = s_j|_{U_i \times_U U_j}
$$
が $\mathcal{F}(U_i \times_U U_j)$ において成り立つならば、
$s_i = s|_{U_i}$ を満たす一意な $s \in \mathcal{F}(U)$ が存在するという意味である。

\begin{remark}
\label{sites-remark-sheaf-condition-empty-covering}
被覆 $\{U_i \to U\}_{i \in I}$ が空の族（すなわち $I = \emptyset$）ならば、
層条件は $\mathcal{F}(U) = \{*\}$ が一元集合であることを意味する。
これは、(\ref{sites-equation-sheaf-condition}) の第二および第三の集合が
集合の圏における空積であり、その圏の終対象、したがって一元集合だからである。
\end{remark}

\begin{example}
\label{sites-example-sheaves-topological}
$X$ を位相空間とする。$X_{Zar}$ を Example
\ref{sites-example-site-topological} で構成したサイトとする。
$X_{Zar}$ 上の層という概念は、Sheaves, Definition
\ref{sheaves-definition-sheaf} で導入した $X$ 上の層という概念と一致する。
\end{example}

\begin{example}
\label{sites-example-topological-wrong}
$X$ を位相空間とする。Example \ref{sites-example-site-topological} のサイト
$X_{Zar}$ と同じであるが、空集合の空被覆を許さないサイト $X'_{Zar}$ を考える。
言い換えると、被覆 $\{\emptyset \to \emptyset\}$ は許すが、
添字集合が空である被覆は許さない。これもサイトを定めることは容易に示せる。
しかし、$X'_{Zar}$ 上の層は $X_{Zar}$ 上の層とは異なると主張する。
極端な例として、$X = \{p\}$ が一元集合である場合を考えよう。
このとき $X'_{Zar}$ の対象は $\emptyset, X$ であり、
$\emptyset$ の任意の被覆は $\{\emptyset \to \emptyset\}$ によって、
$X$ の任意の被覆は $\{X \to X\}$ によって細分される。
したがって、この上の層は、任意に選んだ集合 $\mathcal{F}(\emptyset)$ と
任意に選んだ集合 $\mathcal{F}(X)$、および任意の制限写像
$\mathcal{F}(X) \to \mathcal{F}(\emptyset)$ によって与えられることは明らかである。
ゆえに $X'_{Zar}$ 上の層は、開集合が
$\{\emptyset, \{\eta\}, \{p, \eta\}\}$ である二点空間
$\{\eta, p\}$ 上の通常の層と同じである。一般に、$X'_{Zar}$ 上の層は、
空集合と、開集合 $U \subset X$ に対する
$U \cup \{\eta\}$ という形の集合を開集合とする空間
$X \amalg \{\eta\}$ 上の層と同じである。
\end{example}


\begin{definition}
\label{sites-definition-category-sheaves-sets}
集合の層の圏 {\it $\Sh(\mathcal{C})$} とは、
$\textit{PSh}(\mathcal{C})$ の充満部分圏で、その対象が集合の層であるものをいう。
\end{definition}

\noindent
$\mathcal{A}$ を圏とする。被覆族の添字集合として現れる任意の $I$ について、
$I$ および $I \times I$ で添字付けられた積が $\mathcal{A}$ に存在するならば、
上の Definition \ref{sites-definition-sheaf-sets} は意味をもち、
$\mathcal{A}$ に値をもつ $\mathcal{C}$ 上の層という概念を定める。
$\mathcal{A}$ における図式
$$
\xymatrix{
\mathcal{F}(U) \ar[r]
&
\prod\nolimits_{i\in I}
\mathcal{F}(U_i)
\ar@<1ex>[r]^-{\text{pr}_0^*} \ar@<-1ex>[r]_-{\text{pr}_1^*}
&
\prod\nolimits_{(i_0, i_1) \in I \times I}
\mathcal{F}(U_{i_0} \times_U U_{i_1})
}
$$
が等化子図式であるための必要十分条件は、$\mathcal{A}$ の任意の対象 $X$ に対し、
集合の図式
$$
\xymatrix{
\Mor_\mathcal{A}(X, \mathcal{F}(U)) \ar[r]
&
\prod
\Mor_\mathcal{A}(X, \mathcal{F}(U_i))
\ar@<1ex>[r]^-{\text{pr}_0^*} \ar@<-1ex>[r]_-{\text{pr}_1^*}
&
\prod
\Mor_\mathcal{A}(X, \mathcal{F}(U_{i_0} \times_U U_{i_1}))
}
$$
が等化子図式であることである。

\medskip\noindent
$\mathcal{A}$ を任意の圏とする。$\mathcal{F}$ を $\mathcal{A}$ に値をもつ前層とする。
任意の対象 $X\in \Ob(\mathcal{A})$ を選ぶ。このとき規則
$$
\mathcal{F}_X(U) = \Mor_\mathcal{A}(X, \mathcal{F}(U)).
$$
で定められる集合の前層 $\mathcal{F}_X$ が得られる。
上の議論から、すべての前層 $\mathcal{F}_X$ が集合の層であることを要求すれば、
適切な定義が得られる。

\begin{definition}
\label{sites-definition-sheaf}
$\mathcal{C}$ をサイト、$\mathcal{A}$ を圏とし、$\mathcal{F}$ を
$\mathcal{A}$ に値をもつ $\mathcal{C}$ 上の前層とする。
$\mathcal{A}$ のすべての対象 $X$ に対して、上で定義した集合の前層
$\mathcal{F}_X$ が層であるとき、$\mathcal{F}$ は{\it 層}であるという。
\end{definition}












\section{終域を固定した射の族}
\label{sites-section-refinements}

\noindent
この節では、終域が同じ射の族に関するいくつかの概念を導入する。

\begin{definition}
\label{sites-definition-morphism-coverings}
$\mathcal{C}$ を圏とする。$\mathcal{U} = \{U_i \to U\}_{i\in I}$ を
$\mathcal{C}$ における終域を固定した射の族とする。
$\mathcal{V} = \{V_j \to V\}_{j\in J}$ を別のそのような族とする。
\begin{enumerate}
\item
{\it $\mathcal{C}$ の終域を固定した射の族の $\mathcal{U}$ から
$\mathcal{V}$ への射}、または単に{\it $\mathcal{U}$ から $\mathcal{V}$ への射}
とは、射 $U \to V$、集合の写像 $\alpha : I \to J$、および各 $i\in I$
に対する射 $U_i \to V_{\alpha(i)}$ で、図式
$$
\xymatrix{
U_i \ar[r] \ar[d]
&
V_{\alpha(i)} \ar[d]
\\
U \ar[r]
&
V
}
$$
が可換となるものからなるデータをいう。
\item 特別な場合として、$U = V$ かつ $U \to V$ が恒等写像ならば、
$\mathcal{U}$ を族 $\mathcal{V}$ の{\it 細分}という。
\end{enumerate}
\end{definition}

\noindent
自明だが重要な注意として、$\mathcal{V} = \{V_j \to V\}_{j \in J}$ が
{\it 空の射の族}、すなわち $J = \emptyset$ ならば、
$I \not = \emptyset$ である族 $\mathcal{U} = \{U_i \to V\}_{i \in I}$ が
$\mathcal{V}$ を細分することはない！

\begin{definition}
\label{sites-definition-combinatorial-tautological}
$\mathcal{C}$ を圏とする。
$\mathcal{U} = \{\varphi_i : U_i \to U\}_{i\in I}$ と
$\mathcal{V} = \{\psi_j : V_j \to U\}_{j\in J}$ を、
終域を固定した二つの射の族とする。
\begin{enumerate}
\item 写像 $\alpha : I \to J$ と $\beta : J\to I$ で
$\varphi_i = \psi_{\alpha(i)}$ および $\psi_j = \varphi_{\beta(j)}$ を
満たすものが存在するとき、$\mathcal{U}$ と $\mathcal{V}$ は
{\it 組合せ的に同値}であるという。
\item 写像 $\alpha : I \to J$ と $\beta : J\to I$、およびすべての
$i\in I$ と $j \in J$ に対して可換図式
$$
\xymatrix{
U_i \ar[rd] \ar[rr] & &
V_{\alpha(i)} \ar[ld] & &
V_j \ar[rd] \ar[rr] & &
U_{\beta(j)} \ar[ld] \\
&
U & & & &
U &
}
$$
で水平の矢が同型であるものが存在するとき、$\mathcal{U}$ と
$\mathcal{V}$ は{\it 自明に同値}であるという。
\end{enumerate}
\end{definition}

\begin{lemma}
\label{sites-lemma-tautological-combinatorial}
$\mathcal{C}$ を圏とする。
$\mathcal{U} = \{\varphi_i : U_i \to U\}_{i\in I}$ と
$\mathcal{V} = \{\psi_j : V_j \to U\}_{j\in J}$ を、
同じ固定終域をもつ二つの射の族とする。
\begin{enumerate}
\item $\mathcal{U}$ と $\mathcal{V}$ が組合せ的に同値ならば、
自明に同値である。
\item $\mathcal{U}$ と $\mathcal{V}$ が自明に同値ならば、
$\mathcal{U}$ は $\mathcal{V}$ の細分であり、$\mathcal{V}$ は
$\mathcal{U}$ の細分である。
\item 「組合せ的に同値である」という関係は、終域を固定した射の族全体の上の
同値関係である。
\item 「自明に同値である」という関係は、終域を固定した射の族全体の上の
同値関係である。
\item 「$\mathcal{U}$ が $\mathcal{V}$ を細分し、$\mathcal{V}$ が
$\mathcal{U}$ を細分する」という関係は、終域を固定した射の族全体の上の
同値関係である。
\end{enumerate}
\end{lemma}

\begin{proof}
省略する。
\end{proof}

\noindent
以下の補題では、圏 $\mathcal{C}$、$\mathcal{C}$ 上の前層 $\mathcal{F}$、および
すべてのファイバー積 $U_i \times_U U_{i'}$ が存在する族
$\mathcal{U} = \{U_i \to U\}_{i\in I}$ が与えられたとき、図式
(\ref{sites-equation-sheaf-condition}) が等化子図式であることを、
{\it $\mathcal{U}$ に関する $\mathcal{F}$ の層条件}が成り立つという。

\begin{lemma}
\label{sites-lemma-tautological-same-sheaf}
$\mathcal{C}$ を圏とする。
$\mathcal{U} = \{\varphi_i : U_i \to U\}_{i\in I}$ と
$\mathcal{V} = \{\psi_j : V_j \to U\}_{j\in J}$ を、
同じ固定終域をもつ二つの射の族とする。ファイバー積
$U_i \times_U U_{i'}$ と $V_j \times_U V_{j'}$ が存在すると仮定する。
$\mathcal{U}$ と $\mathcal{V}$ が自明に同値ならば、任意の
$\mathcal{C}$ 上の前層 $\mathcal{F}$ に対して、$\mathcal{U}$ に関する
$\mathcal{F}$ の層条件と $\mathcal{V}$ に関する $\mathcal{F}$ の層条件は同値である。
\end{lemma}

\begin{proof}
まず、$\varphi : A \to B$ が圏 $\mathcal{C}$ の同型ならば、
$\varphi^* : \mathcal{F}(B) \to \mathcal{F}(A)$ も同型であることに注意する。
$\beta : J \to I$ を写像とし、仮定によって存在する $U$ 上の同型を
$\chi_j : V_j \to U_{\beta(j)}$ とする。$\mathcal{V}$ に関する層条件が
$\mathcal{U}$ に関する層条件を含意することを示そう。
すべての対 $(i, i') \in I \times I$ に対して
$$
s_i|_{U_i \times_U U_{i'}} = s_{i'}|_{U_i \times_U U_{i'}}
$$
が $\mathcal{F}(U_i \times_U U_{i'})$ において成り立つような切断
$s_i \in \mathcal{F}(U_i)$ が与えられたとする。
$s_j = \chi_j^*s_{\beta(j)}$ と定める。任意の対
$(j, j') \in J \times J$ に対して射
$\chi_j \times_{\text{id}_U} \chi_{j'} : V_j \times_U V_{j'} \to
U_{\beta(j)} \times_U U_{\beta(j')}$ も同型である。
したがって構造の移送により
$$
s_j|_{V_j \times_U V_{j'}} = s_{j'}|_{V_j \times_U V_{j'}}
$$
も成り立つ。$\mathcal{V}$ に関する層条件により、すべての $j \in J$ に対して
$s|_{V_j} = s_j$ を満たす一意な $s$ が存在する。
証明の最初の注意により、すべての $i \in \Im(\beta)$ に対しても
$s|_{U_i} = s_i$ となる。$i \in I$, $i \not \in \Im(\beta)$ と仮定する。
このような $i$ に対し、$U$ 上の同型
$U_i \to V_{\alpha(i)} \to U_{\beta(\alpha(i))}$ がある。
これは射 $U_i \to U_i \times_U U_{\beta(\alpha(i))}$ で射影の切断となるものを
与える。$s_i$ と $s_{\beta(\alpha(i))}$ はファイバー積上で同じ元に制限されるので、
$s_{\beta(\alpha(i))}$ を $U_i \to U_{\beta(\alpha(i))}$ によって
引き戻すと $s_i$ になる。したがって、求めるとおり $s_i = s|_{U_i}$ も成り立つ。
\end{proof}

\begin{lemma}
\label{sites-lemma-compare-separated-presheaf-condition}
$\mathcal{C}$ を圏とする。
$\mathcal{V} = \{V_j \to U\}_{j \in J} \to
\mathcal{U} = \{U_i \to U\}_{i \in I}$ を、
$\text{id} : U \to U$、$\alpha : J \to I$、および
$f_j : V_j \to U_{\alpha(j)}$ によって与えられる、$\mathcal{C}$ の
終域を固定した射の族の射とする。$\mathcal{F}$ を $\mathcal{C}$ 上の前層とする。
$\mathcal{F}(U) \to \prod_{j \in J} \mathcal{F}(V_j)$ が単射ならば、
$\mathcal{F}(U) \to \prod_{i \in I} \mathcal{F}(U_i)$ も単射である。
\end{lemma}

\begin{proof}
省略する。
\end{proof}

\begin{lemma}
\label{sites-lemma-compare-sheaf-condition}
$\mathcal{C}$ を圏とする。
$\mathcal{V} = \{V_j \to U\}_{j \in J} \to
\mathcal{U} = \{U_i \to U\}_{i \in I}$ を、
$\text{id} : U \to U$、$\alpha : J \to I$、および
$f_j : V_j \to U_{\alpha(j)}$ によって与えられる、$\mathcal{C}$ の
終域を固定した射の族の射とする。$\mathcal{F}$ を $\mathcal{C}$ 上の前層とする。
次を仮定する。
\begin{enumerate}
\item ファイバー積 $U_i \times_U U_{i'}$, $U_i \times_U V_j$,
$V_j \times_U V_{j'}$ が存在する。
\item $\mathcal{F}$ は $\mathcal{V}$ に関する層条件を満たす。
\item すべての $i \in I$ に対して写像
$\mathcal{F}(U_i) \to \prod_{j \in J} \mathcal{F}(V_j \times_U U_i)$
は単射である。
\end{enumerate}
このとき $\mathcal{F}$ は $\mathcal{U}$ に関する層条件を満たす。
\end{lemma}

\begin{proof}
Lemma \ref{sites-lemma-compare-separated-presheaf-condition} により、写像
$\mathcal{F}(U) \to \prod \mathcal{F}(U_i)$ は単射である。
すべての $i, i' \in I$ に対して
$s_i|_{U_i \times_U U_{i'}} = s_{i'}|_{U_i \times_U U_{i'}}$
を満たす $s_i \in \mathcal{F}(U_i)$ が与えられたとする。
$s_j = f_j^*(s_{\alpha(j)}) \in \mathcal{F}(V_j)$ と置く。
射 $f_j$ は $U$ 上の射なので、射影写像を介して $f_j, f_{j'}$ と両立する
誘導射
% P01-E0061
$f_{jj'} : V_j \times_U V_{j'} \to
U_{\alpha(i)} \times_U U_{\alpha(i')}$
が得られる。したがって
$$
s_j|_{V_j \times_U V_{j'}}
= f_{jj'}^*(s_{\alpha(j)}|_{U_{\alpha(j)} \times_U U_{\alpha(j')}})
= f_{jj'}^*(s_{\alpha(j')}|_{U_{\alpha(j)} \times_U U_{\alpha(j')}})
= s_{j'}|_{V_j \times_U V_{j'}}
$$
がすべての $j, j' \in J$ に対して成り立つ。
ゆえに、$\mathcal{V}$ に関する $\mathcal{F}$ の層条件により、
各 $V_j$ 上で $s_j$ に制限される切断 $s \in \mathcal{F}(U)$ が得られる。
任意の $i \in I$ に対し、$s$ が $U_i$ 上で $s_i$ に制限されることを示せばよい。
$\mathcal{F}$ は (3) を満たすので、すべての $j \in J$ に対して
$s$ と $s_i$ が $U_i \times_U V_j$ 上で同じ元に制限されることを示せば十分である。
これを見るには
$$
s|_{U_i \times_U V_j} = s_j|_{U_i \times_U V_j} =
(\text{id} \times f_j)^*s_{\alpha(j)}|_{U_i \times_U U_{\alpha(j)}} =
(\text{id} \times f_j)^*s_i|_{U_i \times_U U_{\alpha(j)}} =
s_i|_{U_i \times_U V_j}
$$
を用いればよい。
\end{proof}

\begin{lemma}
\label{sites-lemma-refine-same-topology}
$\mathcal{C}$ を圏とする。$\text{Cov}_i$, $i = 1, 2$ を、
それぞれ $\mathcal{C}$ 上にサイトの構造を定める、終域を固定した射の族の
二つの集合とする。
\begin{enumerate}
\item すべての $\mathcal{U} \in \text{Cov}_1$ が、ある
$\mathcal{V} \in \text{Cov}_2$ と自明に同値ならば、
$\Sh(\mathcal{C}, \text{Cov}_2) \subset
\Sh(\mathcal{C}, \text{Cov}_1)$ である。
さらに、すべての $\mathcal{U} \in \text{Cov}_2$ も、ある
$\mathcal{V} \in \text{Cov}_1$ と自明に同値ならば、
% P01-E0062
二つの層の圏は等しい。
\item 各 $\mathcal{U} \in \text{Cov}_1$ に対して
$\mathcal{V} \in \text{Cov}_2$ が存在し、$\mathcal{V}$ が
$\mathcal{U}$ を細分すると仮定する。この場合
$\Sh(\mathcal{C}, \text{Cov}_2) \subset
\Sh(\mathcal{C}, \text{Cov}_1)$ である。
さらに、すべての $\mathcal{U} \in \text{Cov}_2$ に対して
$\mathcal{V} \in \text{Cov}_1$ が存在し、$\mathcal{V}$ が
$\mathcal{U}$ を細分するならば、
二つの層の圏は等しい。
\end{enumerate}
\end{lemma}

\begin{proof}
(1) は Lemma \ref{sites-lemma-tautological-same-sheaf} と定義から直ちに従う。

\medskip\noindent
(2) を証明する。$\mathcal{F}$ をサイト
$(\mathcal{C}, \text{Cov}_2)$ に対する集合の層とする。
$\mathcal{U} \in \text{Cov}_1$ とし、
$\mathcal{U} = \{U_i \to U\}_{i \in I}$ と書く。
仮定により、$\mathcal{U}$ の細分 $\mathcal{V} \in \text{Cov}_2$ を選べる。
$\mathcal{V} = \{V_j \to U\}_{j \in J}$ と書き、細分は
$\alpha : J \to I$ と $f_j : V_j \to U_{\alpha(j)}$ で与えられるとする。
$\{V_j \times_U U_i \to U_i\}_{j \in J}$ も
$\text{Cov}_2$ に属するので、$\mathcal{F}$ は $\mathcal{V}$ およびこれらの
被覆に関する層条件を満たす。ゆえに Lemma
\ref{sites-lemma-compare-sheaf-condition} により、$\mathcal{F}$ は
$\mathcal{U}$ に関する層条件を満たす。
\end{proof}

\begin{lemma}
\label{sites-lemma-choice-set-coverings-immaterial}
$\mathcal{C}$ を圏とする。$\text{Cov}(\mathcal{C})$ を、Definition
\ref{sites-definition-site} の条件 (1), (2), (3) を満たす被覆の真の類とする。
$\text{Cov}_1, \text{Cov}_2 \subset \text{Cov}(\mathcal{C})$ を、いずれも
$\text{Cov}(\mathcal{C})$ の部分集合であり、$\mathcal{C}$ に
サイトの構造を与えるものとする。
すべての被覆 $\mathcal{U} \in \text{Cov}(\mathcal{C})$ が
$\text{Cov}_1$ の被覆と組合せ的に同値であり、かつ
$\text{Cov}_2$ の被覆とも組合せ的に同値ならば、
$\Sh(\mathcal{C}, \text{Cov}_1) =
\Sh(\mathcal{C}, \text{Cov}_2)$ である。
\end{lemma}

\begin{proof}
仮定により、すべての被覆
$\mathcal{U} \in \text{Cov}_1 \subset \text{Cov}(\mathcal{C})$ が
$\text{Cov}_2$ の元と組合せ的に同値であり、$\text{Cov}_1$ と
$\text{Cov}_2$ の役割を逆にしても同様であるから、上の Lemmas
\ref{sites-lemma-refine-same-topology} と
\ref{sites-lemma-tautological-combinatorial} から明らかである。
\end{proof}
















\section{G-集合の例}
\label{sites-section-example-sheaf-G-sets}

\noindent
例として、Example \ref{sites-example-site-on-group} のサイト $\mathcal{T}_G$ を考える。
$\mathcal{T}_G$ 上の層の圏を記述しよう。答えは $\mathcal{T}_G$ の定義における
選択に依存しないことが分かる。実際、証明ではサイト $\mathcal{T}_G$ の
次の性質だけを用いる。
\begin{enumerate}
\item[(a)] $\mathcal{T}_G$ は $G\textit{-Sets}$ の充満部分圏である。
\item[(b)] $\mathcal{T}_G$ は $G$-集合 ${}_GG$ を含む。
\item[(c)] $\mathcal{T}_G$ はファイバー積をもち、それらは
$G\textit{-Sets}$ におけるものと同じである。
\item[(d)] $U \in \Ob(\mathcal{T}_G)$ と $G$-不変部分集合
$O \subset U$ が与えられると、
% P01-E1264
$O$ と同型な $\mathcal{T}_G$ の対象が存在する。
\item[(e)] $U, U_i \in \Ob(\mathcal{T}_G)$ である任意の全射的な写像の族
$\{U_i \to U\}_{i \in I}$ は、$\mathcal{T}_G$ の被覆と組合せ的に同値である。
\end{enumerate}
これらの性質は Sets, Lemmas \ref{sets-lemma-what-is-in-it-G-sets} と
\ref{sets-lemma-coverings-site} によって成り立つ。

\medskip\noindent
写像
$$
\Hom_G({}_GG, {}_GG)
\longrightarrow
G^{opp},
\varphi
\longmapsto
\varphi(1)
$$
は群の同型であることに注意せよ。逆写像は $g \in G$ を写像
$R_g : s \mapsto sg$（すなわち右乗法）へ送る。
$R_{g_1g_2} = R_{g_2} \circ R_{g_1}$ であるから、反対群が必要である。

\medskip\noindent
このことから、$\mathcal{T}_G$ 上の任意の前層 $\mathcal{F}$ に対して、
値 $\mathcal{F}({}_GG)$ は次のように $G$-集合の構造を受け継ぐ。
$g \in G$ と $s \in \mathcal{F}({}_GG)$ に対して $g \cdot s$ を
$\mathcal{F}(R_g)(s)$ と定める。これは左作用である。実際、
$$
(g_1g_2) \cdot s  = \mathcal{F}(R_{g_1g_2})(s) =
\mathcal{F}(R_{g_2}\circ R_{g_1})(s) =
\mathcal{F}(R_{g_1})( \mathcal{F}(R_{g_2})(s)) =
g_1 \cdot (g_2 \cdot s).
$$
ここでは $\mathcal{F}$ が反変であることを用いた。
$\mathcal{F} \to \mathcal{G}$ が $\mathcal{T}_G$ 上の集合の前層の射ならば、
写像 $\mathcal{F}({}_GG) \to \mathcal{G}({}_GG)$ が得られ、これは今定義した
$G$-作用と両立することに注意せよ。以上により関手
$$
\textit{PSh}(\mathcal{T}_G)
\longrightarrow
G\textit{-Sets}, \quad
\mathcal{F}
\longmapsto
\mathcal{F}({}_GG).
$$
を構成した。この構成には、表現可能前層 $h_U$ が $U$ と標準的に同型なものへ
写るという好ましい性質があることの確認を読者に委ねる。
公式では $h_U({}_GG) = \Hom_G({}_GG, U) \cong U$ である。

\medskip\noindent
$S$ を $G$-集合とする。公式\footnote{これは $S$ によって定義される
表現可能前層のように見えるかもしれない。しかし $S$ は、十分大きな
$G$-集合の集合として選ばれた $\mathcal{T}_G$ の対象であるとは限らないので、
そうでない場合もある。}
$$
\mathcal{F}_S(U)
=
\Mor_{G\textit{-Sets}}(U, S).
$$
によって前層 $\mathcal{F}_S$ を定める。これは明らかに前層である。
一方、$\{U_i \to U\}_{i\in I}$ を $\mathcal{T}_G$ の被覆とする。
このとき $\coprod_i U_i \to U$ は全射である。したがって写像
$$
\mathcal{F}_S(U)
=
\Mor_{G\textit{-Sets}}(U, S)
\longrightarrow
\prod \mathcal{F}_S(U_i)
=
\prod \Mor_{G\textit{-Sets}}(U_i, S)
$$
が単射であることは明らかである。また、すべての図式
$$
\xymatrix{
U_i \times_U U_j \ar[d] \ar[r]
&
U_j \ar[d]^{s_j}
\\
U_i \ar[r]^{s_i}
&
S
}
$$
が可換となる $G$-同変写像の族 $s_i : U_i \to S$ が与えられたとする。
このとき $s_i$ が合成 $U_i \to U \to S$ となる一意な $G$-同変写像
$s : U \to S$ がある。実際、$U_i \to U$ のもとで $u$ に写る元
$u_i \in U_i$ が存在するような任意の $i\in I$ を選び、
$s(u) = s_i(u_i)$ と定めればよい。上の図式の可換性は、この構成が
well-defined であることをちょうど意味する。以上により関手
$$
G\textit{-Sets}
\longrightarrow
\Sh(\mathcal{T}_G), \quad
S
\longmapsto
\mathcal{F}_S
.
$$
を構成した。

\medskip\noindent
これで、圏と関手の次の図式が得られた。
$$
\xymatrix{
\textit{PSh}(\mathcal{T}_G) \ar[rr]^{\mathcal{F} \mapsto \mathcal{F}({}_GG)}
&
&
G\textit{-Sets} \ar[ld]_{S \mapsto \mathcal{F}_S}
\\
&
\Sh(\mathcal{T}_G) \ar[lu]
&
}
$$
定義から $\mathcal{F}_S({}_GG) = \Mor_G({}_GG, S) = S$ であることは
直ちに分かる。最後の等号は $1$ における評価による。
これで次の命題はほとんど証明された。

\begin{proposition}
\label{sites-proposition-sheaves-on-group}
関手 $\mathcal{F} \mapsto \mathcal{F}({}_GG)$ と
$S \mapsto \mathcal{F}_S$ は、$\Sh(\mathcal{T}_G)$ と
$G\textit{-Sets}$ の間の互いに擬逆な同値を定める。
\end{proposition}

\begin{proof}
関手を一方の順序で合成すると恒等関手に同型となることは、すでに見た。
もう一方の順序では、任意の層 $\mathcal{H}$ に対して層の自然な写像
$$
can :
\mathcal{H}
\longrightarrow
\mathcal{F}_{\mathcal{H}({}_GG)}.
$$
がある。実際、$\mathcal{T}_G$ の任意の対象 $U$ に対して、$can_U$ を写像
$$
\begin{matrix}
\mathcal{H}(U)
&
\longrightarrow
&
\mathcal{F}_{\mathcal{H}({}_GG)}(U)
=
\Mor_G(U, \mathcal{H}({}_GG))
\\
s
&
\longmapsto
&
(u \mapsto \alpha_u^*s).
\end{matrix}
$$
とする。ここで $\alpha_u : {}_GG \to U$ は写像
$\alpha_u(g) = gu$ であり、$\alpha_u^* : \mathcal{H}(U)
\to \mathcal{H}({}_GG)$ は引き戻し写像である。
自明だが紛らわしい確認により、これは実際に前層の写像である。
$can$ が同型であることを示さなければならない。すべての
$U \in \Ob(\mathcal{T}_G)$ に対して $can_U$ が同型であることを示す。
$U = {}_GG$ の場合は（重要だが容易なので）読者に委ねる。
$\mathcal{T}_G$ の一般の対象 $U$ は $G$-軌道の非交和
$U = \coprod_{i\in I} O_i$ である。射の族
$\{O_i \to U\}_{i \in I}$ は $\mathcal{T}_G$ の被覆と自明に同値である
（この節の初めに列挙した $\mathcal{T}_G$ の性質による）。したがって Lemma
\ref{sites-lemma-tautological-same-sheaf} により、層 $\mathcal{H}$ は
$\{O_i \to U\}_{i \in I}$ に関する層条件を満たす。
この被覆に対する層条件から
$\mathcal{H}(U) = \prod_i \mathcal{H}(O_i)$ が従う。
したがって、$U$ が一つの $G$-軌道からなる場合に $can_U$ が同型であることを
示せば十分である。$u \in U$ とし、その安定化部分群を $H \subset G$ とする。
明らかに $\Mor_G(U, \mathcal{H}({}_GG)) = \mathcal{H}({}_GG)^H$ は
$H$-不変元の部分集合に等しい。一方、$g \mapsto gu$ によって与えられる被覆
$\{{}_GG \to U\}$ を考える（これも $\mathcal{T}_G$ のある被覆と
組合せ的に同値であるにすぎず、やはりこの相違は重要でない）。
ファイバー積 $({}_GG)\times_U ({}_GG)$ は
$\{(g, gh), g\in G, h\in H\} \cong \coprod_{h\in H} {}_GG$ に等しい。
したがって、この被覆に対する層条件は
$$
\xymatrix{
\mathcal{H}(U) \ar[r]
&
\mathcal{H}({}_GG)
\ar@<1ex>[r]^-{\text{pr}_0^*} \ar@<-1ex>[r]_-{\text{pr}_1^*}
&
\prod_{h \in H}
\mathcal{H}({}_GG).
}
$$
となる。因子 $\mathcal{H}({}_GG)$ への二つの写像 $\text{pr}_i^*$ は、
$h$ による乗法だけ異なる。このことと、$U = {}_GG$ に対して
$can$ が同型であることから結論が従う。
\end{proof}





















\section{層化}
\label{sites-section-sheafification}

\noindent
層化を定義するため、被覆の第0
% P01-E0063
{\v C}ech コホモロジー群とその関手性を調べる。

\medskip\noindent
$\mathcal{F}$ を $\mathcal{C}$ 上の集合の前層とし、
$\mathcal{U} = \{U_i \to U\}_{i \in I}$ を $\mathcal{C}$ の被覆とする。
等化子を表すために $\mathcal{F}(\mathcal{U})$ という記法を用いる。
$$
H^0(\mathcal{U}, \mathcal{F})
=
\{
(s_i)_{i\in I} \in \prod\nolimits_i \mathcal{F}(U_i)
\mid
s_i|_{U_i \times_U U_j} = s_j|_{U_i \times_U U_j}
\ \forall i, j \in I
\}.
$$
後で見るように、これは被覆 $\mathcal{U}$ に関する $U$ 上の
$\mathcal{F}$ の第0 {\v C}ech コホモロジーである。
小さな注意として、すべての射 $U_i \to U$ が表現可能ならば
$H^0(\mathcal{U}, \mathcal{F})$ を定義できる。すなわち、
$\mathcal{U}$ はサイトの被覆である必要はない。
標準写像 $\mathcal{F}(U) \to H^0(\mathcal{U}, \mathcal{F})$ がある。
被覆の射 $\mathcal{U} \to \mathcal{V}$ が可換図式
$$
\xymatrix{
& U_i \ar[rr] & & V_{\alpha(i)} \\
U_i \times_U U_j \ar[rr] \ar[ur] \ar[dr] & &
V_{\alpha(i)} \times_V V_{\alpha(j)} \ar[ur] \ar[dr] & \\
& U_j \ar[rr] & & V_{\alpha(j)}
}.
$$
を誘導することは明らかである。これはさらに写像
$H^0(\mathcal{V}, \mathcal{F}) \to
H^0(\mathcal{U}, \mathcal{F})$ を生じ、この写像は
$\mathcal{F}(V) \to \mathcal{F}(U)$ と両立する。

\medskip\noindent
構成により、前層 $\mathcal{F}$ が層であるための必要十分条件は、
$\mathcal{C}$ のすべての被覆 $\mathcal{U}$ に対して自然な写像
$\mathcal{F}(U) \to H^0(\mathcal{U}, \mathcal{F})$ が全単射であることである。
この概念を用いて、層の極限に関する次の簡単な補題を証明する。

\begin{lemma}
\label{sites-lemma-limit-sheaf}
\begin{slogan}
層の図式の極限は存在し、前層としての極限と一致する
\end{slogan}
$\mathcal{F} : \mathcal{I} \to \Sh(\mathcal{C})$ を図式とする。
このとき $\lim_\mathcal{I} \mathcal{F}$ は存在し、前層の圏における極限に等しい。
\end{lemma}

\begin{proof}
$\lim_i \mathcal{F}_i$ を前層としての極限とする。
これが層であることを示せば、層の圏における極限でもあることが直ちに従う。
層の性質を証明するため、$\mathcal{V} = \{V_j \to V\}_{j\in J}$ を被覆とする。
$(s_j)_{j\in J}$ を $H^0(\mathcal{V}, \lim_i \mathcal{F}_i)$ の元とする。
射影写像を用いると、$H^0(\mathcal{V}, \mathcal{F}_i)$ の元
$(s_{j, i})_{j\in J}$ が得られる。$\mathcal{F}_i$ の層の性質により、
$s_{j, i} = s_i|_{V_j}$ を満たす一意な
$s_i \in \mathcal{F}_i(V)$ が存在する。
$\phi : i \to i'$ を添字圏の射とする。
$\mathcal{F}(\phi) : \mathcal{F}_i \to \mathcal{F}_{i'}$ が
$s_i$ を $s_{i'}$ へ写すことを示したい。すべての $j$ に対して切断
% P01-E1265
$s_{i, j}$ と $s_{i', j}$ についてこれは成り立つことが分かっているので、
$\mathcal{F}_{i'}$ の層の性質によって主張が従う。
これで $(\lim_i \mathcal{F}_i)(V)$ の元
$s = (s_i)_{i \in \Ob(\mathcal{I})}$ が得られた。
この元が $s_j = s|_{V_j}$ という所要の性質をもつことの確認を読者に委ねる。
\end{proof}

\begin{example}
\label{sites-example-singleton-sheaf}
特別な例は空図式にわたる極限である。これは（前）層の圏の終対象を与える。
各 $\mathcal{C}$ の対象 $U$ に一元集合を対応させ、制限写像を一意な写像とする
前層であり、しかもこの前層は層である。この層をしばしば $*$ と書く。
\end{example}

\noindent
$\mathcal{J}_U$ を $U$ のすべての被覆からなる圏とする。
言い換えると、$\mathcal{J}_U$ の対象は $\mathcal{C}$ における $U$ の被覆であり、
射は細分である。サイトに関する本章の規約により、これは実際に圏である。
すなわち、対象と射の集まりはいずれも集合をなす。
$\{\text{id}_U\}$ が対象なので $\Ob(\mathcal{J}_U)$ は空でない。
上の注意により、構成 $\mathcal{U} \mapsto H^0(\mathcal{U}, \mathcal{F})$ は
$\mathcal{J}_U$ 上の反変関手である。次で定める。
$$
\mathcal{F}^{+}(U)
=
\colim_{\mathcal{J}_U^{opp}}
H^0(\mathcal{U}, \mathcal{F})
$$
極限と余極限については Categories, Section
\ref{categories-section-limits} を参照せよ。後で
$\mathcal{F}^{+}(U)$ が $U$ 上の $\mathcal{F}$ の第0
{\v C}ech コホモロジーであることを見る。

\medskip\noindent
余極限の構造をさらに述べる前に、集合の集まり
$\mathcal{F}^{+}(U)$, $U \in \Ob(\mathcal{C})$ を前層にする。
$V \to U$ を $\mathcal{C}$ の射とする。サイトの公理により関手\footnote{この構成は
実際にはファイバー積 $U_i \times_U V$ の選択を伴い、したがって選択公理を用いる。
得られる写像は選択に依存しない。以下を参照せよ。}
$$
\mathcal{J}_U
\longrightarrow
\mathcal{J}_V, \quad
\{U_i \to U\}
\longmapsto
\{U_i \times_U V \to V\}.
$$
がある。射影写像は関手的な被覆の射
$\{U_i \times_U V \to V\} \to \{U_i \to U\}$ を与え、したがって上の構成により
集合の関手的な写像
$H^0(\{U_i \to U\}, \mathcal{F}) \to
H^0(\{U_i \times_U V \to V\}, \mathcal{F})$ を与える。
言い換えると、関手
$H^0(-, \mathcal{F}) : \mathcal{J}_U^{opp} \to \textit{Sets}$ から合成
$\mathcal{J}_U^{opp} \to \mathcal{J}_V^{opp}
\xrightarrow{H^0(-, \mathcal{F})} \textit{Sets}$ への変換がある。
したがって余極限の一般論により標準写像
$\mathcal{F}^+(U) \to \mathcal{F}^+(V)$ が得られる。
上の集合 $\mathcal{F}^+(U)$ の記述で言えば、この写像は
$s = (s_i) \in H^0(\{U_i \to U\}, \mathcal{F})$ に対応する元を
$(s_i|_{V \times_U U_i})
\in H^0(\{U_i \times_U V \to V\}, \mathcal{F})$ に対応する元へ写す。

\begin{lemma}
\label{sites-lemma-plus-presheaf}
上の構成は前層 $\mathcal{F}^+$ と、前層の標準写像
$\mathcal{F} \to \mathcal{F}^+$ を定める。
\end{lemma}

\begin{proof}
示すべきことは、射 $W \to V \to U$ が与えられたとき、合成
$\mathcal{F}^+(U) \to \mathcal{F}^+(V) \to \mathcal{F}^+(W)$ が
写像 $\mathcal{F}^+(U) \to \mathcal{F}^+(W)$ と等しいことだけである。
これは、被覆 $\{U_i \to U\}$ と
$s = (s_i) \in H^0(\{U_i \to U\}, \mathcal{F})$ が与えられたとき、標準的に
$W \times_U U_i \cong W \times_V (V \times_U U_i)$ であり、
$s_i|_{W \times_U U_i}$ がこの同型によって
$(s_i|_{V \times_U U_i})|_{W \times_V (V \times_U U_i)}$ に対応することを
直接確かめれば示せる。
\end{proof}

\noindent
より間接的には、Lemma \ref{sites-lemma-independent-refinement} により、
上のような元 $s$ を、$W$ の任意の被覆から $\{U_i \to U\}$ への
任意の射によって引き戻してよく、常に $\mathcal{F}^+(W)$ の同じ元が得られる。

\begin{lemma}
\label{sites-lemma-plus-functorial}
対応 $\mathcal{F} \mapsto
(\mathcal{F} \to \mathcal{F}^+)$ は関手である。
\end{lemma}

\begin{proof}
これを証明する代わりに、何を証明すべきかを正確に述べる。
$\mathcal{F} \to \mathcal{G}$ を前層の写像とし、次の図式の可換性を証明せよ。
$$
\xymatrix{
\mathcal{F} \ar[r] \ar[d]
&
\mathcal{F}^{+} \ar[d]
\\
\mathcal{G} \ar[r]
&
\mathcal{G}^{+}
}
$$
\end{proof}

\noindent
次の二つの補題は、上の余極限が有向集合にわたる余極限であることを含意する。

\begin{lemma}
\label{sites-lemma-common-refinement}
サイト $\mathcal{C}$ のある対象 $U$ の二つの被覆
$\{U_i \to U\}$ と $\{V_j \to U\}$ が与えられたとき、
それらの共通細分となる被覆が存在する。
\end{lemma}

\begin{proof}
$\mathcal{C}$ はサイトなので、すべての $i$ に対して族
$\{V_j \times_U U_i \to U_i\}_j$ は被覆である。
さらに別の公理により $\{V_j \times_U U_i \to U\}_{i, j}$ は
$U$ の被覆である。この被覆が与えられた二つの被覆を細分することは明らかである。
\end{proof}

\begin{lemma}
\label{sites-lemma-independent-refinement}
同じ射 $U \to V$ を誘導する被覆の射
$f, g: \mathcal{U} \to \mathcal{V}$ は、同じ写像
$H^0(\mathcal{V}, \mathcal{F}) \to  H^0(\mathcal{U}, \mathcal{F})$ を誘導する。
\end{lemma}

\begin{proof}
$\mathcal{U} = \{U_i \to U\}_{i\in I}$ および
$\mathcal{V} = \{V_j \to V\}_{j\in J}$ とする。
射 $f$ は写像 $U\to V$、写像 $\alpha : I \to J$、および写像
$f_i : U_i \to V_{\alpha(i)}$ からなる。同様に、$g$ は写像
$\beta : I \to J$ と写像 $g_i : U_i \to V_{\beta(i)}$ を定める。
$f$ と $g$ は同じ写像 $U\to V$ を誘導するので、すべての $i\in I$ に対して図式
$$
\xymatrix{
&
V_{\alpha(i)} \ar[dr]
\\
U_i \ar[ur]^{f_i} \ar[dr]_{g_i}
&
&
V
\\
&
V_{\beta(i)} \ar[ur]
}
$$
は可換である。したがって $f$ と $g$ はファイバー積を経由して
$$
\xymatrix{
&
V_{\alpha(i)}
\\
U_i \ar[r]^-\varphi \ar[ur]^{f_i} \ar[dr]_{g_i}
&
V_{\alpha(i)} \times_V V_{\beta(i)} \ar[u]_{\text{pr}_1} \ar[d]^{\text{pr}_2}
\\
&
V_{\beta(i)}.
}
$$
と分解する。ここで $s = (s_j)_j \in H^0(\mathcal{V}, \mathcal{F})$ とする。
するとすべての $i\in I$ に対して
$$
(f^*s)_i
=
f_i^*(s_{\alpha(i)})
=
\varphi^*\text{pr}_1^*(s_{\alpha(i)})
=
\varphi^*\text{pr}_2^*(s_{\beta(i)})
=
g_i^*(s_{\beta(i)})
=
(g^*s)_i,
$$
である。中央の等号は $H^0(\mathcal{V}, \mathcal{F})$ の定義による。
これにより $f$ と $g$ が誘導する写像
$H^0(\mathcal{V}, \mathcal{F}) \to H^0(\mathcal{U}, \mathcal{F})$
は等しいことが分かる。
\end{proof}

\begin{remark}
\label{sites-remark-both-refine-same-H0}
特に、この補題により、$\mathcal{U}$ が $\mathcal{V}$ の細分で、かつ
$\mathcal{V}$ が $\mathcal{U}$ の細分ならば、標準的な同一視
$H^0(\mathcal{U}, \mathcal{F}) =
H^0(\mathcal{V}, \mathcal{F})$ がある。
\end{remark}

\noindent
これら二つの補題と $\mathcal{J}_U$ が空でないことから、図式
$H^0(-, \mathcal{F}) : \mathcal{J}_U^{opp}
\to \textit{Sets}$ はフィルター付きである。Categories, Definition
\ref{categories-definition-directed} を参照せよ。
したがって Categories, Section
\ref{categories-section-directed-colimits} により、余極限
$\mathcal{F}^{+}(U)$ は次の簡明な方法で記述できる。
すなわち、集合 $\mathcal{F}^{+}(U)$ の各元は、$U$ のある被覆
$\mathcal{U}$ に対する元 $s \in H^0(\mathcal{U}, \mathcal{F})$ から生じる。
第二の元 $s' \in H^0(\mathcal{U}', \mathcal{F})$ が与えられたとき、
$s$ と $s'$ が余極限の同じ元を定めるための必要十分条件は、$U$ の被覆
$\mathcal{V}$ と細分 $f : \mathcal{V} \to \mathcal{U}$ および
$f' : \mathcal{V} \to \mathcal{U}'$ で、
$H^0(\mathcal{V}, \mathcal{F})$ において
$f^*s = (f')^*s'$ を満たすものが存在することである。
自明な被覆 $\{\text{id}_U\}$ は $\mathcal{J}_U$ の対象なので、
標準写像 $\mathcal{F}(U) \to \mathcal{F}^+(U)$ が得られる。

\begin{lemma}
\label{sites-lemma-plus-surjective}
写像 $\theta : \mathcal{F} \to \mathcal{F}^+$ は次の性質をもつ。
$\mathcal{C}$ の任意の対象 $U$ と任意の切断
$s \in \mathcal{F}^+(U)$ に対して、被覆 $\{U_i \to U\}$ が存在し、
$s|_{U_i}$ は写像 $\theta : \mathcal{F}(U_i)
\to \mathcal{F}^{+}(U_i)$ の像に属する。
\end{lemma}

\begin{proof}
$s$ が元 $(s_i) \in H^0(\{U_i \to U\}, \mathcal{F})$ から生じるような
被覆 $\{U_i \to U\}$ を取る。Lemma
\ref{sites-lemma-independent-refinement} により、$s$ の引き戻しを
$\mathcal{F}^+(U_i)$ の元として計算するため、被覆
$\{U_i \to U_i\}$ と（明らかな）被覆の射
$\{U_i \to U_i\} \to \{U_i \to U\}$ を用いてよい。
実際、この被覆を用いると、$s$ の $U_i$ への制限としてちょうど
$\theta(s_i)$ が得られる。
\end{proof}

\begin{definition}
\label{sites-definition-separated}
サイト $\mathcal{C}$ 上の集合の前層 $\mathcal{F}$ が
{\it 分離的}であるとは、すべての被覆
% P01-E0064
$\{U_i \rightarrow U\}$ に対して写像
$\mathcal{F}(U) \to \prod \mathcal{F}(U_i)$ が単射であることをいう。
\end{definition}

\begin{theorem}
\label{sites-theorem-plus}
上の $\mathcal{F}$ について次が成り立つ。
\begin{enumerate}
\item
\label{sites-item-sep}
前層 $\mathcal{F}^+$ は分離的である。
\item
\label{sites-item-sheaf}
$\mathcal{F}$ が分離的ならば $\mathcal{F}^+$ は層であり、
前層の写像 $\mathcal{F} \to \mathcal{F}^+$ は単射である。
\item
\label{sites-item-plus-iso}
$\mathcal{F}$ が層ならば、$\mathcal{F} \to \mathcal{F}^+$ は同型である。
\item
\label{sites-item-plusplus}
前層 $\mathcal{F}^{++}$ は常に層である。
\end{enumerate}
\end{theorem}

\begin{proof}
(\ref{sites-item-sep}) を証明する。
$s, s' \in \mathcal{F}^+(U)$ とし、すべての $i$ に対して
$s|_{U_i} = s'|_{U_i}$ となるような被覆 $\{U_i \to U\}$ が存在すると仮定する。
いま $U$ の被覆が三つある。すなわち、上の被覆 $\{U_i \to U\}$、Lemma
\ref{sites-lemma-plus-surjective} における $s$ の被覆 $\mathcal{U}$、および
$s'$ に対する同様の被覆 $\mathcal{U}'$ である。Lemma
\ref{sites-lemma-common-refinement} により、共通細分 $\{W_j \to U\}$ を取れる。
これは、$s|_{W_j} = \theta(s_j)$ を満たす
$s_j, s'_j \in \mathcal{F}(W_j)$ があり、$s'|_{W_j}$ に対しても同様で、さらに
$\theta(s_j) = \theta(s'_j)$ であることを意味する。
最後の等式は、$s_j|_{W_{jk}} = s'_j|_{W_{jk}}$ となるような被覆
$\{W_{jk} \to W_j\}$ が存在することを意味する。
$\{W_{jk} \to U\}$ は被覆なので、求めるとおり $s, s'$ は
$H^0(\{W_{jk} \to U\}, \mathcal{F})$ の同じ元に写る。

\medskip\noindent
(\ref{sites-item-sheaf}) を証明する。すべての写像
$\mathcal{F}(U) \to H^0(\mathcal{U}, \mathcal{F})$ が単射なので、
$\mathcal{F} \to \mathcal{F}^+$ が単射であることは明らかである。
また、$\mathcal{U} \to \mathcal{U}'$ が細分ならば
$H^0(\mathcal{U}', \mathcal{F})
\to H^0(\mathcal{U}, \mathcal{F})$ が単射であることも明らかである。
ここで $\{U_i \to U\}$ を被覆とし、すべての $i, i' \in I$ に対して層条件
$s_i|_{U_i \times_U U_{i'}} = s_{i'}|_{U_i \times_U U_{i'}}$
を満たす $\mathcal{F}^+(U_i)$ の元の族 $(s_i)$ を取る。
Lemma \ref{sites-lemma-plus-surjective} のように、$s_i|_{U_{ij}}$ が
一意な元 $s_{ij} \in \mathcal{F}(U_{ij})$ の像となるような被覆
$\{U_{ij} \to U_i\}$ を選ぶ。層条件から、$s_{ij}$ と $s_{i'j'}$ は
$U_{ij} \times_U U_{i'j'}$ 上で一致する。実際、それは
$U_i \times_U U_{i'}$ へ写り、そこで等式が成り立つからである。
したがって $(s_{ij}) \in H^0(\{U_{ij} \to U\}, \mathcal{F})$ は
元 $s \in \mathcal{F}^+(U)$ を生じる。
$s|_{U_i} = s_i$ であることの確認を読者に委ねる。

\medskip\noindent
(\ref{sites-item-plus-iso}) を証明する。これは定義から直ちに従う。
実際、層の性質は、$U$ の各被覆 $\mathcal{U}$ に対して、すべての写像
% P01-E0065
$\mathcal{F} \to H^0(\mathcal{U}, \mathcal{F})$ が全単射であることを
ちょうど意味する。

\medskip\noindent
主張 (\ref{sites-item-plusplus}) はこれで明らかである。
\end{proof}

\begin{definition}
\label{sites-definition-associated-sheaf}
$\mathcal{C}$ をサイトとし、$\mathcal{F}$ を $\mathcal{C}$ 上の集合の前層とする。
層 $\mathcal{F}^\# := \mathcal{F}^{++}$ と標準写像
$\mathcal{F} \to \mathcal{F}^\#$ の組を、{\it $\mathcal{F}$ に付随する層}という。
\end{definition}

\begin{proposition}
\label{sites-proposition-sheafification-adjoint}
標準写像 $\mathcal{F} \to \mathcal{F}^\#$ は次の普遍性をもつ。
$\mathcal{G}$ が集合の層である任意の写像
$\mathcal{F} \to \mathcal{G}$ に対して、一意な写像
$\mathcal{F}^\# \to \mathcal{G}$ が存在し、その合成
$\mathcal{F} \to \mathcal{F}^\# \to \mathcal{G}$ は与えられた写像に等しい。
\end{proposition}

\begin{proof}
Lemma \ref{sites-lemma-plus-functorial} により可換図式
$$
\xymatrix{
\mathcal{F} \ar[r] \ar[d]
&
\mathcal{F}^{+} \ar[r] \ar[d]
&
\mathcal{F}^{++} \ar[d]
\\
\mathcal{G} \ar[r]
&
\mathcal{G}^{+} \ar[r]
&
\mathcal{G}^{++}
}
$$
が得られ、Theorem \ref{sites-theorem-plus} により下の水平写像は同型である。
一意性は、$\mathcal{F}^\#$ の各切断が局所的には $\mathcal{F}$ の切断から
生じることを述べる Lemma \ref{sites-lemma-plus-surjective} から従う。
\end{proof}

\noindent
この結果から、関手 $\mathcal{F}
\mapsto (\mathcal{F} \to \mathcal{F}^\#)$ が関手の一意な同型を除いて
一意であることは明らかである。実際、包含関手を一時的に
$i : \Sh(\mathcal{C}) \to \textit{PSh}(\mathcal{C})$ と書く。
上の結果が実際に述べるのは
$$
\Mor_{\textit{PSh}(\mathcal{C})}(\mathcal{F}, i(\mathcal{G}))
=
\Mor_{\Sh(\mathcal{C})}(\mathcal{F}^\#, \mathcal{G}).
$$
である。言い換えると、層化関手は包含関手 $i$ の左随伴である。

\begin{example}
\label{sites-example-constant-sheaf}
$\mathcal{C}$ をサイトとし、$E$ を集合とする。
{\it 値 $E$ の定数層 $\underline{E}$} とは、値 $E$ の定値関手で与えられる
前層の層化である。この層は、$\mathcal{C}$ 上の任意の層 $\mathcal{F}$ に対する規則
$\Mor_{\Sh(\mathcal{C})}(\underline{E}, \mathcal{F}) =
\Mor_{\textit{Sets}}(E, \Gamma(\mathcal{C}, \mathcal{F}))$
によって特徴付けられる。ここで記法は Definition
\ref{sites-definition-global-sections} のものである。
\end{example}

\noindent
この節を二つの補題で終える。

\begin{lemma}
\label{sites-lemma-colimit-sheaves}
\begin{slogan}
層の圏における余極限は、前層の圏における余極限の層化に等しい。
\end{slogan}
$\mathcal{F} : \mathcal{I} \to \Sh(\mathcal{C})$ を図式とする。
このとき $\colim_\mathcal{I} \mathcal{F}$ は存在し、
前層の圏における余極限の層化である。
\end{lemma}

\begin{proof}
層化関手は左随伴なので、すべての余極限と可換する。Categories, Lemma
\ref{categories-lemma-adjoint-exact} を参照せよ。
したがって $\textit{PSh}(\mathcal{C})$ は余極限をもつので、
$\Sh(\mathcal{C})$ も余極限をもつことが従う
（それらは前層における余極限の層化である）。
\end{proof}

\begin{lemma}
\label{sites-lemma-sheafification-exact}
関手 $\textit{PSh}(\mathcal{C}) \to \Sh(\mathcal{C})$,
$\mathcal{F} \mapsto \mathcal{F}^\#$ は完全である。
\end{lemma}

\begin{proof}
これは左随伴なので右完全である。Categories, Lemma
\ref{categories-lemma-exact-adjoint} を参照せよ。
一方、Lemmas \ref{sites-lemma-common-refinement} と Lemma
\ref{sites-lemma-independent-refinement} により、$\mathcal{F}^+$ の構成における
余極限は、実際には有向集合 $\Ob(\mathcal{J}_U)$ にわたる。
ここで $\mathcal{U} \geq \mathcal{U}'$ であることと
$\mathcal{U}$ が $\mathcal{U}'$ の細分であることは同値である。
したがって Categories, Lemma
\ref{categories-lemma-directed-commutes} により、
$\mathcal{F} \to \mathcal{F}^+$ は有限極限と可換する
（前層から前層への関手として）。最後に Lemma
\ref{sites-lemma-limit-sheaf} を用いて結論を得る。
\end{proof}

\begin{lemma}
\label{sites-lemma-sections-sheafification}
$\mathcal{C}$ をサイトとする。$\mathcal{F}$ を $\mathcal{C}$ 上の集合の前層とする。
$\theta^2 : \mathcal{F} \to \mathcal{F}^\#$ を、$\mathcal{F}$ からその層化への
標準写像とする。$U$ を $\mathcal{C}$ の対象とし、
$s \in \mathcal{F}^\#(U)$ とする。このとき被覆 $\{U_i \to U\}$ と切断
$s_i \in \mathcal{F}(U_i)$ が存在し、次を満たす。
\begin{enumerate}
\item $s|_{U_i} = \theta^2(s_i)$ である。
\item すべての $i, j$ に対して $\mathcal{C}$ の被覆
$\{U_{ijk} \to U_i \times_U U_j\}$ が存在し、$s_i$ と $s_j$ の
各 $U_{ijk}$ への引き戻しは一致する。
\end{enumerate}
逆に、任意の被覆 $\{U_i \to U\}$ と (2) を満たす元
$s_i \in \mathcal{F}(U_i)$ が与えられると、(1) を満たす一意な切断
$s \in \mathcal{F}^\#(U)$ が存在する。
\end{lemma}

\begin{proof}
省略する。
\end{proof}

\begin{lemma}
\label{sites-lemma-isomorphism-sheafifications}
$\mathcal{C}$ をサイトとし、$\mathcal{F} \to \mathcal{G}$ を
$\mathcal{C}$ 上の集合の前層の写像とする。
$\mathcal{F}(U) \to \mathcal{G}(U)$ が全単射である
$U \in \Ob(\mathcal{C})$ の集合を $\mathcal{B}$ と書く。
$\mathcal{C}$ の各対象が $\mathcal{B}$ の元による被覆をもつならば、
$\mathcal{F}^\# \to \mathcal{G}^\#$ は同型である。
\end{lemma}

\begin{proof}
$U \in \Ob(\mathcal{C})$ とする。
$\mathcal{F}^\#(U) \to \mathcal{G}^\#(U)$ が全射であることを示そう。
このため、Lemma \ref{sites-lemma-sections-sheafification} を断りなく用いる。
任意の $s \in \mathcal{G}^\#(U)$ に対して被覆
$\{U_i \to U\}$ と切断 $s_i \in \mathcal{G}(U_i)$ が存在し、次を満たす。
\begin{enumerate}
\item $s|_{U_i}$ は $\mathcal{G} \to \mathcal{G}^\#$ による $s_i$ の像である。
\item すべての $i, j$ に対して
% P01-E0066
$\mathcal{D}$ の被覆 $\{U_{ijk} \to U_i \times_U U_j\}$ が存在し、
$s_i$ と $s_j$ の各 $U_{ijk}$ への引き戻しは一致する。
\end{enumerate}
仮定により、各 $i$ に対して $U_{i, a} \in \mathcal{B}$ である被覆
$\{U_{i, a} \to U_i\}$ を選べる。このとき $\{U_{i, a} \to U\}$ は被覆である。
$s_{i, a}$ を $\mathcal{G}(U_{i, a})$ における $s_i$ の像と書くと、
$s_{i, a}$ と $s_{j, b}$ の引き戻しは被覆
$$
\{(U_{i, a} \times_U U_{j, b}) \times_{U_i \times_U U_j} U_{ijk}
\to U_{i, a} \times_U U_{j, b}\}
$$
の各元上で一致する。したがって $U_i \in \mathcal{B}$ と仮定してよい。
議論を繰り返すと、すべての $i, j, k$ に対して
$U_{ijk} \in \mathcal{B}$ とも仮定してよい（詳細は省略する）。
補題の主張における $\mathcal{B}$ の定義により
$\mathcal{F}(U_i) \to \mathcal{G}(U_i)$ は全単射なので、$s_i$ に写る一意な
$t_i \in \mathcal{F}(U_i)$ が得られる。$t_i$ と $t_j$ の各 $U_{ijk}$ への
引き戻しは $\mathcal{F}(U_{ijk})$ において一致する。なぜなら
$U_{ijk} \in \mathcal{B}$ であり、しかも
% P01-E0067
$t_i$ と $t_j$ について一致が成り立つからである。
すると補題により、$t|_{U_i}$ が $t_i$ の像となる一意な
$t \in \mathcal{F}^\#(U)$ が存在する。したがって $t$ は $s$ へ写る。
$\mathcal{F}^\#(U) \to \mathcal{G}^\#(U)$ が単射であることの証明は省略する。
\end{proof}

\section{層の単射写像と全射写像}
\label{sites-section-sheaves-injective}

\begin{definition}
\label{sites-definition-sheaves-injective-surjective}
$\mathcal{C}$ をサイトとし、$\varphi : \mathcal{F}
\to \mathcal{G}$ を集合の層の写像とする。
\begin{enumerate}
\item $\mathcal{C}$ の任意の対象 $U$ に対して写像
$\varphi : \mathcal{F}(U) \to \mathcal{G}(U)$ が単射であるとき、
$\varphi$ は{\it 単射的}であるという。
\item $\mathcal{C}$ の任意の対象 $U$ と任意の切断
$s\in \mathcal{G}(U)$ に対し、被覆 $\{U_i \to U\}$ が存在して、
すべての $i$ について制限 $s|_{U_i}$ が
$\varphi : \mathcal{F}(U_i) \to \mathcal{G}(U_i)$ の像に入るとき、
$\varphi$ は{\it 全射的}であるという。
\end{enumerate}
\end{definition}

\begin{lemma}
\label{sites-lemma-mono-epi-sheaves}
上で定義した単射的写像（それぞれ全射的写像）は、ちょうど圏
$\Sh(\mathcal{C})$ のモノ射（それぞれエピ射）である。
層の写像が同型であるための必要十分条件は、それが単射的かつ全射的であることである。
\end{lemma}

\begin{proof}
省略する。
\end{proof}

\begin{lemma}
\label{sites-lemma-coequalizer-surjection}
$\mathcal{C}$ をサイトとする。$\mathcal{F} \to \mathcal{G}$ を
集合の層の全射とする。このとき図式
$$
\xymatrix{
\mathcal{F} \times_\mathcal{G} \mathcal{F}
\ar@<1ex>[r] \ar@<-1ex>[r]
&
\mathcal{F} \ar[r]
&
\mathcal{G}}
$$
は $\mathcal{G}$ を余等化子として表す。
\end{lemma}

\begin{proof}
$\mathcal{H}$ を集合の層とし、$\varphi : \mathcal{F} \to \mathcal{H}$ を、
二つの写像 $\mathcal{F} \times_\mathcal{G} \mathcal{F} \to \mathcal{F}$ を
等化する層の写像とする。写像 $\mathcal{F} \to \mathcal{G}$ の前層としての像を
$\mathcal{G}' \subset \mathcal{G}$ とする。積
$\mathcal{F} \times_\mathcal{G} \mathcal{F}$ は前層の圏で計算できるので、
これは前層の積 $\mathcal{F} \times_{\mathcal{G}'} \mathcal{F}$ に等しい。
したがって $\varphi$ は前層の一意な写像
$\psi' : \mathcal{G}' \to \mathcal{H}$ を誘導する。
Lemma \ref{sites-lemma-mono-epi-sheaves} により $\mathcal{G}$ は
$\mathcal{G}'$ の層化なので、$\psi'$ は層の写像
$\psi : \mathcal{G} \to \mathcal{H}$ に一意に延長される。
$\varphi$ が $\psi$ と与えられた写像との合成に等しいことの確認は省略する。
\end{proof}














\section{表現可能層}
\label{sites-section-representable-sheaves}

\noindent
$\mathcal{C}$ を圏とする。標準位相とは、すべての表現可能前層が層となるような
最も細かい位相である（これは Definition \ref{sites-definition-canonical-topology} で
形式的に定義されるが、ここでは必要としない）。
この位相は、$\mathcal{C}$ 上のサイトの構造に付随する位相とは限らない。
$\mathcal{C}$ がファイバー積をもつ場合に、この位相を生成する被覆の族を与える。
まず、次の一般的な定義を与える。

\begin{definition}
\label{sites-definition-universal-effective-epimorphisms}
$\mathcal{C}$ を圏とする。族 $\{U_i \to U\}_{i \in I}$ のすべての射
$U_i \to U$ が表現可能であり（Categories, Definition
\ref{categories-definition-representable-morphism} を参照）、任意の
$X\in \Ob(\mathcal{C})$ に対して列
$$
\xymatrix{
\Mor_\mathcal{C}(U, X) \ar[r]
&
\prod\nolimits_{i \in I} \Mor_\mathcal{C}(U_i, X)
\ar@<1ex>[r] \ar@<-1ex>[r]
&
\prod\nolimits_{(i, j) \in I^2} \Mor_\mathcal{C}(U_i \times_U U_j, X)
}
$$
が等化子図式であるとき、この族を{\it 有効エピ族}という。
任意の射 $V \to U$ に対して基底変換
$\{U_i \times_U V \to V\}$ が有効エピ族であるとき、族
$\{U_i \to U\}$ を{\it 普遍有効エピ族}という。
\end{definition}

\noindent
普遍有効エピ族からなる族の類は Definition \ref{sites-definition-site} の公理を満たす。
$\mathcal{C}$ がファイバー積をもつなら、付随する位相は標準位相である。
（この場合にサイトを得るには Sets, Lemma
\ref{sets-lemma-coverings-site} と同様に論じればよい。）

\medskip\noindent
逆に、$\mathcal{C}$ がサイトであり、すべての表現可能前層が層であると仮定する。
すると明らかに、すべての被覆は普遍有効エピ族である。
したがって、サイトの設定では次の定義が「正しい」ものとなる。

\begin{definition}
\label{sites-definition-weaker-than-canonical}
サイト $\mathcal{C}$ 上の位相について、$\mathcal{C}$ のすべての被覆が
普遍有効エピ族であるとき、その位相は{\it 標準位相より弱い}、
または{\it 準標準的}であるという。
\end{definition}

\noindent
表現可能層とは、層でもある表現可能前層のことである。
位相が準標準的でない場合にはこの用語を避けた方がよいかもしれないので、
その場合に限ってこれを形式的に定義する。

\begin{definition}
\label{sites-definition-representable-sheaf}
$\mathcal{C}$ を、その位相が準標準的なサイトとする。
米田埋め込み $h$（Categories, Section
\ref{categories-section-opposite} を参照）は、$\mathcal{C}$ を
$\mathcal{C}$ 上の層の圏の充満部分圏として表す。
この場合、$U \in \Ob(\mathcal{C})$ に対する形 $h_U$ の層を、
$\mathcal{C}$ 上の{\it 表現可能層}という。
記法：$U$ に付随する表現可能層 $h_U$ を
{\it $\underline{U}$} と書くこともある。
\end{definition}

\noindent
定義の状況では、任意の層 $\mathcal{F}$ に対して
$$
\Mor_{\Sh(\mathcal{C})}(h_U, \mathcal{F}) = \mathcal{F}(U)
$$
が成り立つことに注意しよう。実際、これは前層について成り立つ。
式 (\ref{sites-equation-map-representable-into-presheaf}) を参照せよ。
一般には前層 $h_U$ は層ではなく、層を得るにはこれを層化しなければならない。
この場合にも、任意の層 $\mathcal{F}$ に対して
\begin{equation}
\label{sites-equation-map-representable-into-sheaf}
\Mor_{\Sh(\mathcal{C})}(h_U^\#, \mathcal{F}) =
\Mor_{\textit{PSh}(\mathcal{C})}(h_U, \mathcal{F}) =
\mathcal{F}(U)
\end{equation}
が成り立つ。実際、最初の等号は $\#$ の随伴性から従い、二番目は
(\ref{sites-equation-map-representable-into-presheaf}) である。

\begin{lemma}
\label{sites-lemma-covering-surjective-after-sheafification}
\begin{slogan}
被覆は層化後に全射となる。
\end{slogan}
$\mathcal{C}$ をサイトとする。$\{U_i \to U\}_{i \in I}$ がサイト
$\mathcal{C}$ の被覆なら、集合の前層の射
$$
\coprod\nolimits_{i \in I} h_{U_i} \to h_U
$$
は層化後に全射となる。
\end{lemma}

\begin{proof}
上の Lemma \ref{sites-lemma-mono-epi-sheaves} により、
$\coprod\nolimits_{i \in I} h_{U_i}^\# \to h_U^\#$ が
エピ射であることを示せばよい。$\mathcal{F}$ を集合の層とする。
射 $h_U^\# \to \mathcal{F}$ は切断 $s \in \mathcal{F}(U)$ に対応する。
したがって $\Mor(h_U^\#, \mathcal{F})
\to \prod_i \Mor(h_{U_i}^\#, \mathcal{F})$ の単射性は、
$\mathcal{F}$ の層の性質から直ちに従う。
\end{proof}

\noindent
次の補題は、位相が標準位相より弱い場合、あらゆる層がある意味で
表現可能層から組み立てられることを述べる。

\begin{lemma}
\label{sites-lemma-sheaf-coequalizer-representable}
$\mathcal{C}$ をサイトとする。$E \subset \Ob(\mathcal{C})$ を、
$\mathcal{C}$ のすべての対象が $E$ の元による被覆をもつような部分集合とする。
$\mathcal{F}$ を集合の層とする。集合の層の図式
$$
\xymatrix{
\mathcal{F}_1 \ar@<1ex>[r] \ar@<-1ex>[r] &
\mathcal{F}_0 \ar[r] &
\mathcal{F}
}
$$
であって、$\mathcal{F}$ を余等化子として表し、かつ
$\mathcal{F}_i$（$i = 0, 1$）が $U \in E$ に対する形
$h_U^\#$ の層の余積であるものが存在する。
\end{lemma}

\begin{proof}
まず、所望の型のエピ射 $\mathcal{F}_0 \to \mathcal{F}$ が存在することを示す。
実際、単に
$$
\mathcal{F}_0 =
\coprod\nolimits_{U \in E, s \in \mathcal{F}(U)}
(h_U)^\# \longrightarrow \mathcal{F}
$$
とすればよい。ここで $(U, s)$ に対応する成分への制限は、元
$\text{id}_U \in h_U^\#(U)$ を切断 $s \in \mathcal{F}(U)$ へ写す。
Lemma \ref{sites-lemma-mono-epi-sheaves} と $E$ に関する仮定により、これはエピ射である。
$\mathcal{F}_1$ を構成するには、まず
$\mathcal{G} = \mathcal{F}_0 \times_\mathcal{F} \mathcal{F}_0$ とおき、
次に上と同様にエピ射 $\mathcal{F}_1 \to \mathcal{G}$ を構成する。
Lemma \ref{sites-lemma-coequalizer-surjection} を参照せよ。
\end{proof}

\begin{lemma}
\label{sites-lemma-colimit-representable}
$\mathcal{C}$ をサイトとし、$\mathcal{F}$ を $\mathcal{C}$ 上の集合の層とする。
このとき図式 $\mathcal{I} \to \mathcal{C}$、$i \mapsto U_i$ であって
$$
\mathcal{F} = \colim_{i \in \mathcal{I}} h_{U_i}^\#
$$
を満たすものが存在する。さらに、$E \subset \Ob(\mathcal{C})$ が、
$\mathcal{C}$ のすべての対象が $E$ の元による被覆をもつような部分集合なら、
すべての $i \in \Ob(\mathcal{I})$ について $U_i$ は $E$ の元であると仮定できる。
\end{lemma}

\begin{proof}
$\mathcal{I}$ を、対象が $U \in \Ob(\mathcal{C})$ と
$s \in \mathcal{F}(U)$ の対 $(U, s)$ であり、射
$(U, s) \to (U', s')$ が $f^*s' = s$ を満たす $\mathcal{C}$ の射
$f : U \to U'$ であるような圏とする。$\mathcal{I}$ の各対象 $(U, s)$ に対し、
元 $s$ は米田の補題により前層の写像 $s : h_U \to \mathcal{F}$ を定める。
したがって層化の普遍性により、写像 $h_U^\# \to \mathcal{F}$ が得られる。
これらの写像が圏 $\mathcal{I}$ の射と両立することは直ちに分かる。
したがって前層の写像 $\colim_{(U, s)} h_U \to \mathcal{F}$
（ここで余極限は前層の圏でとる）と、層の写像
$\colim_{(U, s)} (h_U)^\# \to \mathcal{F}$
（ここで余極限は層の圏でとる）が得られる。
層化は包含関手 $\Sh(\mathcal{C}) \to \textit{PSh}(\mathcal{C})$ の
左随伴なので（Proposition \ref{sites-proposition-sheafification-adjoint}）、
Categories, Lemma \ref{categories-lemma-adjoint-exact} により
$\colim (h_U)^\# = (\colim h_U)^\#$ である。
したがって $\colim_{(U, s)} h_U \to \mathcal{F}$ が前層の同型であることを
示せば十分である。このため、$\mathcal{C}$ の任意の対象 $X$ に対し、写像
$\colim_{(U, s)} h_U(X) \to \mathcal{F}(X)$ が全単射であることを示す。
実際、この写像には、元 $t \in \mathcal{F}(X)$ を、
$(X, t)$ と $\text{id}_X \in h_X(X)$ に対応する集合
$\colim_{(U, s)} h_U(X)$ の元へ送る逆写像がある。

\medskip\noindent
最後の主張の証明は省略する。
\end{proof}

\section{連続関手}
\label{sites-section-continuous-functors}

\begin{definition}
\label{sites-definition-continuous}
$\mathcal{C}$ と $\mathcal{D}$ をサイトとする。
関手 $u : \mathcal{C} \to \mathcal{D}$ が{\it 連続}であるとは、
任意の $\{V_i \to V\}_{i\in I} \in \text{Cov}(\mathcal{C})$ に対して
次が成り立つことをいう。
\begin{enumerate}
\item $\{u(V_i) \to u(V)\}_{i\in I}$ は $\text{Cov}(\mathcal{D})$ に属する。
\item $\mathcal{C}$ の任意の射 $T \to V$ に対して、射
$u(T \times_V V_i) \to u(T) \times_{u(V)} u(V_i)$ は同型である。
\end{enumerate}
\end{definition}

\noindent
上のような関手 $u$ と $\mathcal{D}$ 上の集合の前層 $\mathcal{F}$ が
与えられたとき、$u^p\mathcal{F}$ を単に前層 $\mathcal{F} \circ u$、
すなわち $\mathcal{C}$ の任意の対象 $V$ に対して
$$
u^p\mathcal{F} (V) = \mathcal{F}(u(V))
$$
と定義したことを思い出そう。

\begin{lemma}
\label{sites-lemma-pushforward-sheaf}
$\mathcal{C}$ と $\mathcal{D}$ をサイトとし、
$u : \mathcal{C} \to \mathcal{D}$ を連続関手とする。
$\mathcal{F}$ が $\mathcal{D}$ 上の層なら、$u^p\mathcal{F}$ も層である。
\end{lemma}

\begin{proof}
$\{V_i \to V\}$ を被覆とする。仮定により
$\{u(V_i) \to u(V)\}$ は $\mathcal{D}$ の被覆であり、
$u(V_i \times_V V_j) = u(V_i)\times_{u(V)}u(V_j)$ である。
したがって $u^p\mathcal{F}$ と被覆 $\{V_i \to V\}$ に対する層条件は、
$\mathcal{F}$ と被覆 $\{u(V_i) \to u(V)\}$ に対する層条件とまったく同じである。
\end{proof}

\noindent
混乱を避けるため、集合の層の部分圏に制限した関手 $u^p$ を、ときに
$$
u^s :
\Sh(\mathcal{D})
\longrightarrow
\Sh(\mathcal{C})
$$
と書く。$u^p$ は左随伴
$u_p : \textit{PSh}(\mathcal{C}) \to \textit{PSh}(\mathcal{D})$ を
もつことを思い出そう。Section \ref{sites-section-functoriality-PSh} を参照せよ。

\begin{lemma}
\label{sites-lemma-adjoint-sheaves}
Lemma \ref{sites-lemma-pushforward-sheaf} の状況で、関手
$u_s : \mathcal{G} \mapsto (u_p \mathcal{G})^\#$ は $u^s$ の左随伴である。
\end{lemma}

\begin{proof}
Lemma \ref{sites-lemma-adjoints-u} と Proposition
\ref{sites-proposition-sheafification-adjoint} から直ちに従う。
\end{proof}

\noindent
次は技術的な補題である。

\begin{lemma}
\label{sites-lemma-technical-up}
Lemma \ref{sites-lemma-pushforward-sheaf} の状況で、$\mathcal{C}$ 上の任意の前層
$\mathcal{G}$ に対して
$(u_p\mathcal{G})^\# = (u_p(\mathcal{G}^\#))^\#$ が成り立つ。
\end{lemma}

\begin{proof}
$\mathcal{D}$ 上の任意の層 $\mathcal{F}$ に対して
\begin{eqnarray*}
\Mor_{\Sh(\mathcal{D})}(u_s(\mathcal{G}^\#), \mathcal{F})
& = &
\Mor_{\Sh(\mathcal{C})}(\mathcal{G}^\#, u^s\mathcal{F}) \\
& = &
\Mor_{\textit{PSh}(\mathcal{C})}(\mathcal{G}^\#, u^p\mathcal{F}) \\
& = &
\Mor_{\textit{PSh}(\mathcal{C})}(\mathcal{G}, u^p\mathcal{F}) \\
& = &
\Mor_{\textit{PSh}(\mathcal{D})}(u_p\mathcal{G}, \mathcal{F}) \\
& = &
\Mor_{\Sh(\mathcal{D})}((u_p\mathcal{G})^\#, \mathcal{F})
\end{eqnarray*}
であり、結果は米田の補題から従う。
\end{proof}

\begin{lemma}
\label{sites-lemma-pullback-representable-sheaf}
$u : \mathcal{C} \to \mathcal{D}$ をサイト間の連続関手とする。
$\mathcal{C}$ の任意の対象 $U$ に対して $u_sh_U^\# = h_{u(U)}^\#$ である。
\end{lemma}

\begin{proof}
Lemmas \ref{sites-lemma-pullback-representable-presheaf} と
\ref{sites-lemma-technical-up} から従う。
\end{proof}


\begin{remark}
\label{sites-remark-quasi-continuous}
（初読時には飛ばしてよい。）
$\mathcal{C}$ と $\mathcal{D}$ をサイトとする。
写像族が自明に同値であるという Definition
\ref{sites-definition-combinatorial-tautological} の定義を用いて、
連続性を定める条件を（わずかに）弱めよう。
$u : \mathcal{C} \to \mathcal{D}$ を関手とする。
任意の $\mathcal{V} = \{V_i \to V\}_{i\in I}
\in \text{Cov}(\mathcal{C})$ に対して次が成り立つとき、
$u$ を{\it 準連続}という。
\begin{enumerate}
\item[(1')] 写像族 $\{u(V_i) \to u(V)\}_{i\in I}$ は
$\text{Cov}(\mathcal{D})$ のある元と自明に同値である。
\item[(2)] $\mathcal{C}$ の任意の射 $T \to V$ に対して、射
$u(T \times_V V_i) \to u(T) \times_{u(V)} u(V_i)$ は同型である。
\end{enumerate}
$u$ が準連続の場合にも Lemmas \ref{sites-lemma-pushforward-sheaf} と
\ref{sites-lemma-adjoint-sheaves} が成り立つことを示そう。

\medskip\noindent
まず、第一の条件により射 $u(V_i) \to u(V)$ は表現可能な射と同型なので、
これらの射は表現可能である。特に、族
$u(\mathcal{V}) = \{u(V_i) \to u(V)\}_{i\in I}$ は、
$\mathcal{D}$ 上の任意の前層 $\mathcal{F}$ に対する
% P01-E0063
第0 {\v C}ech コホモロジー群 $H^0(u(\mathcal{V}), \mathcal{F})$ を定める。
$\mathcal{U} = \{U_j \to u(V)\}_{j \in J}$ を、
$\{u(V_i) \to u(V)\}_{i \in I}$ と自明に同値な
$\text{Cov}(\mathcal{D})$ の元とする。$u(\mathcal{V})$ は
$\mathcal{U}$ の細分であり、逆も成り立つことに注意しよう。
したがって Remark \ref{sites-remark-both-refine-same-H0} により
$H^0(u(\mathcal{V}), \mathcal{F}) = H^0(\mathcal{U}, \mathcal{F})$ である。
特に $\mathcal{F}$ が層なら、第0 {\v C}ech コホモロジー群によって
表された層の性質から
$\mathcal{F}(u(V)) = H^0(u(\mathcal{V}), \mathcal{F})$ である。
したがって $\mathcal{F}$ が層なら $u^p\mathcal{F}$ も層である。
実際、$H^0(\mathcal{V}, u^p\mathcal{F}) =
H^0(u(\mathcal{V}), \mathcal{F})$ であり、右辺は今見たように
$\mathcal{F}(u(V)) = u^p\mathcal{F}(V)$ に等しい。
よって Lemma \ref{sites-lemma-pushforward-sheaf} が成り立つ。
Lemma \ref{sites-lemma-adjoint-sheaves} は直ちに従う。
\end{remark}

\section{サイトの射}
\label{sites-section-morphism-sites}

\begin{definition}
\label{sites-definition-morphism-sites}
$\mathcal{C}$ と $\mathcal{D}$ をサイトとする。
{\it サイトの射} $f : \mathcal{D} \to \mathcal{C}$ とは、
$u_s$ が完全となるような連続関手
$u : \mathcal{C} \to \mathcal{D}$ によって与えられるものをいう。
\end{definition}

\noindent
関手 $u$ の向きは射 $f$ と{\it 反対}であることに注意しよう。
$f \leftrightarrow u$ がサイトの射なら、$f^{-1} = u_s$ および
$f_* = u^s$ と書く。関手 $f^{-1}$ を{\it 逆像関手}、
関手 $f_*$ を{\it 順像関手}という。位相の場合と同様に、随伴性
$$
\Mor_{\Sh(\mathcal{D})}(f^{-1}\mathcal{G}, \mathcal{F})
=
\Mor_{\Sh(\mathcal{C})}(\mathcal{G}, f_*\mathcal{F})
$$
が成り立つ。この定義の動機は次の例にある。

\begin{example}
\label{sites-example-continuous-map}
$f : X  \to Y$ を位相空間の連続写像とする。サイト $X_{Zar}$ と
$Y_{Zar}$ があったことを思い出そう。Example
\ref{sites-example-site-topological} を参照せよ。関手
$u : Y_{Zar} \to X_{Zar}$、$V \mapsto f^{-1}(V)$ を考える。
開被覆の逆像は開被覆なので、この関手は明らかに連続である。
（厳密には、これは $X_{Zar}$ と $Y_{Zar}$ の被覆の集合をどう選んだかに
依存する。しかし、いずれにせよこの関手は準連続である。
Remark \ref{sites-remark-quasi-continuous} を参照せよ。）
関手 $u^s$ が位相における通常の順像関手 $f_*$ と等しいことは容易に確認できる。
したがって $u_s$ は随伴であり、通常の位相的逆像関手 $f^{-1}$ も随伴なので、
標準同型 $f^{-1} \cong u_s$ が得られる。$f^{-1}$ は完全だから、
$u_s$ も完全である。ゆえに $u$ はサイトの射
$f : X_{Zar} \to Y_{Zar}$ を定める。関手 $u_s, u^s$ はいずれにせよ
通常の概念と一致することをすでに見たので、この射も $f$ と書いてよい。
\end{example}

\begin{example}
\label{sites-example-finer-topology}
$\mathcal{C}$ を圏とする。終域を固定した射の族の二つの集合
$$
\text{Cov}(\mathcal{C}) \supset \text{Cov}'(\mathcal{C})
$$
があり、それぞれが $\mathcal{C}$ をサイトにするとする。
$\text{Cov}(\mathcal{C})$ に対応するサイトを $\mathcal{C}_\tau$、
$\text{Cov}'(\mathcal{C})$ に対応するサイトを $\mathcal{C}_{\tau'}$ と書く。
$\mathcal{C}$ 上の恒等関手はサイトの射
$$
\epsilon : \mathcal{C}_\tau \longrightarrow \mathcal{C}_{\tau'}
$$
を定めると主張する。実際、すべての $\tau'$-被覆は $\tau$-被覆なので、
$\text{id} : \mathcal{C}_{\tau'} \to \mathcal{C}_\tau$ は連続である。
したがって関手 $\epsilon_* = \text{id}^s$ は台となる前層上の恒等関手である。
ゆえに $\epsilon_*$ の左随伴 $\epsilon^{-1}$ は、$\tau'$-層
$\mathcal{F}$ を $\mathcal{F}$ の $\tau$-層化へ送る
（層化の普遍性による）。$\tau'$-層の有限極限は前層の有限極限と一致し
（Lemma \ref{sites-lemma-limit-sheaf}）、$\tau$-層化は有限極限と可換する
（Lemma \ref{sites-lemma-sheafification-exact}）。したがって
$\epsilon^{-1}$ は左完全である。また $\epsilon^{-1}$ は左随伴なので、
右完全でもある（Categories, Lemma \ref{categories-lemma-exact-adjoint}）。
ゆえに $\epsilon^{-1}$ は完全であり、Definition
\ref{sites-definition-morphism-sites} のすべての条件を確認した。
\end{example}

\begin{lemma}
\label{sites-lemma-composition-morphisms-sites}
\begin{slogan}
サイト関手の合成はサイト関手の連続性を保つ。
\end{slogan}
$\mathcal{C}_i$（$i = 1, 2, 3$）をサイトとする。
$u : \mathcal{C}_2 \to \mathcal{C}_1$ と
$v : \mathcal{C}_3 \to \mathcal{C}_2$ を、サイトの射を誘導する連続関手とする。
このとき関手 $u \circ v : \mathcal{C}_3 \to \mathcal{C}_1$ は連続であり、
サイトの射 $\mathcal{C}_1 \to \mathcal{C}_3$ を定める。
\end{lemma}

\begin{proof}
$u \circ v$ が連続関手であることは定義から直ちに分かる。
さらに明らかに $(u \circ v)^p = v^p \circ u^p$ であり、したがって
$(u \circ v)^s = v^s \circ u^s$ である。ゆえに関手
$(u \circ v)_s$ と $u_s \circ v_s$ はともに $(u \circ v)^s$ の左随伴である。
したがって $(u \circ v)_s \cong u_s \circ v_s$ であり、
$(u \circ v)_s$ は完全関手の合成として完全である。
\end{proof}

\begin{definition}
\label{sites-definition-composition-morphisms-sites}
$\mathcal{C}_i$（$i = 1, 2, 3$）をサイトとする。
$f : \mathcal{C}_1 \to \mathcal{C}_2$ と
$g : \mathcal{C}_2 \to \mathcal{C}_3$ を、それぞれ連続関手
$u : \mathcal{C}_2 \to \mathcal{C}_1$ と
$v : \mathcal{C}_3 \to \mathcal{C}_2$ によって与えられるサイトの射とする。
{\it 合成} $g \circ f$ とは、関手 $u \circ v$ に対応するサイトの射をいう。
\end{definition}

\noindent
この状況では $(g \circ f)_* = g_* \circ f_*$ および
$(g \circ f)^{-1} = f^{-1} \circ g^{-1}$ である
（Lemma \ref{sites-lemma-composition-morphisms-sites} の証明を参照せよ）。

\begin{lemma}
\label{sites-lemma-directed-morphism}
$\mathcal{C}$ と $\mathcal{D}$ をサイトとし、
$u : \mathcal{C} \to \mathcal{D}$ を連続とする。
Section \ref{sites-section-functoriality-PSh} のすべての圏
$(\mathcal{I}_V^u)^{opp}$ がフィルター圏であると仮定する。
このとき $u$ はサイトの射 $\mathcal{D} \to \mathcal{C}$ を定める。
言い換えると、$u_s$ は完全である。
\end{lemma}

\begin{proof}
$u_s$ は $u^s$ の左随伴なので、$u_s$ は右完全である。Categories, Lemma
\ref{categories-lemma-exact-adjoint} を参照せよ。
したがって $u_s$ が左完全であることを示せば十分である。
言い換えると、$u_s$ が有限極限と可換することを示せばよい。
圏 $\mathcal{I}_Y^{opp}$ はフィルター圏なので、Categories, Lemma
\ref{categories-lemma-directed-commutes} により $u_p$ は有限極限と可換する
（これには $\textit{PSh}$ における極限の記述も用いる。
Section \ref{sites-section-limits-colimits-PSh} を参照せよ）。
また層化も有限極限と可換するので（Lemma
\ref{sites-lemma-sheafification-exact}）、$u_s = \# \circ u_p$ から結論が従う。
\end{proof}

\begin{proposition}
\label{sites-proposition-get-morphism}
$\mathcal{C}$ と $\mathcal{D}$ をサイトとし、
$u : \mathcal{C} \to \mathcal{D}$ を連続とする。
さらに次を仮定する。
\begin{enumerate}
\item 圏 $\mathcal{C}$ は終対象 $X$ をもち、$u(X)$ は
$\mathcal{D}$ の終対象である。
\item 圏 $\mathcal{C}$ はファイバー積をもち、$u$ はそれらと可換する。
\end{enumerate}
このとき $u$ はサイトの射 $\mathcal{D} \to \mathcal{C}$ を定める。
言い換えると、$u_s$ は完全である。
\end{proposition}

\begin{proof}
Lemmas \ref{sites-lemma-directed} と \ref{sites-lemma-directed-morphism} から従う。
\end{proof}

\begin{remark}
\label{sites-remark-explain-left-exact}
上の Proposition \ref{sites-proposition-get-morphism} の条件は、$u$ が左完全、
すなわち有限極限と可換するという条件と同値である。Categories, Lemmas
\ref{categories-lemma-finite-limits-exist} と
\ref{categories-lemma-characterize-left-exact} を参照せよ。
例においては終対象とファイバー積の方が幾何学的意味をもちやすいと思われるので、
この形で述べる方が自然に見える。
\end{remark}

\noindent
Lemma \ref{sites-lemma-continuous-with-continuous-left-adjoint} では、
連続関手がサイトの射を生じることを示す別の方法を与える。

\begin{remark}
\label{sites-remark-quasi-continuous-morphism-sites}
（初読時には飛ばしてよい。）
$\mathcal{C}$ と $\mathcal{D}$ をサイトとする。Definition
\ref{sites-definition-morphism-sites} と同様に、{\it サイトの準射
$f : \mathcal{D} \to \mathcal{C}$} とは、$u_s$ が完全となるような準連続関手
$u : \mathcal{C} \to \mathcal{D}$（Remark
\ref{sites-remark-quasi-continuous} を参照）によって与えられるものをいう。
この設定における Proposition \ref{sites-proposition-get-morphism} の類似は、
「連続」を「準連続」に、「射」を「準射」に置き換えることで得られる。
証明は文字どおり同じである。
\end{remark}

\noindent
Definition \ref{sites-definition-morphism-sites} において、$u_s$ が完全であるという
条件は省略できない。例えば、$u$ が連続であることだけを仮定すると、
次の補題の結論は成り立つとは限らない。

\begin{lemma}
\label{sites-lemma-morphism-of-sites-covering}
$f : \mathcal{D} \to \mathcal{C}$ を、関手
$u : \mathcal{C} \to \mathcal{D}$ によって与えられるサイトの射とする。
$\mathcal{D}$ の任意の対象 $V$ に対して被覆 $\{V_j \to V\}$ が存在し、
各 $j$ について、$\mathcal{C}$ のある対象 $U_j$ への射
$V_j \to u(U_j)$ が存在する。
\end{lemma}

\begin{proof}
$f^{-1} = u_s$ は完全なので $f^{-1}* = *$ である。ここで $*$ は層の圏の
終対象を表す（Example \ref{sites-example-singleton-sheaf}）。
$f^{-1}* = u_s*$ は $u_p*$ の層化なので、被覆 $\{V_j \to V\}$ が存在して
$(u_p*)(V_j)$ は空でない。$(u_p*)(V_j)$ は、対象が射
$V_j \to u(U)$ である圏 $\mathcal{I}^u_{V_j}$ にわたる余極限だから、
補題が従う。
\end{proof}



























\section{トポス}
\label{sites-section-topoi}

\noindent
ここでは本章の目的に適したトポスの定義を与える。
すなわち、トポスとはサイト上の層の圏である。
トポスを指定するにはサイトを指定すればよい。
トポスとサイトの本質的な違いは射の定義にある。
実際、台となるサイトの射から来ないトポスの射が数多く存在することが分かる。

\begin{definition}[トポス]
\label{sites-definition-topos}
{\it トポス}とは、サイト $\mathcal{C}$ 上の層の圏
$\Sh(\mathcal{C})$ のことである。
\begin{enumerate}
\item $\mathcal{C}$ と $\mathcal{D}$ をサイトとする。
$\Sh(\mathcal{D})$ から $\Sh(\mathcal{C})$ への
{\it トポスの射} $f$ とは、関手の対
$f_* : \Sh(\mathcal{D}) \to \Sh(\mathcal{C})$ と
$f^{-1} : \Sh(\mathcal{C}) \to \Sh(\mathcal{D})$ であって、
次を満たすものをいう。
\begin{enumerate}
\item 二変数関手的に
$$
\Mor_{\Sh(\mathcal{D})}(f^{-1}\mathcal{G}, \mathcal{F})
=
\Mor_{\Sh(\mathcal{C})}(\mathcal{G}, f_*\mathcal{F})
$$
が成り立つ。
\item 関手 $f^{-1}$ は有限極限と可換する。すなわち左完全である。
\end{enumerate}
\item $\mathcal{C}$、$\mathcal{D}$、$\mathcal{E}$ をサイトとする。
トポスの射 $f :\Sh(\mathcal{D}) \to \Sh(\mathcal{C})$ と
$g :\Sh(\mathcal{E}) \to \Sh(\mathcal{D})$ が与えられたとき、
{\it 合成 $f\circ g$} とは、関手
$(f \circ g)_* = f_* \circ g_*$ と
$(f \circ g)^{-1} = g^{-1} \circ f^{-1}$ によって定義されるトポスの射をいう。
\end{enumerate}
\end{definition}

\noindent
$\alpha : \mathcal{S}_1 \to \mathcal{S}_2$ が（場合によっては「大きな」）
圏の同値であるとする。$\mathcal{S}_1$ と $\mathcal{S}_2$ がトポスなら、
$f_* = \alpha$ とおき、$f^{-1}$ を $\alpha$ の準逆とすると、
トポスの射 $f : \mathcal{S}_1 \to \mathcal{S}_2$ が得られる。
さらに、この射はトポスの $2$-圏における同値である
（Section \ref{sites-section-2-category} を参照）。
したがって $\mathcal{S}$ がサイト上の層の圏と同値なら
（必ずしもサイト上の層の圏そのものに等しくなくても）、
「$\mathcal{S}$ はトポスである」と言うことには意味がある。
以下では、この記法の濫用をときどき用いる。

\medskip\noindent
{\it 空トポス}とは、$\mathcal{C}$ が空圏であるサイト
$\mathcal{C}$ 上の層のトポスである。このサイトを $\emptyset$ と書くことがある。
これは一意な層をもつサイトである（$\emptyset$ は対象をもたないため）。
したがって $\Sh(\emptyset)$ は、対象を一つ、射を一つだけもつ圏と同値である。

\medskip\noindent
{\it 点トポス}とは、対象を一つ $pt$、射を一つ $\text{id}_{pt}$ だけもち、
被覆が $\{\text{id}_{pt}\}$ だけであるサイト $\mathcal{C}$ 上の層のトポスである。
このサイトを単に $pt$ と書く。層の圏 $ = $ 前層の圏 $ = $ 集合の圏であることは
明らかである。式で書けば $\Sh(pt) = \textit{Sets}$ である。

\medskip\noindent
$\mathcal{C}$ と $\mathcal{D}$ をサイトとし、
$f : \Sh(\mathcal{D}) \to \Sh(\mathcal{C})$ をトポスの射とする。
$f_*$ はすべての極限と可換し、$f^{-1}$ はすべての余極限と可換することに
注意しよう。Categories, Lemma \ref{categories-lemma-adjoint-exact} を参照せよ。
特に、上の定義における $f^{-1}$ に関する条件は、$f^{-1}$ が完全であることを
保証する。トポスの射は Lemma \ref{sites-lemma-cocontinuous-morphism-topoi} または
次の補題を用いて構成されることが多い。

\begin{lemma}
\label{sites-lemma-morphism-sites-topoi}
関手 $u : \mathcal{C} \to \mathcal{D}$ に対応するサイトの射
$f : \mathcal{D} \to \mathcal{C}$ が与えられたとき、関手の対
$(f^{-1} = u_s, f_* = u^s)$ はトポスの射
$\Sh(\mathcal{D}) \to \Sh(\mathcal{C})$ である。
\end{lemma}

\begin{proof}
Definition \ref{sites-definition-morphism-sites} から明らかである。
\end{proof}

\begin{remark}
\label{sites-remark-pt-topos}
トポス $\Sh(pt)$ を生じるサイトは数多くある。有用な例を一つ挙げよう。
$S$ を、少なくとも一つの空でない元を含む（集合からなる）集合とする。
$\mathcal{S}$ を、対象が $S$ の元であり、射が任意の集合写像である圏とする。
$\mathcal{S}$ がファイバー積をもつと仮定する。
例えば $S = \mathcal{P}(\text{infinite set})$ が任意の無限集合の冪集合なら、
この仮定は成り立つ（集合論の演習）。全射的な写像族を被覆と宣言して
$\mathcal{S}$ をサイトにする（Sets, Section
\ref{sets-section-coverings-site} と同様に、被覆族の適切で十分大きな集合を選ぶ）。
$\Sh(\mathcal{S})$ は集合の圏と同値であると主張する。

\medskip\noindent
まず $S$ が一点集合 $e \in S$ を含む場合にこれを証明する。
この場合、関手
\begin{equation}
\label{sites-equation-sheaves-pt-sets}
i^{-1}\mathcal{F} = \mathcal{F}(e), \quad
i_*E = (U \mapsto \Mor_{\textit{Sets}}(U, E))
\end{equation}
によって与えられるトポスの同値
$i : \Sh(pt) \to \Sh(\mathcal{S})$ がある。
実際、$\mathcal{F}$ を $\mathcal{S}$ 上の層とする。
任意の $U \in \Ob(\mathcal{S}) = S$ に対し、被覆
$\{\varphi_u : e \to U\}_{u \in U}$ が得られる。ここで $\varphi_u$ は
$e$ の唯一の元を $u \in U$ へ写す。この場合、層条件は
$\mathcal{F}(U) = \prod_{u \in U} \mathcal{F}(e)$ を含意する。
言い換えると、
$\mathcal{F}(U) = \Mor_{\textit{Sets}}(U, \mathcal{F}(e))$ である。
さらに、この規則は制限写像と両立する。したがって関手
$$
i_* :
\textit{Sets} = \Sh(pt)
\longrightarrow
\Sh(\mathcal{S}), \quad
E \longmapsto (U \mapsto \Mor_{\textit{Sets}}(U, E))
$$
は圏の同値であり、その逆は上で与えた関手 $i^{-1}$ である。

\medskip\noindent
$\mathcal{S}$ が一点集合を含まない場合にも、上で定義した関手 $i_*$ には意味がある。
この場合にも同値であることを示すため、任意の空でない
$\tilde e \in S$ と、像が一点集合である写像
$\varphi : \tilde e \to \tilde e$ を選ぶ。任意の層 $\mathcal{F}$ に対して
$$
\mathcal{F}(e) :=
\Im(
\mathcal{F}(\varphi) :
\mathcal{F}(\tilde e)
\longrightarrow
\mathcal{F}(\tilde e)
)
$$
とおき、これが $i_*$ の準逆であることを示せばよい。詳細は省略する。
\end{remark}

\begin{remark}
\label{sites-remark-morphism-topoi-big}
（トポスの射に関する集合論的問題。初読時には飛ばしてよい。）
上で定義したトポスの射は集合ではなく類である。
言い換えると、それは数学的対象ではなく数式によって与えられる。
与えられた二つのトポス間の射全体を考えることはできるが、
これを数学的対象として導入するのは得策ではない。
一方、$\mathcal{C}$ と $\mathcal{D}$ が与えられたサイトであるとし、関手
$\Phi : \mathcal{C} \to \Sh(\mathcal{D})$ を考える。
これは集合、すなわち数学的対象である。$\Phi$ について順に次を問える。
\begin{enumerate}
\item $\mathcal{D}$ 上の層 $\mathcal{F}$ が与えられたとき、規則
$U \mapsto \Mor_{\Sh(\mathcal{D})}(\Phi(U), \mathcal{F})$ は
$\mathcal{C}$ 上の層を定めるか。定めるなら、これは関手
$\Phi_* : \Sh(\mathcal{D}) \to \Sh(\mathcal{C})$ を定める。
\item $\Phi_*$ は左随伴をもつか。もつなら、その左随伴を $\Phi^{-1}$ と書く。
\item $\Phi^{-1}$ は完全か。
\end{enumerate}
最後の問いにも「はい」と答えられるなら、トポスの射
$(\Phi_*, \Phi^{-1})$ が得られる。さらに、任意のトポスの射
$(f_*, f^{-1})$ に対して $\Phi(U) = f^{-1}(h_U^\#)$ とおけば、
上のような関手 $\Phi$ が得られ、随伴性と両立する同型
$f_* \cong \Phi_*$ と $f^{-1} \cong \Phi^{-1}$ をもつ。
要するに、トポスの射の代わりに $\Phi$ 全体を用いることにより、
(a) トポスの射という概念を数学的対象で置き換え、(b) $\Phi$ 全体は
（類の集まりではなく）類をなす。もちろん、さらに多くのことが言える。
例えば、上の条件 (2) と (3) の意味をより正確に調べることができる。
Section \ref{sites-section-points} では、トポスの点の場合にこれを行う。
\end{remark}

\begin{remark}
\label{sites-remark-quasi-continuous-morphism-topoi}
（初読時には飛ばしてよい。）
$\mathcal{C}$ と $\mathcal{D}$ をサイトとする。
サイトの準射 $f : \mathcal{D} \to \mathcal{C}$
（Remark \ref{sites-remark-quasi-continuous-morphism-sites} を参照）は、
Lemma \ref{sites-lemma-morphism-sites-topoi} とまったく同じ方法で、
$\Sh(\mathcal{D})$ から $\Sh(\mathcal{C})$ へのトポスの射 $f$ を生じる。
\end{remark}

\section{G-集合と射}
\label{sites-section-G-sets-morphisms}

\noindent
$\varphi : G \to H$ を群の準同型とする。Example
\ref{sites-example-site-on-group} と Section
\ref{sites-section-example-sheaf-G-sets} のように（適切な）サイト
$\mathcal{T}_G$ と $\mathcal{T}_H$ を選ぶ。
$u : \mathcal{T}_H \to \mathcal{T}_G$ を、$H$-集合 $U$ に対し、
台集合は同じであるが、$g \in G$ と $\xi \in U$ に対する $G$-作用を
$g \cdot \xi = \varphi(g)\xi$ によって定めた $G$-集合 $U_\varphi$ を
対応させる関手とする。$u$ が有限極限と可換し、連続であることは明らかである
\footnote{集合論的注意：まず $\mathcal{T}_H$ を選ぶ。次に、
$u(\mathcal{T}_H)$ を含み、$\mathcal{T}_H$ の各被覆が
$\mathcal{T}_G$ の被覆に対応するように $\mathcal{T}_G$ を選ぶ。
これは Sets, Lemmas
\ref{sets-lemma-sets-with-group-action}、
\ref{sets-lemma-what-is-in-it-G-sets}、および
\ref{sets-lemma-coverings-site} により可能である。}。
Proposition \ref{sites-proposition-get-morphism} と Lemma
\ref{sites-lemma-morphism-sites-topoi} を適用すると、$\varphi$ に付随するトポスの射
$$
f : \Sh(\mathcal{T}_G) \longrightarrow \Sh(\mathcal{T}_H)
$$
が得られる。Proposition \ref{sites-proposition-sheaves-on-group} を用いると、
随伴関手の対
$$
f_* : G\textit{-Sets} \longrightarrow H\textit{-Sets}, \quad
f^{-1} : H\textit{-Sets} \longrightarrow G\textit{-Sets}.
$$
が得られる。この場合にこれらの関手が何であるかを求めよう。

\medskip\noindent
まず $f_*$ の公式を求める。$G$-集合 $S$ に対応する
$\mathcal{T}_G$ 上の層 $\mathcal{F}_S$ は、規則
$\mathcal{F}_S(U) = \Mor_G(U, S)$ によって与えられることを思い出そう。
一方、$\mathcal{T}_H$ 上の層 $\mathcal{G}$ に対応する $H$-集合は、
規則 $\mathcal{G}({}_HH)$ によって与えられる。したがって
$$
f_*S = \Mor_{G\textit{-Sets}}(({}_HH)_\varphi, S).
$$
である。これをもう少し具体的に書けば
% P01-E0068
$$
f_*S = \{ a : H \to S \mid a(gh) = ga(h) \}
$$
となり、任意の元 $a \in f_*S$ に対する左 $H$-作用は
$(h \cdot a)(h') = a(h'h)$ によって与えられる。

\medskip\noindent
次に $f^{-1}$ を明示的に計算する。$\mathcal{T}_G$ と
$\mathcal{T}_H$ 上の位相は準標準的なので、すべての表現可能前層は層である。
さらに、$\mathcal{T}_H$ の対象 $V$ が与えられると、
$f^{-1}h_V$ は $h_{u(V)}$ に等しい（Lemma
\ref{sites-lemma-pullback-representable-sheaf} を参照）。
したがって表現可能層に対して $f^{-1}S = S_\varphi$ である。
$\mathcal{T}_H$ 上のすべての層は表現可能層の余積なので、これは一般にも成り立つ。
ゆえに任意の $H$-集合 $T$ に対して
$$
f^{-1}T = T_\varphi.
$$
である。$f^{-1}$ と $f_*$ の随伴性は、上の $f_*S$ に対する公式
$$
\Mor_{G\textit{-Sets}}(T_\varphi, S) =
\Mor_{H\textit{-Sets}}(T, f_*S)
$$
によって示される。これは直接証明できる。さらに、これにより
$(f^{-1}, f_*)$ が随伴対をなし、$f^{-1}$ が完全であることは明らかである。
したがって上の構成とは別に、トポスの射
$f : G\textit{-Sets} \to H\textit{-Sets}$ をこの方法で直接定義することもできる。













\section{準コンパクト対象と余極限}
\label{sites-section-quasi-compact}

\noindent
位相空間の場合と同じ言葉を使えるように、次の用語を導入する。

\begin{definition}
\label{sites-definition-quasi-compact}
$\mathcal{C}$ をサイトとする。$\mathcal{C}$ の対象 $U$ が
{\it 準コンパクト}であるとは、$\mathcal{C}$ の被覆
$\mathcal{U} = \{U_i \to U\}_{i \in I}$ が与えられたとき、別の被覆
$\mathcal{V} = \{V_j \to U\}_{j \in J}$ と、$\text{id} : U \to U$、
$\alpha : J \to I$、および $V_j \to U_{\alpha(j)}$ によって与えられる
終域を固定した射の族の射 $\mathcal{V} \to \mathcal{U}$
（Definition \ref{sites-definition-morphism-coverings} を参照）が存在し、
$\alpha$ の像が $I$ の有限部分集合となることをいう。
\end{definition}

\noindent
もちろん、通常の概念から $U$ が準コンパクトであることが従う。

\begin{lemma}
\label{sites-lemma-conclude-quasi-compact}
$\mathcal{C}$ をサイトとし、$U$ を $\mathcal{C}$ の対象とする。
次の条件を考える。
\begin{enumerate}
\item $U$ は準コンパクトである。
\item $\mathcal{C}$ の任意の被覆 $\{U_i \to U\}_{i \in I}$ に対して、
$\mathcal{U}$ を細分する $\mathcal{C}$ の有限被覆
$\{V_j \to U\}_{j = 1, \ldots, m}$ が存在する。
\item $\mathcal{C}$ の任意の被覆 $\{U_i \to U\}_{i \in I}$ に対して、
$\{U_i \to U\}_{i \in I'}$ が $\mathcal{C}$ の被覆となる有限部分集合
$I' \subset I$ が存在する。
\end{enumerate}
このとき常に (3) $\Rightarrow$ (2) $\Rightarrow$ (1) であるが、
逆の含意は一般には成り立たない。
\end{lemma}

\begin{proof}
含意は定義から直ちに従う。$X = [0, 1] \subset \mathbf{R}$ を
（通常の $\epsilon$-$\delta$ 位相をもつ）位相空間とする。
$\mathcal{C}$ を $X$ の開部分空間を対象、包含を射とし、通常の開被覆を
備えた圏とする（Example \ref{sites-example-site-topological} と比較せよ）。
ただし、$\{U \to U\}$ の形の被覆を除くすべての有限被覆を排除するように、
$\mathcal{C}$ の被覆の概念を変更する。これが Definition
\ref{sites-definition-site} の意味でサイトになることは容易に分かる。
このサイトでは対象 $X = U = [0, 1]$ は Definition
\ref{sites-definition-quasi-compact} の意味で準コンパクトであるが、$U$ は (2) を満たさない。
(2) は成り立つが (3) は成り立たない例を作ることは読者に任せる。
\end{proof}

\noindent
準コンパクト対象のトポス論的な意味は次のとおりである。

\begin{lemma}
\label{sites-lemma-quasi-compact}
$\mathcal{C}$ をサイトとし、$U$ を $\mathcal{C}$ の対象とする。
次は同値である。
\begin{enumerate}
\item $U$ は準コンパクトである。
\item 層の任意の全射 $\coprod_{i \in I} \mathcal{F}_i \to h_U^\#$ に対して、
$\coprod_{i \in J} \mathcal{F}_i \to h_U^\#$ が全射となる有限部分集合
$J \subset I$ が存在する。
\end{enumerate}
\end{lemma}

\begin{proof}
(1) を仮定し、$\coprod_{i \in I} \mathcal{F}_i \to h_U^\#$ を全射とする。
このとき $\text{id}_U$ は $U$ 上の $h_U^\#$ の切断である。
したがって被覆 $\{U_a \to U\}_{a \in A}$ が存在し、各 $a \in A$ に対して
$U_a$ 上の $\coprod_{i \in I} \mathcal{F}_i$ の切断 $s_a$ で
$\text{id}_U$ へ写るものが存在する。前層の余積を層化して層の余積を構成すること
（Lemma \ref{sites-lemma-colimit-sheaves}）により、各 $a$ に対して被覆
$\{U_{ab} \to U_a\}_{b \in B_a}$ が存在し、すべての $b \in B_a$ に対して
$\iota(b) \in I$ と、$U_{ab}$ 上の $\mathcal{F}_{\iota(b)}$ の切断
$s_{b}$ で $\text{id}_U|_{U_{ab}}$ へ写るものが存在する。
したがって被覆 $\{U_a \to U\}_{a \in A}$ を
$\{U_{ab} \to U\}_{a \in A, b \in B_a}$ で置き換えることにより、
写像 $\iota : A \to I$ と、各 $a \in A$ に対して $U_a$ 上の
$\mathcal{F}_{\iota(a)}$ の切断 $s_a$ で $\text{id}_U$ へ写るものが
存在すると仮定できる。$U$ は準コンパクトなので、被覆
$\{V_c \to U\}_{c \in C}$、像が有限な写像 $\alpha : C \to A$、および
$U$ 上の射 $V_c \to U_{\alpha(c)}$ が存在する。このとき
$J = \Im(\iota \circ \alpha) \subset I$ が求めるものである。
実際、$\coprod_{c \in C} h_{V_c}^\# \to h_U^\#$ は全射であり
（Lemma \ref{sites-lemma-covering-surjective-after-sheafification}）、
$\coprod_{i \in J} \mathcal{F}_i \to h_U^\#$ を経由する。
（ここでは、合成
% P01-E0069
$h_{V_c}^\# \to h_{U_{\alpha(c)}}
\xrightarrow{s_{\alpha(c)}} \mathcal{F}_{\iota(\alpha(c))} \to h_U^\#$
が、射 $V_c \to U$ から来る写像 $h_{V_c}^\# \to h_U^\#$ であることを
用いた。これは $s_{\alpha(c)}$ が $\text{id}_U|_{U_{\alpha(c)}}$ へ
写ることによる。）

\medskip\noindent
(2) を仮定する。$\{U_i \to U\}_{i \in I}$ を被覆とする。
Lemma \ref{sites-lemma-covering-surjective-after-sheafification} により
$\coprod_{i \in I} h_{U_i}^\# \to h_U^\#$ は全射である。
したがって $\coprod_{j \in J} h_{U_j}^\# \to h_U^\#$ が全射となる
有限部分集合 $J \subset I$ が得られる。上と同様に論じると、
$\mathcal{C}$ における $U$ の被覆 $\{V_c \to U\}_{c \in C}$ と写像
$\iota : C \to J$ が得られ、$\text{id}_U$ は $V_c$ 上の
% P01-E0070
$h_{U_{\iota(c)}}^\#$ の切断 $s_c$ へ持ち上がる。
被覆をさらに細分すると、$\text{id}_U$ へ写る
$s_c \in h_{U_{\iota(c)}}(V_c)$ と仮定できる。
このとき $s_c : V_c \to U_{\iota(c)}$ は $U$ 上の射であり、結論を得る。
\end{proof}

\noindent
上の補題は次の定義を動機づける。

\begin{definition}
\label{sites-definition-quasi-compact-topos}
トポス $\Sh(\mathcal{C})$ の対象 $\mathcal{F}$ が{\it 準コンパクト}であるとは、
$\Sh(\mathcal{C})$ の任意の全射
$\coprod_{i \in I} \mathcal{F}_i \to \mathcal{F}$ に対して、
$\coprod_{i \in J} \mathcal{F}_i \to \mathcal{F}$ が全射となる有限部分集合
$J \subset I$ が存在することをいう。トポス $\Sh(\mathcal{C})$ が
{\it 準コンパクト}であるとは、その終対象 $*$ が準コンパクト対象であることをいう。
\end{definition}

\noindent
Lemma \ref{sites-lemma-quasi-compact} により、サイト $\mathcal{C}$ が終対象 $X$ を
もつなら、$\Sh(\mathcal{C})$ が準コンパクトであることと $X$ が
準コンパクトであることは同値である。

\begin{lemma}
\label{sites-lemma-image-of-quasi-compact}
\begin{slogan}
層の全射は準コンパクト性を伝える。
\end{slogan}
$\mathcal{C}$ をサイトとする。
\begin{enumerate}
\item $U \to V$ が $\mathcal{C}$ の射で、$h_U^\# \to h_V^\#$ が全射かつ
$U$ が準コンパクトなら、$V$ は準コンパクトである。
\item $\mathcal{F} \to \mathcal{G}$ が集合の層の全射で、$\mathcal{F}$ が
準コンパクトなら、$\mathcal{G}$ は準コンパクトである。
\end{enumerate}
\end{lemma}

\begin{proof}
省略する。
\end{proof}

\begin{lemma}
\label{sites-lemma-coproduct-quasi-compact}
\begin{slogan}
% P01-E0071
層の有限余積は準コンパクト性を保つ。
\end{slogan}
$\mathcal{C}$ をサイトとする。$n \geq 1$ であり、
$\mathcal{F}_1, \ldots, \mathcal{F}_n$ が $\mathcal{C}$ 上の準コンパクト層なら、
$\coprod_{i = 1, \ldots, n} \mathcal{F}_i$ は準コンパクトである。
\end{lemma}

\begin{proof}
省略する。
\end{proof}

\noindent
次の二つの補題は、サイトに対する Sheaves, Lemma
\ref{sheaves-lemma-directed-colimits-sections} の類似である。

\begin{lemma}
\label{sites-lemma-directed-colimits-sections}
$\mathcal{C}$ をサイトとする。
$\mathcal{I} \to \Sh(\mathcal{C})$、$i \mapsto \mathcal{F}_i$ を
集合の層のフィルター図式とする。$U \in \Ob(\mathcal{C})$ とし、標準写像
$$
\Psi :
\colim_i \mathcal{F}_i(U)
\longrightarrow
\left(\colim_i \mathcal{F}_i\right)(U)
$$
を考える。上で導入した用語を用いると、次が成り立つ。
\begin{enumerate}
\item すべての遷移写像が単射なら、任意の $U$ に対して $\Psi$ は単射である。
\item $U$ が準コンパクトなら、$\Psi$ は単射である。
\item $U$ が準コンパクトで、すべての遷移写像が単射なら、$\Psi$ は同型である。
\item $U$ が被覆 $\{U_j \to U\}_{j \in J}$ の共終系をもち、
$J$ が有限で、すべての $j, j' \in J$ に対して
$U_j \times_U U_{j'}$ が準コンパクトなら、$\Psi$ は全単射である。
\end{enumerate}
\end{lemma}

\begin{proof}
すべての遷移写像が単射であると仮定する。この場合、前層
$\mathcal{F}' : V \mapsto \colim_i \mathcal{F}_i(V)$ は分離的である
（Definition \ref{sites-definition-separated} を参照）。Lemma
\ref{sites-lemma-colimit-sheaves} により
$(\mathcal{F}')^\# = \colim_i \mathcal{F}_i$ である。
Theorem \ref{sites-theorem-plus} により
$\mathcal{F}' \to (\mathcal{F}')^\#$ は単射である。これで (1) が従う。

\medskip\noindent
$U$ が準コンパクトであると仮定する。$s \in \mathcal{F}_i(U)$ と
$s' \in \mathcal{F}_{i'}(U)$ が左辺の元を定め、それらの $\Psi$ による像が
等しいとする。これは、被覆 $\{U_a \to U\}_{a \in A}$ を選び、
各 $a \in A$ に対して $i_a \in I$、$i_a \geq i$、$i_a \geq i'$ を選んで
$\varphi_{ii_a}(s) = \varphi_{i'i_a}(s')$ とできることを意味する。
$U$ は準コンパクトなので、被覆 $\{V_b \to U\}_{b \in B}$、
像が有限な写像 $\alpha : B \to A$、および $U$ 上の射
$V_b \to U_{\alpha(b)}$ を選べる。$\alpha$ の像は有限なので、すべての
$i_{\alpha(b)}$ のすべてに対して $\geq$ となる $i''\in I$ を選ぶ。すると、すべての $b \in B$ に対して
% P01-E0072
$\varphi_{ii''}(s)$ と $\varphi_{i'i''}(s)$ は $V_b$ 上で一致し、したがって
$\varphi_{ii''}(s) = \varphi_{i'i''}(s)$ である。これで (2) が従う。

\medskip\noindent
$U$ が準コンパクトで、すべての遷移写像が単射であると仮定する。
$s$ を $\Psi$ の終域の元とする。被覆 $\{U_a \to U\}_{a \in A}$ が存在し、
各 $a \in A$ に対して添字 $i_a \in I$ と切断
$s_a \in \mathcal{F}_{i_a}(U_a)$ が存在して、すべての $a \in A$ について
$s|_{U_a}$ は $s_a$ から来る。$U$ は準コンパクトなので、被覆
$\{V_b \to U\}_{b \in B}$、像が有限な写像 $\alpha : B \to A$、および
$U$ 上の射 $V_b \to U_{\alpha(b)}$ を選べる。$\alpha$ の像は有限なので、
すべての $i_{\alpha(b)}$ に対して $\geq$ となる $i \in I$ を選ぶ。
(1) により、切断
$s_b = \varphi_{i_{\alpha(b)} i}(s_{\alpha(b)})|_{V_b}$ は
$V_b \times_U V_{b'}$ 上で一致する。したがってこれらは、$\Psi$ により $s$ へ
写る切断 $s' \in \mathcal{F}_i(U)$ に貼り合わされる。これで (3) が従う。

\medskip\noindent
(4) の仮定を置く。Lemma \ref{sites-lemma-conclude-quasi-compact} により対象 $U$ は
準コンパクトなので、(2) により $\Psi$ は単射である。全射性を示すため、
$s$ を $\Psi$ の終域の元とする。仮定により、すべての
$1 \leq j, j' \leq m$ について $U_j \times_U U_{j'}$ が準コンパクトである
有限被覆 $\{U_j \to U\}_{j = 1, \ldots, m}$ が存在し、
各 $j$ に対して添字 $i_j \in I$ と $s_j \in \mathcal{F}_{i_j}(U_j)$ が存在して、
すべての $j$ について $s|_{U_j}$ は $s_j$ の像である。
$U_j \times_U U_{j'}$ は準コンパクトなので (2) を適用でき、
$i_{jj'} \in I$、$i_{jj'} \geq i_j$、$i_{jj'} \geq i_{j'}$ が存在して、
$\varphi_{i_ji_{jj'}}(s_j)$ と $\varphi_{i_{j'}i_{jj'}}(s_{j'})$ は
$U_j \times_U U_{j'}$ 上で一致する。すべての $i_{jj'}$ 以上の添字
$i \in I$ を選ぶ。すると $\mathcal{F}_i$ の切断
$\varphi_{i_ji}(s_j)$ は $U$ 上の $\mathcal{F}_i$ の切断へ貼り合わされる。
この切断は所望どおり元 $s$ へ写る。
\end{proof}

\begin{lemma}
\label{sites-lemma-directed-colimits-global-sections}
$\mathcal{C}$ をサイトとする。
$\mathcal{I} \to \Sh(\mathcal{C})$、$i \mapsto \mathcal{F}_i$ を
集合の層のフィルター図式とする。標準写像
$$
\Psi :
\colim_i \Gamma(\mathcal{C}, \mathcal{F}_i)
\longrightarrow
\Gamma(\mathcal{C}, \colim_i \mathcal{F}_i)
$$
を考える。次が成り立つ。
\begin{enumerate}
\item すべての遷移写像が単射なら、$\Psi$ は単射である。
\item $\Sh(\mathcal{C})$ が準コンパクトなら、$\Psi$ は単射である。
\item $\Sh(\mathcal{C})$ が準コンパクトで、すべての遷移写像が単射なら、
$\Psi$ は同型である。
\item 次の性質をもつ集合 $S \subset \Ob(\Sh(\mathcal{C}))$ が存在すると仮定する。
\begin{enumerate}
\item 任意の全射 $\mathcal{F} \to *$ に対して、$\mathcal{K} \in S$ と写像
$\mathcal{K} \to \mathcal{F}$ が存在し、$\mathcal{K} \to *$ は全射である。
\item $\mathcal{K} \in S$ に対して、積 $\mathcal{K} \times \mathcal{K}$ は
準コンパクトである。
\end{enumerate}
このとき $\Psi$ は全単射である。
\end{enumerate}
\end{lemma}

\begin{proof}
(1) の証明。すべての遷移写像が単射であると仮定する。
この場合、前層 $\mathcal{F}' : V \mapsto \colim_i \mathcal{F}_i(V)$ は
分離的である（Definition \ref{sites-definition-separated} を参照）。
Lemma \ref{sites-lemma-colimit-sheaves} により
$(\mathcal{F}')^\# = \colim_i \mathcal{F}_i$ である。
Theorem \ref{sites-theorem-plus} により
$\mathcal{F}' \to (\mathcal{F}')^\#$ は単射である。これで (1) が従う。

\medskip\noindent
(2) の証明。$\Sh(\mathcal{C})$ が準コンパクトであると仮定する。
$\Sh(\mathcal{C})$ のすべての $\mathcal{F}$ に対して
$\Gamma(\mathcal{C}, \mathcal{F}) = \Mor(*, \mathcal{F})$ であることを
思い出そう。$a_i, b_i : * \to \mathcal{F}_i$ とし、$i' \geq i$ に対して
$a_{i'}, b_{i'} : * \to \mathcal{F}_{i'}$ を、系の遷移写像との合成とする。
$a = \colim_{i' \geq i} a_{i'}$ とおき、$b$ についても同様におく。
$i' \geq i$ に対して
$$
E_{i'} = \text{Equalizer}(a_{i'}, b_{i'}) \subset *
\quad\text{and}\quad
E = \text{Equalizer}(a, b) \subset *
$$
と書く。Categories, Lemma \ref{categories-lemma-directed-commutes} により
$E = \colim_{i' \geq i} E_{i'}$ である。したがって
$\coprod_{i' \geq i} E_{i'} \to E$ は層の全射である。
よって $E = *$、すなわち $a = b$ なら、$*$ は準コンパクトなので、
ある $i' \geq i$ に対して $E_{i'} = *$ である。
したがって、ある $i' \geq i$ に対して $a_{i'} = b_{i'}$ となる。
これで (2) が従う。

\medskip\noindent
(3) の証明。$\Sh(\mathcal{C})$ が準コンパクトで、すべての遷移写像が
単射であると仮定する。$a : * \to \colim \mathcal{F}_i$ を写像とする。
このとき $E_i = a^{-1}(\mathcal{F}_i) \subset *$ は部分層であり、
$\colim E_i = *$ である（上の参照による）。したがって、ある $i$ について
$E_i = *$ となり、所望どおり $a$ の像は $\mathcal{F}_i$ に含まれる。

\medskip\noindent
(4) の証明。$S \subset \Ob(\Sh(\mathcal{C}))$ が (4)(a), (b) を満たすとする。
(4)(a) を $\text{id} : * \to *$ に適用すると、$\mathcal{K} \in S$ で
$\mathcal{K} \to *$ が全射となるものが存在する。写像
$\mathcal{K} \times \mathcal{K} \to \mathcal{K} \to *$ は全射である。
(4)(b) と Lemma \ref{sites-lemma-image-of-quasi-compact} により、
$\mathcal{K}$ と $\Sh(\mathcal{C})$ は準コンパクトである。
したがって (2) により $\Psi$ は単射である。
$\mathcal{F} = \colim \mathcal{F}_i$ とおく。
$s : * \to \mathcal{F}$ を余極限の大域切断とする。
$\coprod \mathcal{F}_i \to \mathcal{F}$ は全射なので、射影
$$
\coprod\nolimits_{i \in I} * \times_{s, \mathcal{F}} \mathcal{F}_i \to *
$$
は全射である。(4)(a) により、$\mathcal{K} \in S$ と写像
$\mathcal{K} \to \coprod_{i \in I} * \times_{s, \mathcal{F}} \mathcal{F}_i$ が
得られ、$\mathcal{K} \to *$ は全射である。$\mathcal{K}$ は準コンパクトなので、
ある有限部分集合 $I' \subset I$ に対する分解
$\mathcal{K} \to \coprod_{i' \in I'} * \times_{s, \mathcal{F}} \mathcal{F}_{i'}$
が得られる。$i \in I$ を有限部分集合 $I'$ の上界とする。
遷移写像は写像
$\coprod_{i' \in I'} \mathcal{F}_{i'} \to \mathcal{F}_i$ を定める。
これはさらに写像
$\mathcal{K} \to * \times_{s, \mathcal{F}} \mathcal{F}_i$ を生じる。
言い換えると、$\mathcal{K} \in S$、全射 $\mathcal{K} \to *$、および可換図式
$$
\xymatrix{
\mathcal{K} \times \mathcal{K} \ar@<1ex>[r] \ar@<-1ex>[r] &
\mathcal{K} \ar[d] \ar[r] & {*} \ar[d]^s \\
& \mathcal{F}_i \ar[r] & \mathcal{F} \ar@{=}[r] & \colim \mathcal{F}_i
}
$$
が得られる。この図式の上段が余等化子であることに注意しよう。
したがって $i$ を大きくした後、誘導される二つの写像
$a_i, b_i : \mathcal{K} \times \mathcal{K} \to \mathcal{F}_i$ が
等しいことを示せば十分である。
% P01-E0073
これは次の段落で (2) の証明とまったく同じ議論によって示されるので、
読者には残りの証明を飛ばすことを勧める。

\medskip\noindent
$i' \geq i$ に対して、$a_{i'}, b_{i'} :
\mathcal{K} \times \mathcal{K} \to \mathcal{F}_{i'}$ を、$a_i, b_i$ と
系の遷移写像との合成とする。
$a = \colim_{i' \geq i} a_{i'} :
\mathcal{K} \times \mathcal{K} \to \mathcal{F}$ とおき、$b$ についても同様におく。
上の図式の可換性により $a = b$ である。$i' \geq i$ に対して
$$
E_{i'} = \text{Equalizer}(a_{i'}, b_{i'}) \subset
\mathcal{K} \times \mathcal{K}
\quad\text{and}\quad
E = \text{Equalizer}(a, b) \subset
\mathcal{K} \times \mathcal{K}
$$
と書く。Categories, Lemma \ref{categories-lemma-directed-commutes} により
$E = \colim_{i' \geq i} E_{i'}$ である。したがって
$\coprod_{i' \geq i} E_{i'} \to E$ は層の全射である。
$a = b$ なので $E = \mathcal{K} \times \mathcal{K}$ である。
(4)(b) により $\mathcal{K} \times \mathcal{K}$ は準コンパクトだから、
ある $i' \geq i$ に対して
$E_{i'} = \mathcal{K} \times \mathcal{K}$ となる。
したがって、ある $i' \geq i$ に対して $a_{i'} = b_{i'}$ である。
\end{proof}

\begin{remark}
\label{sites-remark-stronger-conditions}
$\mathcal{C}$ をサイトとする。Lemma
\ref{sites-lemma-directed-colimits-global-sections} の (4) の仮定を保証する方法は
いくつかある。以下に幾つか挙げる。
\begin{enumerate}
\item 次の性質をもつ集合 $\mathcal{B} \subset \Ob(\mathcal{C})$ が
存在すると仮定する。
\begin{enumerate}
\item 任意の全射 $\mathcal{F} \to *$ に対して、$m \geq 0$ と
$U_1, \ldots, U_m \in \mathcal{B}$ が存在し、$\mathcal{F}(U_j)$ は空でなく、
$\coprod h_{U_j}^\# \to *$ は全射である。
\item $U, U' \in \mathcal{B}$ に対して、層
$h_U^\# \times h_{U'}^\#$ は準コンパクトである。
\end{enumerate}
\item 次の性質をもつ集合 $\mathcal{B} \subset \Ob(\mathcal{C})$ が
存在すると仮定する。
\begin{enumerate}
\item $m \geq 0$ と $U_1, \ldots, U_m \in \mathcal{B}$ が存在し、
$\coprod h_{U_j}^\# \to *$ は全射である。
\item $U \in \mathcal{B}$ に対して、$U$ の任意の被覆は、
$U_j \in \mathcal{B}$ である有限被覆
$\{U_j \to U\}_{j = 1, \ldots, m}$ によって細分される。
\item $U, U' \in \mathcal{B}$ に対して、$m \geq 0$、
$U_1, \ldots, U_m \in \mathcal{B}$、射 $U_j \to U$ と $U_j \to U'$ が存在し、
$\coprod h_{U_j}^\# \to h_U^\# \times h_{U'}^\#$ は全射である。
\end{enumerate}
\item 次を仮定する。
\begin{enumerate}
\item $\Sh(\mathcal{C})$ は準コンパクトである。
\item $\mathcal{C}$ のすべての対象は、準コンパクト対象を成員とする被覆をもつ。
\item $U$ と $U'$ が準コンパクトなら、層
$h_U^\# \times h_{U'}^\#$ は準コンパクトである。
\end{enumerate}
\end{enumerate}
(1) と (2) の場合、$S \subset \Ob(\Sh(\mathcal{C}))$ を、
$U \in \mathcal{B}$ に対する層 $h_U^\#$ の有限余積からなる集合とする。
(3) の場合、$S \subset \Ob(\Sh(\mathcal{C}))$ を、準コンパクトな
$U \in \Ob(\mathcal{C})$ に対する層 $h_U^\#$ の有限余積からなる集合とする。
\end{remark}

\noindent
後で、余極限について起こりうることに次のような評価が必要となる。

\begin{lemma}
\label{sites-lemma-colimit-over-ordinal-sections}
$\mathcal{C}$ をサイトとし、$\beta$ を順序数とする。
$\beta \to \Sh(\mathcal{C})$、$\alpha \mapsto \mathcal{F}_\alpha$ を
$\beta$ 上の層の系とする。$U \in \Ob(\mathcal{C})$ に対して標準写像
$$
\colim_{\alpha < \beta} \mathcal{F}_\alpha(U)
\longrightarrow
\left(\colim_{\alpha < \beta} \mathcal{F}_\alpha \right)(U)
$$
を考える。$\beta$ の共終数が十分大きければ、この写像はすべての $U$ に対して
全単射である。
\end{lemma}

\begin{proof}
左辺は、前層の圏でとった余極限 $\mathcal{F}_{\colim}$ の $U$ における値である。
Section \ref{sites-section-limits-colimits-PSh} を参照せよ。
$\colim_{\alpha < \beta} \mathcal{F}_\alpha$ は
$\mathcal{F}_{\colim}$ の層化 $\mathcal{F}_{\colim}^\#$ であることを思い出そう。
Lemma \ref{sites-lemma-colimit-sheaves} を参照せよ。
$\mathcal{U} = \{U_i \to U\}_{i \in I}$ を、$\mathcal{C}$ の被覆の集合
$\text{Cov}(\mathcal{C})$ の元とする。$\beta$ の共終数が $I$ の濃度より
大きいなら、
$$
H^0(\mathcal{U}, \mathcal{F}_{\colim}) =
\colim H^0(\mathcal{U}, \mathcal{F}_\alpha) =
\colim \mathcal{F}_\alpha(U) = \mathcal{F}_{\colim}(U)
$$
と主張する。第二、第三の等号は明らかである。第一の等号について、
$s = (s_i) \in H^0(\mathcal{U}, \mathcal{F}_{\colim})$ とする。
すると各 $i$ に対して、元 $s_i$ は、ある $\alpha_i < \beta$ に対する元
$s_{i, \alpha_i} \in \mathcal{F}_{\alpha_i}(U_i)$ から来る。
共終数に関する仮定により、
% P01-E0074
$i$ に依存しない $\alpha_i = \alpha$ を選べる。
すると、ある $\alpha_{i, j} < \beta$ に対して、$s_i$ と $s_j$ は
$\mathcal{F}_{\alpha_{i, j}}(U_i \times_U U_j)$ の同じ元へ写る。
% P01-E0075
$I \times I$ の濃度も $\beta$ の共終数より小さいので、$\alpha$ を大きくした後、
すべての $i, j$ について $\alpha_{i, j} = \alpha$ と仮定できる。
これにより自然写像
$\colim H^0(\mathcal{U}, \mathcal{F}_\alpha)
\to H^0(\mathcal{U}, \mathcal{F}_{\colim})$ が全射であることが分かる。
非常によく似た議論により、それは単射でもある。
特に $\mathcal{F}_{\colim}$ は $\mathcal{U}$ に対する層条件を満たす。
したがって $\beta$ の共終数が、被覆の添字集合 $I$ の濃度全体の上限より
大きければ、結論を得る。
\end{proof}






\section{サイトの余極限}
\label{sites-section-colimit-sites}

\noindent
サイトがサイトの逆系の極限である場合に、Lemma
\ref{sites-lemma-directed-colimits-sections} の類似が必要となる。
簡単のため、被覆の添字集合が有限である場合にのみ構成を説明する。

\begin{situation}
\label{sites-situation-inverse-limit-sites}
ここでは次が与えられている。
\begin{enumerate}
\item 余フィルター添字圏 $\mathcal{I}$。
\item 各 $i \in \Ob(\mathcal{I})$ に対するサイト $\mathcal{C}_i$。
$\mathcal{C}_i$ のすべての被覆の添字集合は有限である。
\item $\mathcal{I}$ の射 $a : i \to j$ に対し、連続関手
$u_a : \mathcal{C}_j \to \mathcal{C}_i$ によって与えられるサイトの射
$f_a : \mathcal{C}_i \to \mathcal{C}_j$。
\end{enumerate}
さらに、$\mathcal{I}$ において $c = a \circ b$ なら
$f_a \circ f_b = f_c$ とする。
\end{situation}

\begin{lemma}
\label{sites-lemma-colimit-sites}
Situation \ref{sites-situation-inverse-limit-sites} では、サイト
$(\mathcal{C}, \text{Cov}(\mathcal{C}))$ を次のように構成できる。
\begin{enumerate}
\item 圏として $\mathcal{C} = \colim \mathcal{C}_i$ とする。
\item $\text{Cov}(\mathcal{C})$ を、$u_i : \mathcal{C}_i \to \mathcal{C}$ による
$\text{Cov}(\mathcal{C}_i)$ の像の和集合とする。
\end{enumerate}
\end{lemma}

\begin{proof}
サイトの射の合成の定義により、$\mathcal{I}$ において
$c = a \circ b$ なら $u_b \circ u_a = u_c$ である。
式 $\mathcal{C} = \colim \mathcal{C}_i$ は、
$\Ob(\mathcal{C}) = \colim \Ob(\mathcal{C}_i)$ および
$\text{Arrows}(\mathcal{C}) = \colim \text{Arrows}(\mathcal{C}_i)$ を意味する。
始域、終域、および合成は $\text{Arrows}(\mathcal{C}_i)$ 上の始域、終域、
および合成から継承される。このようにして圏を得る。
明らかな関手を $u_i : \mathcal{C}_i \to \mathcal{C}$ と書く。
$\mathcal{C}$ の任意の有限図式は、ある $i$ に対して
$\mathcal{C}_i$ の図式の像となることに注意しよう。

\medskip\noindent
$\{U^t \to U\}$ を $\mathcal{C}$ の被覆とする。まず、
$V \to U$ が $\mathcal{C}$ の射なら $U^t \times_U V$ が存在することを示す。
上の注意と被覆の定義により、ある $i$、$\mathcal{C}_i$ の被覆
$\{U_i^t \to U_i\}$、および射 $V_i \to U_i$ を選んで、$u_i$ による像を
与えられたデータとできる。$U^t \times_U V$ は
$u_i$ による $U^t_i \times_{U_i} V_i$ の像であると主張する。
実際、$\mathcal{I}$ の任意の $a : j \to i$ に対して関手 $u_a$ は連続なので、
$u_a(U^t_i \times_{U_i} V_i) = u_a(U^t_i) \times_{u_a(U_i)} u_a(V_i)$ である。
特に、望むなら $i$ を $j$ で置き換えられる。
したがって $W$ が $\mathcal{C}$ の別の対象なら、$W = u_i(W_i)$ と仮定でき、
\begin{align*}
& \Mor_\mathcal{C}(W, u_i(U^t_i \times_{U_i} V_i)) \\
& =
\colim_{a : j \to i}
\Mor_{\mathcal{C}_j}(u_a(W_i), u_a(U^t_i \times_{U_i} V_i)) \\
& =
\colim_{a : j \to i}
\Mor_{\mathcal{C}_j}(u_a(W_i), u_a(U^t_i))
\times_{\Mor_{\mathcal{C}_j}(u_a(W_i), u_a(U_i))}
\Mor_{\mathcal{C}_j}(u_a(W_i), u_a(V_i)) \\
& =
\Mor_\mathcal{C}(W, U^t)
\times_{\Mor_\mathcal{C}(W, U)}
\Mor_\mathcal{C}(W, V)
\end{align*}
となる。これはフィルター余極限が有限極限と可換することによる
（Categories, Lemma \ref{categories-lemma-directed-commutes}）。
また $\{U^t \times_U V \to V\}$ が $\mathcal{C}$ の被覆であることも従う。
これで Definition \ref{sites-definition-site} の公理 (3) が成り立つことが分かる。

\medskip\noindent
Definition \ref{sites-definition-site} の公理 (2) を確認するため、
$\{U^t \to U\}_{t \in T}$ を $\mathcal{C}$ の被覆とし、各 $t$ に対して
$\{U^{ts} \to U^t\}$ を $\mathcal{C}$ の被覆とする。
ある $i$ と $\mathcal{C}_i$ の被覆 $\{U^t_i \to U_i\}_{t \in T}$ を選んで、
$u_i$ による像を $\{U^t \to U\}$ とできる。$T$ は{\bf 有限}なので、
$\mathcal{I}$ の射 $a : j \to i$ と $\mathcal{C}_j$ の被覆
$\{U^{ts}_j \to u_a(U^t_i)\}$ を選んで、$u_j$ による像を
$\{U^{ts} \to U^t\}$ とできる。サイト $\mathcal{C}_j$ に公理 (2) を
適用すると、$\{U^{ts} \to U\}$ は $\mathcal{C}$ の被覆である。

\medskip\noindent
Definition \ref{sites-definition-site} の公理 (1) の証明は省略する。
\end{proof}

\begin{lemma}
\label{sites-lemma-compute-pullback-to-limit}
Situation \ref{sites-situation-inverse-limit-sites} において、
$u_i : \mathcal{C}_i \to \mathcal{C}$ を Lemma
\ref{sites-lemma-colimit-sites} で構成したものとする。このとき $u_i$ は
サイトの射 $f_i : \mathcal{C} \to \mathcal{C}_i$ を定める。
$U_i \in \Ob(\mathcal{C}_i)$ と $\mathcal{C}_i$ 上の層 $\mathcal{F}$ に対して
\begin{equation}
\label{sites-equation-compute-pullback-to-limit}
f_i^{-1}\mathcal{F}(u_i(U_i)) =
\colim_{a : j \to i} f_a^{-1}\mathcal{F}(u_a(U_i))
\end{equation}
である。
\end{lemma}

\begin{proof}
Lemma \ref{sites-lemma-colimit-sites} の証明の議論から、関手 $u_i$ が連続であることは
直ちに分かる。証明を終えるには $f_i^{-1} := u_{i, s}$ が完全関手
$\Sh(\mathcal{C}_i) \to \Sh(\mathcal{C})$ であることを示さなければならない。
実際、$f_i^{-1}$ は左随伴として右完全なので
（Categories, Lemma \ref{categories-lemma-exact-adjoint}）、
左完全であることを示せば十分である。まず
(\ref{sites-equation-compute-pullback-to-limit}) を証明し、その後で完全性を導く。

\medskip\noindent
$\mathcal{C}$ の任意の対象 $V$ に対して、
% P01-E0076
射 $a : j \to i$ と、$V = u_j(V_j)$ を満たす対象
% P01-E0077
$V_j \in \Ob(\mathcal{C})$ を選べる。このとき
$$
\mathcal{G}(V) = \colim_{b : k \to j} f_{a \circ b}^{-1}\mathcal{F}(u_b(V_j))
$$
とおける。余極限の値 $\mathcal{G}(V)$ は、
% P01-E0078
$b : j \to i$ と $u_j(V_j) = V$ を満たす対象 $V_j$ の選択に依存しない。
確認は省略する。さらに $\alpha : V \to V'$ が $\mathcal{C}$ の射なら、
$b : j \to i$ と $u_j(\alpha_j) = \alpha$ を満たす射
$\alpha_j : V_j \to V'_j$ を選べる。各射
$u_b(\alpha_j) : u_b(V_j) \to u_b(V'_j)$ に沿う制限を用いると、
これは写像 $\mathcal{G}(V') \to \mathcal{G}(V)$ を誘導する。
確認により $\mathcal{G}$ は前層である（省略する）。
実際、$\mathcal{G}$ は層条件を満たす。すなわち、$\mathcal{C}$ の任意の被覆
$\mathcal{U} = \{U^t \to U\}$ は有限段階から来る。
ある $\mathcal{I}$ の $a : j \to i$ に対して
$\mathcal{U}_j = \{U^t_j \to U_j\}$ が $u_j$ により $\mathcal{U}$ へ
写るとする。このとき
% P01-E0079
$$
H^0(\mathcal{U}, \mathcal{G}) =
\colim_{b : k \to j} H^0(u_b(\mathcal{U}_j), f_{b \circ a}^{-1}\mathcal{F}) =
\colim_{b : k \to j} f_{b \circ a}^{-1}\mathcal{F}(u_b(U_j)) =
\mathcal{G}(U)
$$
となり、所望の結論を得る。第一の等号はフィルター余極限が有限極限と可換する
ことから従う（Categories, Lemma \ref{categories-lemma-directed-commutes}）。
構成により $\mathcal{G}(U)$ は (\ref{sites-equation-compute-pullback-to-limit}) の
右辺によって与えられる。したがって $\mathcal{G}$ が
$f_i^{-1}\mathcal{F}$ に等しいことを示せば
(\ref{sites-equation-compute-pullback-to-limit}) が従う。

\medskip\noindent
この段落では $\mathcal{G}$ が $f_i^{-1}\mathcal{F}$ と標準的に同型であることを
確認する。この段落は飛ばすことを強く勧める。
確認のため、$\mathcal{C}$ 上の層 $\mathcal{H}$ に関して関手的な全単射
$\Mor_{\Sh(\mathcal{C})}(\mathcal{G}, \mathcal{H}) =
\Mor_{\Sh(\mathcal{C}_i)}(\mathcal{F}, f_{i, *}\mathcal{H})$ が存在することを
示さなければならない。ここで $f_{i, *} = u_i^p$ である。
写像 $\mathcal{G} \to \mathcal{H}$ は、すべての $a : j \to i$、
$b : k \to j$、および $V_j \in \Ob(\mathcal{C}_j)$ に対する両立する写像の系
$$
\varphi_{a, b, V_j} :
f_{a \circ b}^{-1}\mathcal{F}(u_b(V_j))
\longrightarrow
\mathcal{H}(u_j(V_j))
$$
と同じものである。両立性により、写像 $\varphi_{a, b, V_j}$ は
$\varphi_{a \circ b, \text{id}, u_b(V_j)}$ に等しくなければならない。
$a : j \to i$ が与えられると、写像族 $\varphi_{a, \text{id}, V_j}$ は
層の写像 $\varphi_a : f_a^{-1}\mathcal{F} \to f_{j, *}\mathcal{H}$ に対応する。
% P01-E0080
$\varphi_{a, \text{id}, u_a(V_i)}$ と
$\varphi_{\text{id}, \text{id}, V_i}$ の間の両立性により、$\varphi_a$ は
$$
\Mor_{\Sh(\mathcal{C}_j)}(f_a^{-1}\mathcal{F}, f_{j, *}\mathcal{H}) =
\Mor_{\Sh(\mathcal{C}_i)}(\mathcal{F}, f_{a, *}f_{j, *}\mathcal{H}) =
\Mor_{\Sh(\mathcal{C}_i)}(\mathcal{F}, f_{i, *}\mathcal{H})
$$
による写像 $\varphi_{id}$ の随伴となる。したがって最終的に、写像の系
$\varphi_{a, b, V_j}$ 全体は写像
$\varphi_{id} : \mathcal{F} \to f_{i, *}\mathcal{H}$ によって決定される。
逆に、そのような写像 $\psi : \mathcal{F} \to f_{i, *}\mathcal{H}$ が
与えられれば、今の議論を逆にたどって写像族
$\varphi_{a, b, V_j}$ を構成できる。これで
$\mathcal{G} = f_i^{-1}\mathcal{F}$ の証明が終わる。

\medskip\noindent
(\ref{sites-equation-compute-pullback-to-limit}) が成り立つと仮定する。
層の有限極限は前層の圏で計算され（Lemma \ref{sites-lemma-limit-sheaf}）、
関手 $f_a^{-1}$ は有限極限と可換し、フィルター余極限も有限極限と可換するので、
関手 $\mathcal{F} \mapsto f_i^{-1}\mathcal{F}(U)$ は有限極限と可換する。
$\mathcal{C}$ の一般の対象 $V$ に対して
$\mathcal{F} \mapsto f_i^{-1}\mathcal{F}(V)$ が有限極限と可換することを見るには、
上で与えた $f_i^{-1}\mathcal{F}(V) = \mathcal{G}(V)$ の公式を用いて
同じように論じればよい。したがって $f_i^{-1}$ は左完全であり、
補題の証明が完了する。
\end{proof}


\begin{lemma}
\label{sites-lemma-colimit}
Situation \ref{sites-situation-inverse-limit-sites} において、次が与えられていると仮定する。
\begin{enumerate}
\item すべての $i \in \Ob(\mathcal{I})$ に対して
$\mathcal{C}_i$ 上の層 $\mathcal{F}_i$。
\item $a : j \to i$ に対して、$\mathcal{C}_j$ 上の層の写像
$\varphi_a : f_a^{-1}\mathcal{F}_i \to \mathcal{F}_j$。
\end{enumerate}
$c = a \circ b$ のとき
$\varphi_c = \varphi_b \circ f_b^{-1}\varphi_a$ とする。
Lemma \ref{sites-lemma-colimit-sites} のサイト $\mathcal{C}$ 上で
$\mathcal{F} = \colim f_i^{-1}\mathcal{F}_i$ とおく。
$i \in \Ob(\mathcal{I})$ と $X_i \in \text{Ob}(\mathcal{C}_i)$ に対して
$$
\colim_{a : j \to i} \mathcal{F}_j(u_a(X_i)) = \mathcal{F}(u_i(X_i))
$$
である。
\end{lemma}

\begin{proof}
形式的な議論により
% P01-E0081
$$
\colim_{a : j \to i} \mathcal{F}_i(u_a(X_i)) =
\colim_{a : j \to i} \colim_{b : k \to j}
f_b^{-1}\mathcal{F}_j(u_{a \circ b}(X_i))
$$
である。(\ref{sites-equation-compute-pullback-to-limit}) により、内側の余極限は
$f_j^{-1}\mathcal{F}_j(u_i(X_i))$ に等しい。したがって Lemma
\ref{sites-lemma-directed-colimits-sections} により結論を得る。
\end{proof}

\begin{lemma}
\label{sites-lemma-colimit-push-pull}
Situation \ref{sites-situation-inverse-limit-sites} において、
$\mathcal{C}$ 上の層 $\mathcal{F}$ が与えられていると仮定する。このとき
$$
\mathcal{F} = \colim f_i^{-1}f_{i, *}\mathcal{F}
$$
である。ここで遷移写像は $a : j \to i$ に対する
$f_j^{-1}\varphi_a$ であり、$\varphi_a :
f_a^{-1}f_{i, *}\mathcal{F} \to f_{j, *}\mathcal{F}$ は標準写像で、
Lemma \ref{sites-lemma-colimit} と同様のコサイクル条件を満たす。
\end{lemma}

\begin{proof}
射
$$
\varphi_a : f_a^{-1}f_{i, *}\mathcal{F} \to f_{j, *}\mathcal{F}
$$
として、恒等写像
$$
f_{i, *}\mathcal{F} \to f_{a, *}f_{j, *}\mathcal{F}
$$
の随伴を選ぶ。したがって $\varphi_a$ は随伴
$(f_a^{-1}, f_{a, *})$ の余単位である。すべての $a : j \to i$ と
% P01-E0082
$b : k \to i$ に対し、合成 $c = a \circ b$ があるとき
$\varphi_c = \varphi_b \circ f_b^{-1}\varphi_a$ であることを
証明しなければならない。これは Categories, Lemma
\ref{categories-lemma-compose-counits} から従う。
最後に、随伴によって与えられる写像
$$
\colim_{i \in I} f_i^{-1}f_{i, *}\mathcal{F}
\longrightarrow
\mathcal{F}
$$
が同型であることを示さなければならない。
$\mathcal{C}$ の対象 $U$ に対して、写像
% P01-E0083
$$
(\colim_{i \in I} f_i^{-1}\mathcal{F}_i)(U) \to \mathcal{F}(U)
$$
が全単射であることを示す必要がある。ある $i$ と $\mathcal{C}_i$ の対象
$U_i$ を $u_i(U_i) = U$ となるように選ぶ。このとき Lemma
\ref{sites-lemma-colimit} により左辺は
$$
(\colim_{i \in I} f_i^{-1}\mathcal{F}_i)(U) =
\colim_{a : j \to i} f_{j, *}\mathcal{F}(u_a(U_i))
$$
に等しい。$u_j(u_a(U_i)) = U$ なので、定義により、すべての
$a : j \to i$ に対して
$f_{j, *}\mathcal{F}(u_a(U_i)) = \mathcal{F}(U)$ である。
したがって余極限の値は $\mathcal{F}(U)$ であり、証明が完了する。
\end{proof}












\section{前層のさらなる関手性}
\label{sites-section-more-functoriality-PSh}

\noindent
この節では Section \ref{sites-section-functoriality-PSh} の内容を再び扱う。
$u : \mathcal{C} \to \mathcal{D}$ を圏の間の関手とする。
$$
u^p :
\textit{PSh}(\mathcal{D})
\longrightarrow
\textit{PSh}(\mathcal{C})
$$
が、$\mathcal{D}$ 上の $\mathcal{G}$ に前層
$u^p\mathcal{G} = \mathcal{G} \circ u$ を対応させる関手であったことを
思い出そう。この関手は左随伴（すなわち $u_p$）をもつだけでなく、
右随伴ももつことが分かる。

\medskip\noindent
実際、任意の $V \in \Ob(\mathcal{D})$ に対して圏
${}_V\mathcal{I} = {}_V^u\mathcal{I}$ を定義する。
その対象は対 $(U, \psi : u(U) \to V)$ である。
ここで射の向きは、Section \ref{sites-section-functoriality-PSh} で圏
$\mathcal{I}_V^u$ を定義する際に用いた射の向きと反対であることに注意しよう。
射 $(U, \psi) \to (U', \psi')$ は、
$\psi = \psi' \circ u(\alpha)$ を満たす射 $\alpha : U \to U'$ によって
与えられる。さらに、$\mathcal{C}$ 上の任意の集合の前層 $\mathcal{F}$ に対して、
規則 ${}_V\mathcal{F}(U, \psi) = \mathcal{F}(U)$ で定義される関手
${}_V\mathcal{F} : {}_V\mathcal{I}^{opp} \to \textit{Sets}$ を導入する。
次のように定義する。
$$
{}_pu(\mathcal{F})(V) := \lim_{{}_V\mathcal{I}^{opp}} {}_V\mathcal{F}
$$
極限なので、${}_V\mathcal{I}$ の各対象 $(U, \psi)$ に対して射影
$c(\psi) : {}_pu(\mathcal{F})(V) \to \mathcal{F}(U)$ がある。
実際、
$$
{}_pu(\mathcal{F})(V)
=
\left\{
\begin{matrix}
\text{collections }
s_{(U, \psi)} \in \mathcal{F}(U) \\
\forall \beta : (U_1, \psi_1) \to (U_2, \psi_2)
\text{ in }{}_V\mathcal{I} \\
\text{ we have } \beta^*s_{(U_2, \psi_2)} = s_{(U_1, \psi_1)}
\end{matrix}
\right\}
$$
であり、対応は $s \mapsto s_{(U, \psi)} = c(\psi)(s)$ によって与えられる。
$\mathcal{D}$ の任意の射 $V' \to V$ に付随する制限写像
${}_pu(\mathcal{F})(V) \to {}_pu(\mathcal{F})(V')$ の定義は読者に任せる。
得られる前層を ${}_pu\mathcal{F}$ と書く。

\begin{lemma}
\label{sites-lemma-recover-pu}
標準写像 ${}_pu\mathcal{F}(u(U)) \to \mathcal{F}(U)$ があり、
これは制限写像と両立する。
\end{lemma}

\begin{proof}
これは上の射影 $c(\text{id}_{u(U)})$ にほかならない。
\end{proof}

\noindent
前層の任意の写像 $\mathcal{F} \to \mathcal{F}'$ は、関手間の両立する写像の系
${}_V\mathcal{F} \to {}_V\mathcal{F}'$ を生じ、したがって前層の写像
${}_pu\mathcal{F} \to {}_pu\mathcal{F}'$ を生じることに注意しよう。
言い換えると、関手
$$
{}_pu :
\textit{PSh}(\mathcal{C})
\longrightarrow
\textit{PSh}(\mathcal{D})
$$
を定義した。

\begin{lemma}
\label{sites-lemma-adjoints-pu}
関手 ${}_pu$ は関手 $u^p$ の右随伴である。言い換えると、公式
$$
\Mor_{\textit{PSh}(\mathcal{C})}(u^p\mathcal{G}, \mathcal{F})
=
\Mor_{\textit{PSh}(\mathcal{D})}(\mathcal{G}, {}_pu\mathcal{F})
$$
は $\mathcal{F}$ と $\mathcal{G}$ に関して二変数関手的に成り立つ。
\end{lemma}

\begin{proof}
これは Lemma \ref{sites-lemma-adjoints-u} の証明とまったく同じ方法で証明される。
Lemma \ref{sites-lemma-recover-pu} の写像
$u^p{}_pu \mathcal{F} \to \mathcal{F}$ が、右辺から左辺へ移る際に
用いられる写像であることに注意しよう。

\medskip\noindent
別の方法として、$\mathcal{C}$ 上の集合の前層 $\mathcal{F}$ を
$\textit{Sets}^{opp}$ に値をとる $\mathcal{C}^{opp}$ 上の前層
$\mathcal{F}'$ とみなし、$\mathcal{D}$ 上でも同様にみなす。
$({}_pu \mathcal{F})' = u_p(\mathcal{F}')$ および
$(u^p\mathcal{G})' = u^p(\mathcal{G}')$ を確認せよ。
Remark \ref{sites-remark-functoriality-presheaves-values} により、
$\textit{Sets}^{opp}$ に値をとる前層について $u_p$ と $u^p$ の随伴性がある。
すると結果は形式的に従う。
\end{proof}

\noindent
したがって圏の関手 $u : \mathcal{C} \to \mathcal{D}$ が与えられると、
前層の圏の間の関手の列
$$
u_p, u^p, {}_pu
$$
が得られ、連続する各対では前者が後者の左随伴である。

\begin{lemma}
\label{sites-lemma-adjoint-functors}
$u : \mathcal{C} \to \mathcal{D}$ と $v : \mathcal{D} \to \mathcal{C}$ を
圏の関手とし、$v$ が $u$ の右随伴であると仮定する。このとき次が成り立つ。
\begin{enumerate}
\item $\mathcal{D}$ の任意の $V$ に対して $u^ph_V = h_{v(V)}$ である。
\item 圏 $\mathcal{I}^v_U$ は始対象をもつ。
\item 圏 ${}_V^u\mathcal{I}$ は終対象をもつ。
\item ${}_pu = v^p$ である。
\item $u^p = v_p$ である。
\end{enumerate}
\end{lemma}

\begin{proof}
(1) の証明。$V$ を $\mathcal{D}$ の対象とする。仮定により
$u^ph_V(U) = \Mor_\mathcal{D}(u(U), V) = \Mor_\mathcal{C}(U, v(V))$ なので、
$u^ph_V = h_{v(V)}$ である。

\medskip\noindent
(2) の証明。$U$ を $\mathcal{C}$ の対象とする。
$\eta : U \to v(u(U))$ を写像 $\text{id} : u(U) \to u(U)$ の随伴とする。
このとき $(u(U), \eta)$ は $\mathcal{I}_U^v$ の始対象であると主張する。
実際、$\mathcal{I}_U^v$ の対象 $(V, \phi : U \to v(V))$ が与えられると、
射 $\phi$ は写像 $\psi : u(U) \to V$ の随伴であり、この写像は射
$(u(U), \eta) \to (V, \phi)$ を定める。

\medskip\noindent
(3) の証明。$V$ を $\mathcal{D}$ の対象とする。
$\xi : u(v(V)) \to V$ を写像 $\text{id} : v(V) \to v(V)$ の随伴とする。
このとき $(v(V), \xi)$ は ${}_V^u\mathcal{I}$ の終対象であると主張する。
実際、${}_V^u\mathcal{I}$ の対象 $(U, \psi : u(U) \to V)$ が与えられると、
射 $\psi$ は写像 $\phi : U \to v(V)$ の随伴であり、この写像は射
$(U, \psi) \to (v(V), \xi)$ を定める。

\medskip\noindent
したがって $\mathcal{C}$ 上の任意の前層 $\mathcal{F}$ に対して
\begin{eqnarray*}
v^p\mathcal{F}(V)
& = &
\mathcal{F}(v(V)) \\
& = &
\Mor_{\textit{PSh}(\mathcal{C})}(h_{v(V)}, \mathcal{F}) \\
& = &
\Mor_{\textit{PSh}(\mathcal{C})}(u^ph_V, \mathcal{F}) \\
& = &
\Mor_{\textit{PSh}(\mathcal{D})}(h_V, {}_pu\mathcal{F}) \\
& = &
{}_pu\mathcal{F}(V)
\end{eqnarray*}
であり、(4) が従う。(5) は随伴関手の一意性から従う。
\end{proof}





\section{余連続関手}
\label{sites-section-cocontinuous-functors}

\noindent
トポスの射を構成する別の方法がある。これには、次のように定義される
サイト間の余連続関手を用いる。

\begin{definition}
\label{sites-definition-cocontinuous}
$\mathcal{C}$ と $\mathcal{D}$ をサイトとし、
$u : \mathcal{C} \to \mathcal{D}$ を関手とする。
任意の $U \in \Ob(\mathcal{C})$ と $\mathcal{D}$ の任意の被覆
$\{V_j \to u(U)\}_{j \in J}$ に対して、$\mathcal{C}$ の被覆
$\{U_i \to U\}_{i\in I}$ が存在し、写像族
$\{u(U_i) \to u(U)\}_{i \in I}$ が被覆
$\{V_j \to u(U)\}_{j \in J}$ を細分するとき、関手 $u$ は
{\it 余連続}であるという。
\end{definition}

\noindent
$\{u(U_i) \to u(U)\}_{i \in I}$ は一般にはサイト $\mathcal{D}$ の被覆では
{\it ない}ことに注意しよう。

\begin{lemma}
\label{sites-lemma-pu-sheaf}
$\mathcal{C}$ と $\mathcal{D}$ をサイトとし、
$u : \mathcal{C} \to \mathcal{D}$ を余連続とする。
$\mathcal{F}$ を $\mathcal{C}$ 上の層とする。このとき
${}_pu\mathcal{F}$ は $\mathcal{D}$ 上の層であり、これを
${}_su\mathcal{F}$ と書く。
\end{lemma}

\begin{proof}
$\{V_j \to V\}_{j \in J}$ をサイト $\mathcal{D}$ の被覆とする。
図式
$$
\xymatrix{
&\ \phantom{}_pu\mathcal{F}(V) \ar[r] &
\prod {}_pu\mathcal{F}(V_j) \ar@<1ex>[r] \ar@<-1ex>[r] &
\prod {}_pu\mathcal{F}(V_j \times_V V_{j'}) &
}
$$
が等化子図式であることを示さなければならない。
${}_pu$ は $u^p$ の右随伴なので
$$
{}_pu\mathcal{F}(V) =
\Mor_{\textit{PSh}(\mathcal{D})}(h_V, {}_pu\mathcal{F}) =
\Mor_{\textit{PSh}(\mathcal{C})}(u^ph_V, \mathcal{F}) =
\Mor_{\Sh(\mathcal{C})}((u^ph_V)^\#, \mathcal{F})
$$
である。したがって
\begin{equation}
\label{sites-equation-coequalizer}
\xymatrix{
\coprod u^p h_{V_j \times_V V_{j'}} \ar@<1ex>[r] \ar@<-1ex>[r] &
\coprod u^p h_{V_j} \ar[r] &
u^p h_V
}
\end{equation}
が層化後に余等化子図式となることを示せば十分である。
（層の圏における余積は前層の圏における余積の層化であることを思い出そう。
Lemma \ref{sites-lemma-colimit-sheaves} を参照せよ。）

\medskip\noindent
まず (\ref{sites-equation-coequalizer}) の第二の射が層化後に全射となることを示す。
このため Lemma \ref{sites-lemma-mono-epi-sheaves} を用いる。
したがって、$U$ 上の $u^ph_V$ の切断 $s$ が、$U$ のある被覆の成員上で
$\coprod u^p h_{V_j}$ の切断へ持ち上がることを示せば十分である。
$s$ は射 $s : u(U) \to V$ であることに注意しよう。
すると $\{V_j \times_{V, s} u(U) \to u(U)\}$ は $\mathcal{D}$ の被覆である。
$u$ は余連続なので、被覆 $\{U_i \to U\}$ が存在し、
$\{u(U_i) \to u(U)\}$ は
$\{V_j \times_{V, s} u(U) \to u(U)\}$ を細分する。
これは、各制限 $s|_{U_i} : u(U_i) \to V$ が、ある $j$ に対する射
$s_i : u(U_i) \to V_j$ を経由することを意味する。
すなわち、所望どおり $s|_{U_i}$ は
$u^ph_{V_j}(U_i) \to u^ph_V(U_i)$ の像に入る。

\medskip\noindent
$s, s' \in (\coprod u^ph_{V_j})^\#(U)$ が
$(u^ph_V)^\#(U)$ の同じ元へ写るとする。補題の証明を終えるため、
$U$ をある被覆の成員で置き換えた後、$s, s'$ が
(\ref{sites-equation-coequalizer}) の二つの写像による
$\coprod u^p h_{V_j \times_V V_{j'}}$ の同じ切断の像であることを示す。
まず $U$ を被覆の成員で置き換え、
$s \in u^ph_{V_j}(U)$ および $s' \in u^ph_{V_{j'}}(U)$ と仮定してよい。
もう一度そのように置き換えると、$s$ と $s'$ は層化においてではなく
$u^ph_V(U)$ において同じ像をもつと仮定できる。したがって
$s : u(U) \to V_j$ と $s' : u(U) \to V_{j'}$ は $\mathcal{D}$ の射であり、図式
$$
\xymatrix{
u(U) \ar[r]_{s'} \ar[d]_s & V_{j'} \ar[d] \\
V_j \ar[r] & V
}
$$
は可換である。ゆえに $t = (s, s') : u(U) \to V_j \times_V V_{j'}$、
すなわち切断 $t \in u^ph_{V_j \times_V V_{j'}}(U)$ が得られ、
これは所望どおり $s, s'$ へ写る。
\end{proof}

\begin{lemma}
\label{sites-lemma-exact-cocontinuous}
$\mathcal{C}$ と $\mathcal{D}$ をサイトとし、
$u : \mathcal{C} \to \mathcal{D}$ を余連続とする。関手
$\Sh(\mathcal{D}) \to \Sh(\mathcal{C})$、
$\mathcal{G} \mapsto (u^p\mathcal{G})^\#$ は、上の Lemma
\ref{sites-lemma-pu-sheaf} で導入した関手 ${}_su$ の左随伴である。
さらに、この関手は完全である。
\end{lemma}

\begin{proof}
随伴性を次のように証明する。
\begin{eqnarray*}
\Mor_{\Sh(\mathcal{C})}
((u^p\mathcal{G})^\#, \mathcal{F})
& = &
\Mor_{\textit{PSh}(\mathcal{C})}
(u^p\mathcal{G}, \mathcal{F}) \\
& = &
\Mor_{\textit{PSh}(\mathcal{D})}
(\mathcal{G}, {}_pu\mathcal{F}) \\
& = &
\Mor_{\Sh(\mathcal{D})}
(\mathcal{G}, {}_su\mathcal{F}).
\end{eqnarray*}
したがってこれは左随伴であり、ゆえに右完全である。Categories, Lemma
\ref{categories-lemma-exact-adjoint} を参照せよ。
層化が左完全であることはすでに見た。Lemma
\ref{sites-lemma-sheafification-exact} を参照せよ。
さらに包含 $i : \Sh(\mathcal{D}) \to \textit{PSh}(\mathcal{D})$ は Lemma
\ref{sites-lemma-limit-sheaf} により左完全である。最後に、関手 $u^p$ は右随伴
（すなわち $u_p$ の右随伴）なので左完全である。
したがってこの関手は左完全関手の合成
${}^\# \circ u^p \circ i$ であり、ゆえに左完全である。
\end{proof}

\noindent
この節を技術的な補題で終える。

\begin{lemma}
\label{sites-lemma-technical-pu}
% P01-E0084
Lemma \ref{sites-lemma-exact-cocontinuous} の状況とする。
$\mathcal{D}$ 上の任意の前層 $\mathcal{G}$ に対して
$(u^p\mathcal{G})^\# = (u^p(\mathcal{G}^\#))^\#$ である。
\end{lemma}

\begin{proof}
$\mathcal{C}$ 上の任意の層 $\mathcal{F}$ に対して
\begin{eqnarray*}
\Mor_{\Sh(\mathcal{C})}((u^p(\mathcal{G}^\#))^\#, \mathcal{F})
& = &
\Mor_{\Sh(\mathcal{D})}(\mathcal{G}^\#, {}_su\mathcal{F}) \\
& = &
\Mor_{\Sh(\mathcal{D})}(\mathcal{G}^\#, {}_pu\mathcal{F}) \\
& = &
\Mor_{\textit{PSh}(\mathcal{D})}(\mathcal{G}, {}_pu\mathcal{F}) \\
& = &
\Mor_{\textit{PSh}(\mathcal{C})}(u^p\mathcal{G}, \mathcal{F}) \\
& = &
\Mor_{\Sh(\mathcal{C})}((u^p\mathcal{G})^\#, \mathcal{F})
\end{eqnarray*}
であり、結果は米田の補題から従う。
\end{proof}

\begin{remark}
\label{sites-remark-cartesian-cocontinuous}
$u : \mathcal{C} \to \mathcal{D}$ を圏の間の関手とする。
$\mathcal{D}$ の射 $g : u(U) \to V$ と $f : W \to V$ が与えられたとき、関手
$$
\mathcal{C}^{opp} \longrightarrow \textit{Sets},\quad
T \longmapsto
\Mor_\mathcal{C}(T, U)
\times_{\Mor_\mathcal{D}(u(T), V)}
\Mor_\mathcal{D}(u(T), W)
$$
を考えられる。この関手が表現可能なら、対応する $\mathcal{C}$ の対象を
$U \times_{g, V, f} W$ と書く。$\mathcal{C}$ と $\mathcal{D}$ をサイトとする。
次の性質 $P$ を考える。$\mathcal{D}$ の任意の被覆
$\{f_j : V_j \to V\}$ と任意の射 $g : u(U) \to V$ に対して、
% P01-E0085
\begin{enumerate}
\item すべての $i$ に対して $U \times_{g, V, f_i} V_i$ が存在する。
\item $\{U \times_{g, V, f_i} V_i \to U\}$ は $\mathcal{C}$ の被覆である。
\end{enumerate}
連続関手の定義との類似に注意してほしい。$u$ が $P$ をもつなら、
$u$ は余連続である（詳細は省略する）。以下で出会う余連続関手の多くは
$P$ を満たす。
\end{remark}

\section{余連続関手とトポスの射}
\label{sites-section-cocontinuous-morphism-topoi}

\noindent
以上から、余連続関手 $u$ は $u$ と同じ向きのトポスの射を
与えることが明らかである。したがって、これは
定義 \ref{sites-definition-morphism-sites}、
命題 \ref{sites-proposition-get-morphism}、および
補題 \ref{sites-lemma-morphism-sites-topoi} のように、連続な $u$ に
（ある条件のもとで）対応するトポスの射とは逆向きである。

\begin{lemma}
\label{sites-lemma-cocontinuous-morphism-topoi}
$\mathcal{C}$ と $\mathcal{D}$ をサイトとする。
$u : \mathcal{C} \to \mathcal{D}$ を余連続とする。
関手 $g_* = {}_su$ と $g^{-1} = (u^p\ )^\#$ は、
$\Sh(\mathcal{C})$ から $\Sh(\mathcal{D})$ へのトポスの射
$g$ を定める。
\end{lemma}

\begin{proof}
これはちょうど補題 \ref{sites-lemma-exact-cocontinuous} の内容である。
\end{proof}

\begin{lemma}
\label{sites-lemma-composition-cocontinuous}
\begin{slogan}
サイト関手の合成は、サイト関手の余連続性を保つ。
\end{slogan}
$u : \mathcal{C} \to \mathcal{D}$ と
$v : \mathcal{D} \to \mathcal{E}$ を余連続関手とする。このとき
$v \circ u$ は余連続であり、$h = g \circ f$ が成り立つ。ここで
$f : \Sh(\mathcal{C}) \to \Sh(\mathcal{D})$、
resp.\ $g : \Sh(\mathcal{D}) \to \Sh(\mathcal{E})$、
resp.\ $h : \Sh(\mathcal{C}) \to \Sh(\mathcal{E})$ はそれぞれ
$u$、resp.\ $v$、resp.\ $v \circ u$ に対応するトポスの射である。
\end{lemma}

\begin{proof}
$U \in \Ob(\mathcal{C})$ とする。
% P01-E0086
$\{E_i \to v(u(U))\}$ を $\mathcal{E}$ における $U$ の被覆とする。
仮定により、$\mathcal{D}$ における被覆 $\{D_j \to u(U)\}$ で、
$\{v(D_j) \to v(u(U))\}$ が $\{E_i \to v(u(U))\}$ の細分となるものが存在する。
さらに仮定により、$\mathcal{C}$ における被覆 $\{C_l \to U\}$ で、
$\{u(C_l) \to u(U)\}$ が $\{D_j \to u(U)\}$ の細分となるものが存在する。
すると $\{v(u(C_l)) \to v(u(U))\}$ は確かに被覆
$\{E_i \to v(u(U))\}$ の細分である。これにより $v \circ u$ が
余連続であることが証明された。最後の主張を証明するには
${}_sv \circ {}_su = {}_s(v \circ u)$ を示せば十分である。
補題 \ref{sites-lemma-pu-sheaf} により、
${}_pv \circ {}_pu = {}_p(v \circ u)$ を示せば十分である。
${}_pu$、resp.\ ${}_pv$、resp.\ ${}_p(v \circ u)$ は
$u^p$、resp.\ $v^p$、resp.\ $(v \circ u)^p$ の右随伴なので、
$u^p \circ v^p = (v \circ u)^p$ を示せば十分である。
これは定義から直ちに従う。
\end{proof}

\begin{example}
\label{sites-example-open-immersion-cocontinuous}
$X$ を位相空間とする。
$j : U  \to X$ を開部分空間の包含とする。
例 \ref{sites-example-site-topological} で見たサイト $X_{Zar}$ と
$U_{Zar}$ を思い出そう。
$j$ に対応する関手 $u : X_{Zar} \to U_{Zar}$ は連続であり、
サイトの射 $U_{Zar} \to X_{Zar}$ を誘導した。
例 \ref{sites-example-continuous-map} を参照せよ。
これはトポスの射 $(j_*, j^{-1})$ も与える。
次に関手
$v : U_{Zar} \to X_{Zar}$、$V \mapsto v(V) = V$
を考える（同じ開集合を、今度は $X_{Zar}$ の対象とみなすだけである）。
この関手は余連続である。実際、$X$ における開被覆
$v(V) = \bigcup_{j \in J} W_j$ があれば、各 $W_j$ は $U$ の
部分集合でなければならず、したがって $v(V_j)$ の形である。
そして $V = \bigcup_{j \in J} V_j$ は自明に $U$ の開被覆である。
以上と補題 \ref{sites-lemma-cocontinuous-morphism-topoi} から、$v$ に対応する
トポスの射
$$
\Sh(U) \longrightarrow \Sh(X)
$$
が存在し、${}_sv$ と $(v^p\ )^\#$ によって与えられる。
実際には $(v^p\ )^\# = j^{-1}$ かつ ${}_sv = j_*$、すなわち
これは上で得たものと同じトポスの射であると主張する。
これを見る最も簡単な方法は、おそらく $X$ 上の任意の層
$\mathcal{G}$ に対して
$v^p\mathcal{G}(V) = \mathcal{G}(V)$ であり、これは
Sheaves, Lemma \ref{sheaves-lemma-j-pullback} によれば
$j^{-1}\mathcal{G}$ の記述になっていることに気づくことである
（したがってこの場合、層化は不要である）。${}_sv$ と $j_*$ の等式は
随伴関手の一意性から従う（直接計算することもできる）。
\end{example}

\begin{example}
\label{sites-example-open-map-cocontinuous}
これは例 \ref{sites-example-open-immersion-cocontinuous} のわずかな一般化である。
$f : X \to Y$ を位相空間の連続写像とし、$f$ は開写像であると仮定する。
例 \ref{sites-example-site-topological} で見たサイト $X_{Zar}$ と
$Y_{Zar}$ を思い出そう。
$f$ に対応する関手 $u : Y_{Zar} \to X_{Zar}$ は連続であり、
サイトの射 $X_{Zar} \to Y_{Zar}$ を誘導した。
例 \ref{sites-example-continuous-map} を参照せよ。
これはトポスの射 $(f_*, f^{-1})$ も与える。
次に関手
$v : X_{Zar} \to Y_{Zar}$、$U \mapsto v(U) = f(U)$ を考える。
この関手は余連続である。実際、$Y$ における開被覆
$f(U) = \bigcup_{j \in J} V_j$ があれば、$U_j = f^{-1}(V_j) \cap U$
とおくことで開被覆 $U = \bigcup U_j$ を得て、
$f(U) = \bigcup f(U_j)$ は $f(U) = \bigcup V_j$ の細分となる。
以上と補題 \ref{sites-lemma-cocontinuous-morphism-topoi} から、$v$ に対応する
トポスの射
$$
\Sh(X) \longrightarrow \Sh(Y)
$$
が存在し、${}_sv$ と $(v^p\ )^\#$ によって与えられる。
実際には $(v^p\ )^\# = f^{-1}$ かつ ${}_sv = f_*$、すなわち
これは上で得たものと同じトポスの射であると主張する。
$Y$ 上の任意の層 $\mathcal{G}$ に対して
$v^p\mathcal{G}(U) = \mathcal{G}(f(U))$ である。
一方、$u_p\mathcal{G}(U) = \colim_{f(U) \subset V} \mathcal{G}(V)
= \mathcal{G}(f(U))$ と計算できる。実際、明らかに
$(f(U), U \subset f^{-1}(f(U)))$ は、
Section \ref{sites-section-functoriality-PSh} の圏 $\mathcal{I}_U^u$ の
始対象である。したがって $u_p = v^p$ であり、
$f^{-1} = u_s = (v^p\ )^\#$ を得る。
${}_sv$ と $f_*$ の等式は随伴関手の一意性から従う
（直接計算することもできる）。
\end{example}

\noindent
最初の例 \ref{sites-example-open-immersion-cocontinuous} では、関手 $v$ は
連続でもある。しかし二つ目の例 \ref{sites-example-open-map-cocontinuous} では、
定義 \ref{sites-definition-continuous} の条件 (2) が成り立たないことがあるため、
一般には連続でない。したがって次の補題は最初の例には適用できるが、
二つ目には適用できない。

\begin{lemma}
\label{sites-lemma-when-shriek}
\begin{slogan}
連続かつ余連続なサイト関手に対応するトポスの射では、層化は冗長である。
\end{slogan}
$\mathcal{C}$ と $\mathcal{D}$ をサイトとする。
$u : \mathcal{C} \to \mathcal{D}$ を関手とする。次を仮定する。
\begin{enumerate}
\item[(a)] $u$ は余連続である。
\item[(b)] $u$ は連続である。
\end{enumerate}
$g : \Sh(\mathcal{C}) \to \Sh(\mathcal{D})$ を対応するトポスの射とする。
このとき次が成り立つ。
\begin{enumerate}
\item 公式 $g^{-1} = (u^p\ )^\#$ における層化は不要である。
すなわち $g^{-1}(\mathcal{G})(U) = \mathcal{G}(u(U))$ である。
\item $g^{-1}$ は左随伴 $g_{!} = (u_p\ )^\#$ をもつ。
\item $g^{-1}$ は任意の極限および余極限と可換である。
\end{enumerate}
\end{lemma}

\begin{proof}
補題 \ref{sites-lemma-pushforward-sheaf} により、$\mathcal{D}$ 上の任意の層
$\mathcal{G}$ に対して前層 $u^p\mathcal{G}$ は $\mathcal{C}$ 上の層である
（そして $u^s\mathcal{G}$ と書かれる）。補題 \ref{sites-lemma-adjoint-sheaves} により、
関手 $\mathcal{F} \mapsto (u_p\mathcal{F})^\#$ は $u^s$ の左随伴である。
極限と余極限についての主張は、Categories, Section
\ref{categories-section-adjoint} の議論から従う。
\end{proof}

\noindent
補題 \ref{sites-lemma-when-shriek} の状況では、随伴関手の列
$$
g_{!}, \ g^{-1}, \ g_*.
$$
を得る。一般に関手 $g_!$ は完全ではない。実際、一般には
$\Sh(\mathcal{C})$ の終対象を $\Sh(\mathcal{D})$ の終対象に移さない。
Sheaves, Remark \ref{sheaves-remark-j-shriek-not-exact} を参照せよ。
一方、例 \ref{sites-example-open-immersion-cocontinuous} の位相的状況では、
関手 $j_!$ はアーベル層上で完全である。
Modules, Lemma \ref{modules-lemma-j-shriek-exact} を参照せよ。
次の補題はこれをサイトの場合に一般化する。

\begin{lemma}
\label{sites-lemma-preserve-equalizers}
$\mathcal{C}$ と $\mathcal{D}$ をサイトとする。
$u : \mathcal{C} \to \mathcal{D}$ を関手とする。次を仮定する。
\begin{enumerate}
\item[(a)] $u$ は余連続である。
\item[(b)] $u$ は連続である。
\item[(c)] $\mathcal{C}$ にファイバー積と等化子が存在し、
$u$ はそれらと可換である。
\end{enumerate}
このとき上の関手 $g_!$ はファイバー積および等化子と可換である
（より一般に、有限連結極限と可換である）。
\end{lemma}

\begin{proof}
(a)、(b)、(c) を仮定する。
$g_! = (u_p\ )^\#$ である。層の極限は、前層としてとった対応する極限に
等しいことを思い出そう（補題 \ref{sites-lemma-limit-sheaf}）。また、層化は
有限極限と可換である（補題 \ref{sites-lemma-sheafification-exact}）。したがって、
$u_p$ がファイバー積および等化子と可換であることを示せば十分である。
そのためには、Section \ref{sites-section-functoriality-PSh} の圏
$(\mathcal{I}_V^u)^{opp}$ にわたる余極限がファイバー積および等化子と
可換であれば十分である。これは補題 \ref{sites-lemma-almost-directed} と
Categories, Lemma
\ref{categories-lemma-almost-directed-commutes-equalizers} から従う。
\end{proof}

\noindent
次の補題は、位相空間の開埋め込みに対応する射にさらによく似た場合を扱う。

\begin{lemma}
\label{sites-lemma-back-and-forth}
$\mathcal{C}$ と $\mathcal{D}$ をサイトとする。
$u : \mathcal{C} \to \mathcal{D}$ を関手とする。次を仮定する。
\begin{enumerate}
\item[(a)] $u$ は余連続である。
\item[(b)] $u$ は連続である。
\item[(c)] $u$ は充満忠実である。
\end{enumerate}
上の $g_!, g^{-1}, g_*$ に対し、すべての $\mathcal{C}$ 上の層
$\mathcal{F}$ について、標準写像 $\mathcal{F} \to g^{-1}g_!\mathcal{F}$ と
$g^{-1}g_*\mathcal{F} \to \mathcal{F}$ は同型である。
\end{lemma}

\begin{proof}
$X$ を $\mathcal{C}$ の対象とする。
補題 \ref{sites-lemma-pu-sheaf} と \ref{sites-lemma-when-shriek} で、関手
$g^{-1} = (u^p\ )^\#$ と $g_{*} = ({}_pu\ )^\#$ には
層化が不要であることを見た。
$(g^{-1}g_{*}\mathcal{F})(X) = g_{*}\mathcal{F}(u(X))
= \lim \mathcal{F}(Y)$ と計算できる。ここで極限は、対
$(Y, u(Y) \to u(X))$ の圏にわたる。この圏では、射
$u(Y) \to u(X)$ が $\mathcal{C}$ の射 $\alpha$ による
$u(\alpha)$ の形であることは要求されていない。仮定 (c) により、
それらは自動的に $\mathcal{C}$ の射から来る。したがって極限は
$(X, u(\text{id}_X))$ における値、すなわち $\mathcal{F}(X)$ である。
これにより $g^{-1}g_{*}\mathcal{F} = \mathcal{F}$ が証明された。

\medskip\noindent
一方、$(g^{-1}g_{!}\mathcal{F})(X) =
g_{!}\mathcal{F}(u(X)) = (u_p\mathcal{F})^\#(u(X))$ であり、
$u_p\mathcal{F}(u(X)) = \colim \mathcal{F}(Y)$ である。
ここで余極限は、対 $(Y, u(X) \to u(Y))$ の圏にわたる。この圏では、射
$u(X) \to u(Y)$ が $\mathcal{C}$ の射 $\alpha$ による
$u(\alpha)$ の形であることは要求されていない。仮定 (c) により、
それらは自動的に $\mathcal{C}$ の射から来る。したがって余極限は
$(X, u(\text{id}_X))$ における値、すなわち $\mathcal{F}(X)$ である。
ゆえにすべての $X \in \Ob(\mathcal{C})$ について
$u_p\mathcal{F}(u(X)) = \mathcal{F}(X)$ である。
$u$ は余連続かつ連続なので、$\mathcal{D}$ における $u(X)$ の任意の被覆は、
$\mathcal{C}$ における被覆 $\{X_i \to X\}$ から得られる
$\mathcal{D}$ の被覆 (!) $\{u(X_i) \to u(X)\}$ によって細分できる。
したがって $(u_p\mathcal{F})^+(u(X)) = \mathcal{F}(X)$ でもある。
実際、$(u_p\mathcal{F})^+$ の $u(X)$ における値を定める余極限では、
上の被覆 $\{u(X_i) \to u(X)\}$ の余終系に制限できる。
よって $\mathcal{C}$ のすべての対象 $X$ についても
$(u_p\mathcal{F})^+(u(X)) = \mathcal{F}(X)$ である。
この議論をもう一度繰り返すと、$\mathcal{C}$ のすべての対象 $X$ について
$(u_p\mathcal{F})^\#(u(X)) = \mathcal{F}(X)$ が得られる。
これは求める等式 $g^{-1}g_!\mathcal{F} = \mathcal{F}$ を与える。
\end{proof}

\noindent
最後に、対応する位相的な例をもたない場合を挙げる。この補題は、
同じ位相を保ったままスキームの「部分宇宙」を拡大するときに何が起こるかを
見るために用いる。補題の状況では、トポスの射
$g : \Sh(\mathcal{C}) \to \Sh(\mathcal{D})$ は $\Sh(\mathcal{C})$ を
$\Sh(\mathcal{D})$ の部分トポスと同一視し
（Section \ref{sites-section-subtopoi}）、さらにこの埋め込みは
レトラクションをもつ。

\begin{lemma}
\label{sites-lemma-bigger-site}
$\mathcal{C}$ と $\mathcal{D}$ をサイトとする。
$u : \mathcal{C} \to \mathcal{D}$ を関手とする。次を仮定する。
\begin{enumerate}
\item[(a)] $u$ は余連続である。
\item[(b)] $u$ は連続である。
\item[(c)] $u$ は充満忠実である。
\item[(d)] $\mathcal{C}$ にファイバー積が存在し、$u$ はそれらと可換である。
\item[(e)] 終対象 $e_\mathcal{C} \in \Ob(\mathcal{C})$ と
$e_\mathcal{D} \in \Ob(\mathcal{D})$ が存在して、
$u(e_\mathcal{C}) = e_\mathcal{D}$ である。
\end{enumerate}
$g_!, g^{-1}, g_*$ を上のものとする。このとき $u$ は、
$f_* = g^{-1}$、$f^{-1} = g_!$ を満たすサイトの射
$f : \mathcal{D} \to \mathcal{C}$ を定める。合成
$$
\xymatrix{
\Sh(\mathcal{C}) \ar[r]^g &
\Sh(\mathcal{D}) \ar[r]^f &
\Sh(\mathcal{C})
}
$$
はトポス $\Sh(\mathcal{C})$ の恒等射と同型である。さらに、関手
$f^{-1}$ は充満忠実である。
\end{lemma}

\begin{proof}
仮定により、関手 $u$ は命題 \ref{sites-proposition-get-morphism} の仮定を
満たす。したがって $u$ はサイトの射を定め、ゆえに補題
\ref{sites-lemma-morphism-sites-topoi} のようなトポスの射 $f$ を定める。
等式 $f_* = g^{-1}$ と $f^{-1} = g_!$ は、引用した補題および
補題 \ref{sites-lemma-when-shriek} から明らかである。
補題 \ref{sites-lemma-back-and-forth} により、
$f_* \circ g_* = g^{-1} \circ g_* \cong \text{id}$ および
$g^{-1} \circ f^{-1} = g^{-1} \circ g_! \cong \text{id}$ が成り立つ。

\medskip\noindent
$f^{-1}$ が充満忠実であることをなお示さなければならない。
$\mathcal{F}, \mathcal{G} \in \Ob(\Sh(\mathcal{C}))$ とする。
写像
$$
\Mor_{\Sh(\mathcal{C})}(\mathcal{F}, \mathcal{G})
\longrightarrow
\Mor_{\Sh(\mathcal{D})}(f^{-1}\mathcal{F}, f^{-1}\mathcal{G})
$$
が全単射であることを示す必要がある。しかし右辺は
\begin{align*}
\Mor_{\Sh(\mathcal{D})}(f^{-1}\mathcal{F}, f^{-1}\mathcal{G})
& =
\Mor_{\Sh(\mathcal{C})}(\mathcal{F}, f_*f^{-1}\mathcal{G}) \\
& =
\Mor_{\Sh(\mathcal{C})}(\mathcal{F}, g^{-1}f^{-1}\mathcal{G}) \\
& =
\Mor_{\Sh(\mathcal{C})}(\mathcal{F}, \mathcal{G})
\end{align*}
に等しい（最初の等式は随伴による）。これで望む主張が証明された。
\end{proof}

\begin{example}
\label{sites-example-closed-map-cocontinuous-false}
$X$ を位相空間とする。
$i : Z  \to X$ を部分集合の包含とする（誘導位相を入れる）。
関手 $u : X_{Zar} \to Z_{Zar}$、$U \mapsto u(U) = Z \cap U$ を考える。
一見すると、この関手も余連続であるように見えるかもしれない。
結局、$Z$ は誘導位相をもつのだから、
% P01-E0087
$U\cap Z$ の任意の被覆は $X$ における $U$ の被覆から来るはずでは
ないだろうか。そうではない！ もし $U \cap Z = \emptyset$ ならどうなるか。
この場合、空被覆は $U \cap Z$ の被覆であり、空被覆を細分できるのは
空被覆だけである。したがって
% P01-E0088
$u$ が余連続 $\Rightarrow$ $X$ の空でない任意の開集合 $U$ は
$Z$ と空でない共通部分をもつ、という結論を得る。
しかしこれは十分ではない。例えば、$X = \mathbf{R}$ を通常の位相をもつ
実数直線とし、$Z = \mathbf{R} \setminus \{0\}$ とする。このとき $Z$ には
開被覆 $Z = \{x < 0\} \cup \bigcup_n \{1/n < x\}$ があるが、これは
$X$ のいかなる開被覆の制限によっても細分できない。
\end{example}

\section{右随伴をもつ余連続関手}
\label{sites-section-cocontinuous-adjoint}

\noindent
$\mathcal{C}$ と $\mathcal{D}$ をサイトとする。台となる圏の関手
$u : \mathcal{C} \to \mathcal{D}$ と $v : \mathcal{D} \to \mathcal{C}$ があり、
$v$ は $u$ の右随伴であるとする。この場合、$v$ が連続ならば $u$ は
余連続である（補題 \ref{sites-lemma-continuous-with-continuous-left-adjoint}）。
$u$ が余連続ならば、多くの場合（ただし常にではない。例
\ref{sites-example-not-equivalent} を参照）、$v$ は連続である。そしてその場合、
$v$ はサイトの射を定め、それに対応するトポスの射は $u$ が定めるものと
同じである。

\begin{lemma}
\label{sites-lemma-have-functor-other-way}
$\mathcal{C}$ と $\mathcal{D}$ をサイトとし、
$u : \mathcal{C} \to \mathcal{D}$ と $v : \mathcal{D} \to \mathcal{C}$ を
関手とする。$u$ は余連続であり、$v$ は $u$ の右随伴であると仮定する。
$g : \Sh(\mathcal{C}) \to \Sh(\mathcal{D})$ を $u$ に対応するトポスの射とする。
補題 \ref{sites-lemma-cocontinuous-morphism-topoi} を参照せよ。このとき、
\begin{enumerate}
\item $\mathcal{C}$ 上の層 $\mathcal{F}$ に対して、層
$g_*\mathcal{F}$ は前層 $v^p\mathcal{F}$ に等しい。言い換えると、
$(g_*\mathcal{F})(V) = \mathcal{F}(v(V))$ である。
\item $\mathcal{D}$ 上の層 $\mathcal{G}$ に対して、
$g^{-1}\mathcal{G} = (v_p\mathcal{G})^\#$ である。
\end{enumerate}
\end{lemma}

\begin{proof}
(1) の $\mathcal{F}$ に対して
$$
g_*\mathcal{F} =
{}_su\mathcal{F} =
{}_pu\mathcal{F} =
v^p\mathcal{F} =
\mathcal{F} \circ v
$$
である。最初の等式は補題 \ref{sites-lemma-cocontinuous-morphism-topoi}、
二つ目は補題 \ref{sites-lemma-pu-sheaf}、三つ目は補題
\ref{sites-lemma-adjoint-functors} による。最後の等式は Section
\ref{sites-section-functoriality-PSh} における $v^p$ の定義である。
これで (1) が証明された。(2) の $\mathcal{G}$ に対して
$$
g^{-1}\mathcal{G} = (u^p\mathcal{G})^\# = (v_p\mathcal{G})^\#
$$
である。最初の等式は補題 \ref{sites-lemma-cocontinuous-morphism-topoi}、
二つ目は補題 \ref{sites-lemma-adjoint-functors} による。
\end{proof}

\begin{lemma}
\label{sites-lemma-have-functor-other-way-morphism}
記法と仮定を補題 \ref{sites-lemma-have-functor-other-way} のものとする。
さらに $v$ が連続ならば、$v$ はサイトの射
$f : \mathcal{C} \to \mathcal{D}$ を定め、それに対応するトポスの射は
$g$ に等しい。
\end{lemma}

\begin{proof}
補題 \ref{sites-lemma-have-functor-other-way} の結果を断りなく用いる。
$v$ が補題の主張にあるサイトの射 $f$ を定めることを証明するには、
$v_s$ が完全関手であることを示さなければならない
（定義 \ref{sites-definition-morphism-sites} を参照）。しかし
$v_s\mathcal{G} = (v_p\mathcal{G})^\# = g^{-1}\mathcal{G}$ なので、
これは $g$ がトポスの射であることから従う。すると
$f^{-1} = v_s = g^{-1}$ であり、トポスの射として $f = g$ を得る。
\end{proof}

\begin{example}
\label{sites-example-finer-topology-bis}
この例は例 \ref{sites-example-finer-topology} の議論を続けるものであり、
そこから記法 $\mathcal{C}, \tau, \tau', \epsilon$ を借用する。
恒等関手 $v : \mathcal{C}_{\tau'} \to \mathcal{C}_\tau$ は連続関手であり、
恒等関手 $u : \mathcal{C}_\tau \to \mathcal{C}_{\tau'}$ は余連続関手である。
さらに $u$ は $v$ の左随伴である。したがって補題
\ref{sites-lemma-have-functor-other-way} と
\ref{sites-lemma-have-functor-other-way-morphism} の結果を適用でき、
$v$ はサイトの射、すなわち
$$
\epsilon : \mathcal{C}_\tau \longrightarrow \mathcal{C}_{\tau'}
$$
を定める。その対応するトポスの射は、余連続関手 $u$ に対応する
トポスの射と同じである。
\end{example}

\begin{lemma}
\label{sites-lemma-continuous-with-continuous-left-adjoint}
$\mathcal{C}$ と $\mathcal{D}$ をサイトとし、
$v : \mathcal{D} \to \mathcal{C}$ を連続関手とする。
$v$ が左随伴 $u : \mathcal{C} \to \mathcal{D}$ をもつと仮定する。このとき、
\begin{enumerate}
\item $u$ は余連続である。
\item 補題 \ref{sites-lemma-have-functor-other-way} と
\ref{sites-lemma-have-functor-other-way-morphism} の結果が成り立つ。
\end{enumerate}
特に、$v$ はサイトの射 $f : \mathcal{C} \to \mathcal{D}$ を定める。
\end{lemma}

\begin{proof}
$U$ を $\mathcal{C}$ の対象とし、$\{V_j \to u(U)\}_{j \in J}$ を
$\mathcal{D}$ における被覆とする。このとき
% P01-E0089
$\{v(V_j) \to v(u(U))\}_{i \in I}$ は $\mathcal{C}$ における被覆である。
随伴写像 $U \to v(u(U))$ に沿って基底変換すると、
$W_j = v(V_j) \times_{v(u(U))} U$ を満たす被覆
$\{W_j \to U\}_{j \in J}$ を得る。第一射影を
$p_j : W_j \to v(V_j)$ と書けば、写像
$$
u(W_j) \xrightarrow{u(p_j)} u(v(V_j)) \longrightarrow V_j
$$
を得る。ここで二つ目の矢印は随伴写像である。これは固定された終域をもつ
写像族の射 $\{u(W_j) \to u(U)\} \to \{V_j \to u(U)\}$ を定めるので、
$u$ が実際に余連続であることが分かる。他の主張は補題
\ref{sites-lemma-have-functor-other-way} と
\ref{sites-lemma-have-functor-other-way-morphism} から直ちに従う。
\end{proof}

\begin{example}
\label{sites-example-not-equivalent}
$\mathcal{C}$ と $\mathcal{D}$ をサイトとする。台となる圏の関手
$u : \mathcal{C} \to \mathcal{D}$ と $v : \mathcal{D} \to \mathcal{C}$ があり、
$v$ は $u$ の右随伴であるとする。補題
\ref{sites-lemma-continuous-with-continuous-left-adjoint} により、$v$ が連続ならば、
% P01-E0090
$u$ は余連続である。逆に、$u$ が余連続であっても、$v$ が連続であるとは
結論できない。この現象の例をスキームの大エタール・サイトと
滑らかサイトを用いて与えるが、おそらく初等的な例もあるだろう。
実際、スキーム $S$ とサイト $(\Sch/S)_\etale$ および
$(\Sch/S)_{smooth}$ を考える。これらのサイトは同じ台となる圏をもつと
仮定してよい。Topologies, Remark
\ref{topologies-remark-choice-sites} を参照せよ。
$u = v = \text{id}$ とする。このとき、$(\Sch/S)_\etale$ から
$(\Sch/S)_{smooth}$ への関手としての $u$ は余連続である。
実際、スキームの任意の滑らかな被覆はエタール被覆によって細分できる。
More on Morphisms, Lemma
\ref{more-morphisms-lemma-etale-dominates-smooth} を参照せよ。
逆に、$(\Sch/S)_{smooth}$ から $(\Sch/S)_\etale$ への関手 $v$ は
連続でない。一般に滑らかな被覆はエタール被覆ではないからである。
\end{example}

\section{左随伴をもつ余連続関手}
\label{sites-section-cocontinuous-left-adjoint}

\noindent
余連続関手 $u$ が左随伴 $w$ をもつことがある。

\begin{lemma}
\label{sites-lemma-have-left-adjoint}
$\mathcal{C}$ と $\mathcal{D}$ をサイトとする。
$g : \Sh(\mathcal{C}) \to \Sh(\mathcal{D})$ を、連続かつ余連続な関手
$u : \mathcal{C} \to \mathcal{D}$ に対応するトポスの射とする。
補題 \ref{sites-lemma-cocontinuous-morphism-topoi} と
\ref{sites-lemma-when-shriek} を参照せよ。
\begin{enumerate}
\item $w : \mathcal{D} \to \mathcal{C}$ が $u$ の左随伴ならば、
\begin{enumerate}
\item $g_!\mathcal{F}$ は前層 $w^p\mathcal{F}$ に付随する層である。
\item $g_!$ は完全である。
\end{enumerate}
\item $w$ が連続な左随伴ならば、$g_!$ は左随伴をもつ。
\item $w$ が余連続な左随伴ならば、$g_! = h^{-1}$ かつ
$g^{-1} = h_*$ である。ここで $h : \Sh(\mathcal{D}) \to \Sh(\mathcal{C})$ は
$w$ に対応するトポスの射である。
\end{enumerate}
\end{lemma}

\begin{proof}
$g_!\mathcal{F}$ は $u_p\mathcal{F}$ の層化であることを思い出そう。
したがって補題 \ref{sites-lemma-adjoint-functors} による等式 $u_p = w^p$ から
(1)(a) が従う。

\medskip\noindent
(1)(b) を見るため、$g_!$ はすべての余極限と可換することに注意しよう。
実際、$g_!$ は左随伴である
（Categories, Lemma \ref{categories-lemma-adjoint-exact}）。
$i \mapsto \mathcal{F}_i$ を $\Sh(\mathcal{C})$ の有限図式とする。
すると $\lim \mathcal{F}_i$ は前層の圏で計算される
（補題 \ref{sites-lemma-limit-sheaf}）。$w^p$ は右随伴なので
（補題 \ref{sites-lemma-adjoints-u}）、
$w^p \lim \mathcal{F}_i = \lim w^p\mathcal{F}_i$ である。
層化は完全なので（補題 \ref{sites-lemma-sheafification-exact}）、
(1)(a) によって結論を得る。

\medskip\noindent
$w$ が連続であると仮定する。このとき
$g_! = (w^p\ )^\# = w^s$ であるが、層化は不要であり、左随伴 $w_s$ をもつ。
補題 \ref{sites-lemma-pushforward-sheaf} と \ref{sites-lemma-adjoint-sheaves} を参照せよ。

\medskip\noindent
$w$ が余連続であると仮定する。等式 $g_! = h^{-1}$ は (1)(a) と定義から
従う。等式 $g^{-1} = h_*$ は $g_! = h^{-1}$ と随伴関手の一意性から従う。
あるいは補題 \ref{sites-lemma-have-functor-other-way} から導くこともできる。
\end{proof}

\section{下付きシュリークの存在}
\label{sites-section-lower-shriek}

\noindent
この節では、$f^{-1}$ が左随伴 $f_!$ をもつトポスの射 $f$ のいくつかの
場合を論じる。

\begin{lemma}
\label{sites-lemma-existence-lower-shriek}
$\mathcal{C}$、$\mathcal{D}$ を二つのサイトとする。
$f : \Sh(\mathcal{D}) \to \Sh(\mathcal{C})$ をトポスの射とする。
$E \subset \Ob(\mathcal{D})$ を次を満たす部分集合とする。
\begin{enumerate}
\item $V \in E$ に対して $\mathcal{C}$ 上の層 $\mathcal{G}$ が存在し、
$\Sh(\mathcal{C})$ の $\mathcal{F}$ に関して関手的に
$f^{-1}\mathcal{F}(V) =
\Mor_{\Sh(\mathcal{C})}(\mathcal{G}, \mathcal{F})$ となる。
\item $\mathcal{D}$ のすべての対象は $E$ の対象による被覆をもつ。
\end{enumerate}
このとき $f^{-1}$ は左随伴 $f_!$ をもつ。
\end{lemma}

\begin{proof}
米田の補題（Categories, Lemma \ref{categories-lemma-yoneda}）により、
$V \in E$ に対応する層 $\mathcal{G}_V$ は公式
$f^{-1}\mathcal{F}(V) = \Mor_{\Sh(\mathcal{C})}(\mathcal{G}_V, \mathcal{F})$
によって同型を除いて一意に定まる。公式
$f^{-1}\mathcal{F}(V) = \Mor_{\Sh(\mathcal{D})}(h_V^\#, f^{-1}\mathcal{F})$
を思い出そう。$\Mor(\mathcal{G}_V, \mathcal{G}_V)$ の $\text{id}$ に対応する
写像を $i_V : h_V^\# \to f^{-1}\mathcal{G}_V$ と書く。
(1) の関手性により、この全単射は
$$
\Mor_{\Sh(\mathcal{C})}(\mathcal{G}_V, \mathcal{F}) \to
\Mor_{\Sh(\mathcal{D})}(h_V^\#, f^{-1}\mathcal{F}),\quad
\varphi \mapsto f^{-1}\varphi \circ i_V
$$
で与えられる。任意の $V_1, V_2 \in E$ に対し、標準写像
$$
\Mor_{\Sh(\mathcal{D})}(h^\#_{V_2}, h^\#_{V_1})
\to
\Hom_{\Sh(\mathcal{C})}(\mathcal{G}_{V_2}, \mathcal{G}_{V_1}),\quad
\varphi \mapsto f_!(\varphi)
$$
が存在し、これは
$f^{-1}(f_!(\varphi)) \circ i_{V_2} = i_{V_1} \circ \varphi$ によって
特徴づけられる。$\varphi \mapsto f_!(\varphi)$ は合成と両立することに
注意しよう。これは特徴づけから直接分かる。したがって
$h_V^\# \mapsto \mathcal{G}_V$ と $\varphi \mapsto f_!\varphi$ は、
対象が $h_V^\#$ である $\Sh(\mathcal{D})$ の充満部分圏からの関手である。

\medskip\noindent
$J$ を集合とし、$J \to E$、$j \mapsto V_j$ を写像とする。このとき、
上の全単射の積を用いて、関手的全単射
$$
\Mor_{\Sh(\mathcal{C})}(\coprod \mathcal{G}_{V_j}, \mathcal{F})
\longrightarrow
\Mor_{\Sh(\mathcal{D})}(\coprod h_{V_j}^\#, f^{-1}\mathcal{F})
$$
を得る。したがって関手 $f_!$ を、$V \in E$ に対する $h_V^\#$ の余積を
対象とする $\Sh(\mathcal{D})$ の充満部分圏へ拡張できる。

\medskip\noindent
$\mathcal{D}$ 上の任意の層 $\mathcal{H}$ が与えられたとき、余等化子図式
$$
\xymatrix{
\mathcal{H}_1 \ar@<1ex>[r] \ar@<-1ex>[r] &
\mathcal{H}_0 \ar[r] &
\mathcal{H}
}
$$
を選ぶ。ここで $\mathcal{H}_i = \coprod h_{V_{i, j}}^\#$ は
$V_{i, j} \in E$ を満たす余積である。これは仮定 (2) により可能である。
補題 \ref{sites-lemma-sheaf-coequalizer-representable} を参照せよ
（集合論的問題を懸念する読者のために、補題
\ref{sites-lemma-sheaf-coequalizer-representable} で与えられる構成は
標準的であることに注意しておく）。$f_!(\mathcal{H})$ を、
$$
\xymatrix{
f_!\mathcal{H}_1 \ar@<1ex>[r] \ar@<-1ex>[r] &
f_!\mathcal{H}_0 \ar[r] &
f_!\mathcal{H}
}
$$
が余等化子図式となるような $\mathcal{C}$ 上の層として定める。このとき
\begin{align*}
\Mor(f_!\mathcal{H}, \mathcal{F})
& =
\text{Equalizer}(
\xymatrix{
\Mor(f_!\mathcal{H}_0, \mathcal{F}) \ar@<1ex>[r] \ar@<-1ex>[r] &
\Mor(f_!\mathcal{H}_1, \mathcal{F})
}
) \\
& =
\text{Equalizer}(
\xymatrix{
\Mor(\mathcal{H}_0, f^{-1}\mathcal{F}) \ar@<1ex>[r] \ar@<-1ex>[r] &
\Mor(\mathcal{H}_1, f^{-1}\mathcal{F})
}
) \\
& =
\Hom(\mathcal{H}, f^{-1}\mathcal{F})
\end{align*}
したがって $f_!$ を $\mathcal{D}$ 上の層の圏全体へ拡張できることが分かる。
\end{proof}

\section{局所化}
\label{sites-section-localize}

\noindent
$\mathcal{C}$ をサイトとし、$U \in \Ob(\mathcal{C})$ とする。
$U$ 上の対象の圏 $\mathcal{C}/U$ の定義については
Categories, Example \ref{categories-example-category-over-X} を参照せよ。
$U$ 上の対象の射の族 $\{V_j \to V\}$ が $\mathcal{C}/U$ の被覆であることを、
それが $\mathcal{C}$ の被覆であることと同値であると定めて、
$\mathcal{C}/U$ をサイトにする。忘却関手
$$
j_U : \mathcal{C}/U \longrightarrow \mathcal{C}.
$$
を考える。これは明らかに余連続かつ連続である。したがって補題
\ref{sites-lemma-when-shriek} により、$j_U^{-1}$ と $j_{U*}$ で与えられる
トポスの射
$$
j_U : \Sh(\mathcal{C}/U) \longrightarrow \Sh(\mathcal{C})
$$
と、$j_U^{-1}$ の左随伴である関手 $j_{U!}$ を得る。

\begin{definition}
\label{sites-definition-localize}
$\mathcal{C}$ をサイトとし、$U \in \Ob(\mathcal{C})$ とする。
\begin{enumerate}
\item サイト $\mathcal{C}/U$ を、{\it 対象 $U$ におけるサイト
$\mathcal{C}$ の局所化}という。
\item トポスの射 $j_U : \Sh(\mathcal{C}/U) \to \Sh(\mathcal{C})$ を
{\it 局所化射}という。
\item 関手 $j_{U*}$ を{\it 順像関手}という。
\item $\mathcal{C}$ 上の層 $\mathcal{F}$ に対して、層
$j_U^{-1}\mathcal{F}$ を $\mathcal{C}/U$ への
{\it $\mathcal{F}$ の制限}という。
\item $\mathcal{C}/U$ 上の層 $\mathcal{G}$ に対して、層
$j_{U!}\mathcal{G}$ を{\it 空集合による $\mathcal{G}$ の延長}という。
\end{enumerate}
\end{definition}

\noindent
制限 $j_U^{-1}\mathcal{F}$ は、期待どおり規則
$j_U^{-1}\mathcal{F}(X/U) = \mathcal{F}(X)$ によって定められる層である。
空集合による延長もこの場合には非常に簡単に記述できる。次がその記述である。

\begin{lemma}
\label{sites-lemma-describe-j-shriek}
$\mathcal{C}$ をサイトとし、$U \in \Ob(\mathcal{C})$ とする。
$\mathcal{G}$ を $\mathcal{C}/U$ 上の前層とする。このとき
$j_{U!}(\mathcal{G}^\#)$ は、明らかな制限写像をもつ前層
$$
V
\longmapsto
\coprod\nolimits_{\varphi \in \Mor_\mathcal{C}(V, U)}
\mathcal{G}(V \xrightarrow{\varphi} U)
$$
に付随する層である。
\end{lemma}

\begin{proof}
補題 \ref{sites-lemma-when-shriek} により
$j_{U!}(\mathcal{G}^\#) = ((j_U)_p\mathcal{G}^\#)^\#$ である。
補題 \ref{sites-lemma-technical-up} により、これは
$((j_U)_p\mathcal{G})^\#$ に等しい。したがって、$\mathcal{C}/U$ 上の任意の
前層 $\mathcal{G}$ に対して $(j_U)_p$ が上の公式で与えられることを
示せば十分である。
% P01-E0091
Section \ref{sites-section-functoriality-PSh} の定義により
$$
(j_U)_p\mathcal{G}(V)
=
\colim_{(W/U, V \to W)} \mathcal{G}(W)
$$
である。対 $(W/U, V \to W)$ の圏は、各 $\varphi : V \to U$ に対して
対象 $O_\varphi = (\varphi : V \to U, \text{id} : V \to V)$ をもち、
さらに任意の対象へは、いずれかの $O_\varphi$ から一意な射があることが
明らかである。これで結果が従う。
\end{proof}

\begin{lemma}
\label{sites-lemma-describe-j-shriek-representable}
$\mathcal{C}$ をサイトとし、$U \in \Ob(\mathcal{C})$ とする。
$X/U$ を $\mathcal{C}/U$ の対象とする。このとき
$j_{U!}(h_{X/U}^\#) = h_X^\#$ である。
\end{lemma}

\begin{proof}
$p : X \to U$ を $X$ の構造射と書く。補題
\ref{sites-lemma-describe-j-shriek} により、$j_{U!}(h_{X/U}^\#)$ は前層
$$
V
\longmapsto
\coprod\nolimits_{\varphi \in \Mor_\mathcal{C}(V, U)}
\{\psi : V \to X \mid p \circ \psi = \varphi\}
$$
に付随する層である。これは明らかに $\Mor_\mathcal{C}(V, X)$ と同じである。
したがって補題が従う。
\end{proof}

\noindent
上の二つの補題のいずれからも $j_{U!}(*) = h_U^\#$ である。
したがって $\mathcal{C}/U$ 上の任意の層 $\mathcal{G}$ に対して、層の
標準写像 $j_{U!}\mathcal{G} \to h_U^\#$ がある。これは $j_{U!}$ の
本質像にある層を特徴づける。

\begin{lemma}
\label{sites-lemma-essential-image-j-shriek}
$\mathcal{C}$ をサイトとし、$U \in \Ob(\mathcal{C})$ とする。
関手 $j_{U!}$ は圏の同値
$$
\Sh(\mathcal{C}/U)
\longrightarrow
\Sh(\mathcal{C})/h_U^\#
$$
を与える。
\end{lemma}

\begin{proof}
$\mathcal{C}/U$ の対象を対 $(X, a)$ と書くことにする。ここで $X$ は
$\mathcal{C}$ の対象であり、$a : X \to U$ は $\mathcal{C}$ の射である。
同様に、$\Sh(\mathcal{C})/h_U^\#$ の対象を対
$(\mathcal{F}, \varphi)$ と書く。関手
$\Sh(\mathcal{C}/U) \to \Sh(\mathcal{C})/h_U^\#$ は $\mathcal{G}$ を
対 $(j_{U!}\mathcal{G}, \gamma)$ に送る。ここで $\gamma$ は
$j_{U!}\mathcal{G} \to j_{U!}*$ と同一視 $j_{U!}* = h_U^\#$ の合成である。

\medskip\noindent
$\Sh(\mathcal{C})/h_U^\#$ から $\Sh(\mathcal{C}/U)$ への関手を構成しよう。
$(\mathcal{F}, \varphi)$ が与えられたとする。$\mathcal{C}/U$ の対象
$(X, a)$ に対し、$\varphi$ によって
$a \in \Mor_\mathcal{C}(X, U) = h_U(X)$ の $h_U^\#(X)$ における像へ
写る元 $s \in \mathcal{F}(X)$ の集合を
$\mathcal{F}_\varphi(X, a)$ とする。
$(X, a) \mapsto \mathcal{F}_\varphi(X, a)$ が $\mathcal{C}/U$ 上の層であることは
容易に分かる。明らかに、規則
$(\mathcal{F}, \varphi) \mapsto \mathcal{F}_\varphi$ は関手
$\Sh(\mathcal{C})/h_U^\# \to \Sh(\mathcal{C}/U)$ を定める。

\medskip\noindent
さらに関手
$\textit{PSh}(\mathcal{C})/h_U \to \textit{PSh}(\mathcal{C}/U)$、
$(\mathcal{F}, \varphi) \mapsto \mathcal{F}_\varphi$ を考える。ここで
$\mathcal{F}_\varphi(X, a)$ は、$a \in h_U(X)$ へ写る
$\mathcal{F}(X)$ の元の集合と定める。図式
$$
\xymatrix{
\textit{PSh}(\mathcal{C})/h_U \ar[r] \ar[d] &
\textit{PSh}(\mathcal{C}/U) \ar[d] \\
\Sh(\mathcal{C})/h_U^\# \ar[r] &
\Sh(\mathcal{C}/U)
}
$$
は可換であると主張する。ここで縦の矢印は層化で与えられる。
これを見るには\footnote{別の方法として、前層についてのカルテジアン図式
$$
\vcenter{
\xymatrix{
\mathcal{F}_\varphi \ar[r] \ar[d] & {*} \ar[d] \\
\mathcal{F}|_{\mathcal{C}/U} \ar[r] & h_U|_{\mathcal{C}/U}
}
}
\quad\text{for presheaves and}\quad
\vcenter{
\xymatrix{
\mathcal{F}_\varphi \ar[r] \ar[d] & {*} \ar[d] \\
\mathcal{F}|_{\mathcal{C}/U} \ar[r] & h_U^\#|_{\mathcal{C}/U}
}
}
$$
によって、また層について同様に $\mathcal{F}_\varphi$ を記述し、
$\mathcal{C}/U$ への制限が層化と可換であることを用いてもよい。}、
補題 \ref{sites-lemma-plus-presheaf}、\ref{sites-lemma-plus-functorial} と
定理 \ref{sites-theorem-plus} の関手 $\mathcal{F} \mapsto \mathcal{F}^+$ と
この構成が可換であることを示せば十分である。
$\mathcal{F} \mapsto \mathcal{F}^+$ との可換性は、$(X, a)$ が与えられたとき、
$\mathcal{C}/U$ における $(X, a)$ の被覆の圏と $\mathcal{C}$ における $X$ の
被覆の圏が標準的に同一視されることから従う。

\medskip\noindent
次に、$\textit{PSh}(\mathcal{C}/U) \to \textit{PSh}(\mathcal{C})/h_U$ が
$\mathcal{G}$ を対 $(j_{U!}^{PSh}\mathcal{G}, \gamma)$ に送るものとする。
% P01-E0092
ここで $j_{U!}^{PSh}\mathcal{G}$ は補題
\ref{sites-lemma-describe-j-shriek} の公式で定められる前層であり、$\gamma$ は
% P01-E0093
$j_{U!}^{PSh}\mathcal{G} \to j_{U!}*$ と同一視
$j_{U!}^{PSh}* = h_U$ の合成である（公式から明らかである）。
すると図式
$$
\xymatrix{
\textit{PSh}(\mathcal{C}/U) \ar[r] \ar[d] &
\textit{PSh}(\mathcal{C})/h_U \ar[d] \\
\Sh(\mathcal{C}/U) \ar[r] &
\Sh(\mathcal{C})/h_U^\#
}
$$
が可換であることは直ちに明らかである。ここで縦の矢印は層化である。
以上を合わせると、$\textit{PSh}(\mathcal{C}/U)$ の $\mathcal{G}$ に対する
関手的同型 $(j_{U!}^{PSh}\mathcal{G})_\gamma = \mathcal{G}$ と、
$\textit{PSh}(\mathcal{C})/h_U$ の $(\mathcal{F}, \varphi)$ に対する
$j_{U!}^{PSh}\mathcal{F}_\varphi = \mathcal{F}$ が存在することを
示せば十分である。前層 $(j_{U!}^{PSh}\mathcal{G})_\gamma$ の
$(X, a)$ における値は、写像
$$
\coprod\nolimits_{a' : X \to U} \mathcal{G}(X, a') \to \Mor_\mathcal{C}(X, U)
$$
の $a$ 上のファイバーであり、これは $\mathcal{G}(X, a)$ である。
これで最初の等式が証明された。
% P01-E0094
前層 $j_{U!}^{PSh}\mathcal{F}_\varphi$ の $X$ における値は
$$
\coprod\nolimits_{a : X \to U} \mathcal{F}_\varphi(X, a) =
\mathcal{F}(X)
$$
である。実際、集合の写像 $S \to S'$ が与えられたとき、集合 $S$ は
そのファイバーの非交和である。
\end{proof}

\noindent
補題 \ref{sites-lemma-essential-image-j-shriek} によれば、関手 $j_{U!}$ は合成
$$
\Sh(\mathcal{C}/U) \rightarrow
\Sh(\mathcal{C})/h_U^\# \rightarrow
\Sh(\mathcal{C})
$$
であり、最初の矢印は同値である。

\begin{lemma}
\label{sites-lemma-j-shriek-commutes-equalizers-fibre-products}
$\mathcal{C}$ をサイトとし、$U \in \Ob(\mathcal{C})$ とする。
関手 $j_{U!}$ はファイバー積および等化子と可換する
（より一般に有限連結極限と可換する）。特に、
$\Sh(\mathcal{C}/U)$ において $\mathcal{F} \subset \mathcal{F}'$ ならば、
$j_{U!}\mathcal{F} \subset j_{U!}\mathcal{F}'$ である。
\end{lemma}

\begin{proof}
補題 \ref{sites-lemma-essential-image-j-shriek} と、圏の同値がすべての極限と
可換するという事実により、これは関手
$\Sh(\mathcal{C})/h_U^\# \rightarrow \Sh(\mathcal{C})$ がファイバー積および
等化子と可換するという事実に帰着する。あるいは、層化が完全であることを
用い、補題 \ref{sites-lemma-describe-j-shriek} における $j_{U!}$ の記述から
直接証明できる。（また、$\mathcal{C}$ にファイバー積と等化子がある場合、
結果は補題 \ref{sites-lemma-preserve-equalizers} から従う。）
\end{proof}

\begin{lemma}
\label{sites-lemma-j-shriek-reflects-injections-surjections}
$\mathcal{C}$ をサイトとし、$U \in \Ob(\mathcal{C})$ とする。
関手 $j_{U!}$ は単射と全射を反映する。
\end{lemma}

\begin{proof}
$j_{U!}$ が単射と全射を反映することを示す必要がある。
補題 \ref{sites-lemma-mono-epi-sheaves} を参照せよ。補題
\ref{sites-lemma-essential-image-j-shriek} により、これは関手
$\Sh(\mathcal{C})/h_U^\# \to \Sh(\mathcal{C})$ が単射と全射を反映することに
帰着する。
\end{proof}

\begin{lemma}
\label{sites-lemma-compute-j-shriek-restrict}
$\mathcal{C}$ をサイトとし、$U \in \Ob(\mathcal{C})$ とする。
$\mathcal{C}$ 上の任意の層 $\mathcal{F}$ に対して
$j_{U!}j_U^{-1}\mathcal{F} = \mathcal{F} \times h_U^\#$ である。
\end{lemma}

\begin{proof}
補題 \ref{sites-lemma-describe-j-shriek} における $j_{U!}$ の記述から明らかである。
\end{proof}

\begin{lemma}
\label{sites-lemma-relocalize}
$\mathcal{C}$ をサイトとし、$f : V \to U$ を $\mathcal{C}$ の射とする。
このとき、連続かつ余連続な関手の可換図式
$$
\xymatrix{
\mathcal{C}/V \ar[rd]_{j_V} \ar[rr]_j & &
\mathcal{C}/U \ar[ld]^{j_U} \\
& \mathcal{C} &
}
$$
が存在する。関手 $j : \mathcal{C}/V \to \mathcal{C}/U$、
$(a : W \to V) \mapsto (f \circ a : W \to U)$ は、同一視
$(\mathcal{C}/U)/(V/U) = \mathcal{C}/V$ のもとで関手
$j_{V/U} : (\mathcal{C}/U)/(V/U) \to \mathcal{C}/U$ と同一視される。
さらに $j_{V!} = j_{U!} \circ j_!$、
$j_V^{-1} = j^{-1} \circ j_U^{-1}$、および
$j_{V*} = j_{U*} \circ j_*$ である。
\end{lemma}

\begin{proof}
図式の可換性は直ちに分かる。$j$ と $j_{V/U}$ の一致は定義から従う。
補題 \ref{sites-lemma-composition-cocontinuous} により、次のトポスの射の図式
\begin{equation}
\label{sites-equation-relocalize}
\vcenter{
\xymatrix{
\Sh(\mathcal{C}/V) \ar[rd]_{j_V} \ar[rr]_j & &
\Sh(\mathcal{C}/U) \ar[ld]^{j_U} \\
& \Sh(\mathcal{C}) &
}
}
\end{equation}
は可換である。これにより
$j_V^{-1} = j^{-1} \circ j_U^{-1}$ と $j_{V*} = j_{U*} \circ j_*$ が
証明される。等式 $j_{V!} = j_{U!} \circ j_!$ は随伴性から形式的に従う。
\end{proof}

\begin{lemma}
\label{sites-lemma-relocalize-explicit}
記法 $\mathcal{C}$、$f : V \to U$、$j_U$、$j_V$、$j$ を補題
\ref{sites-lemma-relocalize} のものとする。補題
\ref{sites-lemma-essential-image-j-shriek} の同一視
$\Sh(\mathcal{C}/V) = \Sh(\mathcal{C})/h_V^\#$ および
$\Sh(\mathcal{C}/U) = \Sh(\mathcal{C})/h_U^\#$ のもとで、次が成り立つ。
\begin{enumerate}
\item 関手 $j^{-1}$ は次の記述をもつ。
$$
j^{-1}(\mathcal{H} \xrightarrow{\varphi} h_U^\#)
=
(\mathcal{H} \times_{\varphi, h_U^\#, f} h_V^\# \to h_V^\#).
$$
\item 関手 $j_!$ は次の記述をもつ。
$$
j_!(\mathcal{H} \xrightarrow{\varphi} h_V^\#) =
(\mathcal{H} \xrightarrow{h_f \circ \varphi} h_U^\#)
$$
\end{enumerate}
\end{lemma}

\begin{proof}
(2) を証明する。同一視
$\Sh(\mathcal{C}/V) \to \Sh(\mathcal{C})/h_V^\#$ は $\mathcal{G}$ を
$j_{V!}\mathcal{G} \to j_{V!}(*) = h_V^\#$ に送り、
$\Sh(\mathcal{C}/U) \to \Sh(\mathcal{C})/h_U^\#$ についても同様であることを
思い出そう。したがって $j_!\mathcal{G}$ は
$j_{U!}(j_!\mathcal{G}) \to j_{U!}(*) = h_U^\#$ に写り、補題
\ref{sites-lemma-relocalize} により $j_{U!}j_! = j_{V!}$ なので (2) が従う。

\medskip\noindent
$j^{-1}$ が $j_!$ の右随伴であり、(1) の規則が (2) の規則の右随伴である
ことを用いれば、読者は (1) を証明できる。ここでは直接証明を与える。
$\varphi : \mathcal{H} \to h_U^\#$ が $\Sh(\mathcal{C})/h_U^\#$ の対象であると
する。補題 \ref{sites-lemma-essential-image-j-shriek} の証明により、これは
$\mathcal{C}/U$ 上の層 $\mathcal{H}_\varphi$ で、規則
$$
(a : W \to U)
\longmapsto
\{ s \in \mathcal{H}(W) \mid \varphi(s) = a\}
$$
によって定められるものに対応する。これは $\mathcal{C}/U$ 上の規則である。
$\mathcal{C}/V$ への引き戻し
$j^{-1}\mathcal{H}_\varphi$ は規則
$$
(a : W \to V)
\longmapsto
\{ s \in \mathcal{H}(W) \mid \varphi(s) = f \circ a\}
$$
によって与えられる。これは
% P01-E0095
$j^{-1} = j_{U/V}^{-1}$ を $\mathcal{H}_\varphi$ の $\mathcal{C}/V$ への制限として
記述したことによる。一方、$\Sh(\mathcal{C})/h_V^\#$ の対象
$$
\xymatrix{
\mathcal{H}' = \mathcal{H} \times_{\varphi, h_U^\#, f} h_V^\#
\ar[rr]^-{\varphi'} & & h_V^\#
}
$$
にこの規則を適用すると、$\mathcal{H}'_{\varphi'}$ が得られ、これは
\begin{align*}
(a : W \to V)
\longmapsto
&  \{ s' \in \mathcal{H}'(W) \mid \varphi'(s') = a\} \\
= &
\{ (s, a') \in \mathcal{H}(W) \times h_V^\#(W) \mid
a' = a \text{ and } \varphi(s) = f \circ a'\}
\end{align*}
によって与えられる。これは上の $j^{-1}\mathcal{H}_\varphi$ を記述する規則と
まったく同じである。
\end{proof}

\begin{remark}
\label{sites-remark-localize-presheaves}
局所化と前層。$\mathcal{C}$ を圏とし、$U$ を $\mathcal{C}$ の対象とする。
厳密には、関手 $j_U^{-1}$、$j_{U*}$、$j_{U!}$ は前層について定義されて
いない。しかしもちろん、前層を $\mathcal{C}$ 上のカオス位相に関する層と
みなすことができる（例 \ref{sites-example-indiscrete} を参照）。したがって関手
$$
j_U^{-1} :
\textit{PSh}(\mathcal{C})
\longrightarrow
\textit{PSh}(\mathcal{C}/U)
$$
と関手
$$
j_{U*}, j_{U!} :
\textit{PSh}(\mathcal{C}/U)
\longrightarrow
\textit{PSh}(\mathcal{C})
$$
も得る。これらはそれぞれ $j_U^{-1}$ の右随伴、左随伴である。
補題 \ref{sites-lemma-describe-j-shriek} により、$j_{U!}\mathcal{G}$ は前層
$$
V \longmapsto
\coprod\nolimits_{\varphi \in \Mor_\mathcal{C}(V, U)}
\mathcal{G}(V \xrightarrow{\varphi} U)
$$
である。さらに関手 $j_{U!}$ はファイバー積および等化子と可換する。
\end{remark}

\begin{remark}
\label{sites-remark-localization-cartesian-cocontinuous}
$\mathcal{C}$ をサイトとし、$U \to V$ を $\mathcal{C}$ の射とする。
余連続関手 $\mathcal{C}/U \to \mathcal{C}$ と
$j : \mathcal{C}/U \to \mathcal{C}/V$（補題 \ref{sites-lemma-relocalize}）は、
Remark \ref{sites-remark-cartesian-cocontinuous} の性質 $P$ を満たす。
例えば、対象 $(X/U)$、$(W/V)$、射 $g : j(X/U) \to (W/V)$、および被覆
$\{f_i  : (W_i/V) \to (W/V)\}$ があれば、$(X \times_W W_i/U)$ は
$(X/U) \times_{g, (W/V), f_i} (W_i/V)$ の一つの姿であり、族
$\{(X \times_W W_i/U) \to (X/U)\}$ は $\mathcal{C}/U$ の被覆である。
\end{remark}

\section{層の貼り合わせ}
\label{sites-section-glueing-sheaves}

\noindent
この節は Sheaves, Section \ref{sheaves-section-glueing-sheaves} の類似である。

\begin{lemma}
\label{sites-lemma-glue-maps}
\begin{slogan}
層の射は貼り合わさる。
\end{slogan}
$\mathcal{C}$ をサイトとし、$\{U_i \to U\}$ を $\mathcal{C}$ の被覆とする。
$\mathcal{F}$、$\mathcal{G}$ を $\mathcal{C}$ 上の層とする。層の射の族
$$
\varphi_i :
\mathcal{F}|_{\mathcal{C}/U_i}
\longrightarrow
\mathcal{G}|_{\mathcal{C}/U_i}
$$
が与えられ、すべての $i, j \in I$ について $\varphi_i, \varphi_j$ の制限が
同じ射
$\varphi_{ij} : \mathcal{F}|_{\mathcal{C}/U_i \times_U U_j} \to
\mathcal{G}|_{\mathcal{C}/U_i \times_U U_j}$ であるとする。このとき、
各 $\mathcal{C}/U_i$ への制限が $\varphi_i$ と一致する層の射
$$
\varphi :
\mathcal{F}|_{\mathcal{C}/U}
\longrightarrow
\mathcal{G}|_{\mathcal{C}/U}
$$
が一意に存在する。
\end{lemma}

\begin{proof}
補題で用いる制限は補題 \ref{sites-lemma-relocalize} のものである。
$V/U$ を $\mathcal{C}/U$ の対象とする。$V_i = U_i \times_U V$ とおき、
$\mathcal{V} = \{V_i \to V\}$ と書く。
$(U_i \times_U U_j) \times_U V = V_i \times_V V_j$ であることに注意しよう。
すると
$\mathcal{F}|_{\mathcal{C}/U_i}(V_i/U_i) = \mathcal{F}(V_i)$ および
$\mathcal{F}|_{\mathcal{C}/U_i \times_U U_j}(V_i \times_V V_j/U_i \times_U U_j)
= \mathcal{F}(V_i \times_V V_j)$ であり、$\mathcal{G}$ についても同様である。
したがって $V$ 上の切断に対する $\varphi$ を、図式の点線矢印として
定められる。
$$
\xymatrix{
\mathcal{F}(V) \ar@{=}[r] &
H^0(\mathcal{V}, \mathcal{F}) \ar@{..>}[d] \ar[r] &
\prod \mathcal{F}(V_i)
\ar[d]_{\prod \varphi_i}
\ar@<1ex>[r] \ar@<-1ex>[r] &
\prod \mathcal{F}(V_i \times_V V_j) \ar[d]_{\prod \varphi_{ij}} \\
\mathcal{G}(V) \ar@{=}[r] &
H^0(\mathcal{V}, \mathcal{G}) \ar[r] &
\prod \mathcal{G}(V_i)
\ar@<1ex>[r] \ar@<-1ex>[r] &
\prod \mathcal{G}(V_i \times_V V_j)
}
$$
等号は層条件から来る。記法 $H^0(\mathcal{V}, -)$ については
Section \ref{sites-section-sheafification} を参照せよ。これらの写像が制限写像と
両立することの確認は省略する。
\end{proof}

\noindent
前の補題により、サイト $\mathcal{C}$ 上の二つの層
$\mathcal{F}$、$\mathcal{G}$ が与えられたとき、規則
$$
U \longmapsto
\Mor_{\Sh(\mathcal{C}/U)}(
\mathcal{F}|_{\mathcal{C}/U}, \mathcal{G}|_{\mathcal{C}/U})
$$
は層 $\SheafHom(\mathcal{F},\mathcal{G})$ を定める。これは一種の
{\it 内部 Hom 層}である。集合の層の文脈ではほとんど用いられず、
加群の層の文脈でより一般に用いられる。Modules on Sites, Section
\ref{sites-modules-section-internal-hom} を参照せよ。

\begin{lemma}
\label{sites-lemma-internal-hom-sheaf}
\begin{slogan}
サイト上の層の圏はカルテジアン閉である。
\end{slogan}
$\mathcal{C}$ をサイトとする。$\mathcal{F}$、$\mathcal{G}$、$\mathcal{H}$ を
$\mathcal{C}$ 上の層とする。三つの変数すべてについて関手的な標準全単射
$$
\Mor_{\Sh(\mathcal{C})}(\mathcal{F}\times\mathcal{G},\mathcal{H}) =
\Mor_{\Sh(\mathcal{C})}(\mathcal{F},\SheafHom(\mathcal{G},\mathcal{H}))
$$
が存在する。
\end{lemma}

\begin{proof}
補題は、関手 $-\times\mathcal{G}$ と $\SheafHom(\mathcal{G},-)$ が互いに
随伴であることを述べている。これを示すため、単位と余単位の概念を用いる。
Categories, Section \ref{categories-section-adjoint} を参照せよ。単位
$$
\eta_\mathcal{F} :
\mathcal{F}
\longrightarrow
\SheafHom(\mathcal{G},\mathcal{F}\times\mathcal{G})
$$
は $s \in \mathcal{F}(U)$ を射
$\mathcal{G}|_{\mathcal{C}/U} \to
\mathcal{F}|_{\mathcal{C}/U}\times\mathcal{G}|_{\mathcal{C}/U}$ に送る。
これは $V/U$ 上で
$$
\mathcal{G}(V) \longrightarrow \mathcal{F}(V)\times \mathcal{G}(V), \quad
t \longmapsto (s|_{V},t).
$$
によって与えられる。余単位
$$
\epsilon_{\mathcal{H}} :
\SheafHom(\mathcal{G}, \mathcal{H}) \times \mathcal{G}
\longrightarrow
\mathcal{H}
$$
は評価写像であり、規則
$$
\Mor_{\Sh(\mathcal{C}/U)}(
\mathcal{G}|_{\mathcal{C}/U}, \mathcal{H}|_{\mathcal{C}/U})
\times \mathcal{G}(U)
\longrightarrow
\mathcal{H}(U),\quad
(\varphi, s) \longmapsto \varphi(s).
$$
によって与えられる。すると各
$\varphi : \mathcal{F} \times \mathcal{G} \to \mathcal{H}$ に対し、対応する射
$\mathcal{F} \to \SheafHom(\mathcal{G},\mathcal{H})$ は、各切断
$s \in \mathcal{F}(U)$ を $\mathcal{C}/U$ 上の層の射で、$V/U$ 上の切断に
おいて
$$
\mathcal{G}(V) \longrightarrow \mathcal{H}(V),\quad
t \longmapsto \varphi(s|_V, t).
$$
によって与えられるものへ送ることで定まる。逆に、各
$\psi : \mathcal{F} \to \SheafHom(\mathcal{G}, \mathcal{H})$ に対し、対応する射
$\mathcal{F} \times \mathcal{G} \to \mathcal{H}$ は、対象 $U$ 上の切断において
$$
\mathcal{F}(U) \times \mathcal{G}(U) \longrightarrow \mathcal{H}(U),\quad
(s, t) \longmapsto \psi(s)(t)
$$
によって与えられる。これらの構成が互いに逆であることを示す証明の詳細は
省略する。
\end{proof}

\begin{lemma}
\label{sites-lemma-hom-sheaf-hU}
$\mathcal{C}$ をサイトとし、$U \in \Ob(\mathcal{C})$ とする。
$\Sh(\mathcal{C})$ の $\mathcal{F}$ に対して
$\SheafHom(h_U^\#, \mathcal{F}) = j_*(\mathcal{F}|_{\mathcal{C}/U})$ である。
\end{lemma}

\begin{proof}
これは層の同型を直接構成して示せる。ここでは代わりに次のように論じる。
$\mathcal{G}$ を $\mathcal{C}$ 上の層とする。このとき
\begin{align*}
\Mor(\mathcal{G}, j_*(\mathcal{F}|_{\mathcal{C}/U}))
& =
\Mor(\mathcal{G}|_{\mathcal{C}/U}, \mathcal{F}|_{\mathcal{C}/U}) \\
& =
\Mor(j_!(\mathcal{G}|_{\mathcal{C}/U}), \mathcal{F}) \\
& =
\Mor(\mathcal{G} \times h_U^\#, \mathcal{F}) \\
& =
\Mor(\mathcal{G}, \SheafHom(h_U^\#, \mathcal{F}))
\end{align*}
であり、米田の補題によって結論を得る。ここでは補題
\ref{sites-lemma-internal-hom-sheaf} と
\ref{sites-lemma-compute-j-shriek-restrict} を用いた。
\end{proof}

\noindent
$\mathcal{C}$ をサイトとし、$\{U_i \to U\}_{i \in I}$ を
$\mathcal{C}$ の被覆とする。各 $i \in I$ に対し、$\mathcal{F}_i$ を
$\mathcal{C}/U_i$ 上の集合の層とする。各対 $i, j \in I$ に対し、
$$
\varphi_{ij} :
\mathcal{F}_i|_{\mathcal{C}/U_i \times_U U_j}
\longrightarrow
\mathcal{F}_j|_{\mathcal{C}/U_i \times_U U_j}
$$
を集合の層の同型とする。さらに、添字の任意の三つ組 $i, j, k \in I$ に
対して次の図式が可換であると仮定する。
$$
\xymatrix{
\mathcal{F}_i|_{\mathcal{C}/U_i \times_U U_j \times_U U_k}
\ar[rr]_{\varphi_{ik}}
\ar[rd]_{\varphi_{ij}} & &
\mathcal{F}_k|_{\mathcal{C}/U_i \times_U U_j \times_U U_k} \\
&
\mathcal{F}_j|_{\mathcal{C}/U_i \times_U U_j \times_U U_k}
\ar[ru]_{\varphi_{jk}}
}
$$
このようなデータの族 $(\mathcal{F}_i, \varphi_{ij})$ を、
{\it 被覆 $\{U_i \to U\}_{i \in I}$ に関する集合の層の貼り合わせデータ}
という。

\begin{lemma}
\label{sites-lemma-glue-sheaves}
$\mathcal{C}$ をサイトとし、$\{U_i \to U\}_{i \in I}$ を
$\mathcal{C}$ の被覆とする。被覆 $\{U_i \to U\}_{i \in I}$ に関する集合の
層の任意の貼り合わせデータ $(\mathcal{F}_i, \varphi_{ij})$ に対し、
$\mathcal{C}/U$ 上の集合の層 $\mathcal{F}$ と同型
$$
\varphi_i : \mathcal{F}|_{\mathcal{C}/U_i} \to \mathcal{F}_i
$$
で、図式
$$
\xymatrix{
\mathcal{F}|_{\mathcal{C}/U_i \times_U U_j}
\ar[d]_{\text{id}} \ar[r]_{\varphi_i} &
\mathcal{F}_i|_{\mathcal{C}/U_i \times_U U_j} \ar[d]^{\varphi_{ij}} \\
\mathcal{F}|_{\mathcal{C}/U_i \times_U U_j} \ar[r]^{\varphi_j} &
\mathcal{F}_j|_{\mathcal{C}/U_i \times_U U_j}
}
$$
を可換にするものが存在する。
\end{lemma}

\begin{proof}
$\mathcal{C}/U$ 上の層 $\mathcal{F}$ の構成を記述しよう。
$a : V \to U$ を $\mathcal{C}/U$ の対象とする。このとき
$$
\mathcal{F}(V/U) = \{
(s_i)_{i \in I} \in \prod_{i \in I} \mathcal{F}_i(U_i \times_U V/U_i)
\mid
\varphi_{ij}(s_i|_{U_i \times_U U_j \times_U V})
=
s_j|_{U_i \times_U U_j \times_U V}
\}
$$
と定める。制限写像の構成は省略する。これが層であることの確認も省略する。
同型 $\varphi_i$ の構成と、補題の図式が可換であることの証明も省略する。
\end{proof}

\noindent
$\mathcal{C}$ をサイトとし、$\{U_i \to U\}_{i \in I}$ を
$\mathcal{C}$ の被覆とする。$\mathcal{F}$ を $\mathcal{C}/U$ 上の層とする。
$\mathcal{F}$ には、その制限 $\mathcal{F}|_{\mathcal{C}/U_i}$ と標準同型
$$
\left(\mathcal{F}|_{\mathcal{C}/U_i}\right)|_{\mathcal{C}/U_i \times_U U_j}
=
\left(\mathcal{F}|_{\mathcal{C}/U_j}\right)|_{\mathcal{C}/U_i \times_U U_j}
$$
で与えられる{\it 標準貼り合わせデータ}が付随する。この同型は、関手の合成
$\mathcal{C}/U_i \times_U U_j \to \mathcal{C}/U_i \to \mathcal{C}/U$ と
$\mathcal{C}/U_i \times_U U_j \to \mathcal{C}/U_j \to \mathcal{C}/U$ が
等しいことから来る。

\begin{lemma}
\label{sites-lemma-mapping-property-glue}
$\mathcal{C}$ をサイトとし、$\{U_i \to U\}_{i \in I}$ を
$\mathcal{C}$ の被覆とする。$\mathcal{C}/U$ 上の $\mathcal{F}$ に
標準貼り合わせデータを対応させる関手によって、圏
$\Sh(\mathcal{C}/U)$ は貼り合わせデータの圏と同値である。
\end{lemma}

\begin{proof}
補題 \ref{sites-lemma-glue-maps} で関手が充満忠実であることを見た。
補題 \ref{sites-lemma-glue-sheaves} では、それが本質的全射であることを証明した
（準逆関手を明示的に構成した）。
\end{proof}

\noindent
$\mathcal{C}$ をサイトとする。サイト $\mathcal{C}$ 全体上で層を貼り合わせる
版を論じる。$\mathcal{C}$ の各対象 $U$ に対して、$\mathcal{F}_U$ を
$\mathcal{C}/U$ 上の層とする。$\mathcal{C}$ の各射 $f : V \rightarrow U$ には、
$(a : W \to V) \mapsto (f \circ a : W \to U)$ で与えられる関手
$j_f : \mathcal{C}/V \rightarrow \mathcal{C}/U$ が付随することを思い出そう。
そのような各 $f$ に対して、層の同型
$$
c_f : j_f^{-1} \mathcal{F}_U \to \mathcal{F}_V
$$
をとる。$\mathcal{C}$ の任意の二つの矢印 $f : V \to U$ と
$g : W \to V$ に対して、合成 $c_g \circ j_{g}^{-1} c_f$ が
$c_{f \circ g}$ に等しいと仮定する。このようなデータの族
$(\mathcal{F}_{U}, c_f)$ を、{\it $\mathcal{C}$ 上の集合の層の絶対貼り合わせ
データ}という。絶対貼り合わせデータの射
$(\mathcal{F}_{U}, c_f) \to (\mathcal{G}_{U}, c'_f)$ は、層の射の族
$(\varphi_U)$、$\varphi_U : \mathcal{F}_U \to \mathcal{G}_U$ で、
$\mathcal{C}$ のすべての射 $f : V \to U$ に対して図式
$$
\xymatrix{
j_f^{-1} \mathcal{F}_U \ar[r]_{c_f} \ar[d]_{j_f^{-1} \varphi_U} &
\mathcal{F}_V \ar[d]^{\varphi_V} \\
j_f^{-1} \mathcal{G}_U \ar[r]^{c'_f} &
\mathcal{G}_V
}
$$
が可換となるものによって与えられる。

\medskip\noindent
$\mathcal{C}$ 上の任意の層 $\mathcal{F}$ には、その
{\it 標準絶対貼り合わせデータ} $(\mathcal{F}|_{\mathcal{C}/U},c_f)$ が付随する。
ここで $f : V \to U$ に対する標準同型
$c_f : j_f^{-1} \mathcal{F}|_{\mathcal{C}/U} \to
\mathcal{F}|_{\mathcal{C}/V}$ は、補題 \ref{sites-lemma-relocalize} の関係
$j_V = j_U \circ j_f$ から来る。$\mathcal{C}$ の層の任意の射
$\varphi : \mathcal{F} \to \mathcal{G}$ は、標準絶対貼り合わせデータの射
$(\varphi|_{\mathcal{C}/U})$ を誘導する。

\begin{lemma}
\label{sites-lemma-glue-sheaves-absolute}
$\mathcal{C}$ をサイトとする。$\mathcal{C}$ 上の $\mathcal{F}$ に
標準絶対貼り合わせデータを対応させる関手によって、圏
$\Sh(\mathcal{C})$ は絶対貼り合わせデータの圏と同値である。
\end{lemma}

\begin{proof}
絶対貼り合わせデータ $(\mathcal{F}_U, c_f)$ が与えられたとき、
$\mathcal{F}(U) = \mathcal{F}_U(U)$ とおいて $\mathcal{C}$ 上の層
$\mathcal{F}$ を構成する。ここで
% P01-E0096
$f : V \to U$ に沿う制限は可換図式
$$
\xymatrix{
\mathcal{F}_U(U) \ar[r] \ar@{=}[d] &
\mathcal{F}_U(V) \ar[r]^{c_f} &
\mathcal{F}_V(V) \ar@{=}[d] \\
\mathcal{F}(U) \ar[rr] &  &
\mathcal{F}(V)
}
$$
によって与えられる。両立条件
$c_g \circ j_{g}^{-1} c_f = c_{f \circ g}$ により、$\mathcal{F}$ は前層となる。
またこの条件により、写像 $c_f : \mathcal{F}_U(V) \to \mathcal{F}(V)$ は
各 $U$ に対して同型 $\mathcal{F}_U \to \mathcal{F}|_{\mathcal{C}/U}$ を
定める。各 $\mathcal{F}_U$ は層なので、これにより $\mathcal{F}$ も層である。
ここで構成した関手 $(\mathcal{F}_U, c_f) \mapsto \mathcal{F}$ は、
$\mathcal{C}$ 上の層をその標準貼り合わせデータへ送る関手の準逆である。
これ以上の詳細は省略する。
\end{proof}

\begin{remark}
\label{sites-remark-variant-glue-sheaves-absolute}
補題 \ref{sites-lemma-glue-sheaves-absolute} には、代数幾何で現れる変種がある。
すなわち、$\mathcal{C}$ がすべてのファイバー積をもつサイトであり、
各 $U \in \Ob(\mathcal{C})$ に対して次の性質をもつ充満部分圏
$U_\tau \subset \mathcal{C}/U$ が与えられているとする。
\begin{enumerate}
\item $U/U$ は $U_\tau$ に属する。
\item $U_\tau$ の $V/U$ と $\mathcal{C}$ の被覆 $\{V_j \to V\}$ に対して、
$V_j/U$ は $U_\tau$ に属する。
\item $\mathcal{C}$ の射 $U' \to U$ と $U_\tau$ の $V/U$ に対して、
基底変換 $V' = V \times_U U'$ は $U'_\tau$ に属する。
\end{enumerate}
この状況では $\mathcal{C}$ のすべての $U$ に対して $U_\tau$ はサイトであり、
基底変換関手 $U_\tau \to U'_\tau$ は、$\mathcal{C}$ のすべての射
$f : U' \to U$ に対してサイトの射 $f_\tau : U_\tau \to U'_\tau$ を定める。
このとき得られる貼り合わせの主張は次のようになる。$\mathcal{C}$ 上の層
$\mathcal{F}$ は次のデータによって与えられる。
\begin{enumerate}
\item 各 $U \in \Ob(\mathcal{C})$ に対し、$U_\tau$ 上の層
$\mathcal{F}_U$。
\item $\mathcal{C}$ の各 $f : U' \to U$ に対し、写像
$c_f : f_\tau^{-1}\mathcal{F}_U \to \mathcal{F}_{U'}$。
\end{enumerate}
これらのデータは次の条件を満たす。
\begin{enumerate}
\item[(a)] $\mathcal{C}$ の $f : U' \to U$ と $g : U'' \to U'$ が与えられたとき、
合成 $c_g \circ g_\tau^{-1}c_f$ は $c_{f \circ g}$ に等しい。
\item[(b)] $f : U' \to U$ が $U_\tau$ に属するならば、$c_f$ は同型である。
\end{enumerate}
これが必要になれば、ここで正確に主張して証明する。
（この結果は、すべての写像 $c_f$ が同型であることを要求していないため、
上の主張とはわずかに異なることに注意せよ！）
\end{remark}

\section{局所化についての補足}
\label{sites-section-more-localize}

\noindent
この節では、
% P01-E0097
サイトまたは局所化する対象にいくつかの追加仮定を課した局所化について、
幾つかの補題を証明する。

\begin{lemma}
\label{sites-lemma-describe-j-shriek-good-site}
$\mathcal{C}$ をサイトとし、$U \in \Ob(\mathcal{C})$ とする。
$\mathcal{C}$ 上の位相が準標準的であり
（定義 \ref{sites-definition-weaker-than-canonical} を参照）、
$\mathcal{G}$ が $\mathcal{C}/U$ 上の層ならば、
$$
j_{U!}(\mathcal{G})(V)
=
\coprod\nolimits_{\varphi \in \Mor_\mathcal{C}(V, U)}
\mathcal{G}(V \xrightarrow{\varphi} U),
$$
である。言い換えると、補題 \ref{sites-lemma-describe-j-shriek} における層化は
不要である。
\end{lemma}

\begin{proof}
$\mathcal{V} = \{V_i \to V\}_{i \in I}$ をサイト $\mathcal{C}$ における
$V$ の被覆とする。補題 \ref{sites-lemma-describe-j-shriek} の前層
$\mathcal{H}$ に対する層条件を直接確認する。
$(s_i, \varphi_i)_{i \in I} \in \prod_i \mathcal{H}(V_i)$ とする。
% P01-E0098
これは、$\varphi_i : V_i \to U$ が $\mathcal{C}$ の射であり、
$s_i \in \mathcal{G}(V_i \xrightarrow{\varphi_i} U)$ であることを意味する。
対 $(s_i, \varphi_i)$ の $V_i \times_V V_j$ への制限は
$(s_i|_{V_i \times_V V_j/U}, \text{pr}_1 \circ \varphi_i)$ であり、同様に
対 $(s_j, \varphi_j)$ の $V_i \times_V V_j$ への制限は
$(s_j|_{V_i \times_V V_j/U}, \text{pr}_2 \circ \varphi_j)$ である。
したがって、族 $(s_i, \varphi_i)$ が
$\check{H}^0(\mathcal{V}, \mathcal{H})$ に属するならば、
$\text{pr}_1 \circ \varphi_i = \text{pr}_2 \circ \varphi_j$ である。
$\mathcal{C}$ 上の位相が標準位相より弱いという条件により、一意な射
$\varphi : V \to U$ で、$\varphi_i$ が $V_i \to V$ と $\varphi$ の合成に
なるものが存在する。この時点で $\mathcal{G}$ の層条件により、切断 $s_i$ は
一意な切断 $s \in \mathcal{G}(V \xrightarrow{\varphi} U)$ に貼り合わさる。
ゆえに望むとおり $(s, \varphi) \in \mathcal{H}(V)$ である。
\end{proof}

\begin{lemma}
\label{sites-lemma-localize-given-products}
$\mathcal{C}$ をサイトとし、$U \in \Ob(\mathcal{C})$ とする。
$\mathcal{C}$ は対象の対の積をもつと仮定する。このとき、
\begin{enumerate}
\item 関手 $j_U$ は連続な右随伴、すなわち関手
$v(X) = X \times U / U$ をもつ。
\item 関手 $v$ はサイトの射 $\mathcal{C}/U \to \mathcal{C}$ を定め、
それに対応するトポスの射は
$j_U : \Sh(\mathcal{C}/U) \to \Sh(\mathcal{C})$ に等しい。
\item $j_{U*}\mathcal{F}(X) = \mathcal{F}(X \times U/U)$ である。
\end{enumerate}
\end{lemma}

\begin{proof}
関手 $v$ が $j_U$ の右随伴であるとは、$Y/U$ と $X$ が与えられたとき
$$
\Mor_\mathcal{C}(Y, X)
=
\Mor_{\mathcal{C}/U}(Y/U, X \times U/U)
$$
であることを意味し、これは明らかである。$v$ の連続性を確認するため、
$\{X_i \to X\}$ を $\mathcal{C}$ の被覆とする。サイトの第三公理
（定義 \ref{sites-definition-site}）により、
$$
\{X_i \times_X (X \times U) \to X \times_X (X \times U)\}
=
\{X_i \times U \to X \times U\}
$$
も $\mathcal{C}$ の被覆である。したがって $v$ は連続である。
補題の他の主張は補題 \ref{sites-lemma-have-functor-other-way} と
\ref{sites-lemma-have-functor-other-way-morphism} から従う。
\end{proof}

\begin{lemma}
\label{sites-lemma-relocalize-given-fibre-products}
$\mathcal{C}$ をサイトとし、$U \to V$ を $\mathcal{C}$ の射とする。
$\mathcal{C}$ はファイバー積をもつと仮定する。$j$ を補題
\ref{sites-lemma-relocalize} のものとする。このとき、
\begin{enumerate}
\item 関手 $j : \mathcal{C}/U \to \mathcal{C}/V$ は連続な右随伴、
すなわち関手 $v : (X/V) \mapsto (X \times_V U/U)$ をもつ。
\item 関手 $v$ はサイトの射 $\mathcal{C}/U \to \mathcal{C}/V$ を定め、
それに対応するトポスの射は $j$ に等しい。
\item $j_*\mathcal{F}(X/V) = \mathcal{F}(X \times_V U/U)$ である。
\end{enumerate}
\end{lemma}

\begin{proof}
補題 \ref{sites-lemma-localize-given-products} から従う。実際、補題
\ref{sites-lemma-relocalize} により $j$ は局所化関手とみなせる。
\end{proof}

\noindent
開埋め込みの基本的性質は、順像の制限と空集合による延長の制限が
元の層を返すことである。
% P01-E0099
上の $j_U$ に対応する関手について、これは常に成り立つわけではない。
$U$ が「終対象の部分対象」であるときには成り立つ。

\begin{lemma}
\label{sites-lemma-restrict-back}
$\mathcal{C}$ をサイトとし、$U \in \Ob(\mathcal{C})$ とする。
$\mathcal{C}$ のすべての $X$ から $U$ への射は高々一つであると仮定する。
$\mathcal{F}$ を $\mathcal{C}/U$ 上の層とする。標準写像
$\mathcal{F} \to j_U^{-1}j_{U!}\mathcal{F}$ と
$j_U^{-1}j_{U*}\mathcal{F} \to \mathcal{F}$ は同型である。
\end{lemma}

\begin{proof}
これは補題 \ref{sites-lemma-back-and-forth} の特別な場合である。実際、$U$ に
関する仮定は、局所化関手 $\mathcal{C}/U \to \mathcal{C}$ の
% P01-E0100
充満忠実性と同値である。
\end{proof}

\begin{lemma}
\label{sites-lemma-localize-cartesian-square}
$\mathcal{C}$ をサイトとする。$\mathcal{C}$ の可換図式
$$
\xymatrix{
U' \ar[d] \ar[r] & U \ar[d] \\
V' \ar[r] & V
}
$$
をとる。補題 \ref{sites-lemma-relocalize} の射は、連続かつ余連続な関手の図式と
トポスの図式
$$
\vcenter{
\xymatrix{
\mathcal{C}/U' \ar[d]_{j_{U'/V'}} \ar[r]_{j_{U'/U}} &
\mathcal{C}/U \ar[d]^{j_{U/V}} \\
\mathcal{C}/V' \ar[r]^{j_{V'/V}} & \mathcal{C}/V
}
}
\quad\text{and}\quad
\vcenter{
\xymatrix{
\Sh(\mathcal{C}/U') \ar[d]_{j_{U'/V'}} \ar[r]_{j_{U'/U}} &
\Sh(\mathcal{C}/U) \ar[d]^{j_{U/V}} \\
\Sh(\mathcal{C}/V') \ar[r]^{j_{V'/V}} &
\Sh(\mathcal{C}/V)
}
}
$$
を可換にする。さらに、$\mathcal{C}$ の最初の図式がカルテジアンならば、
$j_{V'/V}^{-1} \circ j_{U/V, *} = j_{U'/V', *} \circ j_{U'/U}^{-1}$ である。
\end{lemma}

\begin{proof}
補題の最初の主張にある左の正方形の可換性は定義から直ちに従う。
補題 \ref{sites-lemma-composition-cocontinuous} により、これはトポスの図式の
可換性を導く。図式がカルテジアンであると仮定する。随伴関手の一意性により、
$j_{V'/V}^{-1} \circ j_{U/V, *} = j_{U'/V', *} \circ j_{U'/U}^{-1}$ を
示すことは、
$j_{U/V}^{-1} \circ j_{V'/V!} = j_{U'/U!} \circ j_{U'/V'}^{-1}$ を
示すことと同値である。補題 \ref{sites-lemma-essential-image-j-shriek} の同一視により、
トポスの図式を
$$
\xymatrix{
\Sh(\mathcal{C})/h_{U'}^\# \ar[d] \ar[r] &
\Sh(\mathcal{C})/h_U^\# \ar[d] \\
\Sh(\mathcal{C})/h_{V'}^\# \ar[r] &
\Sh(\mathcal{C})/h_V^\#
}
$$
と考えられる。また補題 \ref{sites-lemma-relocalize-explicit} により、関手
$j^{-1}$ と $j_!$ の解釈を知っている。したがって、
$\mathcal{F} \to h_{V'}^\#$ が与えられたとき、$h_U^\#$ への写像をもつ層として
$$
\mathcal{F} \times_{h_{V'}^\#} h_{U'}^\# =
\mathcal{F} \times_{h_V^\#} h_U^\#
$$
であることを示す必要がある。これは
$h_{U'} = h_{V'} \times_{h_V} h_U$ であり、層化が完全なので
$$
h_{U'}^\# = h_{V'}^\# \times_{h_V^\#} h_U^\#
$$
でもあることから従う。
\end{proof}

\section{局所化と射}
\label{sites-section-localize-morphisms}

\noindent
次の補題は、
% P01-E0101
局所化とサイトおよびトポスの射との関係を理解するために重要である。

\begin{lemma}
\label{sites-lemma-localize-morphism}
$f : \mathcal{C} \to \mathcal{D}$ を、連続関手
$u : \mathcal{D} \to \mathcal{C}$ に対応するサイトの射とする。
$V \in \Ob(\mathcal{D})$ とし、$U = u(V)$ とおく。このとき関手
$u' : \mathcal{D}/V \to \mathcal{C}/U$、
$V'/V \mapsto u(V')/U$ はサイトの射
$f' : \mathcal{C}/U \to \mathcal{D}/V$ を定める。射 $f'$ はトポスの可換図式
$$
\xymatrix{
\Sh(\mathcal{C}/U) \ar[r]_{j_U} \ar[d]_{f'} &
\Sh(\mathcal{C}) \ar[d]^f \\
\Sh(\mathcal{D}/V) \ar[r]^{j_V} &
\Sh(\mathcal{D}).
}
$$
に入る。補題 \ref{sites-lemma-essential-image-j-shriek} の同一視
$\Sh(\mathcal{C}/U) = \Sh(\mathcal{C})/h_U^\#$ および
$\Sh(\mathcal{D}/V) = \Sh(\mathcal{D})/h_V^\#$ を用いると、関手
$(f')^{-1}$ は規則
$$
(f')^{-1}(\mathcal{H} \xrightarrow{\varphi} h_V^\#)
=
(f^{-1}\mathcal{H} \xrightarrow{f^{-1}\varphi} h_U^\#).
$$
で記述される。最後に、$f'_*j_U^{-1} = j_V^{-1}f_*$ である。
\end{lemma}

\begin{proof}
$u'$ が連続であることは明らかである。したがって関手
$f'_* = (u')^s = (u')^p$（Sections \ref{sites-section-functoriality-PSh} と
\ref{sites-section-continuous-functors} を参照）と、その随伴
$(f')^{-1} = (u')_s = ((u')_p\ )^\#$ を得る。主張
$f'_*j_U^{-1} = j_V^{-1}f_*$ は
$$
(j_V^{-1}f_*\mathcal{F})(V'/V)
= f_*\mathcal{F}(V') = \mathcal{F}(u(V'))
= (j_U^{-1}\mathcal{F})(u(V')/U)
= (f'_*j_U^{-1}\mathcal{F})(V'/V)
$$
から従い、これは前層についてさえ成り立つ。先験的に明らかでないのは、
$(f')^{-1}$ が完全であること、図式が可換であること、および
$(f')^{-1}$ の記述が成り立つことである。

\medskip\noindent
$\mathcal{H}$ を $\mathcal{D}/V$ 上の層とする。
$j_{U!}(f')^{-1}\mathcal{H}$ を計算しよう。
\begin{eqnarray*}
j_{U!}(f')^{-1}\mathcal{H}
& =
((j_U)_p(u'_p\mathcal{H})^\#)^\# \\
& =
((j_U)_pu'_p\mathcal{H})^\# \\
& =
(u_p(j_V)_p\mathcal{H})^\# \\
& =
f^{-1}j_{V!}\mathcal{H}
\end{eqnarray*}
% P01-E0102
最初の等式は定義を展開することによる。二つ目の等式は補題
\ref{sites-lemma-technical-up} による。三つ目の等式は
$u \circ j_V = j_U \circ u'$ による。四つ目の等式は再び補題
\ref{sites-lemma-technical-up} による。上の等式はすべて関手の同型であり、
したがって圏と関手の次の図式が可換であると解釈できる。
$$
\xymatrix{
\Sh(\mathcal{C}/U) \ar[r]_{j_{U!}} &
\Sh(\mathcal{C})/h_U^\# \ar[r] &
\Sh(\mathcal{C}) \\
\Sh(\mathcal{D}/V) \ar[r]^{j_{V!}} \ar[u]^{(f')^{-1}} &
\Sh(\mathcal{D})/h_V^\# \ar[r] \ar[u]^{f^{-1}} &
\Sh(\mathcal{D}) \ar[u]^{f^{-1}}
}
$$
中央の矢印が意味をもつのは
$f^{-1}h_V^\# = (h_{u(V)})^\# = h_U^\#$ だからである。補題
\ref{sites-lemma-pullback-representable-sheaf} を参照せよ。特に、これは補題の主張に
ある $(f')^{-1}$ の記述を証明する。補題
\ref{sites-lemma-essential-image-j-shriek} により左の水平矢印は同値であり、仮定に
より $f^{-1}$ は完全なので、$(f')^{-1} = u'_s$ は完全であると結論する。
実際、これは左随伴なので既に右完全である
（Categories, Lemma \ref{categories-lemma-adjoint-exact}）。したがって、
終対象を終対象へ移し、ファイバー積と可換することだけを示せばよい
（Categories, Lemma \ref{categories-lemma-characterize-left-exact}）。
どちらも誘導関手
$f^{-1} : \Sh(\mathcal{D})/h_V^\# \to \Sh(\mathcal{C})/h_U^\#$ について
明らかである。これにより $f'$ がサイトの射であることが証明された。

\medskip\noindent
なお $(f')^{-1}j_V^{-1} = j_U^{-1}f^{-1}$ を確認しなければならない。
そのために上の公式と補題 \ref{sites-lemma-compute-j-shriek-restrict} の記述を用いる。
これらを合わせると、$\mathcal{D}$ 上の任意の層 $\mathcal{G}$ に対して
$$
j_{U!}(f')^{-1}j_V^{-1}\mathcal{G}
=
f^{-1}j_{V!}j_V^{-1}\mathcal{G}
=
f^{-1}(\mathcal{G} \times h_V^\#)
=
f^{-1}\mathcal{G} \times h_U^\#
=
j_{U!}j_U^{-1}f^{-1}\mathcal{G}.
$$
である。関手 $j_{U!}$ は同値
$\Sh(\mathcal{C}/U) \to \Sh(\mathcal{C})/h_U^\#$ を誘導するので、結論を得る。
\end{proof}

\noindent
次の補題は、上のより一般的な補題 \ref{sites-lemma-localize-morphism} の特別な
場合である。

\begin{lemma}
\label{sites-lemma-localize-morphism-strong}
$\mathcal{C}$、$\mathcal{D}$ をサイトとする。
$u : \mathcal{D} \to \mathcal{C}$ を関手とする。
$V \in \Ob(\mathcal{D})$ とし、$U = u(V)$ とおく。次を仮定する。
\begin{enumerate}
\item $\mathcal{C}$ と $\mathcal{D}$ はすべての有限極限をもつ。
\item $u$ は連続である。
\item $u$ は有限極限と可換する。
\end{enumerate}
サイトの射の可換図式
$$
\xymatrix{
\mathcal{C}/U \ar[r]_{j_U} \ar[d]_{f'} & \mathcal{C} \ar[d]^f \\
\mathcal{D}/V \ar[r]^{j_V} & \mathcal{D}
}
$$
が存在する。ここで右の縦矢印は $u$ に対応し、左の縦矢印は関手
$u' : \mathcal{D}/V \to \mathcal{C}/U$、$V'/V \mapsto u(V')/u(V)$ に対応し、
% P01-E0103
水平矢印は、補題 \ref{sites-lemma-localize-given-products} のように関手
$\mathcal{C} \to \mathcal{C}/U$、$X \mapsto X \times U$ および
$\mathcal{D} \to \mathcal{D}/V$、$Y \mapsto Y \times V$ に対応する。
さらに、対応するトポスの射の図式は補題
\ref{sites-lemma-localize-morphism} の図式に等しい。特に
$f'_*j_U^{-1} = j_V^{-1}f_*$ である。
\end{lemma}

\begin{proof}
$u$ は命題 \ref{sites-proposition-get-morphism} の仮定を満たし、したがってその
命題によりサイトの射 $f : \mathcal{C} \to \mathcal{D}$ を誘導する。
$u$ が示された関手 $u'$ を誘導することは明らかである。この関手も命題
\ref{sites-proposition-get-morphism} の仮定を満たすことは明らかである。
したがってサイトの射 $f' : \mathcal{C}/U \to \mathcal{D}/V$ を得る。
図式は、サイトの射の合成の定義
（定義 \ref{sites-definition-composition-morphisms-sites} を参照）と等式
$$
u(Y \times V) = u(Y) \times u(V) = u(Y) \times U
$$
により可換である。この等式は、補題の図式と反対向きの圏と関手の図式が
可換であることを示す。
\end{proof}

\noindent
この時点で、サイトを局所化すること、再局所化すること、および下側の
サイトの対象においてサイトの射を局所化することができる。これらを
組み合わせると、次の種類の図式を得る。

\begin{lemma}
\label{sites-lemma-relocalize-morphism}
$f : \mathcal{C} \to \mathcal{D}$ を、連続関手
$u : \mathcal{D} \to \mathcal{C}$ に対応するサイトの射とする。
$V \in \Ob(\mathcal{D})$、$U \in \Ob(\mathcal{C})$ とし、
$c : U \to u(V)$ を $\mathcal{C}$ の射とする。トポスの可換図式
$$
\xymatrix{
\Sh(\mathcal{C}/U) \ar[r]_{j_U} \ar[d]_{f_c} &
\Sh(\mathcal{C}) \ar[d]^f \\
\Sh(\mathcal{D}/V) \ar[r]^{j_V} &
\Sh(\mathcal{D}).
}
$$
が存在する。$f_c = f' \circ j_{U/u(V)}$ である。ここで
$f' : \Sh(\mathcal{C}/u(V)) \to \Sh(\mathcal{D}/V)$ は補題
\ref{sites-lemma-localize-morphism} のもの、
$j_{U/u(V)} : \Sh(\mathcal{C}/U) \to \Sh(\mathcal{C}/u(V))$ は補題
\ref{sites-lemma-relocalize} のものである。補題
\ref{sites-lemma-essential-image-j-shriek} の同一視
$\Sh(\mathcal{C}/U) = \Sh(\mathcal{C})/h_U^\#$ および
$\Sh(\mathcal{D}/V) = \Sh(\mathcal{D})/h_V^\#$ を用いると、関手
$(f_c)^{-1}$ は規則
$$
(f_c)^{-1}(\mathcal{H} \xrightarrow{\varphi} h_V^\#)
=
(f^{-1}\mathcal{H} \times_{f^{-1}\varphi, h_{u(V)}^\#, c} h_U^\#
\rightarrow h_U^\#).
$$
で記述される。最後に、射 $b : V' \to V$、$a : U' \to U$、
$c' : U' \to u(V')$ が与えられ、図式
$$
\xymatrix{
U' \ar[r]_-{c'} \ar[d]_a & u(V') \ar[d]^{u(b)} \\
U \ar[r]^-c & u(V)
}
$$
が可換ならば、図式
$$
\xymatrix{
\Sh(\mathcal{C}/U') \ar[r]_{j_{U'/U}} \ar[d]_{f_{c'}} &
\Sh(\mathcal{C}/U) \ar[d]^{f_c} \\
\Sh(\mathcal{D}/V') \ar[r]^{j_{V'/V}} &
\Sh(\mathcal{D}/V).
}
$$
は可換である。
\end{lemma}

\begin{proof}
この補題はそれ自体を証明しており、理論の展開のこの段階で既に知っている
事柄を集めたものに近い。例えば、最初の正方形の可換性は
Diagram (\ref{sites-equation-relocalize}) の可換性と、補題
\ref{sites-lemma-localize-morphism} の図式の可換性から従う。
$f_c^{-1}$ の記述は、補題 \ref{sites-lemma-relocalize-explicit} と補題
\ref{sites-lemma-localize-morphism} を組み合わせることから従う。
最後の正方形の可換性は次の等式から従う。
% P01-E0104
$$
f^{-1}\mathcal{H} \times_{h_{u(V)}^\#, c} h_U^\# \times_{h_U^\#} h_{U'}^\#
=
f^{-1}(\mathcal{H} \times_{h_V^\#} h_{V'}^\#)
\times_{h_{u(V'), c'}^\#} h_{U'}^\#
$$
これは $f^{-1}h_V^\# = h_{u(V)}^\#$ および
$f^{-1}h_{V'}^\# = h_{u(V')}^\#$ を用いれば形式的である。補題
\ref{sites-lemma-pullback-representable-sheaf} を参照せよ。
\end{proof}

\noindent
次の補題では、トポスの射が余連続関手から来る場合に、局所化の別の種類の
関手性を得る。これは補題 \ref{sites-lemma-localize-morphism} の図式とは異なる
種類の図式であり、特に一般には、補題 \ref{sites-lemma-localize-morphism} で見た
等式 $f'_*j_U^{-1} = j_V^{-1}f_*$ は次の補題の状況では成り立たない。

\begin{lemma}
\label{sites-lemma-localize-cocontinuous}
$\mathcal{C}$、$\mathcal{D}$ をサイトとする。
$u : \mathcal{C} \to \mathcal{D}$ を余連続関手とする。
$U$ を $\mathcal{C}$ の対象とし、$V = u(U)$ とおく。可換図式
$$
\xymatrix{
\mathcal{C}/U \ar[r]_{j_U} \ar[d]_{u'} & \mathcal{C} \ar[d]^u \\
\mathcal{D}/V \ar[r]^-{j_V} & \mathcal{D}
}
$$
があり、左の縦矢印は
% P01-E0105
$u' : \mathcal{C}/U \to \mathcal{D}/V$、$U'/U \mapsto V'/V$ である。
このとき $u'$ も余連続であり、トポスの可換図式
$$
\xymatrix{
\Sh(\mathcal{C}/U) \ar[r]_{j_U} \ar[d]_{f'} &
\Sh(\mathcal{C}) \ar[d]^f \\
\Sh(\mathcal{D}/V) \ar[r]^-{j_V} &
\Sh(\mathcal{D})
}
$$
を得る。ここで $f$（resp.\ $f'$）は $u$（resp.\ $u'$）に対応する。
\end{lemma}

\begin{proof}
最初の図式の可換性は明らかである。$u'$ が余連続であることを示せば、
これは二つ目の図式の可換性を導く。

\medskip\noindent
$U'/U$ を $\mathcal{C}/U$ の対象とする。
$\{V_j/V \to u(U')/V\}_{j \in J}$ を $\mathcal{D}/V$ における
$u(U')/V$ の被覆とする。$u$ は余連続なので、被覆
$\{U_i' \to U'\}_{i \in I}$ で、族 $\{u(U_i') \to u(U')\}$ が
$\mathcal{D}$ の被覆 $\{V_j \to u(U')\}$ を細分するものが存在する。
言い換えると、添字集合の写像 $\alpha : I \to J$ と、
% P01-E0106
$U'$ 上の射 $\phi_i : u(U_i') \to V_{\alpha(i)}$ が存在する。
$U_i'$ を合成 $U'_i \to U' \to U$ によって $U$ 上の対象とみなす。
すると $\{U'_i/U \to U'/U\}$ は $\mathcal{C}/U$ の被覆で、
$\{u(U_i')/V \to u(U')/V\}$ が $\{V_j/V \to u(U')/V\}$ を細分する
ものとなる（同じ $\alpha$ と同じ写像 $\phi_i$ を用いる）。したがって
$u' : \mathcal{C}/U \to \mathcal{D}/V$ は余連続である。
\end{proof}

\begin{lemma}
\label{sites-lemma-localize-cocontinuous-downstairs}
$\mathcal{C}$、$\mathcal{D}$ をサイトとする。
$u : \mathcal{C} \to \mathcal{D}$ を余連続関手とする。
$V$ を $\mathcal{D}$ の対象とする。${}^u_V\mathcal{I}$ を Section
\ref{sites-section-more-functoriality-PSh} で導入した圏とする。可換図式
$$
\vcenter{
\xymatrix{
\,_V^u\mathcal{I} \ar[r]_j \ar[d]_{u'} &
\mathcal{C} \ar[d]^u \\
\mathcal{D}/V \ar[r]^-{j_V} &
\mathcal{D}
}
}
\quad\text{where}\quad
\begin{matrix}
j : (U, \psi) \mapsto U \\
u' : (U, \psi) \mapsto (\psi : u(U) \to V)
\end{matrix}
$$
がある。${}^u_V\mathcal{I}$ の射の族
$\{(U_i, \psi_i) \to (U, \psi)\}$ が被覆であることを、
$\{U_i \to U\}$ が $\mathcal{C}$ の被覆であることと同値であると定める。
このとき、
\begin{enumerate}
\item ${}^u_V\mathcal{I}$ はサイトである。
\item $j$ は連続かつ余連続である。
\item $u'$ は余連続である。
\item トポスの可換図式
$$
\xymatrix{
\Sh({}^u_V\mathcal{I}) \ar[r]_j \ar[d]_{f'} &
\Sh(\mathcal{C}) \ar[d]^f \\
\Sh(\mathcal{D}/V) \ar[r]^-{j_V} &
\Sh(\mathcal{D})
}
$$
を得る。ここで $f$（resp.\ $f'$）は $u$（resp.\ $u'$）に対応する。
\item $f'_*j^{-1} = j_V^{-1}f_*$ である。
\end{enumerate}
\end{lemma}

\begin{proof}
(1)、(2)、(3)、(4) は定義と、関手 $j$ がファイバー積と可換することの
直接の帰結である。詳細は省略する。(5) を見るため、$f_*$ は
${}_su = {}_pu$ で与えられることを思い出そう。したがって
$V'/V$ における $j_V^{-1}f_*\mathcal{F}$ の値は、$V'$ における
${}_pu\mathcal{F}$ の値、すなわち圏 ${}^u_{V'}\mathcal{I}$ 上の
$\mathcal{F}$ の値の極限である。明らかに圏の同値
$$
{}^u_{V'}\mathcal{I} \to {}^{u'}_{V'/V}\mathcal{I}
$$
がある。$V'/V$ における $f'_*j^{-1}\mathcal{F}$ の値は、圏
${}^{u'}_{V'/V}\mathcal{I}$ 上の $j^{-1}\mathcal{F}$ の値の極限で与えられ、
${}^u_V\mathcal{I}$ の対象における $j^{-1}\mathcal{F}$ の値は
$\mathcal{F}$ の値にほかならない（$j$ は連続かつ余連続なので、補題
\ref{sites-lemma-when-shriek} による）。したがって (5) は成り立つ。
\end{proof}

\noindent
次の二つの結果はやや異なる性質をもつ。

\begin{lemma}
\label{sites-lemma-special-square-cocontinuous}
サイト $\mathcal{C}', \mathcal{C}, \mathcal{D}', \mathcal{D}$ と関手
$$
\xymatrix{
\mathcal{C}' \ar[r]_{v'} \ar[d]_{u'} &
\mathcal{C} \ar[d]^u \\
\mathcal{D}' \ar[r]^v &
\mathcal{D}
}
$$
が与えられたとする。次を仮定する。
\begin{enumerate}
\item $u$、$u'$、$v$、$v'$ は余連続であり、補題
\ref{sites-lemma-cocontinuous-morphism-topoi} によりトポスの射
$f$、$f'$、$g$、$g'$ を与える。
\item $v \circ u' = u \circ v'$ である。
\item $v$ と $v'$ は余連続であると同時に連続である。
\item $\mathcal{D}'$ の任意の対象 $V'$ に対し、$v$ が与える関手
${}^{u'}_{V'}\mathcal{I} \to {}^{\ \ \ u}_{v(V')}\mathcal{I}$ は共終である。
\end{enumerate}
このとき $f'_* \circ (g')^{-1} = g^{-1} \circ f_*$ および
$g'_! \circ (f')^{-1} = f^{-1} \circ g_!$ である。
\end{lemma}

\begin{proof}
圏 ${}^{u'}_{V'}\mathcal{I}$ と ${}^{\ \ \ u}_{v(V')}\mathcal{I}$ は
Section \ref{sites-section-more-functoriality-PSh} で定義された。条件 (4) の関手は、
${}^{u'}_{V'}\mathcal{I}$ の対象 $\psi : u'(U') \to V'$ を
${}^{\ \ \ u}_{v(V')}\mathcal{I}$ の対象
$v(\psi) : uv'(U') = vu'(U') \to v(V')$ に送る。
$g^{-1}$ は $v^p$ で与えられ（補題 \ref{sites-lemma-when-shriek}）、
$f_*$ は ${}_su = {}_pu$ で与えられることを思い出そう。したがって
$V'$ における $g^{-1}f_*\mathcal{F}$ の値は $v(V')$ における
${}_pu\mathcal{F}$ の値、すなわち極限
$$
\lim_{u(U) \to v(V') \in \Ob({}^{\ \ \ u}_{v(V')}\mathcal{I}^{opp})}
\mathcal{F}(U)
$$
である。同じ議論により、$V'$ における
$f'_*(g')^{-1}\mathcal{F}$ の値は極限
$$
\lim_{u'(U') \to V' \in \Ob({}^{u'}_{V'}\mathcal{I}^{opp})} \mathcal{F}(v'(U'))
$$
で与えられる。したがって仮定 (4) と Categories, Lemma
\ref{categories-lemma-initial} により、これらは一致し、補題の最初の等式が
証明される。二つ目の等式は、随伴の一意性により最初の等式から従う。
\end{proof}

\begin{lemma}
\label{sites-lemma-special-square-continuous}
サイト $\mathcal{C}', \mathcal{C}, \mathcal{D}', \mathcal{D}$ と関手
$$
\xymatrix{
\mathcal{C}' \ar[r]_{v'} &
\mathcal{C} \\
\mathcal{D}' \ar[r]^v \ar[u]^{u'} &
\mathcal{D} \ar[u]_u
}
$$
が与えられたとする。Sections \ref{sites-section-morphism-sites} と
\ref{sites-section-cocontinuous-morphism-topoi} の記法のもとで、次を仮定する。
\begin{enumerate}
\item $u$ と $u'$ は連続であり、サイトの射 $f$ と $f'$ を与える。
\item $v$ と $v'$ は余連続であり、トポスの射 $g$ と $g'$ を与える。
\item $u \circ v = v' \circ u'$ である。
\item $v$ と $v'$ は余連続であると同時に連続である。
\end{enumerate}
このとき\footnote{この一般性では、トポスの射として
$f \circ g'$ が $g \circ f'$ に等しいかどうかは分からない
（前者から後者への標準的な $2$-矢印があるが、同型とは限らない）。}
$f'_* \circ (g')^{-1} = g^{-1} \circ f_*$ および
$g'_! \circ (f')^{-1} = f^{-1} \circ g_!$ である。
\end{lemma}

\begin{proof}
実際、
$$
f'_*(g')^{-1}\mathcal{F} = (u')^p((v')^p\mathcal{F})^\# =
(u')^p(v')^p\mathcal{F}
$$
である。
% P01-E0107
最初の等式は定義により、二つ目は補題 \ref{sites-lemma-when-shriek} による。また
$$
g^{-1}f_*\mathcal{F} = (v^pu^p\mathcal{F})^\# =
((u')^p(v')^p\mathcal{F})^\# = (u')^p(v')^p\mathcal{F}
$$
である。最初の等式は定義により、二つ目は
$u \circ v = v' \circ u'$ により、三つ目は
$(u')^p(v')^p\mathcal{F}$ が層であることを既に見たことによる。これにより
$f'_* \circ (g')^{-1} = g^{-1} \circ f_*$ が証明され、等式
$g'_! \circ (f')^{-1} = f^{-1} \circ g_!$ は左随伴の一意性から従う。
\end{proof}

\section{トポスの射}
\label{sites-section-morphisms-topoi}

\noindent
本節では、任意のトポスの射が、サイトの射から生じるトポスの射と
同値であることを示す。\cite[Expos\'e III, Proposition 4.9.4]{SGA4}
と比較されたい。

\begin{lemma}
\label{sites-lemma-equivalence}
$\mathcal{C}$、$\mathcal{D}$ をサイトとする。
$u : \mathcal{C} \to \mathcal{D}$ を関手とする。次を仮定する。
\begin{enumerate}
\item $u$ は余連続である。
\item $u$ は連続である。
\item $\mathcal{C}$ における $a, b : U' \to U$ が
$u(a) = u(b)$ を満たすとき、$\mathcal{C}$ に被覆
$\{f_i : U'_i \to U'\}$ が存在して
$a \circ f_i = b \circ f_i$ となる。
\item $U', U \in \Ob(\mathcal{C})$ と
$\mathcal{D}$ における射 $c : u(U') \to u(U)$ が与えられると、
$\mathcal{C}$ に被覆 $\{f_i : U_i' \to U'\}$ と射
$c_i : U_i' \to U$ が存在して
$u(c_i) = c \circ u(f_i)$ となる。
\item $V \in \Ob(\mathcal{D})$ が与えられると、$\mathcal{D}$ に
$\{u(U_i) \to V\}_{i \in I}$ という形の $V$ の被覆が存在する。
\end{enumerate}
このとき、補題 \ref{sites-lemma-cocontinuous-morphism-topoi} により
余連続関手 $u$ に付随するトポスの射
$$
g : \Sh(\mathcal{C}) \longrightarrow \Sh(\mathcal{D})
$$
は同値である。
\end{lemma}

\begin{proof}
$u$ が性質 (1) -- (5) を満たすと仮定する。随伴写像
$$
\mathcal{G} \longrightarrow g_*g^{-1}\mathcal{G}
\quad\text{and}\quad
g^{-1}g_*\mathcal{F} \longrightarrow \mathcal{F}
$$
が同型であることを示す。

\medskip\noindent
補題 \ref{sites-lemma-when-shriek} が適用でき、$\mathcal{D}$ 上の任意の層
$\mathcal{G}$ に対して
$g^{-1}\mathcal{G}(U) = \mathcal{G}(u(U))$ であることに注意する。
次に、$\mathcal{F}$ を $\mathcal{C}$ 上の層とし、
$V$ を $\mathcal{D}$ の対象とする。定義により
$g_*\mathcal{F}(V) = \lim_{u(U) \to V} \mathcal{F}(U)$ である。したがって
$$
g^{-1}g_*\mathcal{F}(U) = \lim_{U', u(U') \to u(U)} \mathcal{F}(U')
$$
である。ここで射 $\psi : u(U') \to u(U)$ は
$u(\alpha)$ という形であるとは限らない。
% P01-E0109
このような対 $(U', \psi)$ の圏は終対象、すなわち
$(U, \text{id})$ をもち、これにより極限から
$\mathcal{F}(U)$ への写像が得られる。
$(s_{(U', \psi)})$ をこの極限の元とする。
$s_{(U', \psi)}$ が値
$s_{(U, \text{id})} \in \mathcal{F}(U)$ によって一意に定まることを
示したい。性質 (4) により、任意の $(U', \psi)$ に対して被覆
$\{U'_i \to U'\}$ が存在し、合成
$u(U'_i) \to u(U') \to u(U)$ は、ある $\mathcal{C}$ における
$c_i : U'_i \to U$ に対する $u(c_i)$ という形になる。したがって
$$
s_{(U', \psi)}|_{U'_i} = c_i^*(s_{(U, \text{id})}).
$$
$\mathcal{F}$ は層なので、実際に $s_{(U', \psi)}$ は
$s_{(U, \text{id})}$ によって定まる。これで一意性が示された。
存在について、任意の
$s \in \mathcal{F}(U)$、$\psi : u(U') \to u(U)$、
$\{f_i : U_i' \to U'\}$ および $c_i : U_i' \to U$ が与えられ、
上のように $\psi \circ u(f_i) = u(c_i)$ であると仮定する。
$$
s_{(U', \psi)}|_{U'_i} = c_i^*(s).
$$
を満たす（一意な）元
$s_{(U', \psi)} \in \mathcal{F}(U')$ が存在すると主張する。
実際、図式
$$
\xymatrix{
U_i' \times_{U'} U_j' \ar[r]_-{\text{pr}_2} \ar[d]_{\text{pr}_1} &
U_j' \ar[d]^{c_j} \\
U_i' \ar[r]^{c_i} & U}
$$
は可換とは限らないので、先験的には元
$c_i^*(s)|_{U_i' \times_{U'} U_j'}$ と
$c_j^*(s)|_{U_i' \times_{U'} U_j'}$ が一致するかは明らかでない。
しかし補題の条件 (3) により、被覆
$\{f_{ijk} : U'_{ijk} \to U_i' \times_{U'} U_j'\}_{k \in K_{ij}}$
が存在して
$c_i \circ \text{pr}_1 \circ f_{ijk} = c_j \circ \text{pr}_2 \circ f_{ijk}$
となる。したがって
$$
f_{ijk}^* \left(c_i^*s|_{U_i' \times_{U'} U_j'}\right)
=
f_{ijk}^* \left(c_j^*s|_{U_i' \times_{U'} U_j'}\right)
$$
である。ゆえに $\mathcal{F}$ の層条件から
$c_i^*(s)|_{U_i' \times_{U'} U_j'} = c_j^*(s)|_{U_i' \times_{U'} U_j'}$
であり、さらに $\mathcal{F}$ の層条件から
$s_{(U', \psi)}$ が存在する。一意性により、このように得られる族
$(s_{(U', \psi)})$ は極限の元になる。
% P01-E0108
しかも $s_{(U, \psi)} = s$ である。これで
$g^{-1}g_*\mathcal{F} \to \mathcal{F}$ が同型であることが示された。

\medskip\noindent
$\mathcal{G}$ を $\mathcal{D}$ 上の層とし、
$V$ を $\mathcal{D}$ の対象とする。このとき
$$
g_*g^{-1}\mathcal{G}(V) =
\lim_{U, \psi : u(U) \to V} \mathcal{G}(u(U))
$$
である。前段落により、$V = u(U)$ という形の対象 $V$ 上での層
$g_*g^{-1}\mathcal{G}$ の値は $\mathcal{G}(u(U))$ に等しい。
（形式的には、$g^{-1}g_*g^{-1} \cong g^{-1}$ と、証明の冒頭で与えた
$g^{-1}$ の記述から従う。非形式的には、ここでの極限と上の極限を
比較すればよい。）したがって随伴写像
$\mathcal{G} \to g_*g^{-1}\mathcal{G}$ は、$u(U)$ という形の任意の
対象上の切断について全単射である。公理 (5) により、
$V$ には $u(U)$ という形の対象による被覆が存在するので、
随伴写像は同型であることが容易に分かる。
\end{proof}

\noindent
補題 \ref{sites-lemma-equivalence} に現れる余連続関手に名前を付けておくと
便利である。

\begin{definition}
\label{sites-definition-special-cocontinuous-functor}
$\mathcal{C}$、$\mathcal{D}$ をサイトとする。
{\it $\mathcal{C}$ から $\mathcal{D}$ への特殊余連続関手 $u$}
とは、補題 \ref{sites-lemma-equivalence} の仮定と結論を満たす
余連続関手 $u : \mathcal{C} \to \mathcal{D}$ のことである。
\end{definition}

\begin{lemma}
\label{sites-lemma-localize-special-cocontinuous}
$\mathcal{C}$、$\mathcal{D}$ をサイトとする。
$u : \mathcal{C} \to \mathcal{D}$ を特殊余連続関手とする。
$\mathcal{C}$ の各対象 $U$ に対し、補題
\ref{sites-lemma-localize-cocontinuous} におけるような可換図式
$$
\xymatrix{
\mathcal{C}/U \ar[r]_{j_U} \ar[d] & \mathcal{C} \ar[d]^u \\
\mathcal{D}/u(U) \ar[r]^-{j_{u(U)}} & \mathcal{D}
}
$$
をもつ。左の垂直矢印は特殊余連続関手である。
したがってトポスの可換図式
$$
\xymatrix{
\Sh(\mathcal{C}/U) \ar[r]_{j_U} \ar[d] &
\Sh(\mathcal{C}) \ar[d]^u \\
\Sh(\mathcal{D}/u(U)) \ar[r]^-{j_{u(U)}} &
\Sh(\mathcal{D})
}
$$
において、垂直矢印は同値である。
\end{lemma}

\begin{proof}
図式の存在と可換性は補題 \ref{sites-lemma-localize-cocontinuous} で
既に示した。誘導された関手
$u : \mathcal{C}/U \to \mathcal{D}/u(U)$ について、
補題 \ref{sites-lemma-equivalence} の仮定 (1) -- (5) を確認しなければ
ならない。これはまったく機械的である。

\medskip\noindent
性質 (1)。これは補題 \ref{sites-lemma-localize-cocontinuous} である。

\medskip\noindent
性質 (2)。$\{U_i'/U \to U'/U\}_{i \in I}$ を
$\mathcal{C}/U$ における $U'/U$ の被覆とする。
$u$ は連続なので、
$\{u(U_i')/u(U) \to u(U')/u(U)\}_{i \in I}$ は
$\mathcal{D}/u(U)$ における $u(U')/u(U)$ の被覆である。
したがって $u : \mathcal{C}/U \to \mathcal{D}/u(U)$ について
(2) が成り立つ。

\medskip\noindent
性質 (3)。$\mathcal{C}/U$ における射
$a, b : U''/U \to U'/U$ が、$\mathcal{D}/u(U)$ において
$u(a) = u(b)$ を満たすとする。$u$ は (3) を満たすので、
$\mathcal{C}$ に被覆 $\{f_i : U''_i \to U''\}$ が存在して
$a \circ f_i = b \circ f_i$ となる。これは
$\mathcal{C}/U$ に、$a \circ f_i = b \circ f_i$ を満たす被覆
$\{f_i : U''_i/U \to U''/U\}$ を与える。したがって
$u : \mathcal{C}/U \to \mathcal{D}/u(U)$ について (3) が成り立つ。

\medskip\noindent
性質 (4)。$U''/U, U'/U \in \Ob(\mathcal{C}/U)$ と、
$\mathcal{D}/u(U)$ における射
$c : u(U'')/u(U) \to u(U')/u(U)$ が与えられたとする。
$u$ は性質 (4) を満たすので、$\mathcal{C}$ に被覆
$\{f_i : U_i'' \to U''\}$ と射 $c_i : U_i'' \to U'$ が存在し、
$u(c_i) = c \circ u(f_i)$ となる。
$U_i'' \to U'' \to U$ という合成により、
$U_i''$ を $U$ 上の対象とみなす。
$c_i$ は $U$ 上の射でないかもしれない。
しかし $u(c_i)$ は $u(U)$ 上の射なので、$u$ に対して
性質 (3) を適用し、被覆 $\{f_{ik} : U''_{ik} \to U''_i\}$ を
見いだして、$c_{ik} = c_i \circ f_{ik} : U''_{ik} \to U'$ が
$U$ 上の射となるようにできる。したがって
$\{f_i \circ f_{ik} : U''_{ik}/U \to U''/U\}$ は
$\mathcal{C}/U$ の被覆である。
% P01-E0110
しかも $u(c_{ik}) = c \circ u(f_{ik})$ である。
ゆえに $u : \mathcal{C}/U \to \mathcal{D}/u(U)$ について
(4) が成り立つ。

\medskip\noindent
性質 (5)。$h : V \to u(U)$ を
$\mathcal{D}/u(U)$ の対象とする。$u$ は性質 (5) を満たすので、
$\mathcal{D}$ に被覆 $\{c_i : u(U_i) \to V\}$ が存在する。
性質 (4) により、被覆 $\{f_{ij} : U_{ij} \to U_i\}$ と射
$c_{ij} : U_{ij} \to U$ を見いだして
$u(c_{ij}) = h \circ c_i \circ u(f_{ij})$ とできる。
したがって $\{u(U_{ij})/u(U) \to V/u(U)\}$ は
$\mathcal{D}/u(U)$ における所望の形の被覆であり、
$u : \mathcal{C}/U \to \mathcal{D}/u(U)$ について
(5) が成り立つ。
\end{proof}

\begin{lemma}
\label{sites-lemma-special-equivalence}
$\mathcal{C}$ をサイトとする。
$\mathcal{C}' \subset \Sh(\mathcal{C})$ を、次を満たす
充満部分圏（対象の集合をもつもの）とする。
\begin{enumerate}
\item すべての $U \in \Ob(\mathcal{C})$ に対して
$h_U^\# \in \Ob(\mathcal{C}')$ である。
\item $\mathcal{C}'$ は $\Sh(\mathcal{C})$ における
ファイバー積で保たれる。
\end{enumerate}
$\mathcal{C}'$ の被覆を、層の全射
$\coprod_{i \in I} \mathcal{F}_i \to \mathcal{F}$ を与える
任意の写像族
$\{\mathcal{F}_i \to \mathcal{F}\}_{i \in I}$ と定める。このとき
\begin{enumerate}
\item $\mathcal{C}'$ はサイトである（被覆の集合を一つ選ぶことについては
Sets, Lemma \ref{sets-lemma-coverings-site} を参照）。
\item $\mathcal{C}'$ 上の表現可能前層は層である
（すなわち $\mathcal{C}'$ の位相は準標準である。定義
\ref{sites-definition-weaker-than-canonical} を参照）。
\item 関手 $v : \mathcal{C} \to \mathcal{C}'$、
$U \mapsto h_U^\#$ は特殊余連続関手であり、したがって同値
$g : \Sh(\mathcal{C}) \to \Sh(\mathcal{C}')$ を誘導する。
\item 任意の $\mathcal{F} \in \Ob(\mathcal{C}')$ に対して
$g^{-1}h_\mathcal{F} = \mathcal{F}$ である。
\item 任意の $U \in \Ob(\mathcal{C})$ に対して
$g_*h_U^\# = h_{v(U)} = h_{h_U^\#}$ である。
\end{enumerate}
\end{lemma}

\begin{proof}
% P01-E0111
注意：上の主張のいくつかは、初めは少し分かりにくく見えるかもしれない。
これは $\mathcal{C}'$ の対象を $\mathcal{C}$ 上の層とも
みなせるからである。補題で記述した $\mathcal{C}'$ の被覆が
定義 \ref{sites-definition-site} の条件を満たすことの証明は省略する。

\medskip\noindent
$\{\mathcal{F}_i \to \mathcal{F}\}$ を層の射の全射族とする。
$\mathcal{G}$ を別の層とする。補題の (2) は
$$
\xymatrix{
\Mor_{\Sh(\mathcal{C})}(
\coprod_{i \in I} \mathcal{F}_i, \mathcal{G})
\ar@<1ex>[r] \ar@<-1ex>[r]
&
\Mor_{\Sh(\mathcal{C})}(
\coprod_{(i_0, i_1) \in I \times I}
\mathcal{F}_{i_0} \times_\mathcal{F} \mathcal{F}_{i_1}, \mathcal{G})
}
$$
の等化子が
$\Mor_{\Sh(\mathcal{C})}(\mathcal{F}, \mathcal{G}).$
であるということである。これは明らかである
（例えば補題 \ref{sites-lemma-coequalizer-surjection} を用いよ）。

\medskip\noindent
(3) を証明するには、補題 \ref{sites-lemma-equivalence} の条件
(1) -- (5) を確認しなければならない。$v$ が余連続であることは、
補題 \ref{sites-lemma-mono-epi-sheaves} における層の全射の記述と同値である。
関手 $v$ は連続である。実際、
$U \mapsto h_U^\#$ はファイバー積と可換であり、被覆を被覆に移す
（補題 \ref{sites-lemma-sheafification-exact} と補題
\ref{sites-lemma-covering-surjective-after-sheafification} を参照）。
補題 \ref{sites-lemma-equivalence} の性質 (3)、(4) は、射
$f : h_{U'}^\# \to h_U^\#$ に関する主張である。
そのような射は $h_U^\#(U')$ の元と同じものである。
したがって (3) と (4) は層化の構成から直ちに従う。
補題 \ref{sites-lemma-equivalence} の性質 (5) は
補題 \ref{sites-lemma-sheaf-coequalizer-representable} である。
% P01-E0112
$v$ に付随するトポスの同値を
$g : \Sh(\mathcal{C}) \to \Sh(\mathcal{C}')$ と書く。
これは補題 \ref{sites-lemma-equivalence} による。

\medskip\noindent
$\mathcal{F}$ を補題の (4) におけるものとする。
任意の $U \in \Ob(\mathcal{C})$ に対して
$$
g^{-1}h_\mathcal{F}(U) = h_\mathcal{F}(v(U))
= \Mor_{\Sh(\mathcal{C})}(h_U^\#, \mathcal{F})
= \mathcal{F}(U)
$$
である。
% P01-E0113
最初の等式は補題 \ref{sites-lemma-when-shriek} から従う。
したがって (4) が成り立つ。

\medskip\noindent
$\mathcal{F} \in \Ob(\mathcal{C}')$ とし、
$U \in \Ob(\mathcal{C})$ とする。このとき
\begin{align*}
g_*h_U^\#(\mathcal{F})
& =
\Mor_{\Sh(\mathcal{C}')}(h_\mathcal{F}, g_*h_U^\#) \\
& =
\Mor_{\Sh(\mathcal{C})}(g^{-1}h_\mathcal{F}, h_U^\#) \\
& =
\Mor_{\Sh(\mathcal{C})}(\mathcal{F}, h_U^\#) \\
& =
\Mor_{\mathcal{C}'}(\mathcal{F}, h_U^\#)
\end{align*}
となり、所望の結論を得る（第三の等式は上で示した）。
\end{proof}

\noindent
これを用いると、任意のトポスを変形して、すべての有限極限をもつ
サイト上にあるようにできる。

\begin{lemma}
\label{sites-lemma-topos-good-site}
$\Sh(\mathcal{C})$ をトポスとする。
$\{\mathcal{F}_i\}_{i \in I}$ を $\mathcal{C}$ 上の層の集合とする。
サイトの特殊余連続関手
$u : \mathcal{C} \to \mathcal{C}'$ により誘導されるトポスの同値
$g : \Sh(\mathcal{C}) \to \Sh(\mathcal{C}')$ で、次を満たすものが
存在する。
\begin{enumerate}
\item $\mathcal{C}'$ は準標準位相をもつ。
\item $\mathcal{C}'$ の射の族 $\{V_j \to V\}$ が
$\mathcal{C}'$ の被覆と（組合せ的に）同値であることと、
$\coprod h_{V_j} \to h_V$ が全射であることは同値である。
\item $\mathcal{C}'$ はファイバー積と終対象をもつ
（すなわち $\mathcal{C}'$ はすべての有限極限をもつ）。
\item $\mathcal{C}'$ 上の表現可能層の任意の部分層は表現可能である。
\item 各 $g_*\mathcal{F}_i$ は表現可能層である。
\end{enumerate}
\end{lemma}

\begin{proof}
充満部分圏
$\mathcal{C}_1 \subset \Sh(\mathcal{C})$ を、すべての
$U \in \Ob(\mathcal{C})$ に対する $h_U^\#$、与えられた層
$\mathcal{F}_i$、および終層 $*$（例
\ref{sites-example-singleton-sheaf} を参照）からなるものとする。
充満部分圏を帰納的に
% P01-E0114
$$
\mathcal{C}_1 \subset \mathcal{C}_2 \subset \mathcal{C}_2 \subset \ldots
\subset \Sh(\mathcal{C})
$$
と定める。すなわち、$\mathcal{C}_n$ が与えられたとき、
$\mathcal{C}_{n + 1}$ を $\mathcal{C}_n$ の対象のすべての
ファイバー積と部分層からなる充満部分圏とする
（$\mathcal{C}_{n + 1}$ は対象の集合をもつことに注意する）。
$\mathcal{C}' = \bigcup_{n \geq 1} \mathcal{C}_n$ とおく。
$\mathcal{C}'$ の被覆とは、$\mathcal{C}$ 上の層の写像として
$\coprod \mathcal{G}_j \to \mathcal{G}$ が全射となるような、
$\mathcal{C}'$ の対象の射の任意の族
$\{\mathcal{G}_j \to \mathcal{G}\}_{j \in J}$ のことである。
関手 $v : \mathcal{C} \to \mathcal{C'}$ は
$U \mapsto h_U^\#$ で与えられる。
補題 \ref{sites-lemma-special-equivalence} を適用すればよい。
\end{proof}

\noindent
以下が本節の目標である。

\begin{lemma}
\label{sites-lemma-morphism-topoi-comes-from-morphism-sites}
\begin{reference}
この主張は
\cite[Proposition 4.9.4. Expos\'e IV]{SGA4}
と密接に関係する。結果全体を得るには、さらに
\cite[Remarque 4.7.4, Expos\'e IV]{SGA4}
も用いるべきである。
\end{reference}
$\mathcal{C}$、$\mathcal{D}$ をサイトとする。
$f : \Sh(\mathcal{C}) \to \Sh(\mathcal{D})$ をトポスの射とする。
このとき、サイト $\mathcal{C}'$ と関手の図式
$$
\xymatrix{
\mathcal{C} \ar[r]_v & \mathcal{C}' & \mathcal{D} \ar[l]^u
}
$$
で次を満たすものが存在する。
\begin{enumerate}
\item 関手 $v$ は特殊余連続関手である。
\item 関手 $u$ はファイバー積と可換であり、連続で、
サイトの射 $\mathcal{C}' \to \mathcal{D}$ を定める。
\item トポスの射 $f$ は、トポスの射の合成
$$
\Sh(\mathcal{C}) \longrightarrow
\Sh(\mathcal{C}') \longrightarrow
\Sh(\mathcal{D})
$$
と一致する。ここで第一の矢印は補題 \ref{sites-lemma-equivalence} により
$v$ から生じ、第二の矢印は補題
\ref{sites-lemma-morphism-sites-topoi} により $u$ から生じる。
\end{enumerate}
\end{lemma}

\begin{proof}
充満部分圏
$\mathcal{C}_1 \subset \Sh(\mathcal{C})$ を、すべての
$U \in \Ob(\mathcal{C})$ と $V \in \Ob(\mathcal{D})$ に対する
$h_U^\#$ および $f^{-1}h_V^\#$ からなるものとする。
$\mathcal{C}_{n + 1}$ を、$\mathcal{C}_n$ の対象のすべての
ファイバー積からなる充満部分圏とする。
$\mathcal{C}' = \bigcup_{n \geq 1} \mathcal{C}_n$ とおく。
$\mathcal{C}'$ の被覆とは、$\mathcal{C}$ 上の層の写像として
$\coprod_{i \in I} \mathcal{F}_i \to \mathcal{F}$ が全射となるような
任意の族 $\{\mathcal{F}_i \to \mathcal{F}\}_{i \in I}$ のことである。
関手 $v : \mathcal{C} \to \mathcal{C'}$ は
$U \mapsto h_U^\#$ で与えられる。
関手 $u : \mathcal{D} \to \mathcal{C'}$ は
$V \mapsto f^{-1}h_V^\#$ で与えられる。

\medskip\noindent
(1) は補題 \ref{sites-lemma-special-equivalence} から従う。

\medskip\noindent
補題の (2) と (3) を証明する。$V \mapsto h_V^\#$ と
$f^{-1}$ はともにファイバー積と可換なので、関手 $u$ もそうである。
さらに $f^{-1}$ は完全で任意の余極限と可換なので、
被覆を層の射の全射族に移す。したがって $u$ は連続である。
これがサイトの射を定めることを見るには、さらに $u_s$ が
完全であることを示さなければならない。そのため
$g^{-1} \circ u_s = f^{-1}$ を示す。実際、そうすれば
$g^{-1}$ は同値で $f^{-1}$ は完全なので結論が従う。
$g^{-1}$ は $g_*$ の随伴であり、$u_s$ は $u^s$ の随伴であり、
$f^{-1}$ は $f_*$ の随伴なので、
$u^s \circ g_* = f_*$ を証明するだけでも十分である。
$\mathcal{F}$ を $\mathcal{C}$ 上の層とし、
$V$ を $\mathcal{D}$ の対象とする。このとき
\begin{align*}
(u^s g_{\ast}\mathcal{F})(V)
& =
(g_{\ast}\mathcal{F})(f^{-1}h_V^\#) \\
& =
\Mor_{Sh(\mathcal{C}')}(h_{f^{-1}h_V^\#},g_{\ast}\mathcal{F}) \\
& =
\Mor_{Sh(\mathcal{C})}(g^{-1}h_{f^{-1}h_V^\#},\mathcal{F}) \\
& =
\Mor_{Sh(\mathcal{C})}(f^{-1}h_V^\#,\mathcal{F}) \\
& =
\Mor_{Sh(\mathcal{D})}(h_V^\#,f_{\ast}\mathcal{F}) \\
& =
f_{\ast}\mathcal{F}(V)
\end{align*}
である。
% P01-E0115
第一の等式は $u^s = u^p$ から従う。
第二の等式は米田の補題による。
第三の等式は $g^{-1}$ と $g_*$ の随伴性から従う。
第四の等式は補題 \ref{sites-lemma-special-equivalence} (4) による。
第五の等式は $f^{-1}$ と $f_*$ の随伴性から従う。
第六の等式は米田の補題による。
したがって $u^s g_*\mathcal{F} = f_*\mathcal{F}$ であり、
補題の証明が完了する。
\end{proof}

\begin{remark}
\label{sites-remark-morphism-topoi-comes-from-morphism-sites}
記法と仮定は補題 \ref{sites-lemma-morphism-topoi-comes-from-morphism-sites}
のとおりとする。サイト $\mathcal{D}$ が終対象とファイバー積を
もつならば、関手 $u : \mathcal{D} \to \mathcal{C}'$ は
命題 \ref{sites-proposition-get-morphism} のすべての仮定を満たす。
実際、補題で述べた性質に加えて、$u$ は $\mathcal{D}$ の終対象を
$\mathcal{C}'$ の終対象に移す。これは $u$ の構成から明らかである。
したがって、まず
% P01-E0116
補題 \ref{sites-lemma-topos-good-site}
を $\mathcal{D}$ に適用し、次に
補題 \ref{sites-lemma-morphism-topoi-comes-from-morphism-sites}
を、得られたトポスの射
$\Sh(\mathcal{C}) \to \Sh(\mathcal{D}')$
に適用すると、次の主張を得る。
任意のトポスの射
$f : \Sh(\mathcal{C}) \to \Sh(\mathcal{D})$
は可換図式
$$
\xymatrix{
\Sh(\mathcal{C}) \ar[d]_g \ar[r]_f &
\Sh(\mathcal{D}) \ar[d]^e \\
\Sh(\mathcal{C}') \ar[r]^{f'} &
\Sh(\mathcal{D}')
}
$$
に組み込まれ、次の性質を満たす。
\begin{enumerate}
\item 射 $e$ と $g$ は、特殊余連続関手
$\mathcal{C} \to \mathcal{C}'$ および
$\mathcal{D} \to \mathcal{D}'$ により与えられる同値である。
\item サイト $\mathcal{C}'$ と $\mathcal{D}'$ は
ファイバー積と終対象をもち、準標準位相をもつ。
\item 射 $f' : \mathcal{C}' \to \mathcal{D}'$ は、
命題 \ref{sites-proposition-get-morphism} を適用できる関手
$u : \mathcal{D}' \to \mathcal{C}'$ に対応するサイトの射から生じる。
\item $\mathcal{C}$ 上の任意の層の集合 $\mathcal{F}_i$
（それぞれ $\mathcal{D}$ 上の $\mathcal{G}_j$）が与えられたとき、
これらはそれぞれ $\mathcal{C}'$ 上
（それぞれ $\mathcal{D}'$ 上）の表現可能層であると仮定できる。
\end{enumerate}
$\mathcal{C}$ と $\mathcal{D}$ を
$\mathcal{C}'$ と $\mathcal{D}'$ で置き換えることは、しばしば有用である。
\end{remark}

\begin{remark}
\label{sites-remark-equivalence-topoi-comes-from-morphism-sites}
記法と仮定は補題 \ref{sites-lemma-morphism-topoi-comes-from-morphism-sites}
のとおりとする。さらに、もとのトポスの射
$\Sh(\mathcal{C}) \to \Sh(\mathcal{D})$ が同値であると仮定する。
このとき、補題
\ref{sites-lemma-morphism-topoi-comes-from-morphism-sites}
の証明における構成は二つの関手
$$
\mathcal{C} \rightarrow \mathcal{C}' \leftarrow \mathcal{D}
$$
を与え、両者とも特殊余連続関手である。したがってこの場合、
トポスの射を実際に、Remark
\ref{sites-remark-morphism-topoi-comes-from-morphism-sites} におけるような合成
$$
\Sh(\mathcal{C}) \rightarrow
\Sh(\mathcal{C}') =
\Sh(\mathcal{D}') \leftarrow
\Sh(\mathcal{D})
$$
に分解できる。ただし中央の射は恒等射である。
\end{remark}

\section{トポスの局所化}
\label{sites-section-localize-topoi}

\noindent
局所化に関する内容の一部を、一見より一般的なトポスの場合について
繰り返す。実際にはこれは一般化ではない。基礎となるサイトをいつでも
拡大し、サイトの対象で局所化していると仮定できるからである。

\begin{lemma}
\label{sites-lemma-localize-topos}
$\mathcal{C}$ をサイトとし、$\mathcal{F}$ を
$\mathcal{C}$ 上の層とする。このとき圏
$\Sh(\mathcal{C})/\mathcal{F}$ はトポスである。
標準的なトポスの射
$$
j_\mathcal{F} :
\Sh(\mathcal{C})/\mathcal{F}
\longrightarrow
\Sh(\mathcal{C})
$$
が存在し、これは節 \ref{sites-section-localize} における意味で局所化であり、
次を満たす。
\begin{enumerate}
\item 関手 $j_\mathcal{F}^{-1}$ は
$\mathcal{H} \mapsto \mathcal{H} \times \mathcal{F}/\mathcal{F}$ という関手である。
\item 関手 $j_{\mathcal{F}!}$ は忘却関手
$\mathcal{G}/\mathcal{F} \mapsto \mathcal{G}$ である。
\end{enumerate}
\end{lemma}

\begin{proof}
補題 \ref{sites-lemma-topos-good-site} を適用する。これは、
$\mathcal{C}$ が準標準位相をもつサイトであり、ある
$U \in \Ob(\mathcal{C})$ に対して
$\mathcal{F} = h_U = h_U^\#$ であると仮定してよいことを意味する。
したがって節 \ref{sites-section-localize} の内容を適用できる。特に、
合成
$$
\Sh(\mathcal{C}/U)
\to
\Sh(\mathcal{C})/h_U^\#
\to \Sh(\mathcal{C})
$$
が $j_{U!}$ に等しくなるような同値
$\Sh(\mathcal{C}/U) = \Sh(\mathcal{C})/h_U^\#$ が存在する。
補題 \ref{sites-lemma-essential-image-j-shriek} を参照されたい。
% P01-E0117
逆関手を
$a : \Sh(\mathcal{C})/h_U^\# \to \Sh(\mathcal{C}/U)$ と書く。
すると $j_{\mathcal{F}!} = j_{U!} \circ a$、
$j_\mathcal{F}^{-1} = a^{-1} \circ j_U^{-1}$、および
$j_{\mathcal{F}, *} = j_{U, *} \circ a$ である。
$j_{\mathcal{F}!}$ の記述は上から従い、
$j_\mathcal{F}^{-1}$ の記述は補題
\ref{sites-lemma-compute-j-shriek-restrict} から従う。
\end{proof}

\begin{lemma}
\label{sites-lemma-localize-pushforward}
補題 \ref{sites-lemma-localize-topos} の状況において、関手
$j_{\mathcal{F}, *}$ は
% P01-E0118
$\varphi : \mathcal{G} \to \mathcal{F}$ に層
$$
U
\longmapsto
\{\alpha : \mathcal{F}|_U \to \mathcal{G}|_U
\text{ such that } \alpha \text{ is a right inverse to }\varphi|_U \}.
$$
を対応させる関手である。
\end{lemma}

\begin{proof}
任意の $\varphi : \mathcal{G} \to \mathcal{F}$ に対し、対応する
$\Sh(\mathcal{C})/\mathcal{F}$ の対象を
$\mathcal{G}/\mathcal{F}$ と書くことにする。このとき
$$
(j_{\mathcal{F}, *}(\mathcal{G}/\mathcal{F}))(U) =
\Mor_{\Sh(\mathcal{C})}(h_U^\#, j_{\mathcal{F}, *}(\mathcal{G}/\mathcal{F})) =
\Mor_{\Sh(\mathcal{C})/\mathcal{F}}(j_{\mathcal{F}}^{-1}h_U^{\#},
(\mathcal{G}/\mathcal{F})).
$$
である。補題 \ref{sites-lemma-localize-topos} により、この集合は、集合の写像
$$
\Mor_{\Sh(\mathcal{C})}(h_U^\# \times \mathcal{F}, \mathcal{G})
\xrightarrow{\varphi \circ}
\Mor_{\Sh(\mathcal{C})}(h_U^\# \times \mathcal{F}, \mathcal{F}).
$$
のもとでの元 $h_U^\# \times \mathcal{F} \to \mathcal{F}$ の
ファイバーである。補題 \ref{sites-lemma-internal-hom-sheaf} の随伴により
\begin{align*}
\Mor_{\Sh(\mathcal{C})}(h_U^{\#}\times\mathcal{F}, \mathcal{G})
& =
\Mor_{\Sh(\mathcal{C})}(h_U^{\#},\SheafHom(\mathcal{F}, \mathcal{G})) \\
& =
\Mor_{\Sh(\mathcal{C}/U)}(\mathcal{F}|_{\mathcal{C}/U},
\mathcal{G}|_{\mathcal{C}/U}), \\
\Mor_{\Sh(\mathcal{C})}(h_U^{\#} \times \mathcal{F}, \mathcal{F})
& =
\Mor_{\Sh(\mathcal{C})}(h_U^{\#},\SheafHom(\mathcal{F},\mathcal{F})) \\
& =
\Mor_{\Sh(\mathcal{C}/U)}(\mathcal{F}|_{\mathcal{C}/U},
\mathcal{F}|_{\mathcal{C}/U}),
\end{align*}
であり、この随伴のもとで写像 $\varphi\circ$ は写像
$$
\Mor_{\Sh(\mathcal{C}/U)}(\mathcal{F}|_{\mathcal{C}/U},
\mathcal{G}|_{\mathcal{C}/U})
\longrightarrow
\Mor_{\Sh(\mathcal{C}/U)}(\mathcal{F}|_{\mathcal{C}/U},
\mathcal{F}|_{\mathcal{C}/U}),\quad
\psi \longmapsto \varphi|_{\mathcal{C}/U} \circ \psi,
$$
となり、元 $h_U^\# \times \mathcal{F} \to \mathcal{F}$ は
$\text{id}_{\mathcal{F}|_{\mathcal{C}/U}}$ となる。したがって
$(j_{\mathcal{F}, *}\mathcal{G}/\mathcal{F})(U)$ は写像
$$
\Mor_{\Sh(\mathcal{C}/U)}(\mathcal{F}|_{\mathcal{C}/U},
\mathcal{G}|_{\mathcal{C}/U})
\xrightarrow{\varphi|_{\mathcal{C}/U}\circ}
\Mor_{\Sh(\mathcal{C}/U)}(\mathcal{F}|_{\mathcal{C}/U},
\mathcal{F}|_{\mathcal{C}/U}),
$$
のもとでの $\text{id}_{\mathcal{F}|_{\mathcal{C}/U}}$ の
ファイバーと同型である。これは所望の
$$
\{\alpha : \mathcal{F}|_U \to \mathcal{G}|_U
\text{ such that } \alpha \text{ is a right inverse to }\varphi|_U \}
$$
である。
\end{proof}

\begin{lemma}
\label{sites-lemma-localize-topos-site}
$\mathcal{C}$ をサイトとし、$\mathcal{F}$ を
$\mathcal{C}$ 上の層とする。$\mathcal{C}/\mathcal{F}$ を、
$U \in \Ob(\mathcal{C})$ かつ $s \in \mathcal{F}(U)$ である対
$(U, s)$ の圏とする。$\mathcal{C}/\mathcal{F}$ における被覆を、
$\{U_i \to U\}$ が $\mathcal{C}$ の被覆となるような族
$\{(U_i, s_i) \to (U, s)\}$ とする。このとき
$j : \mathcal{C}/\mathcal{F} \to \mathcal{C}$ はサイトの
連続かつ余連続な関手であり、トポスの射
$j : \Sh(\mathcal{C}/\mathcal{F}) \to \Sh(\mathcal{C})$ を誘導する。
実際、$j$ を $j_\mathcal{F}$ に移す同値
$\Sh(\mathcal{C}/\mathcal{F}) =
\Sh(\mathcal{C})/\mathcal{F}$ が存在する。
\end{lemma}

\begin{proof}
$\mathcal{C}/\mathcal{F}$ がサイトであり、$j$ が連続かつ
余連続であることの確認は省略する。補題
\ref{sites-lemma-when-shriek} により、示されたトポスの射 $j$ が存在し、
$j^{-1}\mathcal{G}(U, s) = \mathcal{G}(U)$ であり、
$j^{-1}$ は左随伴 $j_!$ をもつ。
$\mathcal{C}/\mathcal{F}$ 上の射
$\varphi : * \to j^{-1}\mathcal{G}$ は、各対 $(U, s)$ に
制限写像と両立する切断
$\varphi(s) \in \mathcal{G}(U)$ を対応させる規則と同じものである。
したがってこれは、$\mathcal{C}$ 上の射
$\varphi : \mathcal{F} \to \mathcal{G}$ と同じものである。
ゆえに $j_!* = \mathcal{F}$ である。特に、任意の
$\mathcal{H} \in \Sh(\mathcal{C}/\mathcal{F})$ に対して標準写像
$$
j_!\mathcal{H} \to j_!* = \mathcal{F}
$$
が存在する。すなわち、関手
$j'_! : \Sh(\mathcal{C}/\mathcal{F}) \to \Sh(\mathcal{C})/\mathcal{F}$
を得る。この関手の逆は、
$\Sh(\mathcal{C})/\mathcal{F}$ の対象
$\varphi : \mathcal{G} \to \mathcal{F}$ に層
$$
a(\mathcal{G}/\mathcal{F}) :
(U, s) \longmapsto \{t \in \mathcal{G}(U) \mid \varphi(t) = s\}
$$
を対応させる規則である。
$a(\mathcal{G}/\mathcal{F})$ が層であり、$a$ が
$j'_!$ の逆であることの確認は省略する。
\end{proof}

\begin{definition}
\label{sites-definition-localize-topos}
$\mathcal{C}$ をサイトとし、$\mathcal{F}$ を
$\mathcal{C}$ 上の層とする。
\begin{enumerate}
\item トポス $\Sh(\mathcal{C})/\mathcal{F}$ を
{\it $\mathcal{F}$ におけるトポス $\Sh(\mathcal{C})$ の局所化}
と呼ぶ。
\item 補題 \ref{sites-lemma-localize-topos} のトポスの射
$j_\mathcal{F} :
\Sh(\mathcal{C})/\mathcal{F}
\to
\Sh(\mathcal{C})$ を {\it 局所化射} と呼ぶ。
\end{enumerate}
\end{definition}

\noindent
層 $\mathcal{F}$ がサイトのある対象 $U$ に対する $h_U^\#$ に
等しいとき、トポスの局所化が $U$ におけるサイトの局所化上の
層の圏に等しいことを示す。さらに、任意の関手性がこの同一視と
両立することを確認する。

\begin{lemma}
\label{sites-lemma-localize-compare}
$\mathcal{C}$ をサイトとする。
$\mathcal{F} = h_U^\#$ が $\mathcal{C}$ のある対象 $U$ に対して
成り立つとする。このとき、補題
\ref{sites-lemma-localize-topos} で構成した
$j_\mathcal{F} : \Sh(\mathcal{C})/\mathcal{F}
\to \Sh(\mathcal{C})$ は、節 \ref{sites-section-localize} で構成したトポスの射
$j_U : \Sh(\mathcal{C}/U) \to \Sh(\mathcal{C})$
と、補題 \ref{sites-lemma-essential-image-j-shriek} の同一視
$\Sh(\mathcal{C}/U) = \Sh(\mathcal{C})/h_U^\#$ のもとで一致する。
\end{lemma}

\begin{proof}
補題 \ref{sites-lemma-essential-image-j-shriek} において、合成
$\Sh(\mathcal{C}/U) \to \Sh(\mathcal{C})/h_U^\#
\to \Sh(\mathcal{C})$
が $j_{U!}$ であることを見た。関手
$\Sh(\mathcal{C})/h_U^\# \to \Sh(\mathcal{C})$
は補題 \ref{sites-lemma-localize-topos} により
$j_{\mathcal{F}!}$ である。したがって、この同一視のもとで
$j_{\mathcal{F}!} = j_{U!}$ である。ゆえに随伴性から
$j_\mathcal{F}^{-1} = j_U^{-1}$ であり、さらに随伴性から
$j_{\mathcal{F}, *} = j_{U, *}$ である。
\end{proof}

\begin{lemma}
\label{sites-lemma-relocalize-topos}
$\mathcal{C}$ をサイトとする。
$s : \mathcal{G} \to \mathcal{F}$ が $\mathcal{C}$ 上の層の射ならば、
トポスの射の自然な可換図式
$$
\xymatrix{
\Sh(\mathcal{C})/\mathcal{G} \ar[rd]_{j_\mathcal{G}} \ar[rr]_j & &
\Sh(\mathcal{C})/\mathcal{F} \ar[ld]^{j_\mathcal{F}} \\
& \Sh(\mathcal{C}) &
}
$$
が存在する。ここで $j = j_{\mathcal{G}/\mathcal{F}}$ は、対象
$\mathcal{G}/\mathcal{F}$ におけるトポス
$\Sh(\mathcal{C})/\mathcal{F}$ の局所化である。特に
$$
j^{-1}(\mathcal{H} \to \mathcal{F}) =
(\mathcal{H} \times_\mathcal{F} \mathcal{G} \to \mathcal{G})
$$
および
$$
j_!(\mathcal{E} \xrightarrow{e} \mathcal{G}) =
(\mathcal{E} \xrightarrow{s \circ e} \mathcal{F}).
$$
が成り立つ。
\end{lemma}

\begin{proof}
$j^{-1}$ と $j_!$ の記述は、補題
\ref{sites-lemma-localize-topos} におけるこれらの関手の記述から得られる。
関手の等式
$j_{\mathcal{G}!} = j_{\mathcal{F}!} \circ j_!$ は、これらの関手を
忘却関手として記述すれば明らかである。随伴性により、等式
$j_\mathcal{G}^{-1} = j^{-1} \circ j_\mathcal{F}^{-1}$ および
$j_{\mathcal{G}, *} = j_{\mathcal{F}, *} \circ j_*$
も得られる。
\end{proof}

\begin{lemma}
\label{sites-lemma-relocalize-compare}
$\mathcal{C}$ と $s : \mathcal{G} \to \mathcal{F}$ は
補題 \ref{sites-lemma-relocalize-topos} のとおりとする。
$\mathcal{G} = h_V^\#$、$\mathcal{F} = h_U^\#$ であり、
$s : \mathcal{G} \to \mathcal{F}$ が $\mathcal{C}$ の射
$V \to U$ から生じるならば、補題
\ref{sites-lemma-relocalize-topos} の図式は、図式
(\ref{sites-equation-relocalize}) と、同一視
$\Sh(\mathcal{C}/V) = \Sh(\mathcal{C})/h_V^\#$ および
$\Sh(\mathcal{C}/U) = \Sh(\mathcal{C})/h_U^\#$
のもとで同一視される。この同一視は補題
\ref{sites-lemma-essential-image-j-shriek} による。
\end{lemma}

\begin{proof}
これは $j^{-1}$ の記述が一致するためである。
補題 \ref{sites-lemma-relocalize-explicit} および
補題 \ref{sites-lemma-relocalize-topos} を参照されたい。
\end{proof}

\section{局所化とトポスの射}
\label{sites-section-localize-morphisms-topoi}

\noindent
本節は、トポスの射についての節
\ref{sites-section-localize-morphisms} の類似である。

\begin{lemma}
\label{sites-lemma-localize-morphism-topoi}
$f : \Sh(\mathcal{C}) \to \Sh(\mathcal{D})$
をトポスの射とし、$\mathcal{G}$ を $\mathcal{D}$ 上の層とする。
$\mathcal{F} = f^{-1}\mathcal{G}$ とおく。このときトポスの可換図式
$$
\xymatrix{
\Sh(\mathcal{C})/\mathcal{F} \ar[r]_{j_\mathcal{F}} \ar[d]_{f'} &
\Sh(\mathcal{C}) \ar[d]^f \\
\Sh(\mathcal{D})/\mathcal{G} \ar[r]^{j_\mathcal{G}} &
\Sh(\mathcal{D}).
}
$$
が存在する。射 $f'$ は
$$
(f')^{-1}(\mathcal{H} \xrightarrow{\varphi} \mathcal{G})
=
(f^{-1}\mathcal{H} \xrightarrow{f^{-1}\varphi} \mathcal{F})
$$
という性質によって特徴づけられ、さらに
$f'_*j_\mathcal{F}^{-1} = j_\mathcal{G}^{-1}f_*$ である。
\end{lemma}

\begin{proof}
この主張はトポスについてのもので、基礎となるサイトには言及しないので、
サイトを自由に変更してよい。したがって Remark
\ref{sites-remark-morphism-topoi-comes-from-morphism-sites} の議論により、
$f$ は、すべての有限極限をもち準標準位相をもつサイト間の、
命題 \ref{sites-proposition-get-morphism} の仮定を満たす連続関手
$u : \mathcal{D} \to \mathcal{C}$ によって与えられ、さらに
$\mathcal{G} = h_V$ が $\mathcal{D}$ のある対象 $V$ に対して
成り立つと仮定してよい。このとき補題
\ref{sites-lemma-pullback-representable-sheaf} により
$\mathcal{F} = f^{-1}h_V = h_{u(V)}$ である。
% P01-E0119
補題 \ref{sites-lemma-localize-morphism} により、トポスの射の可換図式
$$
\xymatrix{
\Sh(\mathcal{C}/U) \ar[r]_{j_U} \ar[d]_{f'} &
\Sh(\mathcal{C}) \ar[d]^f \\
\Sh(\mathcal{D}/V) \ar[r]^{j_V} &
\Sh(\mathcal{D}),
}
$$
を得て、$f'_*j_U^{-1} = j_V^{-1}f_*$ が成り立つ。
補題 \ref{sites-lemma-localize-compare} により、
$j_\mathcal{F}$ と $j_U$、および $j_\mathcal{G}$ と $j_V$ を
同一視できる。$(f')^{-1}$ の記述は補題
\ref{sites-lemma-localize-morphism} で与えられている。
\end{proof}

\begin{lemma}
\label{sites-lemma-localize-morphism-compare}
$f : \mathcal{C} \to \mathcal{D}$ を、連続関手
$u : \mathcal{D} \to \mathcal{C}$ により与えられるサイトの射とする。
$V$ を $\mathcal{D}$ の対象とし、$U = u(V)$ とおく。
$\mathcal{G} = h_V^\#$ および
$\mathcal{F} = h_U^\# = f^{-1}h_V^\#$ とおく（補題
\ref{sites-lemma-pullback-representable-sheaf} を参照）。
このとき、補題 \ref{sites-lemma-localize-morphism-topoi} の
トポスの射の図式は、補題 \ref{sites-lemma-localize-morphism} の
トポスの射の図式と、補題
\ref{sites-lemma-localize-compare} の同一視
$j_\mathcal{F}= j_U$ および $j_\mathcal{G} = j_V$ のもとで一致する。
\end{lemma}

\begin{proof}
これは完全に自明というわけではない。補題
\ref{sites-lemma-localize-morphism-topoi} の証明で選んだ
$f$ を生じるサイトの射は、この補題で与えられたサイトの射とは
異なるかもしれないからである。しかしどちらの場合も関手
$(f')^{-1}$ は同じ規則で記述される。したがって両者は一致し、
付随するトポスの射も同じである。いくつかの詳細は省略する。
\end{proof}

\begin{lemma}
\label{sites-lemma-relocalize-morphism-topoi}
$f : \Sh(\mathcal{C}) \to \Sh(\mathcal{D})$
をトポスの射とする。
$\mathcal{G} \in \Sh(\mathcal{D})$、
$\mathcal{F} \in \Sh(\mathcal{C})$ とし、
$s : \mathcal{F} \to f^{-1}\mathcal{G}$ を層の射とする。
トポスの可換図式
$$
\xymatrix{
\Sh(\mathcal{C})/\mathcal{F} \ar[r]_{j_\mathcal{F}} \ar[d]_{f_s} &
\Sh(\mathcal{C}) \ar[d]^f \\
\Sh(\mathcal{D})/\mathcal{G} \ar[r]^{j_\mathcal{G}} &
\Sh(\mathcal{D}).
}
$$
が存在する。ここで
$f_s = f' \circ j_{\mathcal{F}/f^{-1}\mathcal{G}}$ であり、
% P01-E0120
補題 \ref{sites-lemma-localize-morphism-topoi} における
$f' :
\Sh(\mathcal{C})/f^{-1}\mathcal{G}
\to
\Sh(\mathcal{D})/\mathcal{F}$
と、補題 \ref{sites-lemma-relocalize-topos} における
$j_{\mathcal{F}/f^{-1}\mathcal{G}} :
\Sh(\mathcal{C})/\mathcal{F}
\to
\Sh(\mathcal{C})/f^{-1}\mathcal{G}$
を用いた。関手 $(f_s)^{-1}$ は規則
$$
(f_s)^{-1}(\mathcal{H} \xrightarrow{\varphi} \mathcal{G})
=
(f^{-1}\mathcal{H} \times_{f^{-1}\varphi, f^{-1}\mathcal{G}, s} \mathcal{F}
\rightarrow \mathcal{F}).
$$
で記述される。最後に、任意の射
$b : \mathcal{G}' \to \mathcal{G}$、
$a : \mathcal{F}' \to \mathcal{F}$、および
$s' : \mathcal{F}' \to f^{-1}\mathcal{G}'$ が与えられ、図式
$$
\xymatrix{
\mathcal{F}' \ar[r]_-{s'} \ar[d]_a &
f^{-1}\mathcal{G}' \ar[d]^{f^{-1}b} \\
\mathcal{F} \ar[r]^-s & f^{-1}\mathcal{G}
}
$$
が可換ならば、図式
$$
\xymatrix{
\Sh(\mathcal{C})/\mathcal{F}'
\ar[r]_{j_{\mathcal{F}'/\mathcal{F}}} \ar[d]_{f_{s'}} &
\Sh(\mathcal{C})/\mathcal{F} \ar[d]^{f_s} \\
\Sh(\mathcal{D})/\mathcal{G}' \ar[r]^{j_{\mathcal{G}'/\mathcal{G}}} &
\Sh(\mathcal{D})/\mathcal{G}.
}
$$
は可換である。
\end{lemma}

\begin{proof}
第一の正方形の可換性は、補題
\ref{sites-lemma-relocalize-topos} の図式の可換性と、
補題 \ref{sites-lemma-localize-morphism-topoi} の図式の可換性から従う。
$f_s^{-1}$ の記述は、補題
\ref{sites-lemma-localize-morphism-topoi} の $(f')^{-1}$ の記述と、
補題 \ref{sites-lemma-relocalize-topos} の
$(j_{\mathcal{F}/f^{-1}\mathcal{G}})^{-1}$ の記述を組み合わせれば従う。
最後の正方形の可換性は、形式的な等式
$$
f^{-1}\mathcal{H} \times_{f^{-1}\mathcal{G}, s} \mathcal{F}
\times_\mathcal{F} \mathcal{F}'
=
f^{-1}(\mathcal{H} \times_\mathcal{G} \mathcal{G}')
\times_{f^{-1}\mathcal{G}', s'} \mathcal{F}'
$$
から従う。
\end{proof}

\begin{lemma}
\label{sites-lemma-relocalize-morphism-compare}
$f : \mathcal{C} \to \mathcal{D}$ を、連続関手
$u : \mathcal{D} \to \mathcal{C}$ により与えられるサイトの射とする。
$V$ を $\mathcal{D}$ の対象とし、$c : U \to u(V)$ を射とする。
% P01-E0121
$\mathcal{G} = h_V^\#$ および
$\mathcal{F} = h_U^\# = f^{-1}h_V^\#$ とおく。
$s : \mathcal{F} \to f^{-1}\mathcal{G}$ を $c$ により誘導される写像とする。
このとき、補題 \ref{sites-lemma-relocalize-morphism} の
トポスの射の図式は、補題
\ref{sites-lemma-relocalize-morphism-topoi} のトポスの射の図式と、
補題 \ref{sites-lemma-localize-compare} の同一視
$j_\mathcal{F} = j_U$ および $j_\mathcal{G} = j_V$ のもとで一致する。
\end{lemma}

\begin{proof}
これは補題 \ref{sites-lemma-relocalize-compare} と
\ref{sites-lemma-localize-morphism-compare} を組み合わせれば従う。
\end{proof}

\section{点}
\label{sites-section-points}

\begin{definition}
\label{sites-definition-point-topos}
$\mathcal{C}$ をサイトとする。
{\it トポス $\Sh(\mathcal{C})$ の点}とは、
$\Sh(pt)$ から $\Sh(\mathcal{C})$ へのトポスの射 $p$ のことである。
\end{definition}

\noindent
サイトの点を関手
$u : \mathcal{C} \to \textit{Sets}$ によって定義する。
後で $u$ がサイトの射を定め、それが $\mathcal{C}$ に付随する
トポスの点を生じることが分かる。補題
\ref{sites-lemma-site-point-morphism} を参照されたい。

\medskip\noindent
$\mathcal{C}$ をサイトとし、$p = u$ を関手
$u : \mathcal{C} \to \textit{Sets}$ とする。
この奇妙な言い方を導入するのは、関手の近傍を語るのは不自然に
思われるからである。そこで「点」$p$ を、関手 $u$ という圏論的データで
与えられる幾何学的なものと考える。
$p$ が実際には $u$ に等しいという事実は問題にならない。
$p$ の {\it 近傍} とは、$U \in \Ob(\mathcal{C})$ かつ
$x \in u(U)$ である対 $(U, x)$ のことである。
{\it 近傍の射} $(V, y) \to (U, x)$ とは、
$u(\alpha)(y) = x$ を満たす $\mathcal{C}$ の射
$\alpha :V \to U$ によって与えられるものをいう。
近傍の圏は「大きな」圏ではないことに注意する。

\medskip\noindent
前層 $\mathcal{F}$ の $p$ における {\it 茎} を
\begin{equation}
\label{sites-equation-stalk}
\mathcal{F}_p = \colim_{\{(U, x)\}^{opp}} \mathcal{F}(U).
\end{equation}
と定める。この余極限は $p$ の近傍の圏の反対圏上で取る。
言い換えると、$\mathcal{F}_p$ の元は三つ組 $(U, x, s)$ によって
与えられる。ここで $(U, x)$ は $p$ の近傍で、
$s \in \mathcal{F}(U)$ である。三つ組の等しさは、上のような
$\alpha$ に対する
$(U, x, s) \sim (V, y, \alpha^*s)$ で生成される同値関係である。

\medskip\noindent
$\varphi : \mathcal{F} \to \mathcal{G}$ が集合の前層の射ならば、
茎の標準写像
$\varphi_p : \mathcal{F}_p \to \mathcal{G}_p$ を得る。したがって
{\it 茎関手}
$$
\textit{PSh}(\mathcal{C}) \longrightarrow \textit{Sets}, \quad
\mathcal{F} \longmapsto \mathcal{F}_p.
$$
を得る。茎関手は任意の関手
$p = u : \mathcal{C} \to \textit{Sets}$ を用いて定義した。
定義が機能するための条件は不要である\footnote{$u(U) = \emptyset$
がすべての $U$ に対して成り立つ場合は避けるべきである。}。
一方、一般には茎関手はよい性質をもたないので、$p$ が実際に点
（下の定義を参照）でない限り、この概念を使わない方がよいだろう。

\begin{definition}
\label{sites-definition-point}
$\mathcal{C}$ をサイトとする。
サイト $\mathcal{C}$ の {\it 点 $p$} とは、次を満たす関手
$u : \mathcal{C} \to \textit{Sets}$ によって与えられるものをいう。
\begin{enumerate}
\item $\mathcal{C}$ の各被覆 $\{U_i \to U\}$ に対して、写像
$\coprod u(U_i) \to u(U)$ は全射である。
\item $\mathcal{C}$ の各被覆 $\{U_i \to U\}$ と各射
$V \to U$ に対して、写像
$u(U_i \times_U V) \to u(U_i) \times_{u(U)} u(V)$ は全単射である。
\item 茎関手 $\Sh(\mathcal{C}) \to \textit{Sets}$、
$\mathcal{F} \mapsto \mathcal{F}_p$ は左完全である。
\end{enumerate}
\end{definition}

\noindent
これらの条件は、定義 \ref{sites-definition-continuous} と定義
\ref{sites-definition-morphism-sites} の条件を模したものなので、
見覚えがあるはずである。(3) から $*_p = \{*\}$ が従うことに
注意する。例 \ref{sites-example-singleton-sheaf} を参照されたい。
したがって少なくともある $U$ について
$u(U) \not = \emptyset$ である（空な余極限は空集合を与える）。
これが実際にトポス $\Sh(\mathcal{C})$ の点を生じることを
下で示す（補題 \ref{sites-lemma-point-site-topos}）。
その前に、一般の関手 $u$ についていくつかの補題を証明する。

\begin{lemma}
\label{sites-lemma-points-recover}
$\mathcal{C}$ をサイトとし、
$p = u : \mathcal{C} \to \textit{Sets}$ を関手とする。
$U \in \Ob(\mathcal{C})$ に対して、関手的な同型
$(h_U)_p = u(U)$ が存在する。
\end{lemma}

\begin{proof}
$(h_U)_p$ の元は三つ組 $(V, y, f)$ によって与えられる。
ここで $V \in \Ob(\mathcal{C})$、$y\in u(V)$、および
$f \in h_U(V) = \Mor_\mathcal{C}(V, U)$ である。
% P01-E0122
このような二つの三つ組 $(V, y, f)$、$(V', y', f')$ は、
$u(\phi)(y) = y'$ かつ $f' \circ \phi = f$ を満たす射
$\phi : V \to V'$ が存在すれば同じ元を定める。一般には、正しい
同値関係を得るため、このような恒等関係の鎖を取らなければならない。
いずれにせよ、各 $(V, y, f)$ は元
$(U, u(f)(y), \text{id}_U)$ と同値である。上のような $\phi$ が
存在すれば、三つ組 $(V, y, f)$ と $(V', y', f')$ は同じ元
$(U, u(f)(y), \text{id}_U) = (U, u(f')(y'), \text{id}_U)$ を定める。
これにより写像
$u(U) \to (h_U)_p$、
$x \mapsto \text{class of }(U, x, \text{id}_U)$
が全単射であることが示された。
\end{proof}

\noindent
$\mathcal{C}$ をサイトとし、
$p = u : \mathcal{C} \to \textit{Sets}$ を関手とする。
節 \ref{sites-section-functoriality-PSh} の構成と同様に、集合 $E$ が
与えられたとき、前層 $u^pE$ を規則
\begin{equation}
\label{sites-equation-skyscraper}
U
\longmapsto
u^pE(U) = \Mor_{\textit{Sets}}(u(U), E) = \text{Map}(u(U), E).
\end{equation}
によって定める。これは関手
$u^p : \textit{Sets} \to \textit{PSh}(\mathcal{C})$、
$E \mapsto u^pE$ を定める。

\begin{lemma}
\label{sites-lemma-adjoint-point-push-stalk}
% P01-E0123
任意の関手 $u : \mathcal{C} \to \textit{Sets}$ に対して、
関手 $u^p$ は前層上の茎関手の右随伴である。
\end{lemma}

\begin{proof}
$\mathcal{F}$ を $\mathcal{C}$ 上の前層とし、$E$ を集合とする。
射 $\mathcal{F} \to u^pE$ は写像の両立する系
$\mathcal{F}(U) \to \text{Map}(u(U), E)$、すなわち写像の両立する系
$\mathcal{F}(U) \times u(U) \to E$ によって与えられる。
一方、$\mathcal{F}_p$ の定義により、写像
$\mathcal{F}_p \to E$ は、各三つ組 $(U, x, \sigma)$ に
$E$ の元を対応させ、同値な三つ組を同じ元に移す規則によって
与えられる。式 (\ref{sites-equation-stalk}) の周りの議論を参照されたい。
これは同じく、写像の両立する系
$\mathcal{F}(U) \times u(U) \to E$ を意味する。
\end{proof}

\noindent
節 \ref{sites-section-continuous-functors} と同様に、次の補題が成り立つ。

\begin{lemma}
\label{sites-lemma-point-pushforward-sheaf}
$\mathcal{C}$ をサイトとし、
$p = u : \mathcal{C} \to \textit{Sets}$ を関手とする。
$\mathcal{C}$ の各被覆 $\{U_i \to U\}$ に対して次を仮定する。
\begin{enumerate}
\item 写像 $\coprod u(U_i) \to u(U)$ は全射である。
\item 写像
$u(U_i \times_U U_j) \to u(U_i) \times_{u(U)} u(U_j)$ は全射である。
\end{enumerate}
このとき次が成り立つ。
\begin{enumerate}
\item 任意の集合 $E$ に対して、前層 $u^pE$ は層である。
これを $u^sE$ と書く。
\item 茎関手 $\Sh(\mathcal{C}) \to \textit{Sets}$ と関手
$u^s: \textit{Sets} \to \Sh(\mathcal{C})$ は随伴である。
\item 集合の各前層 $\mathcal{F}$ に対して
$\mathcal{F}_p = \mathcal{F}^\#_p$ である。
\end{enumerate}
\end{lemma}

\begin{proof}
第一の主張は、層の定義、仮定 (1) と (2)、および
$u^pE$ の定義から直ちに従う。第二の主張は、$u^p$ と茎関手の
随伴性（補題 \ref{sites-lemma-adjoint-point-push-stalk}）を層に制限した
ものにほかならない。第三の主張は、任意の集合 $E$ に対して
$$
\text{Map}(\mathcal{F}_p, E) =
\Mor_{\textit{PSh}(\mathcal{C})}(\mathcal{F}, u^pE) =
\Mor_{\Sh(\mathcal{C})}(\mathcal{F}^\#, u^sE) =
\text{Map}(\mathcal{F}^\#_p, E)
$$
となることから従う。ここでは層化の随伴性を用いた。
\end{proof}

\noindent
特に、$p = u$ が点ならば補題
\ref{sites-lemma-point-pushforward-sheaf} が成り立つ。この場合、
層 $u^sE$ を、$p$ において値 $E$ をもつ「摩天楼」層と考える。

\begin{definition}
\label{sites-definition-pushforward-point}
$p$ を、関手 $u$ によって与えられるサイト $\mathcal{C}$ の点とする。
集合 $E$ に対し、
% P01-E0124
$p_*E = u^sE$ を、上の補題
\ref{sites-lemma-point-pushforward-sheaf} で記述した層として定める。
これを {\it 摩天楼層} と呼ぶこともある。
\end{definition}

\noindent
特に、摩天楼層と茎について次の随伴性をもつ。
$$
\Mor_{\Sh(\mathcal{C})}(\mathcal{F}, p_*E)
=
\text{Map}(\mathcal{F}_p, E)
$$
これは、ときどき用いる記法
$p^{-1}\mathcal{F} = \mathcal{F}_p$ の動機となる。

\begin{lemma}
\label{sites-lemma-point-site-topos}
$\mathcal{C}$ をサイトとする。
\begin{enumerate}
\item $p$ をサイト $\mathcal{C}$ の点とする。このとき上で導入した
関手の対 $(p_*, p^{-1})$ はトポスの射
$\Sh(pt) \to \Sh(\mathcal{C})$ を定める。
\item $p = (p_*, p^{-1})$ をトポス
$\Sh(\mathcal{C})$ の点とする。このとき関手
$u : U \mapsto p^{-1}(h_U^\#)$ はサイト $\mathcal{C}$ の点
$p'$ を生じ、それに付随するトポスの射
$(p'_*, (p')^{-1})$ は $p$ に等しい。
\end{enumerate}
\end{lemma}

\begin{proof}
(1) の証明。上により、関手 $p_*$ と $p^{-1}$ は随伴である。
定義 \ref{sites-definition-point} により、関手 $p^{-1}$ は
完全であることを要求される。したがって定義
\ref{sites-definition-topos} で課した条件はすべて満たされ、(1) が成り立つ。

\medskip\noindent
(2) の証明。$\{U_i \to U\}$ を $\mathcal{C}$ の被覆とする。
すると補題 \ref{sites-lemma-covering-surjective-after-sheafification} により
$\coprod (h_{U_i})^\# \to h_U^\#$ は全射である。
$p^{-1}$ は完全なので（トポスの射の定義による）、
$\coprod u(U_i) \to u(U)$ は全射である。
これで定義 \ref{sites-definition-point} の (1) が示された。
層化は完全である。補題
\ref{sites-lemma-sheafification-exact} を参照されたい。
したがって $U \times_V W$ が $\mathcal{C}$ に存在すれば
$$
h_{U \times_V W}^\# =  h_U^\# \times_{h_V^\#} h_W^\#
$$
である。$p^{-1}$ は完全なので、
$u(U \times_V W) = u(U) \times_{u(V)} u(W)$ であることが分かる。
これで定義 \ref{sites-definition-point} の (2) が示された。
$p' = u$ とし、$\mathcal{F}_{p'}$ を $u$ を用いて式
(\ref{sites-equation-stalk}) で定義した茎関手とする。
標準的な比較写像
$c : \mathcal{F}_{p'} \to \mathcal{F}_p = p^{-1}\mathcal{F}$
が存在する。実際、$\mathcal{F}_{p'}$ の元 $\xi$ を表す三つ組
$(U, x, \sigma)$ が与えられたとき、$\sigma$ を写像
$\sigma : h_U^\# \to \mathcal{F}$ と考え、
$x \in u(U) = p^{-1}(h_U^\#)$ なので
$c(\xi) = p^{-1}(\sigma)(x)$ とおける。補題
\ref{sites-lemma-points-recover} により
$(h_U)_{p'} = u(U)$ である。定義
\ref{sites-definition-point} の条件 (1) と (2) は $p'$ について
成り立つので、補題 \ref{sites-lemma-point-pushforward-sheaf} により
$(h_U^\#)_{p'} = (h_U)_{p'}$ でもある。したがって
$$
(h_U^\#)_{p'} = (h_U)_{p'} = u(U) = p^{-1}(h_U^\#)
$$
である。この全単射は比較写像
$c : (h_U^\#)_{p'} \to p^{-1}(h_U^\#)$ に等しいと主張する
（確認は省略する）。$\mathcal{C}$ 上の任意の層は、
$h_U^\#$ という形の層の余積間の写像の余等化子である。
補題 \ref{sites-lemma-sheaf-coequalizer-representable} を参照されたい。
茎関手 $\mathcal{F} \mapsto \mathcal{F}_{p'}$ と関手 $p^{-1}$ は、
ともに左随伴なので任意の余極限と可換する。
したがって $c$ はすべての層 $\mathcal{F}$ に対して同型である。
ゆえに茎関手 $\mathcal{F} \mapsto \mathcal{F}_{p'}$ は
$p^{-1}$ と同型であり、特に完全である。
これで定義 \ref{sites-definition-point} の条件 (3) が成り立ち、
$p'$ は点であることが示された。最後の主張も既に示されている。
実際、$p^{-1} = (p')^{-1}$ を見た。
\end{proof}

\noindent
実は、次の補題で示すように、点は常にサイトの射に対応する。

\begin{lemma}
\label{sites-lemma-site-point-morphism}
$\mathcal{C}$ をサイトとし、$p$ を関手
$u : \mathcal{C} \to \textit{Sets}$ によって与えられる
$\mathcal{C}$ の点とする。すべての $U \in \Ob(\mathcal{C})$ に
対して $u(U) \subset S_0$ となるような無限集合 $S_0$ を取る。
$\mathcal{S}$ を、Remark \ref{sites-remark-pt-topos} において冪集合
$S = \mathcal{P}(S_0)$ から構成したサイトとする。このとき
\begin{enumerate}
\item 同値 $i : \Sh(pt) \to \Sh(\mathcal{S})$ が存在する。
\item 関手 $u : \mathcal{C} \to \mathcal{S}$ はサイトの射
$f : \mathcal{S} \to \mathcal{C}$ を誘導する。
\item 合成
$$
\Sh(pt) \to
\Sh(\mathcal{S}) \to
\Sh(\mathcal{C})
$$
は補題 \ref{sites-lemma-point-site-topos} のトポスの射
$(p_*, p^{-1})$ である。
\end{enumerate}
\end{lemma}

\begin{proof}
(1) は Remark \ref{sites-remark-pt-topos} で見た。
さらに、この同値は集合 $E$ に、規則
$V \mapsto \Mor_{\textit{Sets}}(V, E)$ で定義される
$\mathcal{S}$ 上の層 $i_*E$ を対応させることを思い出そう。
(2) は、$\mathcal{C}$ の点の定義
（定義 \ref{sites-definition-point}）とサイトの射の定義
（定義 \ref{sites-definition-morphism-sites}）から明らかである。
最後に $f_*i_*E$ を考える。構成により
$$
f_*i_*E(U) = i_*E(u(U)) = \Mor_{\textit{Sets}}(u(U), E)
$$
であり、これは式 (\ref{sites-equation-skyscraper}) により
$p_*E(U)$ に等しい。これで (3) が示された。
\end{proof}

\noindent
位相空間の場合とは異なり、値 $E$ をもつ摩天楼層の茎が
$E$ であるとは限らない。一般には次が成り立つ。

\begin{lemma}
\label{sites-lemma-stalk-skyscraper}
$\mathcal{C}$ をサイトとし、
$p : \Sh(pt) \to \Sh(\mathcal{C})$ を
$\mathcal{C}$ に付随するトポスの点とする。
任意の集合 $E$ に対して標準写像
$$
E \longrightarrow (p_*E)_p \longrightarrow E
$$
が存在し、その合成は $\text{id}_E$ である。
\end{lemma}

\begin{proof}
随伴写像 $(p_*E)_p = p^{-1}p_*E \to E$ は常に存在する。
$E = \{*\}$ のとき、この写像は同型である。実際、$p_*$ と
$p^{-1}$ はともに左完全なので、終対象を終対象に移す。
したがって $e \in E$ が与えられたとき、写像
$i_e : \{*\} \to E$ を考えることができ、これは
$$
\xymatrix{
p^{-1}p_*\{*\} \ar[rr]_{p^{-1}p_*i_e} \ar[d]_{\cong} & &
p^{-1}p_*E \ar[d] \\
\{*\} \ar[rr]^{i_e} & & E
}
$$
を与える。これから所望の写像
$E \to (p_*E)_p = p^{-1}p_*E$ が得られる。
\end{proof}

\begin{lemma}
\label{sites-lemma-skyscraper-functor-exact}
$\mathcal{C}$ をサイトとし、
$p : \Sh(pt) \to \Sh(\mathcal{C})$ を
$\mathcal{C}$ に付随するトポスの点とする。
関手 $p_* : \textit{Sets} \to \Sh(\mathcal{C})$ は次の性質をもつ。
任意の極限と可換し、左完全で、忠実で、全射を全射に移し、
余等化子と可換し、単射を反映し、全射を反映し、同型を反映する。
\end{lemma}

\begin{proof}
$p_*$ は右随伴なので、任意の極限と可換し、左完全である。
$p^{-1}p_*E \to E$ が標準的に分裂する全射であることから、
$p_*$ が忠実であり、単射、全射、および同型を反映することが従う。
補題 \ref{sites-lemma-point-site-topos} により、
$p$ は基礎となるサイト $\mathcal{C}$ の点
$u : \mathcal{C} \to \textit{Sets}$ から生じると仮定してよい。
この場合、層 $p_*E$ は
$$
p_*E(U) = \Mor_{\textit{Sets}}(u(U), E)
$$
で与えられる。式 (\ref{sites-equation-skyscraper}) と定義
\ref{sites-definition-pushforward-point} を参照されたい。
この公式から、$p_*$ が全射を全射に移し、
余等化子を余等化子に移すことが直ちに従う。
\end{proof}

\section{点の構成}
\label{sites-section-construct-points}

\noindent
本節では、サイトから集合の圏への関手が、そのサイトの点を定めるための
判定条件を与える。

\begin{lemma}
\label{sites-lemma-neighbourhoods-cofiltered}
$\mathcal{C}$ をサイトとし、
$p = u : \mathcal{C} \to \textit{Sets}$ を関手とする。
$p$ の近傍の圏が余フィルターならば、茎関手
(\ref{sites-equation-stalk}) は左完全である。
\end{lemma}

\begin{proof}
$\mathcal{I} \to \Sh(\mathcal{C})$、
$i \mapsto \mathcal{F}_i$ を層の有限図式とする。
この系の極限の茎が、茎の極限と一致することを示さなければならない。
$\mathcal{F}$ を、この系の {\it 前層} としての極限とする。
補題 \ref{sites-lemma-limit-sheaf} によれば、これは層であり、層の圏における
極限でもある。したがって
$\mathcal{F}_p = \lim_\mathcal{I} \mathcal{F}_{i, p}$
を示せばよい。さらに $\mathcal{F}$ は簡単な記述をもつことを
思い出そう。節 \ref{sites-section-limits-colimits-PSh} を参照されたい。
よって示すべきことは
$$
\lim_i \colim_{\{(U, x)\}^{opp}} \mathcal{F}_i(U)
=
\colim_{\{(U, x)\}^{opp}} \lim_i \mathcal{F}_i(U).
$$
である。仮定により近傍の圏の反対圏はフィルターなので、
これは Categories, Lemma
\ref{categories-lemma-directed-commutes} から従う。
\end{proof}

\begin{lemma}
\label{sites-lemma-neighbourhoods-directed}
$\mathcal{C}$ をサイトとし、$\mathcal{C}$ は終対象 $X$ と
ファイバー積をもつと仮定する。
$p = u : \mathcal{C} \to \textit{Sets}$ を次を満たす関手とする。
\begin{enumerate}
\item $u(X)$ は一元集合である。
\item 同じ終域をもつ射 $U \to W$ と $V \to W$ の各対に対して、
写像 $u(U \times_W V) \to u(U) \times_{u(W)} u(V)$ は全単射である。
\end{enumerate}
% P01-E0125
このとき $p$ の近傍の圏は余フィルターであり、したがって茎関手
$\Sh(\mathcal{C}) \to \textit{Sets}$、
$\mathcal{F} \to \mathcal{F}_p$ は有限極限と可換する。
\end{lemma}

\begin{proof}
補題 \ref{sites-lemma-directed} との類似に注意されたい。
$\mathcal{C}$ に関する仮定から、$\mathcal{C}$ は有限極限をもつ。
Categories, Lemma \ref{categories-lemma-finite-limits-exist} を
参照されたい。仮定 (1) から、近傍の圏は空でない。
$(U, x)$ と $(V, y)$ を近傍とする。このとき (2) により
$u(U \times V) = u(U \times_X V) =
u(U) \times_{u(X)} u(V) = u(U) \times u(V)$
である。したがって、$(U, x)$ と $(V, y)$ の両方に写る近傍
$(U \times_X V, z)$ が存在する。
$a, b : (V, y) \to (U, x)$ を近傍の圏の二つの射とする。
$W$ を $a, b : V \to U$ の等化子とする。
Categories, Lemma \ref{categories-lemma-finite-limits-exist}
の証明と同様に、$W$ をファイバー積によって
$$
W = (V \times_{a, U, b} V) \times_{(pr_1, pr_2), V \times V, \Delta} V
$$
と書ける。(2) の全単射性により、
$((y, y), y)$ に写る元 $z \in u(W)$ が存在する。
すると $(W, z) \to (V, y)$ は所望のように $a, b$ を等化する。
「したがって」の部分は補題
\ref{sites-lemma-neighbourhoods-cofiltered} である。
\end{proof}

\begin{proposition}
\label{sites-proposition-point-limits}
$\mathcal{C}$ をサイトとし、$\mathcal{C}$ に有限極限が存在すると
仮定する（すなわち $\mathcal{C}$ はファイバー積と終対象をもつ）。
このようなサイト $\mathcal{C}$ の点 $p$ は、次を満たす関手
$u : \mathcal{C} \to \textit{Sets}$ によって与えられる。
\begin{enumerate}
\item $u$ は有限極限と可換する。
\item $\{U_i \to U\}$ が被覆ならば、
$\coprod_i u(U_i) \to u(U)$ は全射である。
\end{enumerate}
\end{proposition}

\begin{proof}
まず $p$ が関手 $u$ によって与えられる点
（定義 \ref{sites-definition-point}）であると仮定する。
条件 (2) は点の定義から直接満たされる。
補題 \ref{sites-lemma-points-recover} により
$(h_U)_p = u(U)$ である。補題
\ref{sites-lemma-point-pushforward-sheaf} により
$(h_U^\#)_p = (h_U)_p$ である。したがって $u$ は関手の合成
$$
\mathcal{C} \xrightarrow{h}
\textit{PSh}(\mathcal{C}) \xrightarrow{{}^\#}
\Sh(\mathcal{C}) \xrightarrow{()_p}
\textit{Sets}
$$
に等しい。これらの関手はすべて左完全なので、$u$ は (1) を満たす。

\medskip\noindent
逆に、$u$ が (1) と (2) を満たすと仮定する。
この場合、$u$ が定義 \ref{sites-definition-point} の最初の二条件を
満たすことは直ちに分かる。また、その茎関手は完全である。
実際、補題 \ref{sites-lemma-point-pushforward-sheaf} により左随伴であり、
補題 \ref{sites-lemma-neighbourhoods-directed} により有限極限と可換する。
\end{proof}

\begin{remark}
\label{sites-remark-improve-proposition-points-limits}
実際、$\mathcal{C}$ をサイトとし、$\mathcal{C}$ は終対象 $X$ と
ファイバー積をもつと仮定する。
$p = u: \mathcal{C} \to \textit{Sets}$ を次を満たす関手とする。
\begin{enumerate}
% P01-E0126
\item $u(X) = \{*\}$ は一元集合である。
\item 同じ終域をもつ射 $U \to W$ と $V \to W$ の各対に対して、
写像 $u(U \times_W V) \to u(U) \times_{u(W)} u(V)$ は全射である。
\item 各被覆 $\{U_i \to U\}$ に対して、
写像 $\coprod u(U_i) \to u(U)$ は全射である。
\end{enumerate}
このとき一般には、$p$ は $\mathcal{C}$ の点では {\bf ない}。
例として、二つの対象 $\{U, X\}$ と、ただ一つの非恒等矢印
$U \to X$ をもつ圏 $\mathcal{C}$ を考える。
$\mathcal{C}$ に自明な位相、すなわち被覆が
$\{U \to U\}$ と $\{X \to X\}$ だけである位相を入れる。
層 $\mathcal{F}$ は前層と同じもので、三つ組
$(A, B, A \to B)$ からなる。すなわち
$A = \mathcal{F}(X)$、$B = \mathcal{F}(U)$ であり、
$A \to B$ は $U \to X$ に対応する制限写像である。
$U \times_X U = U$ なのでファイバー積が存在することに注意する。
$u(X) = \{*\}$、$u(U) = \{*_1, *_2\}$ である関手
$u = p$ を考える。これは (1)、(2)、(3) を満たすが、
対応する茎関手 (\ref{sites-equation-stalk}) は
$$
(A, B, A \to B) \longmapsto B \amalg_A B
$$
であり、完全ではない。実際、
$(\emptyset, \{1\}, \emptyset \to \{1\}) \to
(\{1\}, \{1\}, \{1\} \to \{1\})$
を考えよう。これは層の単射だが、茎関手により集合の非単射
$$
\{1\} \amalg \{1\} \longrightarrow \{1\} \amalg_{\{1\}} \{1\}
$$
に移される。
\end{remark}

\begin{example}
\label{sites-example-point-topological}
$X$ を位相空間とし、$X_{Zar}$ を例
\ref{sites-example-site-topological} のサイトとする。
$x \in X$ を点とし、関手
$$
u : X_{Zar} \longrightarrow \textit{Sets}, \quad
U \mapsto
\left\{
\begin{matrix}
\emptyset & \text{if} & x \not \in U \\
\{*\} & \text{if} & x \in U
\end{matrix}
\right.
$$
を考える。この関手は積およびファイバー積と可換し、
被覆を写像の全射族に移す。したがってサイト $X_{Zar}$ の点
$p$ を得る。茎関手が Sheaves, Section
\ref{sheaves-section-stalks} で定義した $x$ における茎と
一致することは直ちに確かめられる。
\end{example}

\begin{example}
\label{sites-example-point-topology}
$X$ を位相空間とする。トポス $\Sh(X)$ の点は何であろうか。
これを見るため、$X_{Zar}$ を例
\ref{sites-example-site-topological} のサイトとする。
補題 \ref{sites-lemma-point-site-topos} により、
$\Sh(X)$ の点はこのサイトの点に対応する。
$p$ を、関手 $u : X_{Zar} \to \textit{Sets}$ によって
与えられるサイト $X_{Zar}$ の点とする。
命題 \ref{sites-proposition-point-limits} におけるこのような $u$ の
特徴づけを用いる。これから直ちに
$u(\emptyset) = \emptyset$ および
$u(U \cap V) = u(U) \times u(V)$ が従う。特に対角写像により
$u(U) = u(U) \times u(U)$ なので、$u(U)$ は一元集合または空集合である。
さらに $U = \bigcup U_i$ が開被覆ならば
$$
u(U) = \emptyset \Rightarrow \forall i, \ u(U_i) = \emptyset
\quad\text{and}\quad
u(U) \not = \emptyset \Rightarrow \exists i, \ u(U_i) \not = \emptyset.
$$
である。$u(W) = \emptyset$ を満たす一意な最大の開集合
$W \subset X$ が存在することが分かる。これは
$u(U) = \emptyset$ を満たすすべての開集合 $U$ の合併である。
$Z = X \setminus W$ とおく。$Z = Z_1 \cup Z_2$ で
$Z_i \subset Z$ が閉ならば、
$W = (X \setminus Z_1) \cap (X \setminus Z_2)$ である。よって
$\emptyset = u(W) = u(X \setminus Z_1) \times u(X \setminus Z_2)$
なので、$u(X \setminus Z_1) = \emptyset$ または
$u(X \setminus Z_2) = \emptyset$ である。
これは $X \setminus Z_1 = W$ または $X \setminus Z_2 = W$
を意味する。言い換えると、$Z$ は既約である。
ここで $u$ は規則
$$
u : X_{Zar} \longrightarrow \textit{Sets}, \quad
U \mapsto
\left\{
\begin{matrix}
\emptyset & \text{if} & Z \cap U = \emptyset \\
\{*\} & \text{if} & Z \cap U \not = \emptyset
\end{matrix}
\right.
$$
によって記述される。
任意の既約閉集合 $Z \subset X$ に対して、この関手は命題
\ref{sites-proposition-point-limits} の仮定 (1)、(2) を満たし、
したがって点を定めることに注意する。
言い換えると、サイト $X_{Zar}$ の点は $X$ の既約閉部分集合と
一対一に対応する。特に $X$ がソバー位相空間ならば、
$X_{Zar}$ の点と $X$ の点は一対一に対応する。
例 \ref{sites-example-point-topological} を参照されたい。
\end{example}

\begin{example}
\label{sites-example-point-G-sets}
例 \ref{sites-example-site-on-group} および節
\ref{sites-section-example-sheaf-G-sets} で記述したサイト
$\mathcal{T}_G$ を考える。忘却関手
$u : \mathcal{T}_G \to \textit{Sets}$ は積およびファイバー積と
可換し、被覆を全射族に移す。したがって
$\mathcal{T}_G$ の点を定める。
$\Sh(\mathcal{T}_G)$ と $G\textit{-Sets}$ を同一視する。
茎関手
$$
p^{-1} :
\Sh(\mathcal{T}_G) = G\textit{-Sets}
\longrightarrow
\textit{Sets}
$$
は忘却関手である。順像 $p_*$ は、集合 $S$ を
$G$-集合 $\text{Map}(G, S)$ に移す関手
$$
\textit{Sets}
\longrightarrow
\Sh(\mathcal{T}_G) = G\textit{-Sets}
$$
である。ここで作用は $g \cdot \psi = \psi \circ R_g$ であり、
$R_g$ は右乗法である。特に、集合として
$p^{-1}p_*S = \text{Map}(G, S)$ であり、補題
\ref{sites-lemma-stalk-skyscraper} の写像
$S \to \text{Map}(G, S) \to S$ は明らかなものである。
\end{example}

\begin{example}
\label{sites-example-indiscrete-points}
$\mathcal{C}$ をカオス位相を入れた圏とする
（例 \ref{sites-example-indiscrete}）。
$\mathcal{C}$ の各対象 $U_0$ に対して、関手
$u : U \mapsto \Mor_\mathcal{C}(U_0, U)$ は
$\mathcal{C}$ の点 $p$ を定める。
実際、被覆は恒等写像によって与えられるものだけなので、
定義 \ref{sites-definition-point} の条件 (1) と (2) は直ちに成り立つ。
% P01-E0127
条件 (2) が成り立つのは
$\mathcal{F}_p = \mathcal{F}(U_0)$ であり、さらに
% P01-E0128
位相が離散的なので $U_0$ 上の切断を取ることが完全関手だからである。
\end{example}

\section{点とトポスの射}
\label{sites-section-functorial-points}

\noindent
本節では、点とトポスの射についていくつか注意を述べる。

\begin{lemma}
\label{sites-lemma-point-functor}
$u : \mathcal{C} \to \mathcal{D}$ を関手とし、
$v : \mathcal{D} \to \textit{Sets}$ を関手として
$w = v \circ u$ とおく。$v$ と $w$ にそれぞれ付随する茎関手
(\ref{sites-equation-stalk}) をそれぞれ $q$ と $p$ と書く。
このとき $\mathcal{C}$ 上の前層 $\mathcal{F}$ に関手的に
$(u_p\mathcal{F})_q = \mathcal{F}_p$ である。
\end{lemma}

\begin{proof}
これは単純な圏論的事実である。
\begin{align*}
(u_p\mathcal{F})_q
& =
\colim_{(V, y)} \colim_{U, \phi : V \to u(U)} \mathcal{F}(U) \\
& = \colim_{(V, y, U, \phi : V \to u(U))} \mathcal{F}(U) \\
& = \colim_{(U, x)} \mathcal{F}(U) \\
& = \mathcal{F}_p
\end{align*}
である。第一の等式は $u_p$ の定義と茎関手の定義による。
$y \in v(V)$ であることに注意する。第二の等式では単に余極限を
まとめた。第三の等式を見るため、図式圏の関手 $F$ を規則
$$
F((V, y, U, \phi : V \to u(U))) = (U, v(\phi)(y)).
$$
によって定め、Categories, Lemma
\ref{categories-lemma-colimit-constant-connected-fibers} を適用する。
$w(U) = v(u(U))$ なのでこれは意味をもつ。
Categories, Lemma
\ref{categories-lemma-colimit-constant-connected-fibers} の仮定を確認しよう。
$F$ は右逆
$(U, x) \mapsto (u(U), x, U, \text{id} : u(U) \to u(U))$
をもつことに注意する。$x \in w(U) = v(u(U))$ なので、これも
意味をもつ。一方、$(U, x)$ 上のファイバー圏には常に射
$$
(V, y, U, \phi : V \to u(U))
\longrightarrow
(u(U), v(\phi)(y), U, \text{id} : u(U) \to u(U))
$$
が存在し、ファイバー圏が連結であることが分かる。
第四の最後の等式は明らかである。
\end{proof}

\begin{lemma}
\label{sites-lemma-point-morphism-sites}
\begin{slogan}
サイトの写像は点上の写像を定め、引き戻しは茎を保つ。
\end{slogan}
$f : \mathcal{D} \to \mathcal{C}$ を、連続関手
$u : \mathcal{C} \to \mathcal{D}$ によって与えられるサイトの射とする。
$q$ を関手 $v : \mathcal{D} \to \textit{Sets}$ によって
与えられる $\mathcal{D}$ の点とする。定義
\ref{sites-definition-point} を参照されたい。このとき関手
$v \circ u : \mathcal{C} \to \textit{Sets}$ は
$\mathcal{C}$ の点 $p$ を定め、さらに
$\mathcal{C}$ 上の任意の層 $\mathcal{F}$ に対して標準的な同一視
$$
(f^{-1}\mathcal{F})_q = \mathcal{F}_p
$$
が存在する。
\end{lemma}

\begin{proof}[補題 \ref{sites-lemma-point-morphism-sites} の第一の証明]
$u$ は連続で、$v$ は点を定めるので、
$v \circ u$ が定義 \ref{sites-definition-point} の条件 (1) と (2) を
満たすことは直ちに分かる。表示された等式を証明しよう。
$\mathcal{F}$ を $\mathcal{C}$ 上の層とする。このとき
$$
(f^{-1}\mathcal{F})_q = (u_s\mathcal{F})_q =
(u_p \mathcal{F})_q = \mathcal{F}_p
$$
である。第一の等式は $f^{-1} = u_s$ から、第二の等式は補題
\ref{sites-lemma-point-pushforward-sheaf} から、第三の等式は補題
\ref{sites-lemma-point-functor} から従う。これにより、$p$ は定義
\ref{sites-definition-point} の条件 (3) も満たすことが分かる。
これは完全関手の合成だからである。以上で証明が完了する。
\end{proof}

\begin{proof}[補題 \ref{sites-lemma-point-morphism-sites} の第二の証明]
補題 \ref{sites-lemma-site-point-morphism} により、
$(q_*, q^{-1})$ を
$$
\Sh(pt) \xrightarrow{i} \Sh(\mathcal{S})
\xrightarrow{h} \Sh(\mathcal{D})
$$
と分解できる。ここで第二のトポスの射は、関手
$v : \mathcal{D} \to \mathcal{S}$ によって誘導されるサイトの射
$h : \mathcal{S} \to \mathcal{D}$ から得られる
（$\mathcal{S} \subset \textit{Sets}$ は $v$ の像に現れるすべての
対象を含む充満部分圏なので、これは意味をもつ）。補題
\ref{sites-lemma-composition-morphisms-sites} により、合成
$v \circ u : \mathcal{C} \to \mathcal{S}$ はサイトの射
$g : \mathcal{S} \to \mathcal{C}$ を定める。特に関手
$v \circ u : \mathcal{C} \to \mathcal{S}$ は連続である。
$\mathcal{S}$ における被覆の定義、すなわち注意
\ref{sites-remark-pt-topos} を参照すると、これは $v \circ u$ が定義
\ref{sites-definition-point} の条件 (1) と (2) を満たすことを意味する。
一方、$i$ の構成（注意 \ref{sites-remark-pt-topos}）により、
% P01-E0129
$$
g_*i_*E(U) = i_*E(v(u(U)) = \Mor_{\textit{Sets}}(v(u(U)), E)
$$
であることが分かる。これは式
(\ref{sites-equation-skyscraper}) の $(v \circ u)^pE$ に等しい式と
同じであることに注意する。
% P01-E0130
補題 \ref{sites-lemma-point-pushforward-sheaf} により、関手
$g_*i_* = (v \circ u)^p = (v \circ u)^s$ は茎関手
$\mathcal{F} \mapsto \mathcal{F}_q$ の右随伴である。
% P01-E0131
したがって茎関手 $p^{-1}$ は $i^{-1} \circ g^{-1}$ と標準的に
同型である。ゆえにこれは完全であり、$p$ は点であると結論できる。
最後に、構成により $g = f \circ h$ なので、
$p^{-1} = i^{-1} \circ h^{-1} \circ f^{-1} = q^{-1} \circ f^{-1}$
であることが分かる。すなわち、補題の表示式を得る。
\end{proof}

\begin{lemma}
\label{sites-lemma-point-morphism-topoi}
$f : \Sh(\mathcal{D}) \to \Sh(\mathcal{C})$ をトポスの射とする。
$q : \Sh(pt) \to \Sh(\mathcal{D})$ を点とする。このとき
$p = f \circ q$ はトポス $\Sh(\mathcal{C})$ の点であり、
$\mathcal{C}$ 上の任意の層 $\mathcal{F}$ に対して標準的な同一視
$$
(f^{-1}\mathcal{F})_q = \mathcal{F}_p
$$
が存在する。
\end{lemma}

\begin{proof}
これは定義と、トポスの射を合成できるという事実から直ちに従う。
\end{proof}





\section{局所化と点}
\label{sites-section-localize-points}

\noindent
本節では、局所化 $\mathcal{C}/U$ の点が $\mathcal{C}$ の点から
簡単な仕方で構成されることを示す。

\begin{lemma}
\label{sites-lemma-point-localize}
$\mathcal{C}$ をサイトとする。$p$ を
$u : \mathcal{C} \to \textit{Sets}$ によって与えられる
$\mathcal{C}$ の点とする。$U$ を $\mathcal{C}$ の対象とし、
$x \in u(U)$ とする。関手
$$
v : \mathcal{C}/U \longrightarrow \textit{Sets}, \quad
(\varphi : V \to U) \longmapsto \{y \in u(V) \mid u(\varphi)(y) = x\}
$$
はサイト $\mathcal{C}/U$ の点 $q$ を定め、図式
$$
\xymatrix{
& \Sh(pt) \ar[d]^p \ar[ld]_q \\
\Sh(\mathcal{C}/U) \ar[r]^{j_U} &
\Sh(\mathcal{C})
}
$$
は可換である。言い換えると、$\mathcal{C}$ 上の任意の層について
% P01-E0132
$\mathcal{F}_p = (j_U^{-1}\mathcal{F})_q$ である。
\end{lemma}

\begin{proof}
補題 \ref{sites-lemma-site-point-morphism} のように $S$ と
$\mathcal{S}$ を選ぶ。同補題のように
$\Sh(pt) = \Sh(\mathcal{S})$ と同一視し、$u$ によって誘導される
トポスの射を $p = f : \Sh(\mathcal{S}) \to \Sh(\mathcal{C})$
と書くことができる。補題 \ref{sites-lemma-localize-morphism} により、
トポスの可換図式
$$
\xymatrix{
\Sh(\mathcal{S}/u(U)) \ar[r]_-{j_{u(U)}} \ar[d]_{p'} &
\Sh(\mathcal{S}) \ar[d]^p \\
\Sh(\mathcal{C}/U) \ar[r]^{j_U} &
\Sh(\mathcal{C}),
}
$$
を得る。ここで $p'$ は関手
$u' : \mathcal{C}/U \to \mathcal{S}/u(U)$、
$V/U \mapsto u(V)/u(U)$ によって与えられる。
集合 $E$ を、値 $x$ の定値写像 $E \to u(U)$ を備えた集合 $E$ に
移すことによって得られる関手
$j_x : \mathcal{S} \cong \mathcal{S}/x$ を考える。
すると $j_x$ は余連続な充満忠実関手であり、連続な右随伴
$v_x : (\psi : E \to u(U)) \mapsto \psi^{-1}(\{x\})$ をもつ。
$j_{u(U)} \circ j_x = \text{id}_\mathcal{S}$ かつ
$v_x \circ u' = v$ であることに注意する。
これらの観察から、次のトポスの可換図式を得る。
$$
\xymatrix{
\Sh(\mathcal{S}) \ar[rd]^a \ar[rdd]_q
\ar `r[rrr] `d[dd]^p [rrdd] & & & \\
& \Sh(\mathcal{S}/u(U)) \ar[r]_-{j_{u(U)}} \ar[d]^{p'} &
\Sh(\mathcal{S}) \ar[d]^p & \\
& \Sh(\mathcal{C}/U) \ar[r]^{j_U} &
\Sh(\mathcal{C}) &
}
$$
より詳しくは次の通りである。
\begin{enumerate}
\item 射 $a : \Sh(\mathcal{S}) \to \Sh(\mathcal{S}/u(U))$ は、
余連続関手 $j_x$ に付随するトポスの射である。これは連続関手
$v_x$ に付随する射に等しい。補題
\ref{sites-lemma-cocontinuous-morphism-topoi} および節
\ref{sites-section-cocontinuous-adjoint} を参照されたい。
\item $j_{u(U)} \circ j_x = \text{id}_\mathcal{S}$ なので、合成
$p \circ j_{u(U)} \circ a = p$ である。
\item 合成 $p' \circ a$ はトポスの射を与える。さらに、これは
連続関手 $v_x \circ u' = v$ に付随するトポスの射である。
したがって $v$ は実際に $\mathcal{C}/U$ の点 $q$ を定め、
構成により上の図式に収まる。
\end{enumerate}
これで補題の証明は終わる。
\end{proof}

\begin{lemma}
\label{sites-lemma-points-above-point}
$\mathcal{C}$、$p$、$u$、$U$ を補題
\ref{sites-lemma-point-localize} のようにとる。補題
\ref{sites-lemma-point-localize} の構成は、$p$ 上にある
$\mathcal{C}/U$ の点 $q$ と $u(U)$ の元 $x$ との一対一対応を与える。
\end{lemma}

\begin{proof}
$q$ を、関手 $v : \mathcal{C}/U \to \textit{Sets}$ によって
与えられる $\mathcal{C}/U$ の点で、トポスの射として
$j_U \circ q = p$ を満たすものとする。
$\mathcal{C}$ の任意の対象 $V$ に対して
$u(V) = p^{-1}(h_V^\#)$ であることを思い出す。補題
\ref{sites-lemma-point-site-topos} を参照されたい。同様に、
$\mathcal{C}/U$ の任意の対象 $V/U$ に対して
$v(V/U) = q^{-1}(h_{V/U}^\#)$ である。次の二つの図式を考える。
$$
\vcenter{
\xymatrix{
\Mor_{\mathcal{C}/U}(W/U, V/U) \ar[r] \ar[d] &
\Mor_\mathcal{C}(W, V) \ar[d] \\
\Mor_{\mathcal{C}/U}(W/U, U/U) \ar[r] &
\Mor_\mathcal{C}(W, U)
}
}
\quad
\vcenter{
\xymatrix{
h_{V/U}^\# \ar[r] \ar[d] & j_U^{-1}(h_V^\#) \ar[d] \\
h_{U/U}^\# \ar[r] & j_U^{-1}(h_U^\#)
}
}
$$
右の図式は、$W/U$ を左の集合の図式に移す
$\mathcal{C}/U$ 上の前層の図式を層化したものである。
（ここには小さな技術的注意がある。すなわち
$(j_U^{-1}h_V)^\# = j_U^{-1}(h_V^\#)$ であり、$h_U$ についても
同様である。補題 \ref{sites-lemma-technical-pu} を参照されたい。）
左の集合の図式はカルテジアンであることに注意する。
層化は完全なので（補題 \ref{sites-lemma-sheafification-exact}）、
右の図式もカルテジアンであると結論できる。

\medskip\noindent
右の図式に完全関手 $q^{-1}$ を適用すると、集合のカルテジアン図式
$$
\xymatrix{
v(V/U) \ar[r] \ar[d] & u(V) \ar[d] \\
v(U/U) \ar[r] & u(U)
}
$$
を得る。ここでは
% P01-E0133
$q^{-1} \circ j^{-1} = p^{-1}$ を用いた。
$U/U$ は $\mathcal{C}/U$ の終対象なので、$v(U/U)$ は一元集合である。
したがって $v(U/U)$ の $u(U)$ における像は元 $x$ であり、
上の水平写像は所望の全単射
$v(V/U) \to  \{y \in u(V) \mid y \mapsto x\text{ in }u(U)\}$
を与える。
\end{proof}

\begin{lemma}
\label{sites-lemma-stalk-j-shriek}
$\mathcal{C}$ をサイトとする。$p$ を
$u : \mathcal{C} \to \textit{Sets}$ によって与えられる
$\mathcal{C}$ の点とする。$U$ を $\mathcal{C}$ の対象とする。
$\mathcal{C}/U$ 上の任意の層 $\mathcal{G}$ に対して
$$
(j_{U!}\mathcal{G})_p =
\coprod\nolimits_q \mathcal{G}_q
$$
である。ここで余積は、補題 \ref{sites-lemma-point-localize} のように
$x \in u(U)$ に付随する $\mathcal{C}/U$ の点 $q$ 全体にわたる。
\end{lemma}

\begin{proof}
補題 \ref{sites-lemma-describe-j-shriek} における、前層
$V \mapsto
\coprod\nolimits_{\varphi \in \Mor_\mathcal{C}(V, U)}
\mathcal{G}(V/_\varphi U)$
に付随する層としての $j_{U!}\mathcal{G}$ の記述を用いる。
また、$p$ における $j_{U!}\mathcal{G}$ の茎はこの前層の茎に等しい。
補題 \ref{sites-lemma-point-pushforward-sheaf} を参照されたい。
したがって
$$
(j_{U!}\mathcal{G})_p =
\colim_{(V, y)} \coprod\nolimits_{\varphi : V \to U} \mathcal{G}(V/_\varphi U)
$$
であることが分かる。この余極限の各元 $(V, y, \varphi, s)$ に
$x = u(\varphi)(y) \in u(U)$ を対応させることができる。
よって
$$
(j_{U!}\mathcal{G})_p =
\coprod\nolimits_{x \in u(U)}
\colim_{(\varphi : V \to U, y), \ u(\varphi)(y) = x} \mathcal{G}(V/_\varphi U).
$$
を得る。点 $q$ の構成により、これは補題の式に等しい。
\end{proof}

\begin{remark}
\label{sites-remark-not-pushforward}
注意：補題 \ref{sites-lemma-stalk-j-shriek} の結果には
$j_{U, *}$ に対する類似は存在しない。
\end{remark}




\section{トポスの 2-射}
\label{sites-section-2-category}

\noindent
本節では、トポスの $2$-射という概念について簡単に述べる。

\begin{definition}
\label{sites-definition-2-morphism-topoi}
$f, g : \Sh(\mathcal{C}) \to \Sh(\mathcal{D})$ を二つのトポスの射とする。
{\it $f$ から $g$ への 2-射}とは、関手の自然変換
$t : f_* \to g_*$ によって与えられるものをいう。
\end{definition}

\noindent
$t$ を図式的に次のように表すことがある。
$$
\xymatrix{
\Sh(\mathcal{C})
\rrtwocell^f_g{t}
&
&
\Sh(\mathcal{D})
}
$$
$f^{-1}$ は $f_*$ に随伴し、$g^{-1}$ は $g_*$ に随伴するので、
$t$ は関手の自然変換 $t : g^{-1} \to f^{-1}$
（通常は同じ記号で表す）も誘導する。これは図式
$$
\xymatrix{
\Mor_{\Sh(\mathcal{D})}(\mathcal{G}, f_*\mathcal{F})
\ar[d]_{t \circ -} \ar@{=}[r] &
\Mor_{\Sh(\mathcal{C})}(f^{-1}\mathcal{G}, \mathcal{F})
\ar[d]^{- \circ t}
\\
\Mor_{\Sh(\mathcal{D})}(\mathcal{G}, g_*\mathcal{F})
\ar@{=}[r] &
\Mor_{\Sh(\mathcal{C})}(g^{-1}\mathcal{G}, \mathcal{F})
}
$$
が可換であるという条件によって一意に特徴づけられる。
集合論的困難（注意 \ref{sites-remark-morphism-topoi-big}）のため、
トポスの 2-圏は得られない。しかし水平合成と垂直合成を定義し、
Categories, Section \ref{categories-section-2-categories} に挙げられた
狭義 2-圏の公理が成り立つことを示すことはできる。
すなわち、2-射の垂直合成は明らかであり
（関手の自然変換を合成すればよい）、1-射の合成は定義
\ref{sites-definition-topos} で定めた。また
$$
\xymatrix{
\Sh(\mathcal{C})
\rtwocell^f_g{t}
&
\Sh(\mathcal{D})
\rtwocell^{f'}_{g'}{s}
&
\Sh(\mathcal{E})
}
$$
の水平合成は、Categories, Definition
\ref{categories-definition-horizontal-composition} で導入された
関手の自然変換 $s \star t$ によって定める。明示的には
$s \star t$ は
$$
\xymatrix{
f'_*f_*\mathcal{F}
\ar[r]^{f'_*t} &
f'_*g_*\mathcal{F}
\ar[r]^s &
g'_*g_*\mathcal{F}
}
\quad\text{or}\quad
\xymatrix{
f'_*f_*\mathcal{F}
\ar[r]^s &
g'_*f_*\mathcal{F}
\ar[r]^{g'_*t} &
g'_*g_*\mathcal{F}
}
$$
によって与えられる（これらの写像は等しい）。これらの定義は
Categories, Section \ref{categories-section-formal-cat-cat} の定義と
一致するので、Categories, Lemma
\ref{categories-lemma-properties-2-cat-cats} から、この定義のもとで
狭義 2-圏の公理が成り立つ。




\section{点の間の射}
\label{sites-section-morphisms-points}

\begin{lemma}
\label{sites-lemma-maps-u-points}
$\mathcal{C}$ をサイトとする。
$u, u' : \mathcal{C} \to \textit{Sets}$ を二つの関手とし、
$t : u' \to u$ を関手の自然変換とする。このとき茎関手の標準的な
自然変換 $t_{stalk} : \mathcal{F}_{p'} \to \mathcal{F}_p$ を得る。
これは補題 \ref{sites-lemma-points-recover} の同一視を通じて $t$ と一致する。
\end{lemma}

\begin{proof}
省略する。
\end{proof}

\begin{definition}
\label{sites-definition-morphism-points}
$\mathcal{C}$ をサイトとする。$p, p'$ を、関手
$u, u' : \mathcal{C} \to \textit{Sets}$ によって与えられる
$\mathcal{C}$ の点とする。{\it 射 $f : p \to p'$}とは、関手の自然変換
$$
f_u : u' \to u.
$$
によって与えられるものをいう。
\end{definition}

\noindent
関手の自然変換の向きが逆であることに注意する。後で見るように、
射 $f$ を次の図式における一種の $2$-矢印と考えれば、これは自然である。
$$
\xymatrix{
\textit{Sets}
=
\Sh(pt)
\rrtwocell^p_{p'}{f}
&
&
\Sh(\mathcal{C})
}
$$
すなわち、$f_u$ は摩天楼層の構成の間の標準的な関手の自然変換
$p_* \to p'_*$ を誘導することを後で見る。

\medskip\noindent
これはかなり重要な概念であり、ここでより完全に扱うに値する。
望まれる事項を列挙する。
\begin{enumerate}
\item 例 \ref{sites-example-point-G-sets} で述べた
$\mathcal{T}_G$ の点の自己同型を記述する。
\item $\Mor(p, p')$ を $\Mor(p_*, p'_*)$ によって記述する。
\item 位相空間における点の特殊化を記述する。位相空間 $X$ において
$x' \in \overline{\{x\}}$ ならば射 $p \to p'$ が存在することを示す。
ここで $p$（resp.\ $p'$）は $x$（resp.\ $x'$）に付随する
$X_{Zar}$ の点である。
\end{enumerate}




\section{十分多くの点をもつサイト}
\label{sites-section-sites-enough-points}

\begin{definition}
\label{sites-definition-enough-points}
$\mathcal{C}$ をサイトとする。
\begin{enumerate}
\item 点の族 $\{p_i\}_{i\in I}$ が{\it 保守的}であるとは、
すべてのファイバー上で同型 $\mathcal{F}_{p_i} \to \mathcal{G}_{p_i}$ を
% P01-E0134
誘導する任意の層の写像 $\phi : \mathcal{F} \to \mathcal{G}$ が
同型であることをいう。
\item 保守的な点の族が存在するとき、$\mathcal{C}$ は
{\it 十分多くの点をもつ}という。
\end{enumerate}
\end{definition}

\noindent
このとき「完全性」は茎で確認できることが分かる。

\begin{lemma}
\label{sites-lemma-exactness-stalks}
$\mathcal{C}$ をサイトとし、$\{p_i\}_{i\in I}$ を保守的な点の族とする。
このとき次が成り立つ。
\begin{enumerate}
\item 任意の層の写像 $\varphi : \mathcal{F} \to \mathcal{G}$ に対して、
$\forall i, \varphi_{p_i}$ が単射ならば $\varphi$ は単射である。
\item 任意の層の写像 $\varphi : \mathcal{F} \to \mathcal{G}$ に対して、
$\forall i, \varphi_{p_i}$ が全射ならば $\varphi$ は全射である。
\item 層の写像の任意の対
$\varphi_1, \varphi_2 : \mathcal{F} \to \mathcal{G}$ に対して、
$\forall i, \varphi_{1, p_i} = \varphi_{2, p_i}$ ならば
$\varphi_1 = \varphi_2$ である。
\item 有限図式 $\mathcal{G} : \mathcal{J}
\to \Sh(\mathcal{C})$、層 $\mathcal{F}$、および射
$q_j : \mathcal{F} \to \mathcal{G}_j$ が与えられたとき、
$(\mathcal{F}, q_j)$ がこの図式の極限であるための必要十分条件は、
各 $i$ に対して茎 $(\mathcal{F}_{p_i}, (q_j)_{p_i})$ が極限であることである。
\item 有限図式 $\mathcal{F} : \mathcal{J}
\to \Sh(\mathcal{C})$、層 $\mathcal{G}$、および射
$e_j : \mathcal{F}_j \to \mathcal{G}$ が与えられたとき、
$(\mathcal{G}, e_j)$ がこの図式の余極限であるための必要十分条件は、
各 $i$ に対して茎 $(\mathcal{G}_{p_i}, (e_j)_{p_i})$ が
余極限であることである。
\end{enumerate}
\end{lemma}

\begin{proof}
すべての茎関手が任意の有限極限および有限余極限と可換し、したがって
積、ファイバー積などと可換することを繰り返し用いる。また、単射は
モノ射であり、全射はエピ射であることも用いる。
層の写像 $\varphi : \mathcal{F} \to \mathcal{G}$ が単射であるための
必要十分条件は、
$\mathcal{F} \to \mathcal{F} \times_\mathcal{G} \mathcal{F}$ が
同型であることである。これにより (1) を得る。
同様に、$\varphi : \mathcal{F} \to \mathcal{G}$ が全射であるための
必要十分条件は、
$\mathcal{G} \amalg_\mathcal{F} \mathcal{G} \to \mathcal{G}$ が
同型であることである。これにより (2) を得る。
写像 $a, b : \mathcal{F} \to \mathcal{G}$ が等しいための必要十分条件は、
$\text{Equalizer}(a, b : \mathcal{F} \to \mathcal{G}) \to \mathcal{F}$
が同型であることである\footnote{等化子は積とファイバー積によって
表せる。Categories, Lemma
\ref{categories-lemma-finite-limits-exist} の証明を参照されたい。}。
これにより (3) を得る。(4) と (5) は定義およびこの証明の冒頭の
注意から直ちに従う。
\end{proof}

\begin{lemma}
\label{sites-lemma-enough}
$\mathcal{C}$ をサイトとし、$\{(p_i, u_i)\}_{i\in I}$ を点の族とする。
この族が保守的であるための必要十分条件は、任意の層 $\mathcal{F}$、
任意の $U\in \Ob(\mathcal{C})$、および相異なる任意の切断の対
$s, s' \in \mathcal{F}(U)$、$s \not = s'$ に対して、ある $i$ と
$x\in u_i(U)$ が存在して、三つ組 $(U, x, s)$ と $(U, x, s')$ が
$\mathcal{F}_{p_i}$ の相異なる元を定めることである。
\end{lemma}

\begin{proof}
この族が保守的であり、$\mathcal{F}$、$U$、$s, s'$ が補題のように
与えられているとする。切断 $s$, $s'$ は相異なる写像
$a, a' : (h_U)^\# \to \mathcal{F}$ を定める。したがって補題
\ref{sites-lemma-exactness-stalks} により、ある $i$ に対して
$a_{p_i} \not = a'_{p_i}$ である。補題
\ref{sites-lemma-points-recover} および
\ref{sites-lemma-point-pushforward-sheaf} により
$(h_U)^\#_{p_i} = u_i(U)$ であることを思い出す。
したがって $\mathcal{F}_{p_i}$ において
$a_{p_i}(x) \not = a'_{p_i}(x)$ となる $x \in u_i(U)$ が存在する。
定義を展開すると、$(U, x, s)$ と $(U, x, s')$ は補題の主張の
通りであることが分かる。

\medskip\noindent
逆を証明するため、補題の点の存在条件を仮定する。
$\phi : \mathcal{F} \to \mathcal{G}$ をすべての茎で同型となる
層の写像とする。補題 \ref{sites-lemma-mono-epi-sheaves} により、
$\phi$ が単射かつ全射であることを示せばよい。
単射性は仮定から直ちに従う。全射性を示すには、
$\mathcal{G} \amalg_\mathcal{F} \mathcal{G} \to \mathcal{G}$ が
同型であることを示さなければならない
（Categories, Lemma \ref{categories-lemma-characterize-mono-epi}）。
この写像は明らかに全射なので、単射性を確認すれば十分である。
これは仮定により
$\mathcal{G} \amalg_\mathcal{F} \mathcal{G} \to \mathcal{G}$ が
すべての茎で単射であることから従う。
\end{proof}

\noindent
次の補題における点 $q_{i, x}$ は、補題
\ref{sites-lemma-points-above-point} によれば、点 $p_i$ 上にある
$\mathcal{C}/U$ の点全体にほかならない。

\begin{lemma}
\label{sites-lemma-localize-enough}
$\mathcal{C}$ をサイトとする。$U$ を $\mathcal{C}$ の対象とする。
% P01-E0135
$\{(p_i, u_i)\}_{i\in I}$ を $\mathcal{C}$ の点の族とする。
$x \in u_i(U)$ に対して、補題 \ref{sites-lemma-point-localize} で構成した
$\mathcal{C}/U$ の点を $q_{i, x}$ とする。
$\{p_i\}$ が保守的な点の族ならば、
$\{q_{i, x}\}_{i \in I, x \in u_i(U)}$ は
$\mathcal{C}/U$ の保守的な点の族である。特に $\mathcal{C}$ が
十分多くの点をもつならば、各局所化 $\mathcal{C}/U$ もそうである。
\end{lemma}

\begin{proof}
補題 \ref{sites-lemma-essential-image-j-shriek} により、$j_{U!}$ は同値
$j_{U!} : \Sh(\mathcal{C}/U) \to \Sh(\mathcal{C})/h_U^\#$ を
誘導する。また、補題 \ref{sites-lemma-stalk-j-shriek} により
$(j_{U!}\mathcal{G})_{p_i} = \coprod_x \mathcal{G}_{q_{i, x}}$ である。
したがって結果は形式的に従う。
\end{proof}

\noindent
次の補題は、点の存在をサイト上で局所的に確認できることを述べる。

\begin{lemma}
\label{sites-lemma-enough-points-local}
$\mathcal{C}$ をサイトとする。$\{U_i\}_{i \in I}$ を
$\mathcal{C}$ の対象の族とする。次を仮定する。
\begin{enumerate}
\item $\coprod h_{U_i}^\# \to *$ は層の全射である。
\item 各局所化 $\mathcal{C}/U_i$ は十分多くの点をもつ。
\end{enumerate}
このとき $\mathcal{C}$ は十分多くの点をもつ。
\end{lemma}

\begin{proof}
各 $i \in I$ に対して、$\{p_j\}_{j \in J_i}$ を
$\mathcal{C}/U_i$ の保守的な点の族とする。$j \in J_i$ に対して、
% P01-E0136
$q_j : \Sh(pt) \to \Sh(\mathcal{C})$ を、$p_j$ と局所化射
$\Sh(\mathcal{C}/U_i) \to \Sh(\mathcal{C})$ との合成とする。
補題 \ref{sites-lemma-point-morphism-topoi} により $q_j$ は点である。
点の族 $\{q_j\}_{j \in \coprod J_i}$ が保守的であると主張する。
実際、$\mathcal{F} \to \mathcal{G}$ を $\mathcal{C}$ 上の層の写像で、
すべての $j \in \coprod J_i$ に対して
$\mathcal{F}_{q_j} \to \mathcal{G}_{q_j}$ が同型であるものとする。
$W$ を $\mathcal{C}$ の対象とする。仮定 (1) により、被覆
$\{W_a \to W\}$ と射 $W_a \to U_{i(a)}$ が存在する。
補題 \ref{sites-lemma-point-morphism-topoi} により
$(\mathcal{F}|_{\mathcal{C}/U_{i(a)}})_{p_j} = \mathcal{F}_{q_j}$ および
$(\mathcal{G}|_{\mathcal{C}/U_{i(a)}})_{p_j} = \mathcal{G}_{q_j}$ である。
点の族 $\{p_j\}_{j \in J_{i(a)}}$ は保守的なので、
$\mathcal{F}|_{U_{i(a)}} \to \mathcal{G}|_{U_{i(a)}}$ は同型である。
したがって各 $a$ に対して
$\mathcal{F}(W_a) \to \mathcal{G}(W_a)$ は全単射である。
同様に各 $a, b$ に対して
$\mathcal{F}(W_a \times_W W_b) \to \mathcal{G}(W_a \times_W W_b)$
は全単射である。層条件により
$\mathcal{F}(W) \to \mathcal{G}(W)$ は全単射、すなわち
$\mathcal{F} \to \mathcal{G}$ は同型である。
\end{proof}

\begin{lemma}
\label{sites-lemma-check-morphism-sites}
$u : \mathcal{C} \to \mathcal{D}$ をサイトの連続関手とする。
$\{(q_i, v_i)\}_{i\in I}$ を $\mathcal{D}$ の保守的な点の族とする。
各関手 $u_i = v_i \circ u$ が $\mathcal{C}$ の点を定めるならば、
$u$ はサイトの射 $f : \mathcal{D} \to \mathcal{C}$ を定める。
\end{lemma}

\begin{proof}
関手 $u_i$ に対応する $\textit{PSh}(\mathcal{C})$ 上の茎関手
(\ref{sites-equation-stalk}) を $p_i$ と書く。
% P01-E0137
$$
(f^{-1}\mathcal{F})_{q_i} =
(u_s\mathcal{F})_{q_i} =
(u_p\mathcal{F})_{q_i} =
\mathcal{F}_{p_i}
$$
である。第一の等式は $f^{-1} = u_s$ により、第二の等式は補題
\ref{sites-lemma-point-pushforward-sheaf} により、第三の等式は補題
\ref{sites-lemma-point-functor} による。したがって $p_i$ が点ならば、
$f$ によって引き戻してから $q_i$ で茎を取る操作は完全関手である。
点の族 $\{q_i\}$ は保守的なので、これは $f^{-1}$ が完全関手である
ことを意味する。よって定義 \ref{sites-definition-morphism-sites} により、
$f$ はサイトの射であることが分かる。
\end{proof}


\section{点の存在判定条件}
\label{sites-section-criterion-points}

\noindent
本節は \cite[Expos\'e VI]{SGA4} に収録された Deligne の付録に
対応する。実際、ほとんど文字通り同じである。

\medskip\noindent
$\mathcal{C}$ をサイトとする。$(I, \geq)$ を有向集合とし、
$(U_i, f_{ii'})$ を $I$ 上の逆系とする。Categories, Definition
\ref{categories-definition-system-over-poset} を参照されたい。
データ $(I, \geq, U_i, f_{ii'})$ が与えられたとき、
$$
u : \mathcal{C} \longrightarrow \textit{Sets}, \quad
u(V) = \colim_i \Mor_\mathcal{C}(U_i , V)
$$
と定める。節 \ref{sites-section-points} のように、$u$ に付随する茎関手を
$\mathcal{F} \mapsto \mathcal{F}_p$ とする。この特別な場合には、
定義から実際に
$$
\mathcal{F}_p = \colim_i \mathcal{F}(U_i)
$$
であることが直接分かる。$u$ はすべての有限極限
（$\mathcal{C}$ において表現可能なものをいう）と可換することに
注意する。これは各関手
$V \mapsto \Mor_\mathcal{C}(U_i , V)$ がそうだからである。
% P01-E0138
Categories, Lemma \ref{categories-lemma-directed-commutes} を参照されたい。

\medskip\noindent
系 $(I, \geq, U_i, f_{ii'})$ が
$(J, \geq, V_j, g_{jj'})$ の{\it 細分}であるとは、$J \subset I$ であり、
% P01-E0139
$J$ 上の順序が $I$ の順序から誘導され、さらに
$V_j = U_j$、$g_{jj'} = f_{jj'}$ であることをいう
（言い換えると、$J$ 上の逆系は $I$ 上の逆系から誘導される）。
$(I, \geq, U_i, f_{ii'})$ に付随する関手を $u$ とし、
$(J, \geq, V_j, g_{jj'})$ に付随する関手を $u'$ とする。
$u'$ の余極限は $u$ の余極限に現れる系の部分系にわたるので、
これは関手の自然変換
$$
u' \longrightarrow u
$$
を誘導する。特に、補題 \ref{sites-lemma-maps-u-points} により、
付随する茎関手の自然変換
$\mathcal{F}_{p'} \to \mathcal{F}_p$ を得る。

\begin{lemma}
\label{sites-lemma-refine}
$\mathcal{C}$ をサイトとする。
$(J, \geq, V_j, g_{jj'})$ を上のような系とし、付随する関手の対を
$(u', p')$ とする。$\mathcal{F}$ を $\mathcal{C}$ 上の層とする。
$s, s' \in \mathcal{F}_{p'}$ を相異なる元とする。
$\{W_k \to W\}$ を $\mathcal{C}$ の有限被覆とし、
$f \in u'(W)$ とする。このとき $(J, \geq, V_j, g_{jj'})$ の細分
$(I, \geq, U_i, f_{ii'})$ であって、$s, s'$ が
$\mathcal{F}_p$ の相異なる元に移り、かつ $u(W)$ における $f$ の像が
$u(W_k)$ のいずれか一つの像に属するものが存在する。
\end{lemma}

\begin{proof}
$f$ が $f' : V_{j_0} \to W$ によって表されるような
$j_0 \in J$ が存在する。$j \geq j_0$ に対して
% P01-E0140
$V_{j, k} = V_j \times_{f'\circ f_{j j_0}, W} W_k$ とおく。
このとき $\{V_{j, k} \to V_j\}$ はサイト $\mathcal{C}$ の有限被覆である。
したがって
$\mathcal{F}(V_j) \subset \prod_k \mathcal{F}(V_{j, k})$ である。
Categories, Lemma \ref{categories-lemma-directed-commutes} を再び用いると、
$$
\mathcal{F}_{p'} =
\colim_j \mathcal{F}(V_j)
\longrightarrow
\prod\nolimits_k \colim_j \mathcal{F}(V_{j, k})
$$
は単射であることが分かる。よって $s$ と $s'$ が
$\colim_j \mathcal{F}(V_{j, k})$ において相異なる像をもつような
$k$ が存在する。$J_0 = \{j \in J, j \geq j_0\}$ および
$I = J \amalg J_0$ とおく。第二の成分のどの元も第一の成分の
どの元より小さくならないように $I$ を順序づけ、それ以外では
$J$ の順序を用いる。$j \in I$ が第一の成分に属するときは $V_j$ を、
$j \in I$ が第二の成分に属するときは $V_{j, k}$ を用いる。
逆系の遷移写像の
定義は省略する。以上により、$s, s'$ は $\mathcal{F}_p$ において
相異なる像をもつ。また構成により、$f'$ の $V_{j, k}$ への制限は
$W_k$ を経由する。
\end{proof}

\begin{lemma}
\label{sites-lemma-refine-all-at-once}
$\mathcal{C}$ をサイトとする。
$(J, \geq, V_j, g_{jj'})$ を上のような系とし、付随する関手の対を
$(u', p')$ とする。$\mathcal{F}$ を $\mathcal{C}$ 上の層とし、
$s, s' \in \mathcal{F}_{p'}$ を相異なる元とする。このとき
$(J, \geq, V_j, g_{jj'})$ の細分 $(I, \geq, U_i, f_{ii'})$ であって、
$s, s'$ が $\mathcal{F}_p$ の相異なる元に移り、さらにサイト
$\mathcal{C}$ の任意の有限被覆 $\{W_k \to W\}$ と任意の
$f \in u'(W)$ に対して、$u(W)$ における $f$ の像が
$u(W_k)$ のいずれか一つの像に属するものが存在する。
\end{lemma}

\begin{proof}
$E$ を対 $(\{W_k \to W\}, f\in u'(W))$ 全体の集合とする。
次の条件を満たす対
$(E' \subset E, (I, \geq, U_i, f_{ii'}))$ を考える。
\begin{enumerate}
% P01-E0141
\item $(I, \geq, U_i, g_{ii'})$ は
$(J, \geq, V_j, g_{jj'})$ の細分である。
\item $s, s'$ は $\mathcal{F}_p$ の相異なる元に移る。
\item 各対 $(\{W_k \to W\}, f\in u'(W)) \in E'$ に対して、
$u(W)$ における $f$ の像は $u(W_k)$ のいずれか一つの像に属する。
\end{enumerate}
このような対を、第一成分では包含によって、第二成分では細分によって
順序づける。上のような対
$(E' \subset E, (I, \geq, U_i, f_{ii'}))$ 全体の類を
$\mathcal{S}$ と書く。
% P01-E0142
$\mathcal{S}$ はツォルンの補題の仮定を満たすと主張する。
実際、$(E'_a, (I_a, \geq, U_i, f_{ii'}))_{a \in A}$ を
$\mathcal{S}$ の全順序部分集合とする。このとき
$E' = \bigcup_{a \in A} E'_a$ と定め、
$I = \bigcup_{a \in A} I_a$ とおくことができる。対応する対
$(E' , (I, \geq, U_i, f_{ii'}))$ は $\mathcal{S}$ の元であると主張する。
条件 (1) と (3) は明らかである。条件 (2) については
$$
u = \colim_{a \in A} u_a
\text{ and correspondingly }
\mathcal{F}_p = \colim_{a \in A} \mathcal{F}_{p_a}
$$
であることに注意する。この茎における $s, s'$ の像の相異性は、
集合の有向余極限の記述から従う。Categories, Section
\ref{categories-section-directed-colimits} を参照されたい。
この状況では単に
$(E', (I, \ldots)) = \bigcup_{a \in A}(E'_a, (I_a, \ldots))$
と書くことにする。

\medskip\noindent
さて、$\mathcal{S}$ が集合であればツォルンの補題を適用でき、
上の補題 \ref{sites-lemma-refine} と合わせて本補題が容易に証明される。
しかし $\mathcal{S}$ は類なので適用できない。この問題を回避するため、
$E$ 上の整列順序を選ぶ。$e \in E$ に対して
$E'_e = \{e' \in E \mid e' \leq e\}$ とおく。超限再帰により、
$(E'_e, (I_e, \ldots)) \in \mathcal{S}$ であって
$e_1 \leq e_2 \Rightarrow (E'_{e_1}, (I_{e_1}, \ldots))
\leq (E'_{e_2}, (I_{e_2}, \ldots))$ となる対を構成する。
$e \in E$ とし、$e = (\{W_k \to W\}, f\in u'(W))$ と書く。
$e$ が直前元 $e - 1$ をもつならば、補題 \ref{sites-lemma-refine} のように、
系 $e = (\{W_k \to W\}, f\in u'(W))$ に関して
$(I_e, \ldots)$ を $(I_{e - 1}, \ldots)$ の細分とする。
$e$ が直前元をもたないならば、同じ系
$e = (\{W_k \to W\}, f\in u'(W))$ に関して
$(I_e, \ldots)$ を
$\bigcup_{e' < e} (I_{e'}, \ldots)$ の細分とする。
最後に、合併 $\bigcup_{e \in E} I_e$ が補題の問題の解となる。
\end{proof}

\begin{proposition}
\label{sites-proposition-criterion-points}
\begin{reference}
\cite[Expos\'e VI, Appendix by Deligne, Proposition 9.0]{SGA4}
\end{reference}
$\mathcal{C}$ をサイトとする。次を仮定する。
\begin{enumerate}
\item $\mathcal{C}$ に有限極限が存在する。
\item 各被覆 $\{U_i \to U\}_{i \in I}$ は
$\mathcal{C}$ の有限被覆による細分をもつ。
\end{enumerate}
このとき $\mathcal{C}$ は十分多くの点をもつ。
\end{proposition}

\begin{proof}
$\mathcal{C}$ 上の任意の層 $\mathcal{F}$、任意の
$U \in \Ob(\mathcal{C})$、および相異なる任意の切断
$s, s' \in \mathcal{F}(U)$ に対して、$s, s'$ の
$\mathcal{F}_p$ における像が相異なるような点 $p$ が存在することを
示さなければならない。補題 \ref{sites-lemma-enough} を参照されたい。
$J = \{1\}$、$V_1 = U$、$g_{11} = \text{id}_U$ である系
$(J, \geq, V_j, g_{jj'})$ を考える。補題
\ref{sites-lemma-refine-all-at-once} を適用する。同補題により、もとの系の
細分である系 $(I, \geq, U_i, f_{ii'})$ であって
$s_p \not = s'_p$ を満たし、さらにサイト $\mathcal{C}$ の
任意の有限被覆 $\{W_k \to W\}$ に対して写像
$\coprod_k u(W_k) \to u(W)$ が全射となるものを得る。
$\mathcal{C}$ の各被覆は有限被覆によって細分できるので、サイト
$\mathcal{C}$ の{\it 任意の}被覆 $\{W_k \to W\}$ に対して
$\coprod_k u(W_k) \to u(W)$ が全射であると結論できる。
これは $u = p$ が点であることを意味する。命題
\ref{sites-proposition-point-limits}（および $u$ が有限極限と可換することを
保証する本節冒頭の議論）を参照されたい。
\end{proof}

\begin{lemma}
\label{sites-lemma-criterion-points}
$\mathcal{C}$ をサイトとする。$I$ を集合とし、各 $i \in I$ に対して
$U_i$ を $\mathcal{C}$ の対象で、次を満たすものとする。
\begin{enumerate}
% P01-E0143
\item $\coprod h_{U_i}$ は $\Sh(\mathcal{C})$ の終対象へ全射する。
\item $\mathcal{C}/U_i$ は命題
\ref{sites-proposition-criterion-points} の仮定を満たす。
\end{enumerate}
このとき $\mathcal{C}$ は十分多くの点をもつ。
\end{lemma}

\begin{proof}
仮定 (2) と命題により、$\mathcal{C}/U_i$ は十分多くの点をもつ。
$\mathcal{C}/U_i$ の点は、補題 \ref{sites-lemma-point-morphism-sites} の
手続きにより $\mathcal{C}$ の点を与える。したがって次を示せば十分である。
$\phi : \mathcal{F} \to \mathcal{G}$ が $\mathcal{C}$ 上の層の写像で、
各 $i$ に対して $\phi|_{\mathcal{C}/U_i}$ が同型ならば、$\phi$ は
同型である。仮定 (1) により、$\mathcal{C}$ の各対象 $W$ に対して、
被覆 $\{W_j \to W\}_{j \in J}$ が存在し、各 $j \in J$ に対して
ある $i \in I$ と射 $f_j : W_j \to U_i$ が存在する。
このとき写像 $\mathcal{F}(W_j) \to \mathcal{G}(W_j)$ は全単射であり、
同様に
$\mathcal{F}(W_j \times_W W_{j'}) \to \mathcal{G}(W_j \times_W W_{j'})$
も全単射である。層条件から、所望のように
$\mathcal{F}(W) \to \mathcal{G}(W)$ は全単射である。
\end{proof}



\section{弱可縮対象}
\label{sites-section-w-contractible}

\noindent
サイトの対象が{\it 弱可縮}であるとは、次の補題の同値な条件を
満たすことをいう。

\begin{lemma}
\label{sites-lemma-w-contractible}
$\mathcal{C}$ をサイトとし、$U$ を $\mathcal{C}$ の対象とする。
次の条件は同値である。
\begin{enumerate}
\item 各被覆 $\{U_i \to U\}$ に対して、
$\coprod h_{U_i} \to h_U$ の層化の右逆となる層の写像
$h_U^\# \to \coprod h_{U_i}^\#$ が存在する。
\item 集合の層の各全射 $\mathcal{F} \to \mathcal{G}$ に対して、
写像 $\mathcal{F}(U) \to \mathcal{G}(U)$ は全射である。
\end{enumerate}
\end{lemma}

\begin{proof}
(1) を仮定し、$\mathcal{F} \to \mathcal{G}$ を集合の層の全射とする。
$s \in \mathcal{G}(U)$ に対して、被覆 $\{U_i \to U\}$ と、
$s|_{U_i}$ に移る $t_i \in \mathcal{F}(U_i)$ が存在する。定義
\ref{sites-definition-sheaves-injective-surjective} を参照されたい。
(\ref{sites-equation-map-representable-into-sheaf}) を用いて、$t_i$ を写像
$t_i : h_{U_i}^\# \to \mathcal{F}$ と考える。(1) から得られる写像
$h_U^\# \to \coprod h_{U_i}^\#$ を
$\coprod t_i : \coprod h_{U_i}^\# \to \mathcal{F}$ の前に合成すると、
% P01-E0144
$s$ に移る切断 $t \in \mathcal{F}(U)$ を得る。したがって (2) が成り立つ。

\medskip\noindent
(2) を仮定する。$\{U_i \to U\}$ を被覆とする。このとき
$\coprod h_{U_i}^\# \to h_U^\#$ は全射である
（補題 \ref{sites-lemma-covering-surjective-after-sheafification}）。
したがって (2) により、$h_U^\#$ の切断 $\text{id}_U$ に移る
$\coprod h_{U_i}^\#$ の切断 $s$ が存在する。この切断は写像
$h_U^\# \to \coprod h_{U_i}^\#$ に対応し、これは
$\coprod h_{U_i} \to h_U$ の層化の右逆である。これにより (1) を得る。
\end{proof}

\begin{definition}
\label{sites-definition-w-contractible}
$\mathcal{C}$ をサイトとする。
\begin{enumerate}
\item $\mathcal{C}$ の対象 $U$ が補題
\ref{sites-lemma-w-contractible} の同値な条件を満たすとき、その対象は
{\it 弱可縮}であるという。
\item $\mathcal{C}$ の各対象 $U$ が、すべての $i$ に対して $U_i$ が
弱可縮であるような被覆 $\{U_i \to U\}$ をもつとき、サイトは
{\it 十分多くの弱可縮対象をもつ}という。
\item より一般に、$P$ を $\mathcal{C}$ の対象の性質とする。
$\mathcal{C}$ の各対象 $U$ が、すべての $i$ に対して $U_i$ が $P$ を
もつような被覆 $\{U_i \to U\}$ をもつとき、$\mathcal{C}$ は
{\it 十分多くの $P$ 対象をもつ}という。
\end{enumerate}
\end{definition}

\noindent
$\mathbf{A}^1_\mathbf{C}$ の小エタールサイトには弱可縮対象が
一つも存在しない。一方、任意のスキームの小プロエタールサイトは
十分多くの可縮対象をもつ。

\section{順像の完全性}
\label{sites-section-pushforward}

\noindent
$f$ をトポスの射とする。一般に関手 $f_*$ は左完全でしかない。
トポスの射の特定のクラスを取り出すため、この関手には多くの追加条件を
課すことができる。ここではそれらをまとめ、論理的な依存関係のいくつかを
指摘する。次の補題の一部は純粋に圏論的である
（すなわち、トポスの射であることには依存せず、随伴関手の対があれば
十分である）。

\begin{lemma}
\label{sites-lemma-exactness-properties}
$f : \Sh(\mathcal{C}) \to \Sh(\mathcal{D})$ をトポスの射とする。
次の性質（集合の層について）を考える。
\begin{enumerate}
\item $f_*$ は忠実である。
\item $f_*$ は充満忠実である。
\item $\Sh(\mathcal{C})$ のすべての $\mathcal{F}$ に対して
$f^{-1}f_*\mathcal{F} \to \mathcal{F}$ は全射である。
\item $f_*$ は全射を全射に移す。
\item $f_*$ は余等化子と可換する。
\item $f_*$ は押し出しと可換する。
\item $\Sh(\mathcal{C})$ のすべての $\mathcal{F}$ に対して
$f^{-1}f_*\mathcal{F} \to \mathcal{F}$ は同型である。
\item $f_*$ は単射を反映する。
\item $f_*$ は全射を反映する。
\item $f_*$ は全単射を反映する。
\item 任意の全射 $\mathcal{F} \to f^{-1}\mathcal{G}$ に対して、
$f^{-1}\mathcal{G}' \to f^{-1}\mathcal{G}$ が
$\mathcal{F} \to f^{-1}\mathcal{G}$ を経由するような全射
$\mathcal{G}' \to \mathcal{G}$ が存在する。
\end{enumerate}
このとき次の含意が成り立つ。
\begin{enumerate}
\item[(a)] (2) $\Rightarrow$ (1),
\item[(b)] (3) $\Rightarrow$ (1),
\item[(c)] (7) $\Rightarrow$ (1), (2), (3), (8), (9), (10).
\item[(d)] (3) $\Leftrightarrow$ (9),
\item[(e)] (6) $\Rightarrow$ (4) and (5) $\Rightarrow$ (4),
\item[(f)] (4) $\Leftrightarrow$ (11),
\item[(g)] (9) $\Rightarrow$ (8), (10), and
\item[(h)] (2) $\Leftrightarrow$ (7).
\end{enumerate}
これを図示すると
$$
\xymatrix{
(6) \ar@{=>}[rd] & & & & & (9) \ar@{=>}[r] \ar@{=>}[rd] & (8) \\
& (4) \ar@{<=>}[r] & (11) &
(2) \ar@{<=>}[r] &
(7) \ar@{=>}[ru] \ar@{=>}[rd] & & (10) \\
(5) \ar@{=>}[ur] & & & & & (3) \ar@{=>}[r] & (1)
}
$$
となる。
\end{lemma}

\begin{proof}
(a) の証明：これは定義から直ちに従う。

\medskip\noindent
(b) の証明。$a, b : \mathcal{F} \to \mathcal{F}'$ を
$\mathcal{C}$ 上の層の写像とする。$f_*a = f_*b$ ならば
$f^{-1}f_*a = f^{-1}f_*b$ である。可換図式
$$
\xymatrix{
\mathcal{F} \ar@<-1ex>[r] \ar@<1ex>[r] & \mathcal{F}' \\
f^{-1}f_*\mathcal{F} \ar@<-1ex>[r] \ar@<1ex>[r] \ar[u] &
f^{-1}f_*\mathcal{F}' \ar[u]
}
$$
を考える。下の二つの矢印が等しく、垂直の矢印が全射ならば、
上の二つの矢印は等しい。したがって (b) が従う。

\medskip\noindent
(c) の証明。$a : \mathcal{F} \to \mathcal{F}'$ を
$\mathcal{C}$ 上の層の写像とする。可換図式
$$
\xymatrix{
\mathcal{F} \ar[r] & \mathcal{F}' \\
f^{-1}f_*\mathcal{F} \ar[r] \ar[u] &
f^{-1}f_*\mathcal{F}' \ar[u]
}
$$
を考える。(7) が成り立つならば、垂直の矢印は同型である。
したがって $f_*a$ が単射（resp.\ 全射、resp.\ 全単射）ならば、
下の矢印は単射（resp.\ 全射、resp.\ 全単射）であり、ゆえに上の
矢印も単射（resp.\ 全射、resp.\ 全単射）である。よって (7) は
(8)、(9)、(10) を含意する。また (7) が (3) を含意することは明らかである。
含意 (7) $\Rightarrow$ (2), (1) は、後で見る (a) と (h) から従う。

\medskip\noindent
(d) の証明。(3) を仮定する。$a : \mathcal{F} \to \mathcal{F}'$ を
$\mathcal{C}$ 上の層の写像で、$f_*a$ が全射であるものとする。
$f^{-1}$ は完全なので、これは
$f^{-1}f_*a : f^{-1}f_*\mathcal{F} \to f^{-1}f_*\mathcal{F}'$ が
全射であることを意味する。(3) と合わせると $a$ は全射である。
これは (9) が成り立つことを意味する。
(9) を仮定する。$\mathcal{F}$ を $\mathcal{C}$ 上の層とする。
写像 $f^{-1}f_*\mathcal{F} \to \mathcal{F}$ が全射であることを
示さなければならない。
$f_*f^{-1}f_*\mathcal{F} \to f_*\mathcal{F}$ が全射であることを
示せば十分である。実際、片側逆となる標準写像
$f_*\mathcal{F} \to f_*f^{-1}f_*\mathcal{F}$ が存在するので、
これは成り立つ。

\medskip\noindent
(e) の証明。以下、Categories, Lemma
\ref{categories-lemma-characterize-mono-epi} を断りなく用いる。
$\mathcal{F} \to \mathcal{F}'$ が全射ならば、
$\mathcal{F}' \amalg_\mathcal{F} \mathcal{F}' \to \mathcal{F}'$ は
同型である。したがって (6) により
$$
f_*\mathcal{F}' \amalg_{f_*\mathcal{F}} f_*\mathcal{F}' =
f_*(\mathcal{F}' \amalg_\mathcal{F} \mathcal{F}')
\longrightarrow
f_*\mathcal{F}'
$$
も同型である。すると
$f_*\mathcal{F} \to f_*\mathcal{F}'$ は全射である。よって (6) は
(4) を含意する。$\mathcal{F} \to \mathcal{F}'$ が全射ならば、補題
\ref{sites-lemma-coequalizer-surjection} により $\mathcal{F}'$ は二つの射影
$\mathcal{F} \times_{\mathcal{F}'} \mathcal{F} \to \mathcal{F}$ の
余等化子である。したがって (5) が成り立つならば、
$f_*\mathcal{F}'$ は二つの射影
$$
f_*(\mathcal{F} \times_{\mathcal{F}'} \mathcal{F}) =
f_*\mathcal{F} \times_{f_*\mathcal{F}'} f_*\mathcal{F}
\longrightarrow
f_*\mathcal{F}
$$
の余等化子である。これは明らかに
$f_*\mathcal{F} \to f_*\mathcal{F}'$ が全射であることを意味する。
よって (5) も (4) を含意する。

\medskip\noindent
(f) の証明。(4) を仮定する。
$\mathcal{F} \to f^{-1}\mathcal{G}$ を $\mathcal{C}$ 上の層の全射とする。
(4) により $f_*\mathcal{F} \to f_*f^{-1}\mathcal{G}$ は全射である。
$\mathcal{G}'$ をファイバー積
$$
\xymatrix{
f_*\mathcal{F} \ar[r] & f_*f^{-1}\mathcal{G} \\
\mathcal{G}' \ar[u] \ar[r] & \mathcal{G} \ar[u]
}
$$
とする。このとき $\mathcal{G}' \to \mathcal{G}$ も全射である。
可換図式
$$
\xymatrix{
\mathcal{F} \ar[r] & f^{-1}\mathcal{G} \\
f^{-1}f_*\mathcal{F} \ar[r] \ar[u] & f^{-1}f_*f^{-1}\mathcal{G} \ar[u] \\
f^{-1}\mathcal{G}' \ar[u] \ar[r] & f^{-1}\mathcal{G} \ar[u]
}
$$
を考えると、所望の結果が得られる。逆に (11) を仮定する。
% P01-E0145
$a : \mathcal{F} \to \mathcal{F}'$ を $\mathcal{C}$ 上の層の全射とする。
ファイバー積の図式
$$
\xymatrix{
\mathcal{F} \ar[r] & \mathcal{F}' \\
\mathcal{F}'' \ar[u] \ar[r] & f^{-1}f_*\mathcal{F}' \ar[u]
}
$$
を考える。下の水平矢印は全射なので、(11) により全射
$\gamma : \mathcal{G}' \to f_*\mathcal{F}'$ であって、
$f^{-1}\gamma$ が
$\mathcal{F}'' \to f^{-1}f_*\mathcal{F}'$ を経由するものを選べる。
$$
\xymatrix{
& \mathcal{F} \ar[r] & \mathcal{F}' \\
f^{-1}\mathcal{G}' \ar[r] & \mathcal{F}'' \ar[u] \ar[r] &
f^{-1}f_*\mathcal{F}' \ar[u]
}
$$
これを $f_*$ によって順像に移すと、可換図式
$$
\xymatrix{
& f_*\mathcal{F} \ar[r] & f_*\mathcal{F}' \\
f_*f^{-1}\mathcal{G}' \ar[r] & f_*\mathcal{F}'' \ar[u] \ar[r] &
f_*f^{-1}f_*\mathcal{F}' \ar[u] \\
\mathcal{G}' \ar[u] \ar[rr] & & f_*\mathcal{F}' \ar[u]
}
$$
を得る。これにより (4) が成り立つ。

\medskip\noindent
(g) の証明。(9) を仮定する。以下、Categories, Lemma
\ref{categories-lemma-characterize-mono-epi} を断りなく用いる。
$a : \mathcal{F} \to \mathcal{F}'$ を $\mathcal{C}$ 上の層の写像で、
$f_*a$ が単射であるものとする。これは
$f_*\mathcal{F} \to f_*\mathcal{F} \times_{f_*\mathcal{F}'} f_*\mathcal{F} =
f_*(\mathcal{F} \times_{\mathcal{F}'} \mathcal{F})$ が同型であることを
意味する。よって (9) により
$\mathcal{F} \to \mathcal{F} \times_{\mathcal{F}'} \mathcal{F}$ は
全射、すなわち同型である。したがって $a$ は単射であり、(8) が
成り立つ。(10) は明らかに (8) $+$ (9) と同値なので、(g) の証明は
終わる。

\medskip\noindent
(h) の証明。これは Categories, Lemma
\ref{categories-lemma-adjoint-fully-faithful} である。
\end{proof}

\noindent
次は、関手 $f_*$ が全射を全射に移すことを保証するサイトの射の条件である。

\begin{lemma}
\label{sites-lemma-weaker}
$f : \mathcal{D} \to \mathcal{C}$ を、連続関手
$u : \mathcal{C} \to \mathcal{D}$ に付随するサイトの射とする。
$\mathcal{C}$ の任意の対象 $U$ および $\mathcal{D}$ の任意の被覆
$\{V_j \to u(U)\}$ に対して、$\mathcal{C}$ の被覆
$\{U_i \to U\}$ で、層の写像
$$
\coprod h_{u(U_i)}^\# \to h_{u(U)}^\#
$$
が層の写像
$$
\coprod h_{V_j}^\# \to h_{u(U)}^\#.
$$
を経由するものが存在すると仮定する。このとき $f_*$ は層の全射を
層の全射に移す。
\end{lemma}

\begin{proof}
$a : \mathcal{F} \to \mathcal{G}$ を $\mathcal{D}$ 上の層の全射とする。
$U$ を $\mathcal{C}$ の対象とし、
$s \in f_*\mathcal{G}(U) = \mathcal{G}(u(U))$ とする。仮定により、
被覆 $\{V_j \to u(U)\}$ と、$a(s_j) = s|_{V_j}$ を満たす切断
$s_j \in \mathcal{F}(V_j)$ が存在する。切断 $s$、$s_j$ および $a$ は、
層の写像の可換図式
$$
\xymatrix{
\coprod h_{V_j}^\# \ar[r]_-{\coprod s_j} \ar[d] & \mathcal{F} \ar[d]^a \\
h_{u(U)}^\# \ar[r]^s & \mathcal{G}
}
$$
を与えると考えられる。仮定により、上の可換図式を次のように拡張できる
被覆 $\{U_i \to U\}$ が存在する。
$$
\xymatrix{
& \coprod h_{V_j}^\# \ar[r]_-{\coprod s_j} \ar[d] & \mathcal{F} \ar[d]^a \\
\coprod h_{u(U_i)}^\# \ar[r] \ar[ur] &
h_{u(U)}^\# \ar[r]^s & \mathcal{G}
}
$$
$\mathcal{F}$ は層なので、左下隅から右上隅への写像は切断の族
$s_i \in \mathcal{F}(u(U_i))$、すなわち切断
$s_i \in f_*\mathcal{F}(U_i)$ に対応する。図式の可換性から、
$a(s_i)$ は $s$ の $U_i$ への制限に等しい。言い換えると、
$f_*a$ が層の全射であることを示した。
\end{proof}

\begin{example}
\label{sites-example-weaker-not-coequalizer}
$f : \mathcal{D} \to \mathcal{C}$ が補題 \ref{sites-lemma-weaker} の仮定を
満たすとする。このとき一般に、$f_*$ が余等化子または押し出しと
可換するとは限らない。実際、$f$ を位相空間の射
$X = \{1, 2\} \to Y = \{*\}$ に付随するサイトの射とする
（例 \ref{sites-example-continuous-map}）。ここで $Y$ は一元空間であり、
$X = \{1, 2\}$ は二点離散空間である。$X$ 上の層 $\mathcal{F}$ は
集合の対 $(A_1, A_2)$ によって与えられる。このとき
$f_*\mathcal{F}$ は集合 $A_1 \times A_2$ に対応する。したがって
$a = (a_1, a_2), b = (b_1, b_2) : (A_1, A_2) \to (B_1, B_2)$
が $X$ 上の層の写像ならば、$a, b$ の余等化子は $(C_1, C_2)$ である。
ここで $C_i$ は $a_i, b_i$ の余等化子である。一方、
$f_*a, f_*b$ の余等化子は
$$
a_1 \times a_2, b_1 \times b_2 :
A_1 \times A_2 \longrightarrow B_1 \times B_2
$$
の余等化子であり、一般には $C_1 \times C_2$ と異なる。実際、
$A_2 = \emptyset$ ならば $A_1 \times A_2 = \emptyset$ なので、表示した
矢印の余等化子は $B_1 \times B_2$ であるが、一般に
$C_1 \not = B_1$ である。押し出しについても同様の例が成り立つ。
\end{example}

\noindent
次の補題は、サイトの射の関手 $f_*$ が単射と全射を反映するための
判定条件を与える。これは $f_*$ が忠実であり、写像
$f^{-1}f_*\mathcal{F} \to \mathcal{F}$ が常に全射であることも含意する。

\begin{lemma}
\label{sites-lemma-cover-from-below}
$f : \mathcal{D} \to \mathcal{C}$ を、関手
$u : \mathcal{C} \to \mathcal{D}$ によって与えられるサイトの射とする。
$\mathcal{D}$ の各対象 $V$ に対して、$\mathcal{C}$ の対象 $U_i$ と射
$u(U_i) \to V$ で、$\{u(U_i) \to V\}$ が $\mathcal{D}$ の被覆となる
ものが存在すると仮定する。この場合、関手
$f_* : \Sh(\mathcal{D}) \to \Sh(\mathcal{C})$ は単射と全射を反映する。
\end{lemma}

\begin{proof}
% P01-E0146
$\alpha : \mathcal{F} \to \mathcal{G}$ を $\mathcal{D}$ 上の層の写像とする。
仮定により、$\mathcal{D}$ の各対象 $V$ に対して、ある
$U_i \in \Ob(\mathcal{C})$ が存在し、層条件から
$\mathcal{F}(V) \subset \prod \mathcal{F}(u(U_i)) =
\prod f_*\mathcal{F}(U_i)$ を得る。$\mathcal{G}$ についても同様である。
したがって $f_*\alpha$ が単射ならば $\alpha$ が単射であることは
明らかである。言い換えると、$f_*$ は単射を反映する。

\medskip\noindent
$f_*\alpha$ が全射であると仮定する。上のような
$V, U_i, u(U_i) \to V$ と切断 $s \in \mathcal{G}(V)$ に対して、
$s|_{u(U_{ij})}$ が $\mathcal{F}(u(U_{ij}))$ の像に属するような被覆
$\{U_{ij} \to U_i\}$ が存在する。$u$ は連続であり、サイトの公理を
用いると $\{u(U_{ij}) \to V\}$ は被覆なので、$s$ は局所的に像に
属すると結論できる。したがって $\alpha$ は全射である。
言い換えると、$f_*$ は全射を反映する。
\end{proof}

\begin{example}
\label{sites-example-cover-from-below}
補題 \ref{sites-lemma-cover-from-below} の仮定を満たす射
$f : \mathcal{D} \to \mathcal{C}$ を構成する。具体的には、
$\varphi : G \to H$ を有限群の射とする。可算 $G$-集合と可算 $H$-集合、
および共同全射な写像の可算族を被覆とするサイト
$\mathcal{D} = \mathcal{T}_G$ と $\mathcal{C} = \mathcal{T}_H$ を考える
（例 \ref{sites-example-site-on-group}）。節
\ref{sites-section-G-sets-morphisms} で述べた関手を
$u : \mathcal{T}_H \to \mathcal{T}_G$ とし、対応するサイトの射を
$f : \mathcal{T}_G \to \mathcal{T}_H$ とする。
$\varphi$ が単射ならば、各可算 $G$-集合は $G$-集合として可算
$H$-集合の商である（$\varphi$ が単射でなければこれは成り立たない）。
したがって $f$ は補題 \ref{sites-lemma-cover-from-below} の仮定を満たす。
$\mathcal{T}_G$ 上の層 $\mathcal{F}$ が $G$-集合 $S$ に対応するならば、
標準写像
$$
f^{-1}f_*\mathcal{F} \longrightarrow \mathcal{F}
$$
は写像
$$
\text{Map}_G(H, S) \longrightarrow S,\quad
a \longmapsto a(1_H)
$$
に対応する。$\varphi$ が単射だが全射でないならば、この写像は
全射である（補題 \ref{sites-lemma-cover-from-below} によれば、そうなるべきである）が、
一般には単射でない（例えば
$G = \{1\}$、$H = \{1, \sigma\}$、$S = \{1, 2\}$ とする）。
さらに関手 $f_*$ は余等化子または押し出しと可換しない
（$G = \{1\}$ および $H = \{1, \sigma\}$ とする）。
\end{example}

\section{ほとんど余連続な関手}
\label{sites-section-almost-cocontinuous}

\noindent
$\mathcal{C}$ をサイトとする。圏 $\textit{PSh}(\mathcal{C})$ は始対象を
もつ。これは $\mathcal{C}$ の各対象に空集合を対応させる前層である。
この前層を $\emptyset$ と書く。層化の性質から、$\emptyset$ の層化
$\emptyset^\#$ は $\mathcal{C}$ 上の層の圏 $\Sh(\mathcal{C})$ の
始対象である。

\begin{definition}
\label{sites-definition-empty}
$\mathcal{C}$ をサイトとする。$\mathcal{C}$ の対象 $U$ が
{\it 層論的に空}であるとは、$\emptyset^\# \to h_U^\#$ が
層の同型であることをいう。
\end{definition}

\noindent
次の補題はこの概念をより明示的にする。

\begin{lemma}
\label{sites-lemma-characterize-empty}
$\mathcal{C}$ をサイトとし、$U$ を $\mathcal{C}$ の対象とする。
次は同値である。
\begin{enumerate}
\item $U$ は層論的に空である。
\item 各層 $\mathcal{F}$ に対して $\mathcal{F}(U)$ は一元集合である。
\item $\emptyset^\#(U)$ は一元集合である。
\item $\emptyset^\#(U)$ は空でない。
\item 空な族は $\mathcal{C}$ における $U$ の被覆である。
\end{enumerate}
さらに $U$ が層論的に空ならば、$\mathcal{C}$ の任意の射
$U' \to U$ に対して $U'$ も層論的に空である。
\end{lemma}

\begin{proof}
任意の層 $\mathcal{F}$ に対して
$\mathcal{F}(U) = \Mor_{\Sh(\mathcal{C})}(h_U^\#, \mathcal{F})$ である。
したがって (1) と (2) は同値である。(2) が (3) を、(3) が (4) を
含意することは明らかである。$U$ の各被覆が空でない族によって
与えられるならば、プラス構成の定義により $\emptyset^+(U)$ は空である。
$\emptyset$ は分離前層なので $\emptyset^+ = \emptyset^\#$ であることに
注意する。定理 \ref{sites-theorem-plus} を参照されたい。したがって (4) は
(5) を含意する。(5) が成り立つならば、$\mathcal{F}$ の層条件により、
各層 $\mathcal{F}$ に対して $\mathcal{F}(U)$ は一元集合である。
注意 \ref{sites-remark-sheaf-condition-empty-covering} を参照されたい。
よって (5) は (2) を含意し、(1)--(5) は同値である。補題の最後の主張は、
% P01-E0147
定義 \ref{sites-definition-site} の公理 (3) を $U$ の空被覆に適用して従う。
\end{proof}

\begin{definition}
\label{sites-definition-almost-cocontinuous}
$\mathcal{C}$、$\mathcal{D}$ をサイトとし、
$u : \mathcal{C} \to \mathcal{D}$ を関手とする。$\mathcal{C}$ の各対象
$U$ および各被覆 $\{V_j \to u(U)\}_{j \in J}$ に対して、
$\mathcal{C}$ の被覆 $\{U_i \to U\}_{i \in I}$ で、各 $I$ の $i$ が
次の二条件の少なくとも一つを満たすものが存在するとき、$u$ は
{\it ほとんど余連続}であるという。
\begin{enumerate}
\item $u(U_i)$ は層論的に空である。
\item ある $j \in J$ に対して射 $u(U_i) \to u(U)$ は $V_j$ を経由する。
\end{enumerate}
\end{definition}

\noindent
この定義の動機は位相空間の閉埋入 $i : Z \to X$ から来る。
例 \ref{sites-example-closed-map-cocontinuous-false} で論じたように、連続関手
$X_{Zar} \to Z_{Zar}$、$U \mapsto Z \cap U$ は余連続でない。
しかし上の意味ではほとんど余連続である。$i_*$ は集合の層上では
完全でないが、アーベル群の層上では完全であることが知られている。
% P01-E0148
Sheaves, Remark \ref{sheaves-remark-i-star-not-exact} を参照されたい。
これは一般の連続かつほとんど余連続な関手について成り立つ。

\begin{lemma}
\label{sites-lemma-almost-cocontinuous-sheafification}
$\mathcal{C}$、$\mathcal{D}$ をサイトとし、
$u : \mathcal{C} \to \mathcal{D}$ を関手とする。$u$ は連続かつ
ほとんど余連続であると仮定する。$\mathcal{G}$ を $\mathcal{D}$ 上の
前層で、$V$ が層論的に空であるときは常に $\mathcal{G}(V)$ が
一元集合となるものとする。このとき
$(u^p\mathcal{G})^\# = u^p(\mathcal{G}^\#)$ である。
\end{lemma}

\begin{proof}
$U \in \Ob(\mathcal{C})$ とする。
$(u^p\mathcal{G})^\#(U) = u^p(\mathcal{G}^\#)(U)$ を示さなければならない。
$\mathcal{G}^+$ も補題の仮定を満たす前層なので、
$(u^p\mathcal{G})^+(U) = u^p(\mathcal{G}^+)(U)$ を示せば十分である。
$$
u^p(\mathcal{G}^+)(U) =
\mathcal{G}^+(u(U)) =
\colim_\mathcal{V} \check H^0(\mathcal{V}, \mathcal{G})
$$
である。ここで余極限は $\mathcal{D}$ における $u(U)$ の被覆
$\mathcal{V}$ 全体にわたる。一方、
$$
u^p(\mathcal{G})^+(U) =
\colim_\mathcal{U} \check H^0(u(\mathcal{U}), \mathcal{G})
$$
である。ここで余極限は、$\mathcal{C}$ における $U$ の被覆の圏
$\mathcal{U} = \{U_i \to U\}_{i \in I}$ 全体にわたり、
$u(\mathcal{U}) = \{u(U_i) \to u(U)\}_{i \in I}$ である。
$u$ が連続であるという条件は、各 $u(\mathcal{U})$ が被覆であることを
意味する。$I = I_1 \amalg I_2$ と書く。ここで
$$
I_2 = \{i \in I \mid u(U_i)\text{ is sheaf theoretically empty}\}
$$
である。このとき
$u(\mathcal{U})' = \{u(U_i) \to u(U)\}_{i \in I_1}$ は依然として
被覆である。実際、残りの各部分は空な族で被覆できるので、定義
% P01-E0149
\ref{sites-definition-site} の公理 (2) によって取り除ける。さらに
$\check H^0(u(\mathcal{U}), \mathcal{G}) =
\check H^0(u(\mathcal{U})', \mathcal{G})$ が $\mathcal{G}$ に関する
仮定から従う。最後に $u$ がほとんど余連続であるという条件により、
$u(U)$ の各被覆 $\mathcal{V}$ に対して、$u(\mathcal{U})'$ が
$\mathcal{V}$ を細分するような $U$ の被覆 $\mathcal{U}$ が存在する。
したがって上に表示した二つの余極限は同じ値をもつ。
\end{proof}

\begin{lemma}
\label{sites-lemma-continuous-almost-cocontinuous}
$\mathcal{C}$、$\mathcal{D}$ をサイトとし、
$u : \mathcal{C} \to \mathcal{D}$ を関手とする。$u$ は連続かつ
ほとんど余連続であると仮定する。このとき
$u^s = u^p : \Sh(\mathcal{D}) \to \Sh(\mathcal{C})$ は押し出しおよび
余等化子（より一般に有限連結余極限）と可換する。
\end{lemma}

\begin{proof}
$\mathcal{I}$ を有限連結な添字圏とする。図式
$\mathcal{I} \to \Sh(\mathcal{D})$、$i \mapsto \mathcal{G}_i$ をとる。
% P01-E0150
この図式の余極限は前層の圏における余極限の層化であることが知られている。
補題 \ref{sites-lemma-colimit-sheaves} を参照されたい。前層の圏における
余極限を $\colim^{Psh}$ と書く。$\mathcal{I}$ は有限連結なので、
$\colim^{Psh}_i \mathcal{G}_i$ は補題
\ref{sites-lemma-almost-cocontinuous-sheafification} の仮定を満たす前層である
（一元集合の有限連結余極限は一元集合だからである）。よって同補題から
% P01-E0151
\begin{align*}
u^s(\colim_i \mathcal{G}_i) & =
u^s((\colim^{Psh}_i \mathcal{G}_i)^\#) \\
& = (u^p(\colim^{Psh}_i \mathcal{G}_i))^\# \\
& = (\colim^{PSh}_i u^p(\mathcal{G}_i))^\# \\
& = \colim_i u^s(\mathcal{G}_i)
\end{align*}
を得る。これが所望の等式である。
\end{proof}

\begin{lemma}
\label{sites-lemma-morphism-of-sites-almost-cocontinuous}
$f : \mathcal{D} \to \mathcal{C}$ を、連続関手
$u : \mathcal{C} \to \mathcal{D}$ に付随するサイトの射とする。
$u$ がほとんど余連続ならば、$f_*$ は押し出しおよび余等化子
（より一般に有限連結余極限）と可換する。
\end{lemma}

\begin{proof}
これは補題 \ref{sites-lemma-continuous-almost-cocontinuous} の特別な場合である。
\end{proof}








\section{部分トポス}
\label{sites-section-subtopoi}

\noindent
定義は次の通りである。

\begin{definition}
\label{sites-definition-embedding}
$\mathcal{C}$ と $\mathcal{D}$ をサイトとする。トポスの射
$f : \Sh(\mathcal{D}) \to \Sh(\mathcal{C})$ について、$f_*$ が
充満忠実であるとき、これを{\it 埋め込み}という。
\end{definition}

\noindent
補題 \ref{sites-lemma-exactness-properties} によれば、これは
$\mathcal{D}$ 上の各層 $\mathcal{F}$ に対して随伴写像
$f^{-1}f_*\mathcal{F} \to \mathcal{F}$ が同型であることと同値である。

\begin{definition}
\label{sites-definition-subtopos}
$\mathcal{C}$ をサイトとする。厳密充満部分圏
$E \subset \Sh(\mathcal{C})$ が{\it 部分トポス}であるとは、
トポスの埋め込み $f : \Sh(\mathcal{D}) \to \Sh(\mathcal{C})$ が存在して、
$E$ が関手 $f_*$ の本質像に等しいことをいう。
\end{definition}

\noindent
次の補題で構成する部分トポスは、後の定義で「開」と呼ばれる。

\begin{lemma}
\label{sites-lemma-open-subtopos}
$\mathcal{C}$ をサイトとし、$\mathcal{F}$ を $\mathcal{C}$ 上の層とする。
次は同値である。
\begin{enumerate}
\item $\mathcal{F}$ は $\Sh(\mathcal{C})$ の終対象の部分対象である。
\item トポス $\Sh(\mathcal{C})/\mathcal{F}$ は
$\Sh(\mathcal{C})$ の部分トポスである。
\end{enumerate}
\end{lemma}

\begin{proof}
補題 \ref{sites-lemma-localize-topos} で
$\Sh(\mathcal{C})/\mathcal{F}$ がトポスであることを見た。
実際、その証明を思い出そう。まず補題 \ref{sites-lemma-topos-good-site} を
適用すると、$\mathcal{C}$ は劣標準位相、ファイバー積、終対象 $X$、
および $\mathcal{F} = h_U$ を満たす対象 $U$ をもつサイトであると
仮定できる。補題 \ref{sites-lemma-localize-topos} の証明によれば、トポスの射
$j_\mathcal{F} : \Sh(\mathcal{C})/\mathcal{F} \to \Sh(\mathcal{C})$ は、
ある同一視のもとで局所化射
$j_U : \Sh(\mathcal{C}/U) \to \Sh(\mathcal{C})$ に等しい。

\medskip\noindent
(2) を仮定する。これは $\mathcal{C}/U$ 上のすべての層
$\mathcal{G}$ に対して
$j_U^{-1}j_{U, *}\mathcal{G} \to \mathcal{G}$ が同型であることを
意味する。$\mathcal{C}/U$ の任意の対象 $Z/U$ に対して、補題
\ref{sites-lemma-localize-given-products} により
$$
(j_{U, *}h_{Z/U})(U) = \Mor_{\mathcal{C}/U}(U \times_X U/U, Z/U)
$$
である。上の等式で $\mathcal{G} = h_{Z/U}$ とおくと、すべての
$Z/U$ に対して
$$
\Mor_{\mathcal{C}/U}(U \times_X U/U, Z/U) = \Mor_{\mathcal{C}/U}(U, Z/U)
$$
を得る。米田の補題（Categories, Lemma
\ref{categories-lemma-yoneda}）により、これは
$U \times_X U = U$ を意味する。Categories, Lemma
\ref{categories-lemma-characterize-mono-epi} により $U \to X$ は
モノ射である。言い換えると (1) が成り立つ。

\medskip\noindent
(1) を仮定する。このとき補題 \ref{sites-lemma-restrict-back} により
$j_U^{-1} j_{U, *} = \text{id}$ である。
\end{proof}

\begin{definition}
\label{sites-definition-open-subtopos}
$\mathcal{C}$ をサイトとする。厳密充満部分圏
$E \subset \Sh(\mathcal{C})$ が{\it 開部分トポス}であるとは、
$\Sh(\mathcal{C})$ の終対象の部分層 $\mathcal{F}$ が存在して、
$E$ が補題 \ref{sites-lemma-open-subtopos} で述べた部分トポス
$\Sh(\mathcal{C})/\mathcal{F}$ であることをいう。
\end{definition}

\noindent
これは $\Sh(\mathcal{C})$ の開部分トポスの集まりと
$\Sh(\mathcal{C})$ の終対象の部分対象全体の集合との間に全単射があることを
意味する。開部分トポスが与えられると「閉」補集合が存在する。

\begin{lemma}
\label{sites-lemma-closed-subtopos}
$\mathcal{C}$ をサイトとする。$\mathcal{F}$ を
$\Sh(\mathcal{C})$ の終対象 $*$ の部分層とする。
$\mathcal{F} \times \mathcal{G} \to \mathcal{F}$ が同型となる層
$\mathcal{G}$ 全体の充満部分圏は $\Sh(\mathcal{C})$ の部分トポスである。
\end{lemma}

\begin{proof}
補題 \ref{sites-lemma-topos-good-site} を適用して、$\mathcal{C}$ は同補題に
挙げられた性質をもつサイトであると仮定できる。特に
$\mathcal{C}$ は終対象 $X$（したがって $* = h_X$）と、
$\mathcal{F} = h_U$ を満たす対象 $U$ をもつ。

\medskip\noindent
圏としては $\mathcal{D} = \mathcal{C}$ とする。ただし、射の族
% P01-E0152
$\{V_j \to V\}$ が被覆であるとは、
$\{V_i \to V\} \cup \{U \times_X V \to V\}$ が被覆であることとする。
$\mathcal{C}$ の選び方により、これはちょうど
$$
h_{U \times_X V} \amalg \coprod h_{V_i} \longrightarrow h_V
$$
が全射であることを意味する。$\mathcal{D}$ がサイト、すなわち被覆が
定義 \ref{sites-definition-site} の条件 (1)、(2)、(3) を満たすと主張する。
条件 (1) は成り立つ。条件 (2) について、$\{V_i \to V\}$ と
$\{V_{ij} \to V_i\}$ を $\mathcal{D}$ の被覆とする。このとき合成
$$
\coprod \left(
h_{U \times_X V_i} \amalg \coprod h_{V_{ij}}
\right) \longrightarrow
h_{U \times_X V} \amalg \coprod h_{V_i} \longrightarrow h_V
$$
は全射である。各射 $U \times_X V_i \to V$ は
$U \times_X V$ を経由するので、
$$
h_{U \times_X V} \amalg \coprod h_{V_{ij}} \longrightarrow h_V
$$
は全射である。すなわち $\{V_{ij} \to V\}$ は $\mathcal{D}$ における
$V$ の被覆である。条件 (3) も同様に、層の全射の基底変換が全射である
ことから従う。

\medskip\noindent
（恒等）関手 $u : \mathcal{C} \to \mathcal{D}$ は連続であり、
ファイバー積および終対象と可換することに注意する。したがって命題
\ref{sites-proposition-get-morphism} によりサイトの射
$f : \mathcal{D} \to \mathcal{C}$ を得る。$f_*$ は台となる前層上の
恒等関手なので、充満忠実である。証明を終えるには、$f_*$ の本質像が
補題で指定された充満部分圏 $E \subset \Sh(\mathcal{C})$ であることを
示さなければならない。そのため、
$\mathcal{G} \in \Ob(\Sh(\mathcal{C}))$ が $E$ に属するための
必要十分条件は、$\mathcal{C}$ のすべての対象 $V$ に対して
$\mathcal{G}(U \times_X V)$ が一元集合であることだと注意する。
したがって、このような層は $\mathcal{D}$ のすべての被覆について
層条件を満たす（議論は省略する）。逆に $\mathcal{G}$ が
$\mathcal{D}$ のすべての被覆について層条件を満たすならば、
$\mathcal{D}$ において対象 $U \times_X V$ は空被覆で被覆されるので、
$\mathcal{G}(U \times_X V)$ は一元集合である。
\end{proof}

\begin{definition}
\label{sites-definition-closed-subtopos}
$\mathcal{C}$ をサイトとする。厳密充満部分圏
$E \subset \Sh(\mathcal{C})$ が{\it 閉部分トポス}であるとは、
$\Sh(\mathcal{C})$ の終対象の部分層 $\mathcal{F}$ が存在して、
$E$ が補題 \ref{sites-lemma-closed-subtopos} で述べた部分トポスであることをいう。
\end{definition}

\noindent
これで、トポスの閉埋入と開埋入の意味を定義できる。

\begin{definition}
\label{sites-definition-immersion-topoi}
$f : \Sh(\mathcal{D}) \to \Sh(\mathcal{C})$ をトポスの射とする。
\begin{enumerate}
\item $f$ が埋め込みであり、$f_*$ の本質像が開部分トポスであるとき、
$f$ を{\it 開埋入}という。
\item $f$ が埋め込みであり、$f_*$ の本質像が閉部分トポスであるとき、
$f$ を{\it 閉埋入}という。
\end{enumerate}
\end{definition}

\begin{lemma}
\label{sites-lemma-closed-immersion}
$i : \Sh(\mathcal{D}) \to \Sh(\mathcal{C})$ をトポスの閉埋入とする。
このとき $i_*$ は充満忠実であり、全射を全射に移し、余等化子および
押し出しと可換し、単射、全射、および全単射を反映する。
\end{lemma}

\begin{proof}
$\mathcal{F}$ を $\Sh(\mathcal{C})$ の終対象 $*$ の部分層とし、
$E \subset \Sh(\mathcal{C})$ を、
$\mathcal{F} \times \mathcal{G} \to \mathcal{F}$ が同型となる
$\mathcal{G}$ 全体からなる充満部分圏とする。補題
\ref{sites-lemma-closed-subtopos} により、関手 $i_*$ は包含関手
$\iota : E \to \Sh(\mathcal{C})$ と同型である。

\medskip\noindent
$j_{\mathcal{F}} : \Sh(\mathcal{C})/\mathcal{F} \to \Sh(\mathcal{C})$ を
局所化関手とする（補題 \ref{sites-lemma-localize-topos}）。$E$ は
$j_\mathcal{F}^{-1}\mathcal{G} = *$ を満たす層 $\mathcal{G}$ 全体としても
記述できることに注意する。

\medskip\noindent
% P01-E0153
$a, b : \mathcal{G}_1 \to \mathcal{G}_2$ を $E$ の二つの射とする。
$\iota$ が余等化子と可換することを証明するには、
$\Sh(\mathcal{C})$ における $a$, $b$ の余等化子が $E$ に属することを
示せば十分である。二つの射 $* \to *$ の余等化子は $*$ であり、
$j_\mathcal{F}^{-1}$ は完全なので、これは明らかである。
押し出しについても同様である。

\medskip\noindent
したがって $i_*$ は補題 \ref{sites-lemma-exactness-properties} の性質
(5)、(6)、(7) を満たし、ゆえに射 $i$ は同補題で述べたすべての性質、
特に本補題で述べた性質を満たす。
\end{proof}







\section{代数構造の層}
\label{sites-section-sheaves-algebraic-structures}

\noindent
Sheaves, Section \ref{sheaves-section-algebraic-structures} において、
代数構造の型を対 $(\mathcal{A}, s)$ として導入した。ここで
$\mathcal{A}$ は圏であり、
$s : \mathcal{A} \to \textit{Sets}$ は次の性質をもつ関手である。
\begin{enumerate}
\item $s$ は忠実である。
\item $\mathcal{A}$ は極限をもち、$s$ は極限と可換する。
\item $\mathcal{A}$ はフィルター付き余極限をもち、
$s$ はそれらと可換する。
\item $s$ は同型を反映する。
\end{enumerate}
このような代数構造の型について、空間 $X$ 上の
$\mathcal{A}$ に値をとる前層 $\mathcal{F}$ が層であることと、
付随する集合の前層が層であることが同値であることを見た。さらに、
茎の概念を具体的に記述し、連続写像 $f : X \to Y$ が与えられたとき、
代数構造の層に対する互いに随伴な順像関手と逆像関手を定義した。
これらは台となる集合の層に対する順像および逆像と一致する。また、
代数構造の層を基底から空間のすべての開集合へ拡張する操作も
期待どおりに働く。
% P01-E0154

\medskip\noindent
この内容の一部は、サイトと層の設定でもそのまま成り立つ。
$(\mathcal{A}, s)$ を代数構造の型とし、$\mathcal{C}$ をサイトとする。
$\mathcal{C}$ 上の $\mathcal{A}$ に値をとる前層の圏と層の圏を、それぞれ
$\textit{PSh}(\mathcal{C}, \mathcal{A})$、
$\Sh(\mathcal{C}, \mathcal{A})$ と書く。
\begin{itemize}
\item[] ($\alpha$) $\mathcal{A}$ に値をとる前層が層であることと、
その台となる集合の前層が層であることは同値である。
Sheaves, Lemma \ref{sheaves-lemma-sheaves-structure} の証明を参照せよ。
\item[] ($\beta$) $\mathcal{A}$ に値をとる前層 $\mathcal{F}$ が
与えられたとき、前層 ${\mathcal{F}}^\# = (\mathcal{F}^+)^+$ は層である。
実際、層化の過程に現れる余極限はフィルター付きであり、さらに
有向集合にわたる余極限である（Section \ref{sites-section-sheafification}、
特に Lemma \ref{sites-lemma-sheafification-exact} の証明を参照）。また、
$s$ はフィルター付き余極限と可換する。
\item[] ($\gamma$) 次の可換図式を得る。
$$
\xymatrix{
\Sh(\mathcal{C}, \mathcal{A}) \ar@<1ex>[r] \ar[d]_s &
\textit{PSh}(\mathcal{C}, \mathcal{A}) \ar@<1ex>[l]^{{}^\#} \ar[d]^s\\
\Sh(\mathcal{C}) \ar@<1ex>[r] &
\textit{PSh}(\mathcal{C}) \ar@<1ex>[l]
}
$$
\item[] ($\delta$) $\mathcal{F} = \mathcal{F}^\#$ であることと、
$\mathcal{F}$ が代数構造の層であることは同値である。
\item[] ($\epsilon$) 関手 ${}^\#$ は包含関手の左随伴である。
$$
\Mor_{\textit{PSh}(\mathcal{C}, \mathcal{A})}(\mathcal{G}, \mathcal{F})
=
\Mor_{\Sh(\mathcal{C}, \mathcal{A})}(\mathcal{G}^\#, \mathcal{F})
$$
証明は Proposition \ref{sites-proposition-sheafification-adjoint} の証明と同じである。
\item[] ($\zeta$) 関手
$\mathcal{F} \mapsto \mathcal{F}^\#$ は左完全である。
証明は Lemma \ref{sites-lemma-sheafification-exact} の証明と同じである。
\end{itemize}

\begin{definition}
\label{sites-definition-pushforward-algebraic-structures}
$f : \mathcal{D} \to \mathcal{C}$ を、関手
$u : \mathcal{C} \to \mathcal{D}$ によって与えられるサイトの射とする。
代数構造の前層に対する{\it 順像}関手を規則
$u^p\mathcal{F}(U) = \mathcal{F}(uU)$ により定義する。また、
代数構造の層に対しても同じ規則、すなわち
$f_*\mathcal{F}(U) = \mathcal{F}(uU)$ により定義する。
\end{definition}

\noindent
逆像を定義しようとすると問題が生じる。その理由は、Section
\ref{sites-section-functoriality-PSh} において関手 $u_p$ を定義する余極限が
フィルター付きとは限らないことにある。したがって一般には、上の公理だけでは
代数構造の（前）層の逆像を定義するには不十分である。それでも、ほとんどすべての
場合に、次の補題によって代数構造の（前）層の順像と逆像を定義できる。
% P01-E0155

\begin{lemma}
\label{sites-lemma-push-pull-good-case}
関手 $u : \mathcal{C} \to \mathcal{D}$ が Proposition
\ref{sites-proposition-get-morphism} の仮定を満たし、したがってサイトの射
$f : \mathcal{D} \to \mathcal{C}$ を定めるとする。この場合、
逆像関手 $f^{-1}$（それぞれ $u_p$）と順像関手
$f_*$（それぞれ $u^p$）は、代数構造の層（それぞれ前層）の圏上の
随伴関手の対へ拡張される。さらに、これらの関手は台となる集合の層
（それぞれ前層）を取る操作と可換する。
\end{lemma}

\begin{proof}
$f_* = u^p$ は上で定義した。Proposition
\ref{sites-proposition-get-morphism} の証明において、その命題の仮定のもとでは
$u_p$ を定義するために用いるすべての余極限がフィルター付きであることを見た。
したがって、代数構造の型の定義から、代数構造の前層上の関手 $u_p$ を
まったく同じ余極限によって定義できる。$u_p$ と $u^p$ の随伴性は、
Lemma \ref{sites-lemma-adjoints-u} の証明とまったく同じ方法で証明される。
上で行った代数構造の前層の層化に関する議論から、さらに
$f^{-1}(\mathcal{F}) = (u_p\mathcal{F})^\#$ と定義できる。
\end{proof}

\noindent
一般のサイトの射、さらに一般には任意のトポスの射に対して、
逆像と順像を扱う一つの方法を簡単に説明する。

\medskip\noindent
$\mathcal{C}$ をサイトとする。$\mathcal{A} = \textit{Ab}$ の場合、
$\mathcal{C}$ 上のアーベル（前）層を四つ組
$(\mathcal{F}, +, 0, i)$ と考えることができる。ここでデータは次のとおりである。
\begin{enumerate}
\item[(D1)] $\mathcal{F}$ は集合の層である。
\item[(D2)] $+ : \mathcal{F} \times \mathcal{F} \to \mathcal{F}$ は
集合の層の射である。
\item[(D3)] $0 : * \to \mathcal{F}$ は、一点層（Example
\ref{sites-example-singleton-sheaf} を参照）から $\mathcal{F}$ への射である。
\item[(D4)] $i : \mathcal{F} \to \mathcal{F}$ は集合の層の射である。
\end{enumerate}
これらのデータは次の公理を満たさなければならない。
\begin{enumerate}
\item[(A1)] $+$ は結合的かつ可換である。
\item[(A2)] $0$ は $+$ の単位元である。
\item[(A3)] $+ \circ (1, i) = 0 \circ (\mathcal{F} \to *)$。
\end{enumerate}
Sheaves, Lemma \ref{sheaves-lemma-abelian-presheaves} と比較せよ。
$f : \mathcal{D} \to \mathcal{C}$ をサイトの射とする。
$f^{-1}$ は完全なので
$f^{-1}* = *$ かつ
$f^{-1}(\mathcal{F} \times \mathcal{F}) =
f^{-1}\mathcal{F} \times f^{-1}\mathcal{F}$ であることに注意する。
したがって、$f^{-1}\mathcal{F}$ を四つ組
$(f^{-1}\mathcal{F}, f^{-1}+, f^{-1}0, f^{-1}i)$ としてそのまま定義できる。
$f^{-1}$ は有限極限と可換する関手なので、公理は保たれる。最後に、
$f_*$ と $f^{-1}$ が通常どおり随伴することは容易に確認できる。

\medskip\noindent
\cite{SGA4} ではこの方法が用いられている。そこでは
``{\it esp\`ece the structure alg\'ebrique $\ll$d\'efinie
par limites projectives finie$\gg$}'' と呼ばれるものを導入する。
このような esp\`ece に対して、上の方法により随伴関手の対
$f^{-1}$ と $f_*$ を定義できる。これは、実際に現れるほとんどの
代数構造に明らかに適用できる。この構成を形式化する代わりに、
この方法が適用できる代数構造を列挙し、その都度確認することにする。
実際、この方法は任意のトポスの射に対して成り立つ。
% P01-E0156

\begin{proposition}
\label{sites-proposition-functoriality-algebraic-structures-topoi}
\begin{slogan}
トポスの射は代数構造を保つ。
\end{slogan}
$\mathcal{C}$、$\mathcal{D}$ をサイトとする。
$f = (f^{-1}, f_*)$ を
$\Sh(\mathcal{D}) \to \Sh(\mathcal{C})$ なるトポスの射とする。
上で導入した方法は、次の代数構造の型について、代数構造の層上の
随伴関手の対 $(f^{-1}, f_*)$ を与え、台となる集合の層を取る操作と両立する。
\begin{enumerate}
\item 点付き集合、
\item アーベル群、
\item 群、
\item モノイド、
\item 環、
\item 固定された環上の加群、および
\item 固定された体上のリー代数。
\end{enumerate}
さらに、これらの各場合に、上で ($\alpha$)、($\beta$)、($\gamma$)、
($\delta$)、($\epsilon$)、($\zeta$) と記した結果が成り立つ。
\end{proposition}

\begin{proof}
命題の最後の主張は、列挙した各圏が明らかな忘却関手を備えると、
この節の冒頭で説明した意味で実際に代数構造の型となることから直ちに従う。
Sheaves, Lemma \ref{sheaves-lemma-list-algebraic-structures} を参照せよ。

\medskip\noindent
(2) の証明。命題の前の議論で説明したように、アーベル群の層を四つ組
$(\mathcal{F}, +, 0, i)$ と考える。
$(\mathcal{F}, +, 0, i)$ が $\mathcal{C}$ 上にあるならば、その逆像を
$(f^{-1}\mathcal{F}, f^{-1}+, f^{-1}0, f^{-1}i)$ と定義する。
$(\mathcal{G}, +, 0, i)$ が $\mathcal{D}$ 上にあるならば、その順像を
$(f_*\mathcal{G}, f_*+, f_*0, f_*i)$ と定義する。これは
$f_*(\mathcal{G} \times \mathcal{G}) = f_*\mathcal{G} \times f_*\mathcal{G}$
だから可能である。随伴性は、台となる集合に対する関手の随伴性から従う。実際、
$$
\Mor_{\textit{Ab}(\mathcal{C})}
((\mathcal{F}_1, +, 0, i), (\mathcal{F}_2, +, 0, i))
=
\left\{
\begin{matrix}
\varphi \in \Mor_{\Sh(\mathcal{C})}
(\mathcal{F}_1, \mathcal{F}_2) \\
\varphi \text{ is compatible with }+, 0, i
\end{matrix}
\right\}
$$
である。詳細は読者に委ねる。

\medskip\noindent
この方法は環の層に対しても成り立つ。単位元をもつ環の層を、
どの初等代数の教科書にも載っている公理の列を満たす六つ組
$(\mathcal{O}, + , 0, i, \cdot, 1)$ と考えればよい。

\medskip\noindent
点付き集合の層とは対 $(\mathcal{F}, p)$ であり、ここで
$\mathcal{F}$ は集合の層、$p : * \to \mathcal{F}$ は集合の層の写像である。

\medskip\noindent
群の層は、適切な公理をもつ四つ組 $(\mathcal{F}, \cdot, 1, i)$ によって与えられる。

\medskip\noindent
モノイドの層は、適切な公理をもつ対 $(\mathcal{F}, \cdot)$ によって与えられる。
% P01-E0157

\medskip\noindent
$R$ を環とする。$R$-加群の層は五つ組
$(\mathcal{F}, +, 0, i, \{\lambda_r\}_{r \in R})$ によって与えられる。
ここで四つ組 $(\mathcal{F}, +, 0, i)$ は上で述べたアーベル群の層であり、
$\lambda_r : \mathcal{F} \to \mathcal{F}$ は集合の層の射の族であって、
次を満たす。
% P01-E0158
$\lambda_r \circ 0 = 0$、
$\lambda_r \circ + = + \circ (\lambda_r, \lambda_r)$、
$\lambda_{r + r'} =
+ \circ \lambda_r \times \lambda_{r'} \circ (\text{id}, \text{id})$、
$\lambda_{rr'} = \lambda_r \circ \lambda_{r'}$、
$\lambda_1 = \text{id}$、$\lambda_0 = 0 \circ (\mathcal{F} \to *)$。
% P01-E0159
\end{proof}

\noindent
環の層上の加群の層の圏については、Modules on Sites, Section
\ref{sites-modules-section-sheaves-modules} で論じる。

\begin{remark}
\label{sites-remark-no-pullback-presheaves}
$\mathcal{C}$、$\mathcal{D}$ をサイトとする。
$u : \mathcal{D} \to \mathcal{C}$ を、サイトの射
$\mathcal{C} \to \mathcal{D}$ を定める連続関手とする。
アーベル群の場合でさえ、アーベル群の前層に対する逆像関手はまだ
定義していないことに注意する。アーベル群の圏ではすべての余極限が表現できるので、
集合の前層上の $u_p$ を定義するのに用いたものと同じ余極限によって、
アーベル前層上の関手 $u_p^{ab}$ を確かに定義できる。また、アーベル前層の圏上で
$u_p^{ab}$ は $u^p$ の左随伴となる。しかし、台となる集合の前層上で
$u_p^{ab}$ が $u_p$ と常に一致するとは限らない。
\end{remark}
\section{引き戻し写像}
\label{sites-section-pullback}

\noindent
サイト $\mathcal{C}$ が終対象をもたないことがある。この場合、
大域切断関手を次のように定義する。

\begin{definition}
\label{sites-definition-global-sections}
サイト $\mathcal{C}$ 上の集合の前層 $\mathcal{F}$ の
{\it 大域切断}とは、集合
$$
\Gamma(\mathcal{C}, \mathcal{F}) =
\Mor_{\textit{PSh}(\mathcal{C})}(*, \mathcal{F})
$$
の元をいう。ここで $*$ は $\mathcal{C}$ 上の前層の圏の終対象、
すなわち各対象に一点集合を対応させる前層である。
% P01-E0160
\end{definition}

\noindent
もちろん、同じ定義は層にも適用される。大域切断を計算する一つの方法を示す。

\begin{lemma}
\label{sites-lemma-compute-global-sections}
$\mathcal{C}$ をサイトとする。$a, b : V \to U$ を $\mathcal{C}$ の対象とし、
$$
\xymatrix{
h_V^\# \ar@<1ex>[r] \ar@<-1ex>[r] & h_U^\# \ar[r] & {*}
}
$$
が $\Sh(\mathcal{C})$ における余等化子であるとする。このとき
$\Gamma(\mathcal{C}, \mathcal{F})$ は
$a^*, b^* : \mathcal{F}(U) \to \mathcal{F}(V)$ の等化子である。
% P01-E0161
\end{lemma}

\begin{proof}
$\Mor_{\Sh(\mathcal{C})}(h_U^\#, \mathcal{F}) = \mathcal{F}(U)$ なので、
定義から明らかである。
\end{proof}

\noindent
さて、$f : \Sh(\mathcal{D}) \to \Sh(\mathcal{C})$ をトポスの射とする。
このとき、$\mathcal{C}$ 上の任意の層 $\mathcal{F}$ に対して引き戻し写像
$$
f^{-1} :
\Gamma(\mathcal{C}, \mathcal{F})
\longrightarrow
\Gamma(\mathcal{D}, f^{-1}\mathcal{F})
$$
がある。実際、$f^{-1}$ は完全なので $*$ を $*$ に写す。したがって、
$\mathcal{C}$ 上の $\mathcal{F}$ の大域切断 $s$、すなわち層の写像
$s : * \to \mathcal{F}$ は
$f^{-1}s : * = f^{-1}* \to f^{-1}\mathcal{F}$ へ引き戻せる。

\medskip\noindent
これを少し一般化して、$\mathcal{C}$、$\mathcal{D}$ 上の層の対
$\mathcal{F}$、$\mathcal{G}$ と、写像 $f^{-1}\mathcal{F} \to \mathcal{G}$
を考えることができる。このとき上の構成と明らかな写像
$\Gamma(\mathcal{D}, f^{-1}\mathcal{F}) \to \Gamma(\mathcal{D}, \mathcal{G})$
を合成して、写像
$$
\Gamma(\mathcal{C}, \mathcal{F})
\longrightarrow
\Gamma(\mathcal{D}, \mathcal{G})
$$
を得る。この写像も引き戻し写像と呼ばれることがある。

\medskip\noindent
自然に頻繁に現れる、もう少し一般的な構成は次のとおりである。
トポスの射の可換図式
$$
\xymatrix{
\Sh(\mathcal{D}) \ar[rd]_h \ar[rr]_f & &
\Sh(\mathcal{C}) \ar[ld]^g \\
& \Sh(\mathcal{B})
}
$$
があるとする。次に、$\mathcal{C}$ 上の層 $\mathcal{F}$ があるとする。
このとき{\it 引き戻し写像}
$$
f^{-1} : g_*\mathcal{F} \longrightarrow h_*f^{-1}\mathcal{F}
$$
がある。これは、同一視
$g_*f_*f^{-1}\mathcal{F} = h_*f^{-1}\mathcal{F}$ と、標準写像
$\mathcal{F} \to f_*f^{-1}\mathcal{F}$ に $g_*$ を適用して得られる写像から
生じる写像にほかならない。$g$ が恒等射ならば、大域切断上のこの写像は
上の引き戻し写像と一致する。

\medskip\noindent
前段の状況で、$\mathcal{C}$、$\mathcal{D}$ 上の層の対
$\mathcal{F}$、$\mathcal{G}$ と写像
$f^{-1}\mathcal{F} \to \mathcal{G}$ があるとする。このとき、上の引き戻し写像と、
$f^{-1}\mathcal{F} \to \mathcal{G}$ に $h_*$ を適用して得られる写像を合成して、
写像
$$
g_*\mathcal{F} \longrightarrow h_*\mathcal{G}
$$
を得る。$\mathcal{B}$ の対象上の切断に制限すると、上で述べた大域切断上の
「引き戻し写像」が（サイトを適切に選ぶことにより）得られる。

\medskip\noindent
さらに一般に、トポスの可換図式
$$
\xymatrix{
\Sh(\mathcal{D}) \ar[d]_h \ar[r]_f &
\Sh(\mathcal{C}) \ar[d]^g \\
\Sh(\mathcal{B}) \ar[r]^e & \Sh(\mathcal{A})
}
$$
と $\mathcal{C}$ 上の層 $\mathcal{G}$ があるとする。このとき
{\it 基底変換写像}
$$
e^{-1}g_*\mathcal{G} \longrightarrow h_*f^{-1}\mathcal{G}.
$$
がある。この写像は、写像
$g_*\mathcal{G} \to e_*h_*f^{-1}\mathcal{G} = (e \circ h)_*f^{-1}\mathcal{G}$
に随伴する写像であるが、後者はまさに上で述べた引き戻し写像である。

\begin{remark}
\label{sites-remark-compose-base-change}
トポスの可換図式
$$
\xymatrix{
\Sh(\mathcal{B}') \ar[r]_k \ar[d]_{f'} &
\Sh(\mathcal{B}) \ar[d]^f \\
\Sh(\mathcal{C}') \ar[r]^l \ar[d]_{g'} &
\Sh(\mathcal{C}) \ar[d]^g \\
\Sh(\mathcal{D}') \ar[r]^m &
\Sh(\mathcal{D})
}
$$
を考える。このとき、二つの正方形に対する基底変換写像を合成すると、
外側の長方形に対する基底変換写像が得られる。より正確には、合成
\begin{align*}
m^{-1} \circ (g \circ f)_*
& =
m^{-1} \circ g_* \circ f_* \\
& \to g'_* \circ l^{-1} \circ f_* \\
& \to g'_* \circ f'_* \circ k^{-1} \\
& = (g' \circ f')_* \circ k^{-1}
\end{align*}
がその長方形に対する基底変換写像である。確認は省略する。
\end{remark}

\begin{remark}
\label{sites-remark-compose-base-change-horizontal}
環付きトポスの可換図式
$$
\xymatrix{
\Sh(\mathcal{C}'') \ar[r]_{g'} \ar[d]_{f''} &
\Sh(\mathcal{C}') \ar[r]_g \ar[d]_{f'} &
\Sh(\mathcal{C}) \ar[d]^f \\
\Sh(\mathcal{D}'') \ar[r]^{h'} &
\Sh(\mathcal{D}') \ar[r]^h &
\Sh(\mathcal{D})
}
$$
を考える。このとき、二つの正方形に対する基底変換写像を合成すると、
外側の長方形に対する基底変換写像が得られる。より正確には、合成
\begin{align*}
(h \circ h')^{-1} \circ f_*
& =
(h')^{-1} \circ h^{-1} \circ f_* \\
& \to (h')^{-1} \circ f'_* \circ g^{-1} \\
& \to f''_* \circ (g')^{-1} \circ g^{-1} \\
& = f''_* \circ (g \circ g')^{-1}
\end{align*}
がその長方形に対する基底変換写像である。確認は省略する。
% P01-E0162
\end{remark}

\section{SGA4 との比較}
\label{sites-section-notation-comparison}

\noindent
Section \ref{sites-section-functoriality-PSh} の関手 $u^p$ と $u_p$、および
Section \ref{sites-section-continuous-functors} の関手 $u^s$ と $u_s$ に対する
本章の記法は、\cite[pages 14 and 42]{ArtinTopologies} から採った。
これらの選択をすると、Section \ref{sites-section-more-functoriality-PSh} の関手
${}_pu$ と Section \ref{sites-section-cocontinuous-functors} の関手
${}_su$ に対する記法も自然に思われる。この節では、本章の記法を
SGA4 の記法と比較する。

\medskip\noindent
{\bf 前層：} $u : \mathcal{C} \to \mathcal{D}$ を圏の間の関手とする。
関手 $u^p$ は \cite[Exposee I, Section 5]{SGA4} では $u^*$ と書かれる。
関手 $u_p$ は \cite[Exposee I, Proposition 5.1]{SGA4} では $u_!$ と書かれる。
関手 ${}_pu$ は \cite[Exposee I, Proposition 5.1]{SGA4} では
$u_*$ と書かれる。言い換えると、
$$
u_p, u^p, {}_pu\quad(SP)
\quad\text{versus}\quad
u_!, u^*, u_*\quad(SGA4)
$$
である。SGA4 の各箇所では、これらの関手に異なる記法が用いられていることに
注意しなければならない。

\medskip\noindent
{\bf 層と連続関手：}
$\mathcal{C}$ と $\mathcal{D}$ をサイトとし、
$u : \mathcal{C} \to \mathcal{D}$ を連続関手（Definition
\ref{sites-definition-continuous}）とする。関手 $u^s$ は
\cite[Exposee III, 1.11]{SGA4} では $u_s$ と書かれる。
関手 $u_s$ は \cite[Exposee III, Proposition 1.2]{SGA4} では
$u^s$ と書かれる。言い換えると、
$$
u_s, u^s\quad(SP)
\quad\text{versus}\quad
u^s, u_s\quad(SGA4)
$$
である。$u$ がサイトの射 $f : \mathcal{D} \to \mathcal{C}$
（Definition \ref{sites-definition-morphism-sites}）を定めるとき、
付随するトポスの射（Lemma \ref{sites-lemma-morphism-sites-topoi}）は
\cite[Exposee IV, (4.9.1.1)]{SGA4} のものと同じである。

\medskip\noindent
{\bf 層と余連続関手：}
$\mathcal{C}$ と $\mathcal{D}$ をサイトとし、
$u : \mathcal{C} \to \mathcal{D}$ を余連続関手（Definition
\ref{sites-definition-cocontinuous}）とする。関手 ${}_su$
（Lemma \ref{sites-lemma-pu-sheaf}）は
\cite[Exposee III, Proposition 2.3]{SGA4} では $u_*$ と書かれる。
関手 $(u^p\ )^\#$ は \cite[Exposee III, Proposition 2.3]{SGA4} では
$u^*$ と書かれる。言い換えると、
$$
(u^p\ )^\#, {}_su\quad(SP)
\quad\text{versus}\quad
u^*, u_*\quad(SGA4)
$$
である。したがって、Lemma \ref{sites-lemma-cocontinuous-morphism-topoi} で
$u$ に付随するトポスの射は \cite[Exposee IV, 4.7]{SGA4} のものと同じである。

\medskip\noindent
{\bf トポスの射：} $f$ が関手 $(f^{-1}, f_*)$ によって与えられる
トポスの射ならば、関手 $f^{-1}$ は
\cite[Exposee IV, Definition 3.1]{SGA4} では $f^*$ と書かれる。
集合の層、またはさらに一般に代数構造の層の逆像には $f^{-1}$ を用いる
（Section \ref{sites-section-sheaves-algebraic-structures}）。
環付きトポスの射に対する加群の層の引き戻しには $f^*$ を用いる
（Modules on Sites, Definition \ref{sites-modules-definition-pushforward}）。
\section{位相}
\label{sites-section-topologies}

\noindent
この節では、\cite{SGA4} で定義された圏上の位相を定義する。
層とトポスに関するすべての理論をこの言葉で展開できる。この内容の現代的な解説は
\cite{KS} にある。しかし、代数幾何学で最も重要なのは、サイトがその基礎となる圏上に
定める位相の場合である。したがって、初読者には
{\bf 本節および本章の他のすべての節を飛ばすことを強く勧める}！

\begin{definition}
\label{sites-definition-sieve}
$\mathcal{C}$ を圏とし、$U \in \Ob(\mathcal{C})$ とする。
{\it $U$ 上のふるい $S$}とは、部分前層 $S \subset h_U$ のことである。
\end{definition}

\noindent
言い換えると、$U$ 上のふるいは、各対象
$T \in \Ob(\mathcal{C})$ に対して、すべての射 $T \to U$ の集合の部分集合
$S(T)$ を選ぶ。実際、部分集合の族
$S(T) \subset h_U(T) = \Mor_\mathcal{C}(T, U)$ に課される唯一の条件は
次の規則である。
\begin{equation}
\label{sites-equation-property-sieve}
\left.
\begin{matrix}
(\alpha : T \to U) \in S(T) \\
g : T' \to T
\end{matrix}
\right\} \Rightarrow
(\alpha \circ g : T' \to U) \in S(T')
\end{equation}
念頭に置くとよい図式的な見方は、写像 $S \to h_U$ を
「$S$ から $U$ への射」と考えることである。

\begin{lemma}
\label{sites-lemma-sieves-set}
$\mathcal{C}$ を圏とし、$U \in \Ob(\mathcal{C})$ とする。
\begin{enumerate}
\item $U$ 上のふるい全体は集合である。
\item 包含関係はこの集合上の半順序を定める。
\item ふるいの合併と共通部分はふるいである。
\item
\label{sites-item-sieve-generated}
終域が $U$ である $\mathcal{C}$ の射の族
$\{U_i \to U\}_{i\in I}$ が与えられたとき、各 $U_i \to U$ が
$S(U_i)$ に属するような $U$ 上の最小のふるい $S$ が一意に存在する。
\item ふるい $S = h_U$ は最大ふるいである。
\item 空部分前層は最小ふるいである。
\end{enumerate}
\end{lemma}

\begin{proof}
部分前層の定義により、前層 $\mathcal{F}$ のすべての部分前層の集まりは
$\prod_{U \in \Ob(\mathcal{C})} \mathcal{P}(\mathcal{F}(U))$ の部分集合である。
これは集合である。（ここで $\mathcal{P}(A)$ は $A$ の冪集合を表す。）
したがって、$U$ 上のふるい全体は集合である。

\medskip\noindent
半順序を次のように定義する。すべての $T \to U$ に対して
$S(T) \subset S'(T)$ であるとき、かつそのときに限り $S \leq S'$ とする。
記法は $S \subset S'$ とする。

\medskip\noindent
$U$ 上のふるいの族 $S_i$、$i \in I$ が与えられたとき、
すべての $T \in \Ob(\mathcal{C})$ に対して
$(\bigcup S_i)(T) = \bigcup S_i(T)$ と置くことにより
ふるい $\bigcup S_i$ を定義できる。共通部分 $\bigcap S_i$ も同様に定義する。

\medskip\noindent
主張にある $\{U_i \to U\}_{i\in I}$ が与えられたとする。
前層の射 $h_{U_i} \to h_U$ を考え、これらの前層の写像の像
（Definition \ref{sites-definition-image}）の合併として $S$ を定義すればよい。

\medskip\noindent
補題の最後の二つの主張は明らかである。
\end{proof}

\begin{definition}
\label{sites-definition-sieve-generated}
$\mathcal{C}$ を圏とする。終域が $U$ である $\mathcal{C}$ の射の族
$\{f_i : U_i \to U\}_{i\in I}$ が与えられたとする。Lemma
\ref{sites-lemma-sieves-set} の項目 (\ref{sites-item-sieve-generated}) で記述された
$U$ 上のふるい $S$ を、{\it 射 $f_i$ によって生成される $U$ 上のふるい}という。
\end{definition}

\begin{definition}
\label{sites-definition-pullback-sieve}
$\mathcal{C}$ を圏とし、$f : V \to U$ を $\mathcal{C}$ の射とする。
$S \subset h_U$ をふるいとする。$f$ による $S$ の{\it 引き戻し}を、
規則
$$
(\alpha : T \to V) \in (S \times_U V)(T)
\Leftrightarrow
(f \circ \alpha : T \to U) \in S(T)
$$
によって定義される $V$ のふるい $S \times_U V$ と定義する。
\end{definition}

\noindent
これが実際にふるいであることの確認は読者に委ねる
（ヒント：Equation \ref{sites-equation-property-sieve} を用いよ）。
$S \times_U V$ を、$f : V \to U$ による $S$ の{\it 基底変換}と呼ぶこともある。

\begin{lemma}
\label{sites-lemma-pullback-sieve-section}
$\mathcal{C}$ を圏とし、$U \in \Ob(\mathcal{C})$ とする。
$S$ を $U$ 上のふるいとする。$f : V \to U$ が $S$ に属するならば、
$S \times_U V = h_V$ は最大である。
\end{lemma}

\begin{proof}
定義から自明である。
\end{proof}

\begin{definition}
\label{sites-definition-topology}
$\mathcal{C}$ を圏とする。$\mathcal{C}$ 上の{\it 位相}とは、各
$U \in \Ob(\mathcal{C})$ に $U$ 上のすべてのふるいの集合の部分集合
$J(U)$ を対応させ、次の条件を満たす規則のことである。
\begin{enumerate}
\item $\mathcal{C}$ の任意の射 $f : V \to U$ と任意の元
$S \in J(U)$ に対して、引き戻し $S \times_U V$ は $J(V)$ の元である。
\item $S$ と $S'$ を $U \in \Ob(\mathcal{C})$ 上のふるいとする。
$S \in J(U)$ であり、すべての $f \in S(V)$ に対して引き戻し
$S' \times_U V$ が $J(V)$ に属するならば、$S'$ は $J(U)$ に属する。
\item すべての $U \in \Ob(\mathcal{C})$ に対して、最大ふるい
$S = h_U$ は $J(U)$ に属する。
\end{enumerate}
\end{definition}

\noindent
この場合、$J(U)$ に属するふるいを{\it 被覆ふるい}という。

\begin{lemma}
\label{sites-lemma-topology-basic}
$\mathcal{C}$ を圏とし、$J$ を $\mathcal{C}$ 上の位相とする。
$U \in \Ob(\mathcal{C})$ とする。
\begin{enumerate}
\item $J(U)$ の元の有限共通部分は $J(U)$ に属する。
\item $S \in J(U)$ かつ $S' \supset S$ ならば、$S' \in J(U)$ である。
\end{enumerate}
\end{lemma}

\begin{proof}
(1) の証明。空族の共通部分に関する主張は Definition
\ref{sites-definition-topology} の条件 (3) から従う。
$S, S' \in J(U)$ とし、$S'' = S \cap S'$ を考える。
$S(V)$ に属する各 $V \to U$ に対して
$$
S' \times_U V = S'' \times_U V
$$
である。これは $V \to U$ がすでに $S$ に属するからである。
したがって、定義の第二の公理により $S'' \in J(U)$ である。

\medskip\noindent
(2) の証明。$S \in J(U)$ かつ $S' \supset S$ とする。
$S(V)$ に属する各 $V \to U$ に対して、Lemma
\ref{sites-lemma-pullback-sieve-section} により $S' \times_U V = h_V$ である。
したがって第三の公理により $S' \times_U V \in J(V)$ であり、
第二の公理により $S' \in J(U)$ である。
\end{proof}

\begin{definition}
\label{sites-definition-finer}
$\mathcal{C}$ を圏とし、$J$、$J'$ を $\mathcal{C}$ 上の二つの位相とする。
$\mathcal{C}$ のすべての対象 $U$ に対して $J'(U) \subset J(U)$ であるとき、
かつそのときに限り、$J$ は $J'$ より{\it 細かい}または{\it 強い}という。
この場合、$J'$ は $J$ より{\it 粗い}または{\it 弱い}ともいう。
\end{definition}

\noindent
言い換えると、$J'$ の任意の被覆ふるいは $J$ の被覆ふるいである。
$\mathcal{C}$ 上には最も細かい位相が存在する。すなわち、任意のふるいが
被覆ふるいであるような位相であり、これを $\mathcal{C}$ の
{\it 離散位相}という。また、最も粗い位相も存在する。すなわち、すべての対象
$U$ に対して $J(U) = \{h_U\}$ であるような位相である。これを
{\it カオス位相}または{\it 密着位相}という。

\begin{lemma}
\label{sites-lemma-play-with-topologies}
$\mathcal{C}$ を圏とし、$\{J_i\}_{i\in I}$ を位相の集合とする。
\begin{enumerate}
\item 規則 $J(U) = \bigcap J_i(U)$ は $\mathcal{C}$ 上の位相を定める。
\item すべての位相 $J_i$ より細かい最も粗い位相が存在する。
\end{enumerate}
\end{lemma}

\begin{proof}
第一の主張は定義から直ちに従う。第二の主張は、すべての $J_i$ より細かい
位相すべての共通部分を取ることによって従う。
\end{proof}

\noindent
ここで、動機を与えることなく層を定義できる。

\begin{definition}
\label{sites-definition-sheaf-sets-topology}
$\mathcal{C}$ を位相 $J$ を備えた圏とし、$\mathcal{F}$ を
$\mathcal{C}$ 上の集合の前層とする。すべての
$U \in \Ob(\mathcal{C})$ および $U$ のすべての被覆ふるい $S$ に対して、
標準写像
$$
\Mor_{\textit{PSh}(\mathcal{C})}(h_U, \mathcal{F})
\longrightarrow
\Mor_{\textit{PSh}(\mathcal{C})}(S, \mathcal{F})
$$
が全単射であるとき、$\mathcal{F}$ を $\mathcal{C}$ 上の{\it 層}という。
\end{definition}

\noindent
表示式の左辺が $\mathcal{F}(U)$ に等しいことを思い出そう。
言い換えると、$\mathcal{F}$ が層であることは、$U$ 上の $\mathcal{F}$ の
切断を与えることが、$(\alpha : T \to U) \in S(T)$ によって添字づけられた
両立する切断 $s_{T, \alpha} \in \mathcal{F}(T)$ の族を与えることと同じであり、
これが $U$ 上のすべての被覆ふるい $S$ について成り立つことと同値である。

\begin{lemma}
\label{sites-lemma-topology-presheaves-sheaves}
$\mathcal{C}$ を圏とし、$\{ \mathcal{F}_i \}_{i\in I}$ を
$\mathcal{C}$ 上の集合の前層の族とする。各
$U \in \Ob(\mathcal{C})$ に対して、次の性質をもつ $U$ 上のふるい $S$ の集合を
$J(U)$ と書く。すべての射 $V \to U$ に対して、写像
$$
\Mor_{\textit{PSh}(\mathcal{C})}(h_V, \mathcal{F}_i)
\longrightarrow
\Mor_{\textit{PSh}(\mathcal{C})}(S \times_U V, \mathcal{F}_i)
$$
はすべての $i \in I$ について全単射である。このとき $J$ は
$\mathcal{C}$ 上の位相を定める。この位相は、すべての $\mathcal{F}_i$ が
層となるような最も細かい位相である。
% P01-E0163
\end{lemma}

\begin{proof}
まず、$J$ が位相ならば、補題の最後の主張が直ちに従うことに注意する。

\medskip\noindent
各 $i \in I$ と $U \in \Ob(\mathcal{C})$ に対して、次の性質をもつ
$U$ 上のふるい $S$ の集合を $J_i(U)$ とする。すべての射
$V \to U$ に対して、写像
$$
\Mor_{\textit{PSh}(\mathcal{C})}(h_V, \mathcal{F}_i)
\longrightarrow
\Mor_{\textit{PSh}(\mathcal{C})}(S \times_U V, \mathcal{F}_i)
$$
は全単射である。$J(U) = \bigcap J_i(U)$ であることに注意する。
したがって $J_i$ が位相を定めることを証明すれば、Lemma
\ref{sites-lemma-play-with-topologies} により結論を得る。これにより、次の段落で
論じる場合に帰着される。

\medskip\noindent
族が一つの前層 $\mathcal{F}$ だけからなるとき、$J$ が位相であることを証明する。
位相の第一および第三の公理は直ちに確認できる。そこで、対象 $U$ と
$U$ 上のふるい $S, S'$ があり、$S \in J(U)$ で、$S(V)$ に属するすべての
$V \to U$ に対して $S' \times_U V \in J(V)$ であると仮定する。
$S' \in J(U)$ を示さなければならない。言い換えると、任意の
$f : W \to U$ に対して写像
$$
\mathcal{F}(W) =
\Mor_{\textit{PSh}(\mathcal{C})}(h_W, \mathcal{F})
\longrightarrow
\Mor_{\textit{PSh}(\mathcal{C})}(S' \times_U W, \mathcal{F})
$$
が全単射であることを示せばよい。元
$\varphi \in
\Mor_{\textit{PSh}(\mathcal{C})}(S' \times_U W, \mathcal{F})$ を選ぶ。
$\varphi$ に写る切断 $s \in \mathcal{F}(W)$ を構成する。

\medskip\noindent
$\alpha : V \to W$ が $S \times_U W$ の元であるとする。
引き戻しの定義により、合成 $f \circ \alpha : V \to W  \to U$ は
$S$ に属する。したがって、ふるいの対 $S, S'$ に関する仮定により
$S' \times_U V$ は $J(V)$ に属する。ここで、前層の可換図式
$$
\xymatrix{
S' \times_U V \ar[r] \ar[d] & h_V \ar[d] \\
S' \times_U W \ar[r] & h_W
}
$$
がある。$\varphi$ の $S' \times_U V$ への制限は、元
$s_{V, \alpha} \in \mathcal{F}(V)$ に対応する。これは $J$ の定義と、
$S' \times_U V$ が $J(V)$ に属することから分かる。規則
$(V, \alpha) \mapsto s_{V, \alpha}$ が元
$\psi \in
\Mor_{\textit{PSh}(\mathcal{C})}(S \times_U W, \mathcal{F})$ を定めることの
確認は読者に委ねる。$S \in J(U)$ なので、$J$ の定義から直ちに
$\psi$ は $\mathcal{F}(W)$ の元 $s$ に対応する。

\medskip\noindent
構成 $\varphi \mapsto s$ が上に表示した自然な写像の逆であることの確認は
読者に委ねる。
\end{proof}

\begin{definition}
\label{sites-definition-canonical-topology}
$\mathcal{C}$ を圏とする。すべての表現可能前層が層となるような
$\mathcal{C}$ 上の最も細かい位相（Lemma
\ref{sites-lemma-topology-presheaves-sheaves} を参照）を、
$\mathcal{C}$ の{\it 標準位相}という。
\end{definition}
\section{サイトが定める位相}
\label{sites-section-topology-site}

\noindent
$\mathcal{C}$ を圏とし、$\text{Cov}_1(\mathcal{C})$ と
$\text{Cov}_2(\mathcal{C})$ を、$\mathcal{C}$ 上にサイトの構造を定める
被覆の集合とする。この状況では、$\text{Cov}_1(\mathcal{C})$ と
$\text{Cov}_2(\mathcal{C})$ に対する集合の層の圏が同じになることがある。
例えば Lemma \ref{sites-lemma-refine-same-topology} を参照せよ。

\medskip\noindent
$\mathcal{C}$ 上のある位相 $J$ に関する層の圏は、その位相を決定し、
また位相 $J$ によって決定されることが分かる。これは自明でない主張であり、
後で Theorem \ref{sites-theorem-topology-and-topos} において扱う。

\medskip\noindent
ひとまずこの事実を認めれば、サイトが定める位相を研究することには意味がある。

\begin{lemma}
\label{sites-lemma-site-gives-topology}
$\mathcal{C}$ を被覆 $\text{Cov}(\mathcal{C})$ をもつサイトとする。
$\mathcal{C}$ の各対象 $U$ に対して、次の性質をもつ $U$ 上のふるい $S$ の集合を
$J(U)$ とする。被覆
$\{f_i : U_i \to U\}_{i\in I} \in \text{Cov}(\mathcal{C})$ が存在し、
$f_i$ によって生成されるふるい $S'$（Definition
\ref{sites-definition-sieve-generated} を参照）が $S$ に含まれる。
\begin{enumerate}
\item この $J$ は $\mathcal{C}$ 上の位相である。
\item 前層 $\mathcal{F}$ がこの位相に関する層
（Definition \ref{sites-definition-sheaf-sets-topology} を参照）であることと、
サイト上の層（Definition \ref{sites-definition-sheaf-sets} を参照）であることは同値である。
\end{enumerate}
\end{lemma}

\begin{proof}
第一の主張を示すには、サイトの定義（Definition \ref{sites-definition-site}）の公理
(1)、(2)、(3) が、位相の定義（Definition \ref{sites-definition-topology}）の公理
(3)、(2)、(1) をそれぞれ直ちに含意することに注意すればよい。例として、
$J$ が性質 (2) をもつことを証明する。$U$ を $\mathcal{C}$ の対象とし、
$S, S'$ を $U$ 上のふるいで、$S \in J(U)$ かつ $S(V)$ に属するすべての
$V \to U$ に対して $S' \times_U V \in J(V)$ であるものとする。
$J(U)$ の定義により、サイトの被覆 $\{f_i : U_i \to U\}$ を選べて、
$S$ に関して $h_{U_i} \to h_U$ の像が $S$ に含まれる。
% P01-E0164
各 $S'\times_U U_i$ は $J(U_i)$ に属するので、サイトの被覆
$\{U_{ij} \to U_i\}$ が存在して、$h_{U_{ij}} \to h_{U_i}$ は
$S' \times_U U_i$ に含まれる。基底変換の定義により、これは
$h_{U_{ij}} \to h_U$ が部分前層 $S' \subset h_U$ に含まれることを意味する。
サイトに対する公理 (2) により、$\{U_{ij} \to U\}$ は $U$ の被覆である。
したがって、$J$ の定義から $S' \in J(U)$ である。

\medskip\noindent
$\mathcal{F}$ を前層とし、$\mathcal{F}$ が位相 $J$ に関する層であると仮定する。
$\mathcal{F}$ がサイト上の層でもあることを示す。サイトの被覆
$\{f_i : U_i \to U\}_{i\in I}$ をとる。$s_i \in \mathcal{F}(U_i)$ を、
すべての $i, j$ に対して
$s_i|_{U_i \times_U U_j} = s_j|_{U_i \times_U U_j}$ を満たす切断の族とする。
各 $U_i$ 上で $s_i$ に制限される一意な切断 $s \in \mathcal{F}(U)$ が存在することを
示さなければならない。$S \subset h_U$ を $f_i$ によって生成されるふるいとする。
定義により $S \in J(U)$ であることに注意する。$s$ を構成する代わりに、
位相に関する層条件により、元
% P01-E0165
$$
\varphi \in \Mor_{\textit{PSh}(\mathcal{C})}(S, \mathcal{F}).
$$
を構成すれば十分である。ある対象 $T \in \mathcal{U}$ と
$\alpha \in S(T)$ をとる。これは、$\alpha : T \to U$ が、ある
$i\in I$（$1$ より多い場合もある）に対して $f_i$ を経由する射であることを
正確に意味する。そのような添字 $i$ と分解
$\alpha = f_i \circ \alpha_i$ を一つ選び、
$\varphi(\alpha) = \alpha_i^* s_i$ と定義する。
% P01-E0166
$i'$、$\alpha = f_i \circ \alpha_{i'}'$ が第二の選択ならば、
$\alpha_i^* s_i = (\alpha_{i'}')^* s_{i'}$ である。これはまさに、すべての
$i, j$ に対する条件
$s_i|_{U_i \times_U U_j} = s_j|_{U_i \times_U U_j}$ による。
% P01-E0167
したがって $\varphi(\alpha)$ はよく定義される。$s$ を定める $\varphi$ が正しく、
$s$ が各 $s_i$ に制限されることの確認は読者に委ねる。

\medskip\noindent
$\mathcal{F}$ を前層とし、$\mathcal{F}$ がサイト
$(\mathcal{C}, \text{Cov}(\mathcal{C}))$ 上の層であると仮定する。
$\mathcal{F}$ が位相 $J$ に関する層でもあることを示す。
$U$ を $\mathcal{C}$ の対象とし、$S$ を位相 $J$ に関する
$U$ 上の被覆ふるいとする。元
$$
\varphi \in \Mor_{\textit{PSh}(\mathcal{C})}(S, \mathcal{F}).
$$
をとる。$\varphi$ に制限される一意な元
$\mathcal{F}(U) = \Mor_{\textit{PSh}(\mathcal{C})}(h_U, \mathcal{F})$ が
存在することを示さなければならない。定義により、被覆
$\{f_i : U_i \to U\}_{i\in I} \in \text{Cov}(\mathcal{C})$ が存在し、
$f_i : U_i \to U$ は $S(U_i)$ に属する。したがって
$s_i = \varphi(f_i) \in \mathcal{F}(U_i)$ と置ける。
すべての $i, j$ に対して
$s_i|_{U_i \times_U U_j} = s_j|_{U_i \times_U U_j}$ であることは
よい演習である。したがって、サイト
$(\mathcal{C}, \text{Cov}(\mathcal{C}))$ 上の $\mathcal{F}$ に対する
層条件により、所望の切断 $s$ を得る。詳細は読者に委ねる。
\end{proof}

\begin{definition}
\label{sites-definition-topology-associated-site}
$\mathcal{C}$ を被覆 $\text{Cov}(\mathcal{C})$ をもつサイトとする。
$\mathcal{C}$ に{\it 付随する位相}とは、上の Lemma
\ref{sites-lemma-site-gives-topology} で構成した位相 $J$ のことである。
\end{definition}

\noindent
$\mathcal{C}$ を圏とする。$\text{Cov}_1(\mathcal{C})$ と
$\text{Cov}_2(\mathcal{C})$ を、$\mathcal{C}$ 上にサイトの構造を定める
二つの被覆とする。これらが定める位相が同じになることは十分あり得る。
この場合、$\text{Cov}_1(\mathcal{C})$ と
$\text{Cov}_2(\mathcal{C})$ は $\mathcal{C}$ 上で
{\it 同じ位相を定める}という。このとき Lemma
\ref{sites-lemma-site-gives-topology} により層の圏も同じである。
% P01-E0168

\medskip\noindent
通常、重要なのは被覆の集まりが定める位相だけであり、異なる被覆の集合を選べることは
位相を研究するための道具とみなす。

\begin{remark}
\label{sites-remark-enlarge-coverings}
被覆の類を拡大する。明らかに、$\text{Cov}(\mathcal{C})$ が
$\mathcal{C}$ 上にサイトの構造を定めるならば、
$\text{Cov}(\mathcal{C})$ の元と自明に同値な（Definition
\ref{sites-definition-combinatorial-tautological} を参照）終域を固定した射の族の
任意の集合を $\mathcal{C}$ に加えても、位相は変わらない。
% P01-E0169
\end{remark}

\begin{remark}
\label{sites-remark-shrink-coverings}
被覆の類を縮小する。$\mathcal{C}$ をサイトとし、集合
$$
\mathcal{S} = P(\text{Arrows}(\mathcal{C})) \times \Ob(\mathcal{C})
$$
を考える。ここで $P(\text{Arrows}(\mathcal{C}))$ は射の集合の冪集合、
すなわち射の集合すべてからなる集合である。
$\mathcal{S}_\tau \subset \mathcal{S}$ を、次の性質をもつ
$(T, U) \in \mathcal{S}$ からなる部分集合とする。
(a) すべての $\varphi \in T$ の終域は $U$ である。
(b) 族 $\{\varphi\}_{\varphi \in T}$ は、$\text{Cov}(\mathcal{C})$ の
ある被覆と自明に同値である（Definition
\ref{sites-definition-combinatorial-tautological} を参照）。
$\mathcal{S}_\tau$ の元を被覆とみなしても、Definition
\ref{sites-definition-site} で定義したサイトの概念そのものは得られない。
得られる構造 $(\mathcal{C}, \mathcal{S}_\tau)$ は、少し修正された条件を満たす。
修正された条件は次のとおりである。
\begin{enumerate}
\item[(0')] $\text{Cov}(\mathcal{C}) \subset
P(\text{Arrows}(\mathcal{C})) \times \Ob(\mathcal{C})$。
\item[(1')] $V \to U$ が同型ならば
$(\{V \to U\}, U) \in \text{Cov}(\mathcal{C})$ である。
\item[(2')] $(T, U) \in \text{Cov}(\mathcal{C})$ とし、
$T$ に属する各 $f : U' \to U$ に対して
$(T_f, U') \in \text{Cov}(\mathcal{C})$ が与えられたとする。このとき
$T' = \{f \circ f' \mid f \in T,\ f' \in T_f\}$ と置けば、
$(T', U) \in \text{Cov}(\mathcal{C})$ である。
\item[(3')] $(T, U) \in \text{Cov}(\mathcal{C})$ とし、
$g : V \to U$ を $\mathcal{C}$ の射とする。このとき
\begin{enumerate}
\item $T$ に属する $f : U' \to U$ に対して
$U' \times_{f, U, g} V$ が存在する。
\item ファイバー積のある選択に対して
$T' = \{\text{pr}_2 : U' \times_{f, U, g} V \to V \mid f : U' \to U \in T\}$
と置けば、$(T', V) \in \text{Cov}(\mathcal{C})$ である。
\end{enumerate}
\end{enumerate}
上の (0') -- (3') を満たす構造が与えられたとき、
$\text{Cov}(\mathcal{C})$ を適切に拡大すれば（Sets, Section
\ref{sets-section-coverings-site} と比較せよ）サイトが得られることは容易に確認できる。
少なくとも位相の観点からは、この概念と実際のサイトの概念との差は明らかに小さい。
利点は二つある。上の条件 (0') により被覆は自動的に集合をなし、また (0') により
この型の構造すべての集まりも集合をなす。代償として、至るところで
「自明に同値」と書き続けなければならない。
\end{remark}
\section{位相における層化}
\label{sites-section-sheafification-topology}

\noindent
この節では、位相における層化の構成の類似物を説明する。

\medskip\noindent
$\mathcal{C}$ を圏とし、$J$ を $\mathcal{C}$ 上の位相とする。
$\mathcal{F}$ を集合の前層とする。各 $U \in \Ob(\mathcal{C})$ に対して、
余極限
$$
L\mathcal{F}(U)
=
\colim_{S \in J(U)^{opp}}
\Mor_{\textit{PSh}(\mathcal{C})}(S, \mathcal{F})
$$
を定義する。ここで $J(U)$ を、包含関係で順序づけられた半順序集合とみなす。
Lemma \ref{sites-lemma-sieves-set} を参照せよ。この系の遷移写像を次のように定義する。
$S \subset S'$ が $J(U)$ に属するならば、$S \to S'$ は前層の射である。
したがって、自然な制限写像
$$
\Mor_{\textit{PSh}(\mathcal{C})}(S', \mathcal{F})
\longrightarrow
\Mor_{\textit{PSh}(\mathcal{C})}(S, \mathcal{F}).
$$
がある。これにより、$J(U)$ 上の順序を逆にすれば（上付きの ${}^{opp}$ は
これを表す）、$S \mapsto \Mor_{\textit{PSh}(\mathcal{C})}(S, \mathcal{F})$ は
Categories, Definition \ref{categories-definition-system-over-poset} の意味で
有向系となる。特に $h_U \in J(U)$ なので、同一視
$\mathcal{F}(U) = \Mor_{\textit{PSh}(\mathcal{C})}(h_U, \mathcal{F})$ から
生じる標準写像
$$
\ell : \mathcal{F}(U) \longrightarrow L\mathcal{F}(U)
$$
がある。さらに、$U$ 上の任意の被覆ふるいの対 $S, S'$ に対して
$S \cap S'$ も被覆ふるいなので（Lemma \ref{sites-lemma-sieves-set} を参照）、
$L\mathcal{F}(U)$ を定義する余極限は有向である。
% P01-E0170

\medskip\noindent
$f : V \to U$ を $\mathcal{C}$ の射とし、$S \in J(U)$ とする。
可換図式
$$
\xymatrix{
S \times_U V \ar[r] \ar[d] & h_V \ar[d] \\
S \ar[r] & h_U
}
$$
がある。左の縦写像を用いて、標準的な制限写像
$$
\Mor_{\textit{PSh}(\mathcal{C})}(S, \mathcal{F})
\to \Mor_{\textit{PSh}(\mathcal{C})}(S \times_U V, \mathcal{F}).
$$
を得る。基底変換 $S \mapsto S \times_U V$ は順序を保つ写像
$J(U) \to J(V)$ を誘導する。また、制限写像は Categories, Lemma
{categories-lemma-functorial-colimit} の意味で関手の変換を定める。
したがって、自然な制限写像
% P01-E0171
$$
L\mathcal{F}(U) \longrightarrow L\mathcal{F}(V).
$$
を得る。

\begin{lemma}
\label{sites-lemma-L-presheaf}
上の状況では、次が成り立つ。
% P01-E0172
\begin{enumerate}
\item 対応 $U \mapsto L\mathcal{F}(U)$ と上で定義した制限写像は前層をなす。
\item 写像 $\ell$ は貼り合わさって前層の射
$\ell : \mathcal{F} \to L\mathcal{F}$ を与える。
\item 規則 $\mathcal{F} \mapsto (\mathcal{F} \xrightarrow{\ell}
L\mathcal{F})$ は関手である。
\item $\mathcal{F}$ が $\mathcal{G}$ の部分前層ならば、
$L\mathcal{F}$ は $L\mathcal{G}$ の部分前層である。
\item 写像 $\ell : \mathcal{F} \to L\mathcal{F}$ は次の性質をもつ。
各切断 $s \in L\mathcal{F}(U)$ に対して、$U$ 上の被覆ふるい $S$ と元
$\varphi \in \Mor_{\textit{PSh}(\mathcal{C})}(S, \mathcal{F})$ が存在し、
$\ell(\varphi)$ は $s$ の $S$ への制限に等しい。
\end{enumerate}
\end{lemma}

\begin{proof}
省略する。
\end{proof}

\begin{definition}
\label{sites-definition-presheaf-separated-topology}
$\mathcal{C}$ を圏とし、$J$ を $\mathcal{C}$ 上の位相とする。
集合の前層 $\mathcal{F}$ が{\it 分離的}であるとは、各対象 $U$ と
$U$ 上の各被覆ふるい $S$ に対して標準写像
$\mathcal{F}(U) \to \Mor_{\textit{PSh}(\mathcal{C})}(S, \mathcal{F})$ が
単射であることをいう。
\end{definition}

\begin{theorem}
\label{sites-theorem-L-topology}
$\mathcal{C}$ を圏とし、$J$ を $\mathcal{C}$ 上の位相とする。
$\mathcal{F}$ を集合の前層とする。
\begin{enumerate}
\item 前層 $L\mathcal{F}$ は分離的である。
\item $\mathcal{F}$ が分離的ならば、$L\mathcal{F}$ は層であり、
前層の写像 $\mathcal{F} \to L\mathcal{F}$ は単射である。
\item $\mathcal{F}$ が層ならば、$\mathcal{F} \to L\mathcal{F}$ は同型である。
\item 前層 $LL\mathcal{F}$ は常に層である。
\end{enumerate}
\end{theorem}

\begin{proof}
(3) は $L$ の定義と層の定義（Definition
\ref{sites-definition-sheaf-sets-topology}）から自明である。(4) は他の主張から形式的に従う。

\medskip\noindent
(1) の証明を概説する。$S$ を対象 $U$ の被覆ふるいとする。
$\varphi_i \in L\mathcal{F}(U)$、$i = 1, 2$ が
$\Mor_{\textit{PSh}(\mathcal{C})}(S, L\mathcal{F})$ の同じ元に写るとする。
両方の $\varphi_i$ が元
$\varphi_i \in \Mor_{\textit{PSh}(\mathcal{C})}(S', \mathcal{F})$ によって
表されるような $U$ 上の一つの被覆ふるい $S'$ を選べる。
$S$ と $S'$ の両方を、やはり被覆ふるいである $S' \cap S$ に置き換えることで
$S' = S$ と仮定できる。Lemma \ref{sites-lemma-sieves-set} を参照せよ。
$V\in \Ob(\mathcal{C})$ とし、$\alpha : V \to U$ が $S(V)$ に属するとする。
Lemma \ref{sites-lemma-pullback-sieve-section} により $S \times_U V = h_V$ である。
したがって、$V \to U$ による $\varphi_i$ の制限は、$V$ 上の
$\mathcal{F}$ の切断 $s_{i, V, \alpha}$ に対応する。仮定により、
$V$ の被覆ふるい $S_{V, \alpha}$ が存在して、$s_{i, V, \alpha}$ は
$\Mor_{\textit{PSh}(\mathcal{C})}(S_{V, \alpha}, \mathcal{F})$ の同じ元に制限される。
規則
\begin{eqnarray}
\label{sites-equation-S-prime-prime}
(f : T \to U) \in S''(T)
& \Leftrightarrow &
\exists\ V , \ \alpha : V \to U, \ \alpha \in S(V), \nonumber \\
& &
\exists\ g : T \to V, \ g \in S_{V, \alpha}(T), \\
& &
f = \alpha \circ g \nonumber
\end{eqnarray}
によって $U$ 上のふるい $S''$ を定義する。位相の公理 (2) により、
$S''$ は $U$ 上の被覆ふるいである。構成から、$\varphi_1$ と $\varphi_2$ は
所望のとおり $\Mor_{\textit{PSh}(\mathcal{C})}(S'', L\mathcal{F})$ の
同じ元に制限される。

\medskip\noindent
(2) の証明を概説する。$\mathcal{F}$ を位相 $J$ に関して
$\mathcal{C}$ 上の分離的な集合の前層とする。$S$ を $\mathcal{C}$ の対象
$U$ の被覆ふるいとし、$\varphi \in \Mor_\mathcal{C}(S, L\mathcal{F})$ とする。
% P01-E0173
$\varphi$ に制限される元 $s \in L\mathcal{F}(U)$ を見つけなければならない。
$V\in \Ob(\mathcal{C})$ とし、$\alpha : V \to U$ が $S(V)$ に属するとする。
値 $\varphi(\alpha) \in L\mathcal{F}(V)$ は、$V$ の被覆ふるい
$S_{V, \alpha}$ と前層の射
$\varphi_{V, \alpha} : S_{V, \alpha} \to \mathcal{F}$ によって与えられる。
上の証明と同様、Equation (\ref{sites-equation-S-prime-prime}) により
$U$ 上の被覆ふるい $S''$ を定義する。写像
$$
\varphi'' : S'' \longrightarrow \mathcal{F}
$$
を次の簡単な規則で定義する。各 $f : T \to U$、$f \in S''(T)$ に対して、
Equation (\ref{sites-equation-S-prime-prime}) にある $V, \alpha, g$ を選び、
$$
\varphi''(f) = \varphi_{V, \alpha}(g).
$$
と置く。これは $V, \alpha, g$ の選択に依存しないと主張する。
第二のそのような選択 $V', \alpha', g'$ を考える。
% P01-E0174
$T$ 上の次の被覆ふるいの共通部分
$$
(S_{V, \alpha} \times_{V, g} T) \cap (S_{V', \alpha'} \times_{V', g'} T)
$$
への $\varphi_{V, \alpha}$ と $\varphi_{V', \alpha'}$ の制限は一致する。
実際、これらの制限はともに $\varphi$ の（$f$ を通じた）$T$ への制限に対応し、
$\mathcal{F}$ が分離的であることから所望の等式が従う。共通の制限を $\psi$ と書く。
% P01-E0175
選択に依存しないことは
$\varphi_{V, \alpha}(g) = \psi(\text{id}_T) =
\varphi_{V', \alpha'}(g')$ から従う。これで $\varphi''$ は元
$s \in L\mathcal{F}(U)$ を与える。$s$ が $\varphi$ に制限されることの確認は
読者に委ねる。
\end{proof}

\begin{definition}
\label{sites-definition-associated-sheaf-topology}
$\mathcal{C}$ を位相 $J$ を備えた圏とし、$\mathcal{F}$ を
$\mathcal{C}$ 上の集合の前層とする。標準写像
$\mathcal{F} \to \mathcal{F}^\#$ とともに、層
$\mathcal{F}^\# := LL\mathcal{F}$ を、{\it $\mathcal{F}$ に付随する層}という。
\end{definition}

\begin{proposition}
\label{sites-proposition-sheafification-adjoint-topology}
$\mathcal{C}$ を位相を備えた圏とし、$\mathcal{F}$ を
$\mathcal{C}$ 上の集合の前層とする。標準写像
$\mathcal{F} \to \mathcal{F}^\#$ は次の普遍性をもつ。
$\mathcal{G}$ が集合の層であるような任意の写像
$\mathcal{F} \to \mathcal{G}$ に対して、一意な写像
$\mathcal{F}^\# \to \mathcal{G}$ が存在し、
$\mathcal{F} \to \mathcal{F}^\# \to \mathcal{G}$ は与えられた写像に等しい。
\end{proposition}

\begin{proof}
Proposition \ref{sites-proposition-sheafification-adjoint} の証明と同じである。
\end{proof}
\section{位相と層}
\label{sites-section-topology-and-sheaves}

\begin{lemma}
\label{sites-lemma-sieve-sheafification}
$\mathcal{C}$ を位相 $J$ を備えた圏とし、$U$ を $\mathcal{C}$ の対象、
$S$ を $U$ 上のふるいとする。次は同値である。
\begin{enumerate}
\item ふるい $S$ は被覆ふるいである。
\item 写像 $S \to h_U$ の層化 $S^\# \to h_U^\#$ は同型である。
\end{enumerate}
\end{lemma}

\begin{proof}
まず一般的な注意を二つ述べる。$S^\# = LLS$ および
$h_U^\# = LLh_U$ を用いる。特に Lemma \ref{sites-lemma-L-presheaf} により、
$S^\# \to h_U^\#$ は単射である。また、
$\text{id}_U \in h_U(U)$ であることに注意する。したがって、これは
$U$ 上の $Lh_U$ および $h_U^\# = LLh_U$ の切断を定めるが、これらも
$\text{id}_U$ と書く。

\medskip\noindent
$S$ が被覆ふるいであるとする。合成 $h_U \to h_U^\#$ が標準写像となるような
射 $h_U \to S^\#$ を見つければ十分であることは明らかである。そのような写像を
見つけるには、$\text{id}_U$ に制限される切断 $s \in S^\#(U)$ を
見つければよい。しかし $S$ は被覆ふるいなので、元
$\text{id}_S \in \Mor_{\textit{PSh}(\mathcal{C})}(S, S)$ は、
$Lh_U$ において $\text{id}_U$ に制限される $U$ 上の $LS$ の切断を定める。
これでよい。

\medskip\noindent
$S^\# \to h_U^\#$ が同型であるとする。$1 \in S^\#(U)$ を、
$h_U^\#(U)$ の $\text{id}_U$ に対応する元とする。$S^\# = LLS$ なので、
$U$ 上の被覆ふるい $S'$ が存在し、$1$ は
$$
\varphi \in \Mor_{\textit{PSh}(\mathcal{C})}(S', LS).
$$
から生じる。これはさらに、各 $\alpha : V \to U$、
$\alpha\in S'(V)$ に対して $V$ 上の被覆ふるい $S_{V, \alpha}$ が存在し、
$\varphi(\alpha)$ が前層の射 $S_{V, \alpha} \to S$ に対応することを意味する。
言い換えると、$S_{V, \alpha}$ は $S \times_U V$ に含まれる。
位相の第二の公理により、$S$ は被覆ふるいである。
\end{proof}

\begin{theorem}
\label{sites-theorem-topology-and-topos}
$\mathcal{C}$ を圏とし、$J$、$J'$ を $\mathcal{C}$ 上の位相とする。
次は同値である。
\begin{enumerate}
\item $J = J'$。
\item 位相 $J$ に関する層と位相 $J'$ に関する層は同じである。
\end{enumerate}
\end{theorem}

\begin{proof}
$J = J'$ ならば層の概念が同じであることは恒真である。逆に、Lemma
\ref{sites-lemma-sieve-sheafification} は被覆ふるいを層化関手によって特徴づける。
一方、層化関手
$\textit{PSh}(\mathcal{C}) \to \Sh(\mathcal{C}, J)$ は包含関手
$\Sh(\mathcal{C}, J) \to \textit{PSh}(\mathcal{C})$ の左随伴である。
したがって、部分圏 $\Sh(\mathcal{C}, J)$ と
$\Sh(\mathcal{C}, J')$ が同じならば、層化関手も同じであり、ゆえに
被覆ふるいの集まりも同じである。
\end{proof}

\begin{lemma}
\label{sites-lemma-finer-topology}
仮定と記法を Theorem \ref{sites-theorem-topology-and-topos} と同じとする。
このとき $J \subset J'$ であることと、位相 $J'$ に関するすべての層が
位相 $J$ に関する層であることは同値である。
% P01-E0176
\end{lemma}

\begin{proof}
一方向は明らかである。逆方向について
$\Sh(\mathcal{C}, J') \subset \Sh(\mathcal{C}, J)$ と仮定する。
形式的な議論により、$\mathcal{F}$ が集合の前層で、
$\mathcal{F} \to \mathcal{F}^\#$、それぞれ
$\mathcal{F} \to \mathcal{F}^{\#, \prime}$ が $J$、それぞれ $J'$ に関する
層化ならば、標準写像 $\mathcal{F}^\# \to \mathcal{F}^{\#, \prime}$ が存在して、
$\mathcal{F} \to \mathcal{F}^\# \to \mathcal{F}^{\#, \prime}$ は標準写像
$\mathcal{F} \to \mathcal{F}^{\#, \prime}$ に等しい。
もちろん、$\mathcal{F}^\# \to \mathcal{F}^{\#, \prime}$ は、第二の層を
第一の層の位相 $J'$ に関する層化と同一視する。
これを Lemma \ref{sites-lemma-sieve-sheafification} の写像 $S \to h_U$ に適用する。
可換図式
$$
\xymatrix{
S \ar[r] \ar[d] &
S^\# \ar[r] \ar[d] &
S^{\#, \prime} \ar[d] \\
h_U \ar[r] &
h_U^\# \ar[r] &
h_U^{\#, \prime}
}
$$
を得る。$S$ が位相 $J$ に関する被覆ふるいならば、中央の縦写像は
補題により同型である。右の縦写像は中央の写像の $J'$ に関する層化なので、
これも同型である。再び補題により、$S$ は $J'$ に関する被覆ふるいでもある。
\end{proof}

\section{位相と連続関手}
\label{sites-section-topologies-continuous-functors}

\noindent
連続関手が層上の随伴関手の対をどのように与えるか説明する。

\section{点と位相}
\label{sites-section-points-topologies}

\noindent
Section \ref{sites-section-points} から、関手
$p = u : \mathcal{C} \to \textit{Sets}$ が与えられたとき、茎関手
$$
\textit{PSh}(\mathcal{C}) \longrightarrow \textit{Sets},
\mathcal{F} \longmapsto \mathcal{F}_p.
$$
を定義できることを思い出そう。

\begin{definition}
\label{sites-definition-point-topology}
$\mathcal{C}$ を圏とし、$J$ を $\mathcal{C}$ 上の位相とする。
位相の{\it 点 $p$}は、次を満たす関手
$u : \mathcal{C} \to \textit{Sets}$ によって与えられる。
\begin{enumerate}
\item $U$ 上の各被覆ふるい $S$ に対して、写像
$S_p \to (h_U)_p$ は全射である。
\item 茎関手 $\Sh(\mathcal{C}) \to \textit{Sets}$、
$\mathcal{F} \to \mathcal{F}_p$ は完全である。
% P01-E0177
\end{enumerate}
\end{definition}
